% Options for packages loaded elsewhere
\PassOptionsToPackage{unicode}{hyperref}
\PassOptionsToPackage{hyphens}{url}
\PassOptionsToPackage{dvipsnames,svgnames,x11names}{xcolor}
\PassOptionsToPackage{space}{xeCJK}
%
\documentclass[
  11pt,
  oneside]{book}
\usepackage{amsmath,amssymb}
\usepackage{iftex}
\ifPDFTeX
  \usepackage[T1]{fontenc}
  \usepackage[utf8]{inputenc}
  \usepackage{textcomp} % provide euro and other symbols
\else % if luatex or xetex
  \usepackage{unicode-math} % this also loads fontspec
  \defaultfontfeatures{Scale=MatchLowercase}
  \defaultfontfeatures[\rmfamily]{Ligatures=TeX,Scale=1}
\fi
\usepackage{lmodern}
\ifPDFTeX\else
  % xetex/luatex font selection
  \setmainfont[BoldFont=DejaVuSerif-Bold.ttf,ItalicFont=DejaVuSerif-Italic.ttf,BoldItalicFont=DejaVuSerif-BoldItalic.ttf]{DejaVuSerif.ttf}
  \setmonofont[BoldFont=DejaVuSansMono-Bold.ttf,ItalicFont=DejaVuSansMono-Oblique.ttf,BoldItalicFont=DejaVuSansMono-BoldOblique.ttf]{DejaVuSansMono.ttf}
  \ifXeTeX
    \usepackage{xeCJK}
    \setCJKmainfont[BoldFont=NotoSerifCJKsc-Bold.otf,AutoFakeSlant=0.15]{NotoSerifCJKsc-Regular.otf}
          \fi
  \ifLuaTeX
    \usepackage[]{luatexja-fontspec}
    \setmainjfont[BoldFont=NotoSerifCJKsc-Bold.otf,AutoFakeSlant=0.15]{NotoSerifCJKsc-Regular.otf}
  \fi
\fi
% Use upquote if available, for straight quotes in verbatim environments
\IfFileExists{upquote.sty}{\usepackage{upquote}}{}
\IfFileExists{microtype.sty}{% use microtype if available
  \usepackage[]{microtype}
  \UseMicrotypeSet[protrusion]{basicmath} % disable protrusion for tt fonts
}{}
\makeatletter
\@ifundefined{KOMAClassName}{% if non-KOMA class
  \IfFileExists{parskip.sty}{%
    \usepackage{parskip}
  }{% else
    \setlength{\parindent}{0pt}
    \setlength{\parskip}{6pt plus 2pt minus 1pt}}
}{% if KOMA class
  \KOMAoptions{parskip=half}}
\makeatother
\usepackage{xcolor}
\usepackage[a4paper,left=24mm,right=24mm,top=25mm,bottom=25mm]{geometry}
\usepackage{longtable,booktabs,array}
\usepackage{calc} % for calculating minipage widths
% Correct order of tables after \paragraph or \subparagraph
\usepackage{etoolbox}
\makeatletter
\patchcmd\longtable{\par}{\if@noskipsec\mbox{}\fi\par}{}{}
\makeatother
% Allow footnotes in longtable head/foot
\IfFileExists{footnotehyper.sty}{\usepackage{footnotehyper}}{\usepackage{footnote}}
\makesavenoteenv{longtable}
\setlength{\emergencystretch}{3em} % prevent overfull lines
\providecommand{\tightlist}{%
  \setlength{\itemsep}{0pt}\setlength{\parskip}{0pt}}
\setcounter{secnumdepth}{-\maxdimen} % remove section numbering
\ifLuaTeX
\usepackage[bidi=basic]{babel}
\else
\usepackage[bidi=default]{babel}
\fi
\babelprovide[main,import]{english}
\ifPDFTeX
\else
\babelfont{rm}[BoldFont=DejaVuSerif-Bold.ttf,ItalicFont=DejaVuSerif-Italic.ttf,BoldItalicFont=DejaVuSerif-BoldItalic.ttf]{DejaVuSerif.ttf}
\fi
% get rid of language-specific shorthands (see #6817):
\let\LanguageShortHands\languageshorthands
\def\languageshorthands#1{}
% English mathematics edition; shared publication paths remain unchanged.
\usepackage{fvextra}
\usepackage{xurl}
\usepackage{etoolbox}
\usepackage{fancyhdr}
\setCJKmonofont{NotoSansMonoCJKsc-Regular.otf}[AutoFakeBold=1.5]
\setlength{\emergencystretch}{3em}
\setlength{\parindent}{1.8em}
\setlength{\parskip}{0.32em}
\linespread{1.12}
\renewcommand{\contentsname}{Contents}
\renewcommand{\figurename}{Figure}
\renewcommand{\tablename}{Table}
% Reserve enough room for three-digit page numbers in the mathematics contents.
\makeatletter
\renewcommand{\@pnumwidth}{2.4em}
\renewcommand{\@tocrmarg}{3.4em}
% Long translated headings may wrap without forced word spacing.
\patchcmd{\section}{\normalfont\Large\bfseries}{\normalfont\Large\bfseries\raggedright}{}{}
\patchcmd{\subsection}{\normalfont\large\bfseries}{\normalfont\large\bfseries\raggedright}{}{}
\patchcmd{\subsubsection}{\normalfont\normalsize\bfseries}{\normalfont\normalsize\bfseries\raggedright}{}{}
\renewcommand{\@tocrmarg}{3.4em plus 3em}
\patchcmd{\l@chapter}{\rightskip \@pnumwidth}{\rightskip \@pnumwidth plus 3em}{}{}
\makeatother
\DefineVerbatimEnvironment{verbatim}{Verbatim}{breaklines,breakanywhere,fontsize=\small}
\AtBeginEnvironment{longtable}{\small\setlength{\tabcolsep}{3pt}\renewcommand{\arraystretch}{1.15}}
\pagestyle{fancy}
\fancyhf{}
\fancyfoot[C]{\thepage}
\renewcommand{\headrulewidth}{0pt}
\fancypagestyle{plain}{\fancyhf{}\fancyfoot[C]{\thepage}\renewcommand{\headrulewidth}{0pt}}
\Urlmuskip=0mu plus 1mu
\widowpenalty=10000
\clubpenalty=10000
\ifLuaTeX
  \usepackage{selnolig}  % disable illegal ligatures
\fi
\usepackage{bookmark}
\IfFileExists{xurl.sty}{\usepackage{xurl}}{} % add URL line breaks if available
\urlstyle{same}
\hypersetup{
  pdftitle={UEOT Core: A Complete Mathematical Theory},
  pdfauthor={Degui Qian (钱德桂)},
  pdflang={en},
  colorlinks=true,
  linkcolor={MidnightBlue},
  filecolor={Maroon},
  citecolor={Blue},
  urlcolor={MidnightBlue},
  pdfcreator={LaTeX via pandoc}}

\title{UEOT Core: A Complete Mathematical Theory}
\usepackage{etoolbox}
\makeatletter
\providecommand{\subtitle}[1]{% add subtitle to \maketitle
  \apptocmd{\@title}{\par {\large #1 \par}}{}{}
}
\makeatother
\subtitle{From Object and Evolution to Prediction, Realization, Persistence, Composition, Purpose, and Learning}
\author{Degui Qian (钱德桂)}
\date{2026-11-11}

\begin{document}
\frontmatter
\maketitle

{
\hypersetup{linkcolor=MidnightBlue}
\setcounter{tocdepth}{1}
\tableofcontents
}
\mainmatter
\textbf{Edition and reading responsibilities.} This v4.1 bilingual academic edition revises the v4.0 text at main commit \nolinkurl{3d443ec29eb17af1fc00699a74e2d42d73516d81}, dated 2026-11-11. The original Core v3.0, the 106 P-IDs, the Lean sources, and their validation contracts retain their original identities. The edition audit records added reading guides, explicit hypotheses, historical wording, and bibliographic corrections. Automated structural checks and separated review do not replace external mathematical peer review. In the text, ``actual'' means an object genuinely constructed from the declared mathematical source or an executed program record; it does not by itself mean a natural-world experiment.

\textbf{Copyright and contributions.} © 2025--2026 Degui Qian and the respective rights holders. Expression already licensed under CC BY 4.0 remains available under that license; changing language or format does not revoke it. This edition grants no additional blanket license for the new v4.1 English translation or new editorial expression; rights remain with the respective holders under current repository policy and applicable law. This does not restrict historically licensed material. Third-party material retains its own terms. Mathematical ideas, theorems, and methods do not become copyright-exclusive through this notice. Both texts were prepared with native AI editorial assistance; authorship and project attribution remain unchanged. The complete boundary appears in \href{PUBLICATION_RIGHTS.md}{\nolinkurl{PUBLICATION_RIGHTS}}.

\section{Full Navigation}\label{full-navigation}

\begin{itemize}
\item
  \hyperref[core-theory]{Core synthesis: organization, action, and inheritance from a common source}
\item
  \hyperref[ch-0]{0. Scope of the Book and Mathematical Development of the Original Ideas}
\item
  \hyperref[ch-1]{1. Basic Conventions and Reusable Mathematical Tools}
\item
  \hyperref[ch-2]{2. World Processes, Experimental Protocols, and Time}
\item
  \hyperref[ch-3]{3. Factors, Covariance under Representation, and Object Candidates}
\item
  \hyperref[ch-4]{4. Canonical Predictive Identity: Existence, Minimality, Refinement, and Updating}
\item
  \hyperref[ch-5]{5. Predictive Information, Memory Complexity, and Compressibility}
\item
  \hyperref[ch-6]{6. Structured Openness, Environmental Interfaces, and Causal Boundaries}
\item
  \hyperref[ch-7]{7. Dynamical Closure, General Processes, and Multiscale Consistency}
\item
  \hyperref[ch-8]{8. Persistence: Invariant Structures, Viability, and Relative Identity}
\item
  \hyperref[ch-9]{9. Recovery, Perturbations, and Canonical Recovery Potentials}
\item
  \hyperref[ch-10]{10. Quasi-Stationarity, Spectra, and Persistence Windows}
\item
  \hyperref[ch-11]{11. Historical Identity, Transport, and Developmental Laws}
\item
  \hyperref[ch-12]{12. Object Certificates and Causal Integrity Ω}
\item
  \hyperref[ch-13]{13. Realization of Predictive Identity, Redundancy, and Damage Duality}
\item
  \hyperref[ch-14]{14. Resolution, Coarse-Graining, and the Uncertainty Interval for Microscopic Realizations}
\item
  \hyperref[ch-15]{15. Composition, Overlap, and the Formation of Parent Objects}
\item
  \hyperref[ch-16]{16. Certification from Finite Data: Responses, Carriers, and Protocols}
\item
  \hyperref[ch-17]{17. Identifiability of Inverse Objects, Dynamics, and Purposes}
\item
  \hyperref[ch-18]{18. Feasible Futures, Value, and Purpose Bridges}
\item
  \hyperref[ch-19]{19. Optimal Control and GOD}
\item
  \hyperref[ch-20]{20. How Object Quotients Preserve Decisions and Value}
\item
  \hyperref[ch-21]{21. GOA: Existence and Stability of Long-Term Structures}
\item
  \hyperref[ch-22]{22. Path Information Cost, Reachable Action, and Variational Realization}
\item
  \hyperref[ch-23]{23. Π/Φ, Constraint Geometry, and a Rigorous Dual-Drive Formulation}
\item
  \hyperref[ch-24]{24. Multilevel Value, Integrability, and Coordination Mechanisms}
\item
  \hyperref[ch-25]{25. Reproduction, Heredity, and Positive-Operator Evolution}
\item
  \hyperref[ch-26]{26. Transitions in Individuality and Endogenous Purpose Bridges}
\item
  \hyperref[ch-27]{27. Reflexive Structure, Belief Learning, and Agency}
\item
  \hyperref[ch-28]{28. Complete Interfaces, Algorithms, and Composition Contracts}
\item
  \hyperref[ch-29]{29. Application Method: From Models to Testable Certificates}
\item
  \hyperref[ch-30]{30. Complete Examples and Independent Executable Checks}
\item
  \hyperref[ch-31]{31. Unified Assembly: What the Core Can Jointly Establish}
\item
  \hyperref[ch-32]{32. Typed Unification: Eight Kinds of Hierarchy and Four Generating Mechanisms}
\item
  \hyperref[ch-33]{33. Relational Defects, Quantitative Calibration, and Gauge Covariance}
\item
  \hyperref[ch-34]{34. Complete Interfaces for Operational Objects, Purpose, and Conditional Lifecycles}
\item
  \hyperref[ch-35]{35. Prediction, Carriers, Viability, Repair, and Optimal Control on One Finite Source}
\item
  \hyperref[ch-36]{36. Actual Finite Paths, Posterior Decomposition, and Deployment Coverage}
\item
  \hyperref[ch-37]{37. Finite-Support Viability Under Partial Observation: An Equivalence Theorem and an Executable Check}
\item
  \hyperref[ch-38]{38. From Registered Sampling to Actual Joint Risk After Data-Dependent Selection}
\item
  \hyperref[ch-39]{39. Common-source constructions for execution, faults, reproduction, and purpose}
\item
  \hyperref[appendix-a]{Appendix A. Index of Symbols and Types}
\item
  \hyperref[appendix-b]{Appendix B. Register of Proved Results}
\item
  \hyperref[appendix-c]{Appendix C. Sources, Improvements Returned to the Core, and Standard Dependencies}
\item
  \hyperref[appendix-d]{Appendix D. Complete Executable Verification from v3.0 and Historical Output}
\item
  \hyperref[appendix-e]{Appendix E. Conventions for Maintenance, Extension, and Reproduction}
\item
  \hyperref[appendix-f]{Appendix F. Source Correspondences and Evidence Levels of the Maintained Edition}
\end{itemize}

\chapter{Core synthesis: organization, action, and inheritance from a common source}\label{core-theory}

This synthesis provides a self-contained entry to the present Core edition. Its conceptual starting point is the object--evolution problem in Degui Qian's original paper. It retains the completed system of conditional mathematics developed below, while explicitly revising universal assertions that exceeded the proofs. Chapters 0--39, the 106 stable P identifiers, and their formalization mappings remain in place. The synthesis places the mutually compatible main results before the developmental record. A definitive edition fixes the claims, assumptions, counterexamples, and responsibilities of interpretation; it does not establish empirical universality. Empirical support remains \textbf{UNVERIFIED}.

\section{I. Core objects and the stated outcome}\label{i.-core-objects-and-the-stated-outcome}

A system may be a predictive object without being able to repair itself; it may restore a local function without withstanding every form of damage; and a growth-optimal admissible control need not be selected with certainty in a finite population. In this edition, unification consists in giving these relations a common source and connecting them precisely. The starting point is not that everything that exists optimizes the same objective.

Six types of object must be distinguished. Physical states include actual programs, execution components, and materials. Admissible histories contain only the declared sensor readings and past receipts. Finite body types support conditional transitions and counting, whereas permanent identities support individual continuation and actual parent--offspring relations. An evaluation specifies what is compared; lineage weighting specifies how histories are sampled. Finally, normative value concerns what ought to be pursued and is not determined automatically by the sign of population growth.

The original Core has the following disposition. Full-history prediction and C1 minimal sufficiency, the convex hull of allowed actions and realizability on observation fibers, the source of recovery potentials, C5 rank conditions and gauge freedom, the Perron construction for finite primitive matrices, the declared continuous-process classes, and actual lifecycle and Price sources remain proved conditional results. Persistence does not universally imply a nonempty Ω within an arbitrary operation family: a constantly false failure predicate gives an empty family. A nondegenerate example has states \((x,u)\in\mathrm{Bool}^2\), the nonconstant object readout \(q(x,u)=x\), and allowed cuts that alter only \(u\). This operation family cannot destroy the declared readout capability. General two-dimensional drives, a unique global attractor, and unconditional progress are not theorems of this edition; the earlier stronger claims are explicitly narrowed. Bridges to subjective experience, reasons, and norms require independent philosophical arguments, developed in the companion philosophical book, and are not presented as Lean conclusions.

The new positive class must support a stronger explanation of organization: code and the main execution gate can be damaged, while recovery depends on surviving redundant information, a finite trusted primitive, and affordable material resources. Success, faithfulness, fitness, conditional means, and optimality are all derived from the common source described below. General probability and algebra lemmas may retain explicit assumptions; this concrete application derives those assumptions from its source instead of simply supplying them again.

\section{II. One actual source: components, wallet, and census}\label{ii.-one-actual-source-components-wallet-and-census}

A body has two rows of program code, each consisting of an active bit and two backup bits. The healthy type is \(H\) and the damaged type is \(D\). Damage flips the active bit in each row and closes the main execution gate. There are two actually compiled programs: \(\theta_0=(\mathrm{care},\mathrm{care})\) and \(\theta_1=(\mathrm{bloom},\mathrm{care})\). Their two rows correspond to the \(H/D\) diagnosis made before the current repair. Code \(0\) executes care and code \(1\) executes bloom. A finite bootstrap reads the majority of three bits, charges for writing a bit back, and charges for opening the gate. It does not store an answer specifying the correct program. After repair, execution rereads the actual active bit; it does not select an action from the audit label \(\theta\).

The sensor exposes the six code bits, the main gate, and the bootstrap state. The internal program and external request policies may use these observations and the allowed past. Hidden physical labels, audit identities, and world noise that has not yet occurred are unavailable to the controller. External care, bloom, and stop requests use the same transaction, repair, and payment rules as the internal programs; they confer no free permission to rewrite code.

At each census, an existing individual carries one unit of capital and receives an explicit environmental inflow of five units. Correcting one bit or opening the main gate costs one unit; care costs one unit and bloom costs two. Each actual continuing individual and each newborn receives one unit. Before reading noise, the transaction reserves capital for both possible outputs, giving the following request caps.

\begin{longtable}[]{@{}
  >{\raggedright\arraybackslash}p{(\columnwidth - 10\tabcolsep) * \real{0.1364}}
  >{\raggedleft\arraybackslash}p{(\columnwidth - 10\tabcolsep) * \real{0.1818}}
  >{\raggedleft\arraybackslash}p{(\columnwidth - 10\tabcolsep) * \real{0.1818}}
  >{\raggedleft\arraybackslash}p{(\columnwidth - 10\tabcolsep) * \real{0.1818}}
  >{\raggedleft\arraybackslash}p{(\columnwidth - 10\tabcolsep) * \real{0.1818}}
  >{\raggedright\arraybackslash}p{(\columnwidth - 10\tabcolsep) * \real{0.1364}}@{}}
\toprule\noalign{}
\begin{minipage}[b]{\linewidth}\raggedright
Body and request
\end{minipage} & \begin{minipage}[b]{\linewidth}\raggedleft
Repair cost
\end{minipage} & \begin{minipage}[b]{\linewidth}\raggedleft
Action cost
\end{minipage} & \begin{minipage}[b]{\linewidth}\raggedleft
Reserved output capital
\end{minipage} & \begin{minipage}[b]{\linewidth}\raggedleft
Total request cap
\end{minipage} & \begin{minipage}[b]{\linewidth}\raggedright
Actual treatment
\end{minipage} \\
\midrule\noalign{}
\endhead
\bottomrule\noalign{}
\endlastfoot
H / care & 0 & 1 & 2 & 3 & Execute \\
H / bloom & 0 & 2 & 2 & 4 & Execute \\
D / care & 3 & 1 & 2 & 6 & Execute \\
D / bloom & 3 & 2 & 2 & 7 & Reject \\
\end{longtable}

Rejection retains the existing body and its original wallet, returns the five units of environmental inflow, and does not read the current noise. Stop produces an empty population and returns all six units. Unused funds after execution return to the environment, including capital not allocated because of death. For each event, repair cost plus action cost plus actual output capital plus refund equals old capital plus environmental inflow. A newborn does not receive a free wallet reset.

One census tick processes each existing individual exactly once; newborns begin acting at the next tick. Continuation preserves identity, the original parent pointer, and birth time. A newborn receives a distinct identity pointing to its actual parent. Finite types and permanent identities are separate: a continuation edge is not a strict birth relation. This source explicitly uses a census clock, distinct from the single serial dispatch of the earlier Register source.

The primitive parameter is \(p=(s,\delta,b,\gamma,q)\). The survival categories have probabilities \((1-s,\delta,s-\delta)\), corresponding respectively to neither action retaining the individual, only care retaining it, and both actions retaining it. Birth categories have probabilities \((1-b-\gamma,\gamma,b)\), corresponding to neither action producing a newborn, only bloom doing so, and both doing so. Each of the two outputs is damaged with probability \(q\). These independent primitive categories yield 36 local atoms. The previously completed open positive class is

\[0<\delta<s<1,\quad 0<b,\quad b+\gamma<1,\quad
\delta<\gamma,\quad 0<q<1,\quad s+b>1.\]

This class is nonempty and parametrized, rather than an isolated numerical application. For example, \((3/4,1/8,3/8,1/4,1/3)\) belongs to it. Each existing individual contributes at most two successors per tick. Thus, for \(n\) founders, total environmental inflow through \(T\) is at most \(5n(2^T-1)\), in addition to \(n\) units of initial capital. Indefinite exponential growth depends on continued open supply; it does not arise from a closed finite total wallet.

The last parameter condition has a specific provenance. In the subcritical regime, rejection of an over-budget request preserves the original individual and may outperform continuation. The source was not changed to make rejection mean death, and the request was not removed from the comparison class. The positive class uses \(s+b>1\) and includes the unit row of rejection in the same upper bound. The subcritical counterexample is retained: at \((s,\delta,b,\gamma,q)=(1/4,1/8,1/8,3/16,1/3)\), a damaged body's next-step total mean under care is \(3/8\), whereas rejection gives \(1\).

\section{III. Paid regeneration, formation, and limits of recovery}\label{iii.-paid-regeneration-formation-and-limits-of-recovery}

For both programs and both body types, every allowed damage operation preserves the majority code. Actual recovery restores the two active bits and the main gate. A damaged initial body pays 3, leaving exactly 3 for care and the maximum capital allocation to two outputs. The same restored program then executes the row corresponding to the original diagnosis. Actual outputs copy the restored code and subsequently undergo the declared output damage. Hence the offspring's decodable program equals the parent's program as a consequence of the copying and damage rules, rather than an additional certificate of hereditary success.

Include the balance in the repair state, construct the deterministic kernel from the actual paid recovery stage, and apply the existing \nolinkurl{StrongRepairable} and \texttt{repairRank} interfaces. A healthy source state has rank 0 and a damaged source state rank 1. This rank describes the physical repair stage before noise is read. The subsequent survival and birth event can still kill the individual. The conclusion is therefore finite paid recovery, not eventual permanent health or zero extinction probability.

The causal role of recovery is directly testable. Closing the main execution gate makes raw execution return no output. Flipping only the active bit in the H row currently read by a healthy program changes its raw action. When the bootstrap is cut, a damaged state becomes a self-loop outside the target of the repair kernel. The same occurs with a fixed balance of only 2 and no additional supply during this stage; strong repairability then fails. This does not assert that recovery remains impossible after sufficient resources arrive later. Complete code erasure makes the two original programs physically indistinguishable. Without an additional correct copy or a distinguishing history, no deterministic restorer reading only that state can restore both distinct targets. Private randomness likewise cannot reliably reconstruct two opposite answers from an identical input. These are information and execution limits, not merely restatements of a failed certificate.

The formation bridge uses the four actual initial states obtained from the two compiled programs and H/D types, with the ineffective tag fixed to false; p still ranges over the entire positive class. It starts from actual component readouts and the actual newborn-response law. The complete readable physical carrier determines the finite response, so zero deficiency is derived rather than assumed. Physical component regions cover that carrier, and the existing finite minimization and child-cover theorem then produce a nonempty family of formation candidates. This one-birth response family is nonconstant: the healthy second program has a different birth rate from the other three states. It neither separates every pair of program/body states nor equals full-future predictive identity. An actual bootstrap cut changes the damaged body's successor capability. The result asserts neither a unique natural boundary nor that a population process with positive death probability satisfies an earlier structure requiring certain eternal persistence.

\subsection{Minimal representations for a one-birth response and finite futures}\label{task-dependent-representation}

Fix the same parameter \(p\) in the preceding positive class, internal program execution, the six-unit transaction rule, and the ineffective tag false. Compare only the four initial states of a single founder, \(X=\{\theta_0H,\theta_0D,\theta_1H,\theta_1D\}\). Let \(B_t\) be the \textbf{number of strict births produced by the entire actual population} at census \(t\ge1\), retaining death and zero-population outcomes. Let \(R_1\) be the repair cost actually paid in the founder's first transaction. Successors continue to execute the original program faithfully, each carrying one unit of capital and receiving the next environmental inflow under the original rules.

For a specified task, two initial states are equivalent if all records of that task have the same distribution. A minimal representation here means the coarsest sufficient partition of these four initial states, not a uniquely chosen physical carrier.

\begin{longtable}[]{@{}
  >{\raggedright\arraybackslash}p{(\columnwidth - 4\tabcolsep) * \real{0.3000}}
  >{\raggedright\arraybackslash}p{(\columnwidth - 4\tabcolsep) * \real{0.3000}}
  >{\raggedleft\arraybackslash}p{(\columnwidth - 4\tabcolsep) * \real{0.4000}}@{}}
\toprule\noalign{}
\begin{minipage}[b]{\linewidth}\raggedright
Protected task
\end{minipage} & \begin{minipage}[b]{\linewidth}\raggedright
Exact equivalence classes
\end{minipage} & \begin{minipage}[b]{\linewidth}\raggedleft
Number of classes
\end{minipage} \\
\midrule\noalign{}
\endhead
\bottomrule\noalign{}
\endlastfoot
Only \(B_1\) & \(\{\theta_0H,\theta_0D,\theta_1D\}\); \(\{\theta_1H\}\) & 2 \\
For every finite \(T\), the joint records \((B_1,\ldots,B_T)\) & \(\{\theta_0H,\theta_0D\}\); \(\{\theta_1D\}\); \(\{\theta_1H\}\) & 3 \\
The preceding task together with \(R_1\), or the full public observation record & Four singleton classes & 4 \\
\end{longtable}

\textbf{Proof.} In the first period, \(\theta_0H,\theta_0D,\theta_1D\) all actually execute care, so \(B_1\sim\mathrm{Bernoulli}(b)\). State \(\theta_1H\) executes bloom, so \(B_1\sim\mathrm{Bernoulli}(b+\gamma)\). Since \(\gamma>0\) on the positive class, the first row has exactly two classes.

Couple \(\theta_0H\) and \(\theta_0D\) on the same primitive noise tree. Both execute care initially, and their code agrees after repair. For each local noise atom, their actual \nolinkurl{retained/newborn} body slots coincide, and each actual output receives one unit of capital. First-period repair costs, refunds, and the source type in the audited ancestry may differ, but the internal program does not read these differences. Assigning the same subsequent noise to corresponding body slots produces the same slots and birth counts at every step. Thus the entire birth-count record agrees pathwise for every finite \(T\), not merely in expectation.

The other initial states cannot be merged further. State \(\theta_1H\) is already separated in the first period. For \(\theta_1D\) and \(\theta_0D\), both of which initially execute care, the first-period total successor mean is \(s+b\), with healthy and damaged means \((s+b)(1-q)\) and \((s+b)q\). In the next period, the second program has birth rate higher by \(\gamma\) when healthy, whereas both programs have rate \(b\) when damaged. Therefore

\[\begin{aligned}
E_{\theta_0D} B_2&=(s+b)b,\\
E_{\theta_1D} B_2&=(s+b)\{b+(1-q)\gamma\},\\
E_{\theta_1D} B_2-E_{\theta_0D} B_2&=(s+b)(1-q)\gamma>0.
\end{aligned}\]

The second-period birth count is at most two. Different means therefore imply different record laws, proving both sufficiency and coarseness of the three classes in the second row. At the rational parameter above, the two second-period means are \(27/64\) and \(39/64\), differing by \(3/16\).

Finally, for \(\theta_0H\) and \(\theta_0D\), the cost \(R_1\) is identically 0 and 3, respectively. Their caps are 3 and 6, respectively; both are accepted under the same six-unit rule, and payment occurs before noise is read. The third row therefore has four classes. Full public observations also reveal the distinct program codes and H/D types, so they induce the same four singleton classes. Adding repair cost only to the one-step task \((B_1,R_1)\) still leaves the two damaged programs indistinguishable; that smaller task must not be said to yield four classes. \(\blacksquare\)

No sufficient representation can merge initial states with different record laws, whereas encoding the stated classes is itself sufficient. This establishes a complete classification. It concerns \textbf{birth prediction under fixed internal programs}. It does not shrink the original class of externally allowed control policies to preserve an earlier optimality conclusion, or promote a one-step formation response to the full-task minimal state of C1. Source execution, actual successors, and first moments have Lean endpoints. The full 2/3/4 classification was supplied in the finite hardening as a written proof with independent exact checks; it was not registered as a new Lean theorem.

\section{IV. Mechanism-induced order and realizable growth control}\label{iv.-mechanism-induced-order-and-realizable-growth-control}

Let \(\nu=(1-q,q)\), \(k_c=s+b\), and \(k_b=s-\delta+b+\gamma\). Summing the actual continuation and birth records over all 36 atoms gives

\[L_a(i,j)=s_a\nu_j,\qquad B_a(i,j)=b_a\nu_j,\qquad
M_a(i,j)=(s_a+b_a)\nu_j.\]

Here \(s_c=s\), \(s_b=s-\delta\), \(b_c=b\), and \(b_b=b+\gamma\). Stop gives the zero row; actual rejection of D/bloom gives the unit row \(e_D\). Costs enter the process through admissibility, repair, and actual capital allocation, not by subtracting an arbitrary price from the population matrix. Failure and death atoms all contribute to the mean; one must not condition on success first.

The continuation menus are \(A(H)=\{c,b\}\) and \(A(D)=\{c\}\). Sensors actually distinguish H/D, each body's wallet is recorded separately, and every selector can be compiled into the two-row physical program. Row replacement is therefore implementable. Every continuation matrix is strictly positive, so primitivity follows from the source parameters.

\textbf{Finite spectral optimality theorem.} For a general finite implementable family of row actions, if every continuation selector has a primitive matrix, then a selector with maximal spectral radius exists and its right Perron vector satisfies the Bellman upper bound in every row. The proof uses only the existing \nolinkurl{certificateOfPrimitive}: choose a maximal root among the finitely many selectors. If a row could strictly improve the action on the positive right vector, replace that row and pair with the positive left vector of the replacement matrix. This would give a strictly larger spectral radius, a contradiction. The maximal root, positive vector, and optimality are obtained by construction rather than supplied as inputs to the construction.

In the earlier positive class, \(k_b>k_c\), and the explicit program \(\theta_1\) is a maximizer. Its actual matrix and spectral radius are

\[M_* = \begin{pmatrix}k_b(1-q)&k_bq\\k_c(1-q)&k_cq\end{pmatrix},\qquad
\lambda_*=(1-q)k_b+qk_c.\]

The explicitly computable positive right eigenvector is \(r=(k_b,k_c)\), on the same ray as the vector returned by the existing Perron construction for this matrix. The abstract finite maximum may be noncomputable; the actual returned artifact is the encoded and executed H/bloom, D/care program. These are distinct claims.

Every request still satisfies \(m(i,a)r\le\lambda_*r_i\). The zero row of stop is immediate. The unit rejection row of D/bloom uses \(\lambda_*>1\). Arbitrary real-probability request mixtures depending on admissible histories inherit the bound. Private probabilities need not be rational, and individuals' requests need not be independent. Current requests must be chosen before world noise, and hidden fields of the complete audit cannot become control information.

Summing actual successor counts over parents yields the conditional first moment of the actual population, and hence

\[W_t=\lambda_*^{-t}\sum_{u\in\mathcal P_t}r_{i_t(u)}\]

is nonnegative at every finite admissible history node, with nonincreasing next-step conditional mean. Equality holds when the same source executes optimal requests. We refer to this as a \nolinkurl{supermartingale/martingale} property in the finite conditional-kernel sense. Expectations at each finite T are recursively generated by the same actual branching kernel. This application does not additionally claim a Mathlib \texttt{Martingale} instance on a common infinite path space. Conditional laws at finite histories are defined from actual request branches and noise atoms; conditional means are not assumed independently. With \(r_{\min}>0\) and \(r_{\max}<\infty\), a nonempty initial population gives the two-sided bound for optimal expected population size

\[\frac{\lambda_*^t\langle Z_0,r\rangle}{r_{\max}}
\le E N_t\le
\frac{\lambda_*^t\langle Z_0,r\rangle}{r_{\min}}.\]

Every other admissible history policy satisfies the same upper bound. An empty initial state yields only the zero process. The result concerns optimal exponential growth of expected population size. For the same healthy individual, care has one-step identity-survival probability \(s\), whereas bloom has \(s-\delta\). Thus the source already supplies an actual counterexample to identifying growth optimality with optimal one-step survival of the same identity. It does not compare infinite-lineage survival probabilities. One consistent evaluation is successor reproductive value discounted by \(1/\lambda_*\); the conditional \nolinkurl{equality/inequality} is then its optimality principle. This growth theorem does not replace the earlier registered birth-minus-write-cost objective J, arbitrary short-term objectives, or normative values.

\subsection{Exact growth and internal attainment on the strictly positive probability domain}\label{core-growth-classification}

The preceding positive class gave sufficient conditions under which regeneration and growth optimality hold together. We now relax only the parameter assumptions, retaining the same bodies, two internal programs, public observations, six-unit transactions, stop, actual rejection, and faithful inheritance. The basic probability domain is

\[\mathsf{ProbabilityValid}(p):\quad
0<\delta<s<1,\quad b>0,\quad \gamma>0,\quad
b+\gamma<1,\quad 0<q<1.\]

These conditions make all 36 original local noise atoms strictly positive. The earlier Valid condition is exactly the basic domain together with \(\gamma>\delta\) and \(k_c>1\); its definition and theorems remain unchanged. This extension does not enlarge the domain of the selection-feedback constants or cohort contract in Part V.

\textbf{Growth objective and comparison class.} Fix \(p\) in the basic domain and any finite nonempty canonical initial population \(S_0\), whose members may carry different original programs and silent tags. Let \(N_t\) be actual population length at census \(t\), counting all members after continuation and strict birth. The comparison class remains the complete original class of admissible history policies: requests are chosen before current noise, may use the full public receipt history, and may involve real probabilities and correlated joint randomization across members of a generation. Policies are compared at the given parameter; unknown-parameter learning is not additionally solved.

A policy has positive exponential factor \(\beta\) if there are constants independent of \(t\), with \(0<L\le U<\infty\), such that for every \(t\ge0\), \(L\beta^t\le E N_t\le U\beta^t\). This two-sided contract implies \((E N_t)^{1/t}\to\beta\). Arbitrary admissible policies need not admit such a bound; they are compared by \(G(\pi)=\limsup_{t\to\infty}(E_\pi N_t)^{1/t}\), with a positive root of zero taken to be zero. Exponential optimality does not require maximizing population size at every finite horizon.

The complete menu is obtained from actual transactions, including requests other than continuation actions:

\begin{longtable}[]{@{}
  >{\raggedright\arraybackslash}p{(\columnwidth - 6\tabcolsep) * \real{0.2500}}
  >{\raggedright\arraybackslash}p{(\columnwidth - 6\tabcolsep) * \real{0.2500}}
  >{\raggedright\arraybackslash}p{(\columnwidth - 6\tabcolsep) * \real{0.2500}}
  >{\raggedright\arraybackslash}p{(\columnwidth - 6\tabcolsep) * \real{0.2500}}@{}}
\toprule\noalign{}
\begin{minipage}[b]{\linewidth}\raggedright
Current body
\end{minipage} & \begin{minipage}[b]{\linewidth}\raggedright
Mean row for care
\end{minipage} & \begin{minipage}[b]{\linewidth}\raggedright
Mean row for bloom
\end{minipage} & \begin{minipage}[b]{\linewidth}\raggedright
Mean row for stop
\end{minipage} \\
\midrule\noalign{}
\endhead
\bottomrule\noalign{}
\endlastfoot
H & \(k_c(1-q,q)\) & \(k_b(1-q,q)\) & \((0,0)\) \\
D & \(k_c(1-q,q)\) & \((0,1)\), original rejection & \((0,0)\) \\
\end{longtable}

Set \[k_*=\max(k_c,k_b),\qquad
\lambda_0=k_c,\qquad \lambda_1=(1-q)k_b+qk_c,\qquad
\mu=(1-q)k_*+qk_c.\]

Since \(0<q<1\), we have \(\lambda_1-\lambda_0=(1-q)(\gamma-\delta)\) and \(\mu=\max(\lambda_0,\lambda_1)\).

\textbf{Theorem (exact classification of growth over all policies).} On the basic domain and from a nonempty initial state,

\[\boxed{\rho_*=\max(1,\mu)=\max\{1,k_c,(1-q)k_b+qk_c\}.}\]

A deterministic policy reading only the original public diagnosis attains \(\rho_*\). One finite constant gives an expected-population upper bound for every admissible history policy and every horizon. Moreover, \(\rho_*\) is the least positive exponential base that can appear in such a common upper bound. Consequently, \(\sup_\pi G(\pi)=\rho_*\). The positive \(t\)th root of the optimal expected population size at each horizon also converges to this value. These last two assertions follow by squeezing between the common upper bound and the lower bound of a single attaining policy; no unproved interchange of supremum and limit is used.

\textbf{Proof.} At H, choose the action with the larger of \(k_c,k_b\), choosing care in a tie if desired. Construct strictly positive weights \(r=(r_H,r_D)\) in two cases.

If \(\mu\ge1\), take \(r=(k_*,k_c)\) and choose care at D. A continuation row has weighted value equal to its total output rate times \(\mu\). The selected H and D rows therefore equal \(\mu r_H,\mu r_D\), respectively; all other continuation actions are bounded above by these values. The D/bloom rejection row has value \(k_c\le\mu k_c\), and stop has value zero. Hence every original action satisfies the same upper bound.

If \(\mu<1\), put \(d=1-k_*(1-q)\). Since \(d>qk_c>0\), choose

\[r_D=1,\qquad r_H=\frac{k_*q}{d}>0.\]

Request bloom at D, causing actual rejection under the original rules. The average output-type weight is \((1-q)r_H+qr_D=q/d\), so the selected H row equals \(r_H\), and the selected D rejection row equals \(r_D=1\). The D/care weighted value is \(k_cq/d<1\), all other H rows are no greater than the selected row, and stop remains zero. Thus factor \(1\) dominates every action, with equality for the selected rows. This argument neither requires \(k_*<1\) nor changes rejection into death.

Let \(V_t=\sum_{x\in\mathcal P_t}r_{\operatorname{diagnostic}(x)}\). At a given public history, first sum over original noise and then mix with any joint request kernel. Linearity turns the per-member row bounds into \(E[V_{t+1}\mid\text{admissible history}]\le\rho_*V_t\). Equality holds for the attaining policy. This uses the original actual population list and branching kernel, rather than replacing them with an independent mean process. Induction over finite histories and \(r_{\min}N_t\le V_t\le r_{\max}N_t\) give

\[E_\pi N_t\le\frac{V_0}{r_{\min}}\rho_*^t,\qquad
\frac{V_0}{r_{\max}}\rho_*^t
\le E_{\pi_*}N_t\le
\frac{V_0}{r_{\min}}\rho_*^t.\]

Nonemptiness gives \(V_0>0\). If a smaller positive base also provided a common upper bound for all policies, it would have to dominate the attaining policy's positive lower bound. But positive integer powers of the ratio of the bases are unbounded, a contradiction. An empty initial state gives only the zero-population process, to which the least-base and positive-lower-bound conclusions do not apply. This proves the theorem.

\textbf{Necessary and sufficient conditions for the two existing internal programs.} For each \(\theta\), in a nonempty population carrying that program, the public program selector agrees branch by branch with original internal execution. The original inheritance invariant preserves this fact in subsequent generations. Program weights are \(r^\theta_i=k_\theta(i)>0\), and their actual rows satisfy \(m(i,\theta(i))r^\theta=\lambda_\theta r^\theta_i\). The same finite-history argument therefore gives positive two-sided exponential bounds for that internal program. A single founder also satisfies the original actual output-tree formula \(E_{\theta,i}N_t=k_\theta(i)\lambda_\theta^{t-1}\) for \(t\ge1\).

The factor in a positive two-sided exponential bound is unique. Hence

\[\boxed{
\begin{aligned}
\theta_0\text{ attains all-policy optimality}
&\iff \gamma\le\delta\ \text{and}\ k_c\ge1,\\
\theta_1\text{ attains all-policy optimality}
&\iff \gamma\ge\delta\ \text{and}\ \lambda_1\ge1.
\end{aligned}}\]

\begin{longtable}[]{@{}
  >{\raggedright\arraybackslash}p{(\columnwidth - 4\tabcolsep) * \real{0.3333}}
  >{\raggedright\arraybackslash}p{(\columnwidth - 4\tabcolsep) * \real{0.3333}}
  >{\raggedright\arraybackslash}p{(\columnwidth - 4\tabcolsep) * \real{0.3333}}@{}}
\toprule\noalign{}
\begin{minipage}[b]{\linewidth}\raggedright
Parameter region
\end{minipage} & \begin{minipage}[b]{\linewidth}\raggedright
An attaining policy in the complete original menu
\end{minipage} & \begin{minipage}[b]{\linewidth}\raggedright
The two existing internal programs
\end{minipage} \\
\midrule\noalign{}
\endhead
\bottomrule\noalign{}
\endlastfoot
\(\gamma>\delta,\ \lambda_1\ge1\) & H/bloom, D/care & \(\theta_1\) attains; \(\theta_0\) does not \\
\(\gamma<\delta,\ k_c\ge1\) & H/care, D/care & \(\theta_0\) attains; \(\theta_1\) does not \\
\(\gamma=\delta,\ k_c\ge1\) & Both H continuation actions tie; D/care & Both attain \\
\(\max(\lambda_0,\lambda_1)<1\) & Select the larger H row; actual D/bloom rejection & Neither attains factor \(1\) \\
\end{longtable}

For each program, take its own damaged canonical founder. The same body first actually pays three units for repair and then executes care; its original remaining balance pays for the action and reserved outputs. Successors continue faithfully to carry its code. The conditions therefore classify exactly when each program has a joint witness of paid recovery and optimality. The programs have different backup codes and are not the same literal body. Repair itself holds throughout the basic domain; what can fail is growth optimality of continued execution of the specified program.

For example, \(p=(3/5,1/10,3/10,1/2,1/5)\) gives \(k_c=9/10<1\) and \(\lambda_1=\rho_*=61/50>1\). Thus the earlier condition \(s+b>1\) is not necessary for attainment by \(\theta_1\). Starting at D, rejection maximizes one-period population size with value \(1\). The two-period optimum first uses care and then chooses by the final diagnosis, giving \(279/250\). The fixed \(\theta_1\) has two-period value \(549/500\); the difference \(9/500\) comes from rejection versus care at D in the final period. At the critical point \(p=(2/3,1/8,1/4,1/4,1/3)\), we instead have \(\lambda_1=\rho_*=1\). From the same D founder, internal \(\theta_1\) has expected population size \(11/12\) at every positive horizon, whereas a D founder that is always rejected retains population size \(1\) pathwise. Nevertheless, the two have the same exponential factor.

The subcritical conclusion does not exclude every policy that repairs. An admissible policy can first pay for care once on the original D body and then switch to the larger H row/D rejection policy above, still obtaining a positive expected lower bound and factor \(1\). Long-run exponential growth, finite-horizon population size, current identity survival, and the occurrence of repair are therefore distinct evaluation questions. The advantage of rejection depends on this source's explicit identity-retention and refund rules; it is not a general prescription for natural systems.

The formalization has three new modules. \nolinkurl{CoreRegenerativeGrowthClassification} proves the basic domain, source rows, positive weights, and algebraic rate classification. \nolinkurl{CoreRegenerativeGrowthPopulation} uses the original complete Policy, State, population recursion, and actual programs. \nolinkurl{CoreRegenerativeGrowthAttainment} proves the least uniform base, the necessary and sufficient attainment conditions, and the link to repair of the same damaged founder. The Lean growth contract consists of two-sided constant bounds at every finite horizon and the least base. The \nolinkurl{positive-root/limsup} formulations above are written analytical consequences, not separately claimed Lean limit endpoints. Eight rational probes provide exact cross-checks of the transaction translation; they replace neither the general parameter proof nor empirical evidence. Complete source identity, endpoints, and proportionate validation appear in the \href{../formalization/ueot-core/docs/compression/theory_completion/unified_closure/core_growth_classification_v1/README.md}{growth-classification delivery package}.

\section{V. Actual inheritance feedback and marked lineages}\label{v.-actual-inheritance-feedback-and-marked-lineages}

This part continues to use the original Valid positive class and the separately registered cohort contract. Part IV does not automatically extend the domain of the feedback constants. Both programs are inherited faithfully, so the full mean operator on \((\theta,i)\) is block diagonal by program. No free mutation or immigration is introduced to make the full matrix primitive. The spectral radii of the body-type blocks are

\[\lambda_0=k_c,\qquad \lambda_1=(1-q)k_b+qk_c>\lambda_0.\]

For any \(T\ge1\) and initial body type \(i\), the actual program output tree gives the single founder's exact expected descendant count as \(k_\theta(i)\lambda_\theta^{T-1}\). The tree's two branches are the actual \nolinkurl{retained/newborn} slots. Noise slots that produce no individual are not read by the controller; death contributes zero. At horizon \(T\), a single founder yields at most \(2^T\) individuals and at most \(2^T-1\) actual processing events, so the first and second moments are finite.

For a cohort consisting of one healthy \(\theta_0\) and one healthy \(\theta_1\), the lower program's share of normalized means tends to zero. At the fixed parameter point, the ratio of the means is

\[\frac{\mu_{0,T}}{\mu_{1,T}}=\frac9{10}\left(\frac{27}{29}\right)^{T-1}.\]

This is feedback arising from actual inheritance events, but it is not yet a statement about random proportions. The new source class explicitly has \(\theta\in\{0,1\}\), and the selection application uses the lower program's count. This two-program application is not advertised as a newly proved theorem for arbitrary program families with tied maximal blocks. For \(n\ge1\) actual independent identically structured cohorts, let \(X_{\theta,T},X_T\) be a single cohort's counts and let \(\mu_{\theta,T}=E X_{\theta,T}\), \(\mu_T=E X_T>0\). For \(0<\varepsilon\le1\), the actual aggregate counts satisfy

\[\begin{aligned}
&P\!\left(N_T=0\ \text{or}\ \left[N_T>0\ \text{and}\
\left|Z_{\theta,T}/N_T-\mu_{\theta,T}/\mu_T\right|>\varepsilon\right]\right)\\
&\qquad\le\frac{9\{\operatorname{Var}(X_{\theta,T})+\operatorname{Var}(X_T)\}}
{n\varepsilon^2\mu_T^2}.
\end{aligned}\]

The proof bounds both count deviations from their means by \(\varepsilon n\mu_T/3\), which makes the denominator at least \(2n\mu_T/3\). Finite second moments and the centered second-moment identity for independent product atoms give the constant 9. Zero population is explicitly a failure outcome and cannot be removed by a convention for total division at zero. Combining mean-share concentration with this result, first choose a finite T so that the lower program's mean share is at most \(\varepsilon/2\), and then choose n.~The actual optimal-program share is then at least \(1-\varepsilon\) with the specified probability, with constant 36 in the corresponding bound. Here n counts cohorts, hence there are 2n founders; initial capital plus supply through T is at most \(n[2+10(2^T-1)]\). Increasing n requires these resources. If common environmental noise destroys cohort independence, the bound cannot be used unchanged.

A fixed finite population has no resulting guarantee of certain selection of the optimum over infinite time. In this same-source example, the growth-optimal program has first-period extinction probability \(9/64\); the probability that it is extinct while the lower program remains nonzero after that period is \(243/2048\). This does not prove that the lower program survives forever. The joint zero-population probability is \(45/2048\). Extending the high-program proportion by zero at zero population gives expectation \(2095/4096\); conditioning on nonzero population gives \(2095/4006\); normalizing expected counts gives \(10/19\). These are three different quantities. Faithful copying also cannot create a program outside the initial support.

Lineages require a separate weighting. For the same optimal matrix, \(Q_{ij}=M_{*,ij}r_j/(\lambda_*r_i)\) is a row-stochastic kernel. If path tests read \nolinkurl{continuation/birth} marks, use the marked kernel with \(Q^L_{ij}=L_{*,ij}r_j/(\lambda_*r_i)\) and \(Q^B_{ij}=B_{*,ij}r_j/(\lambda_*r_i)\). The finite-path many-to-one identity starts from the actual parent pointer and mark of each successor and follows by induction on the same noise tree:

\[E\!\left[\sum_{u\in\mathcal P_T}r_{i_T(u)}F(\tau_u)\right]
=\lambda_*^T r_{i_0}\,E^Q_{i_0}[F(Y_0,\ldots,Y_T)].\]

An unmarked F reads the body-type path; a marked F requires the marked kernel. A continuation edge preserves an identity across censuses; only a birth edge is a strict parent--offspring edge. Finitely many founders are mixed by their initial reproductive values. Population mean recursion alone does not imply this lineage identity.

The source's normalized mean body composition is the left Perron vector \(l=\nu\). With the unnormalized \(r=(k_b,k_c)\) used above, the weighted lineage stationary law is \(\nu_i k_i/\lambda_*\). Only after the additional normalizations \(\sum_i l_i=1\) and \(\sum_i l_ir_i=1\) may it be abbreviated as \(l_ir_i\). At the fixed parameter, every weighted lineage row is \((20/29,9/29)\), unlike the ordinary body distribution \((2/3,1/3)\). Actual finite-population proportions are controlled by the probability theorem above; the native controller's kernel is determined by actual requests and execution. These are distinct mathematical types. Q is only a statistically weighted lineage law. Its executability by original actions requires the existing convex-hull and observation-fiber criteria separately.

\section{VI. Same-source representations and a common witness}\label{vi.-same-source-representations-and-a-common-witness}

One actual noninjective quotient of this source removes only an unobservable physical tag that has no effect on transactions. It preserves the six bits, main gate, bootstrap, allowed requests, costs, refunds, \nolinkurl{continuation/birth} marks, and actual identities. For every primitive noise atom, execution followed by projection gives the same complete receipt as projection followed by execution. Finite-path compatibility follows inductively: initial permitted observations and earlier receipt histories agree; the same admissible history kernel gives the same request distribution; the same primitive noise produces the same next receipt; and the argument continues. This preserves the joint law of resource events and counts, with the first-moment intertwining identity as a consequence. Machine checking additionally constructs projections of actual population State and Branch, proves equality of branch weights for the same policy and commutation with next, and inducts on finite horizons to obtain exact pushforward for any terminal test. For tag-invariant tests, the full admissible policy classes in both representations attain the same values. Actual population counts, extinction events, and complete payment-receipt histories have direct applications.

Write the original finite type as \((i,\mathrm{tag})\), the quotient as \(q(i,\mathrm{tag})=i\), and \(H_{xj}=\mathbf1_{q(x)=j}\). Tags are faithfully carried; the extended matrix does not silently resample them. We have \(L H=H\bar L\), \(B H=H\bar B\), and \(M H=H\bar M\). Each fixed-tag body block is primitive, whereas the full extended matrix is reducible. Its positive right vector can be \(r\circ q\), obtained from the quotient vector; the primitive-matrix construction is not incorrectly applied to the whole matrix. Both representations deliver exactly the same admissible observation histories. The same history-action kernel therefore lifts in both directions with the same comparison class, a stronger result than a one-way map for one policy.

Deleting the H/D diagnosis would erase the difference between affordable H/bloom and unaffordable D/bloom. Retaining only decoded content while erasing the main gate would make code prediction identical although raw execution succeeds in one state and fails in the other. Ordinary predictive equivalence therefore does not universally preserve organization, permissions, or control value.

The joint witness of regeneration and growth optimality on the original Valid class is precisely the same \(\theta_1\), the same damaged initial state, and the same six-unit ledger. It first pays three units to restore its two active code bits and execution gate, then executes care in the diagnosed row; when healthy later, the same actual code executes bloom. The population continues to inherit that code. A repair result for one program is not combined with optimality of another, and there is no additional free repair wallet.

\section{VII. Boundaries of proof, execution, and interpretation}\label{vii.-boundaries-of-proof-execution-and-interpretation}

This construction mainly removes three kinds of independent assumptions. Maintainability follows from local redundancy and paid repair; behavioral order follows from actual continuation and birth; and finite-population feedback follows from the actual inheritance tree and independent cohorts. Concrete applications use the existing Perron, repair, formation, admissible-kernel, and probability algebra. The contribution is the connection of these sources and constraints, not a UEOT relabeling of standard theorems.

Formal endpoints occur in \nolinkurl{CoreRegenerativeSource}, \nolinkurl{CoreRegenerativeRepair}, \nolinkurl{CoreRegenerativeFormation}, \nolinkurl{FiniteSourceGrowthControl}, \nolinkurl{CoreRegenerativeControlPopulation}, \nolinkurl{SourceCohortFeedback}, \nolinkurl{CoreRegenerativeFeedback}, \nolinkurl{CoreRegenerativePopulation}, \nolinkurl{CoreRegenerativeLineage}, and \nolinkurl{CoreRegenerativeQuotient}. The public common-program entry is \nolinkurl{CoreTheoryClosure}. Same-source execution calls these source definitions. An independent rational checker checks the complete actual exported records; it does not replace parametrized proofs. The \nolinkurl{core_theory_closure_v1} delivery package binds source identity, compilation, axioms, and execution receipts. The two previously disclosed \nolinkurl{native_decide} numerical trust dependencies in PR443 retain their original endpoint scope. A shared import graph does not imply that all new endpoints depend on them.

The positive conclusions are limited to the nonempty source classes and general assumptions stated in each part. The growth classification in Part IV uses the broader basic probability domain; feedback and other applications retain their own earlier contracts. Universal autopoiesis outside these classes, invention of correct code from total erasure without information, unlimited growth from a closed finite wallet, certain optimization by a fixed finite population, universal two-dimensionality in nature, a unique GOA, and moral obligation inferred from growth are not implicit conclusions. Strong claims refuted by counterexamples are publicly withdrawn or narrowed. Normative and subjective arguments retain their own burden of justification, and empirical calibration remains UNVERIFIED. The developmental chapters below retain the precise content and historical status of each stage; the present Core should be read through this synthesis and the matching final delivery receipts.

\chapter{0. Scope of the Book and Mathematical Development of the Original Ideas}\label{ch-0}

\textbf{Chapter question.} What does this mathematical architecture undertake to explain, and which interpretations still require independent evidence?

\textbf{Reading route.} First relate the original research questions to explicit mathematical objects, then distinguish the responsibilities of definitions, proofs, standard theorems, mechanism assumptions, and empirical conclusions. Each connection below must be read under its own hypotheses; completeness primarily means a complete statement of the declared scope.

This book presents the complete general mathematical account of UEOT (Universal Evolutionary Object Theory). It begins with the questions posed in the original paper: how an object persists through changes in its material and state; how several processes constitute an individual; through which interfaces an object interacts with its environment; and how function, purpose, learning, and evolution can be expressed rigorously within one architecture.

The book reconstructs the results that remain valid from Mathematical Foundations I--IV, Unified Mathematics v1.1--v2.1, the general realization interfaces, and subsequent audits as a continuous exposition. Definitions, proofs, assumptions, counterexamples, algorithmic interfaces, and application procedures are stated here. Understanding the main argument does not require first reading an earlier edition or relying on an unspecified ``theorem from the previous version.'' This edition excludes quantum-specific theory, particle selection, and field-theoretic chapters. Mathematics developed in subsequent branches returns to this core only when its proofs do not require those specialized structures.

Here, ``complete'' means \textbf{a complete statement of the mathematical architecture and application interfaces declared in this edition}. It does not mean that every open problem has been solved, or that all laws of nature have been deduced from a few philosophical propositions. Foundational results on measure extension, measurable selection, spectral theory, and related subjects are explicitly identified as standard theorems. Proofs are provided for the principal connecting results of this theory.

The definitive argument is presented first in the seven parts of the opening ``Definitive Core Synthesis'': organizational regeneration, actual population sources, mechanism-induced ordering, admissible control, inheritance feedback, and legitimate representation are connected within a single source. The provenance, legitimate assumptions, and acceptance evidence for the original conditional contracts appear in §§39.70--39.74 and Appendix F.5. Historical source audits and the open status of earlier stages are retained as records of their respective dates; adoption of the actual release is governed by the final receipts.

\section{0.1 Original Concepts and Their Mathematical Realizations}\label{original-concepts-and-their-mathematical-realizations}

\begin{longtable}[]{@{}
  >{\raggedright\arraybackslash}p{(\columnwidth - 4\tabcolsep) * \real{0.3333}}
  >{\raggedright\arraybackslash}p{(\columnwidth - 4\tabcolsep) * \real{0.3333}}
  >{\raggedright\arraybackslash}p{(\columnwidth - 4\tabcolsep) * \real{0.3333}}@{}}
\toprule\noalign{}
\begin{minipage}[b]{\linewidth}\raggedright
Original question or concept
\end{minipage} & \begin{minipage}[b]{\linewidth}\raggedright
Mathematical realization in this edition
\end{minipage} & \begin{minipage}[b]{\linewidth}\raggedright
Essential qualification
\end{minipage} \\
\midrule\noalign{}
\endhead
\bottomrule\noalign{}
\endlastfoot
Being is continuing Becoming & Processes, historical branches, and validated developmental transport & Persistence does not require an unchanging state \\
Ω-loop and relational identity & Predictive and causal closure, \nolinkurl{recovery/invariant} structures, and mechanism-damage hypergraphs & A geometric loop is a particular mechanism \\
Structured openness and specific coupling & Conditional-kernel factorization through an internal factor and an environmental interface & Resource exchange is not equivalent to loss of identity \\
Type I / Type II & Certificate profiles for passive persistence, endogenous maintenance, and adaptive capacity & No binary classification is imposed on all systems \\
Π/Φ dual drives & Independently anchored resource and information constraints; a low-dimensional response hypothesis & Numerical decomposition and two-dimensionality require separate tests \\
GOD & The optimal-action correspondence for a specified evaluation problem & Usually set-valued; not identical to an instantaneous gradient \\
GOA & Long-term sets, occupation distributions, or conditional-survival structures of an optimal closed loop & A unique fixed point is only a special case \\
Child object and parent object & A parental predictive core, cut irreducibility, and fidelity at the child level & Strong correlation alone does not establish a parent object \\
Self-modification and freedom & An augmented structural state, policy expansion, and the capacity to revise objectives & Expansion has costs and conditions of failure \\
Meaning & Functional contribution, behavioral evaluation, and normative bridging expressed separately & Obligations do not follow directly from stability \\
\end{longtable}

The original paper supplies the research questions and architectural choices. Mathematical theorems establish conditional implications. The two are closely related, but philosophical motivation is not itself a proof.

\section{0.2 The Reading Sequence}\label{the-reading-sequence}

The exposition proceeds as follows:

\begin{enumerate}
\def\labelenumi{\arabic{enumi}.}
\tightlist
\item
  Chapters 1--7: general processes, measurable factors, canonical predictive identity, information and interfaces, and dynamical closure.
\item
  Chapters 8--12: persistence and recovery, quasi-stationary structures, historical development, object certificates, and Ω.
\item
  Chapters 13--17: realization, damage, resolution, composition, finite samples, and inverse problems.
\item
  Chapters 18--24: feasible futures, value, control, GOD/GOA, information costs, dual drives, and hierarchical coordination.
\item
  Chapters 25--27: inheritance, selection, transitions in individuality, endogenous purpose, and reflexive learning.
\item
  Chapters 28--31: complete interfaces, application methods, examples, and cross-module assembly theorems.
\item
  Chapters 32--38: typed hierarchies, defect calibration, object and scientific ports, execution from the same source, finite-support survival, and actual joint risk after training-dependent selection.
\item
  Chapter 39: constructions from a common source for execution, faults, reproduction, continuous processes, decision representations, and unknown mechanisms; thematic entry points are given in the chapter guide.
\item
  Appendices: notation, the 106 results and full source index, original numerical programs, historical records, and maintenance conventions.
\end{enumerate}

The \textbf{direction of dependence} is ``world and protocols → candidate representation → \nolinkurl{predictive/causal/historical} certificates → realization and integrity,'' and, after an evaluation mechanism has been added, ``feasible policies → value → GOD → GOA.'' Selection may change controllers and the structure of future objects, but the mechanism model for each generation must first be specified. A logical circle cannot replace feedback unfolding in time.

\section{0.3 Status of Statements}\label{status-of-statements}

\begin{itemize}
\tightlist
\item
  \textbf{D}: a definition fixing a mathematical object or operational semantics.
\item
  \textbf{P}: a result proved in this book; this does not assert that UEOT originated the result.
\item
  \textbf{K}: an invoked standard theorem, with its applicable assumptions and role stated.
\item
  \textbf{H}: a mechanism, domain, or low-dimensionality hypothesis requiring independent specification or testing.
\item
  \textbf{E}: a certificate supportable by actual data or a specified experiment.
\item
  \textbf{N}: a counterexample or nonidentifiability result.
\item
  \textbf{O}: an open problem, excluded from the chain of proved dependencies.
\end{itemize}

Each principal positive result has a stable identifier, such as P-PRED-01. When assumptions prove insufficient, the statement and its dependents must be revised; an untested ``universal principle'' cannot be added to preserve the conclusion.

\section{0.4 From Maintenance Edition 3.1 to This Edition: Self-Contained Theory and Formal Implementation}\label{from-maintenance-edition-3.1-to-this-edition-self-contained-theory-and-formal-implementation}

Maintenance edition 3.1 carries forward all 106 stable P identifiers in Chapters 0--31 and the appendices, retaining their statements, proofs, standard dependencies, algorithms, and complete numerical code. New connections use chapter-local identifiers such as M-32.1. They neither redefine existing P-IDs nor eliminate domain assumptions by introducing an overarching theorem.

Maintenance edition 3.1 added Chapters 32--38 to express, as continuous mathematical exposition, the completed Lean developments on hierarchy typing, compression by four generators, quantitative defect calibration, operational objects and conditional lifecycles, control and repair from the same source, finite-support viability, and joint risk after data-dependent selection. Each chapter supplies the central proofs, counterexamples establishing the need for assumptions, and precise Lean correspondences. The propositions can be checked without first accepting a philosophical interpretation.

Three senses of ``complete'' must be distinguished: the exposition is self-contained; the propositions within the declared scope have proofs; and a concrete world model has experimental support. The first does not imply either of the others. The 106/106 source-endpoint status, CLOSED status of the P0--P12 work packages, and package closure of C1--C7 are also different facts. The current compressed accounting comprises 11 generating endpoints, 89 retained adapters, and 6 retained boundaries; see Chapter 32. The explicit boundary of P12 and the unconfirmed real-world support of the scientific interfaces are addressed separately in Chapter 34.

The 2,371 numerical checks and their outputs in Appendix D are retained as the historical v3.0 record. The Lean checks added at that earlier stage form a separate body of evidence; see §38.8. This v4.1 publication revision did not rerun those historical checks and does not present old program output as a new execution.

\chapter{1. Basic Conventions and Reusable Mathematical Tools}\label{ch-1}

\textbf{Chapter question.} How can a probability or error symbol retain the same meaning throughout the later chapters?

\textbf{Reading route.} Fix the spaces, probability kernels, pushforwards, total variation, and units before using the two recurring estimates: a common readout cannot increase discrepancy, and the span of a bounded function controls its expectation difference. Average error, worst-case error, and sampling confidence remain distinct.

\section{1.1 Spaces, Kernels, and Notation}\label{spaces-kernels-and-notation}

Unless otherwise stated, random variables take values in standard Borel spaces. Where topology, limits, or compactness are needed, Polish spaces, compact metric spaces, or smooth manifolds are specified explicitly. The standard Borel assumption ensures the usual existence of regular conditional probabilities. It does not make arbitrary quotient spaces or selection maps automatically measurable.

A probability kernel is written \[
K:\mathcal X\longrightarrow\mathcal P(\mathcal Y),\qquad
K(x,B)=P(Y\in B\mid x).
\] For a measurable map \(f:\mathcal Y\to\mathcal Z\), the pushforward kernel is \[
(f_\#K)(x,A)=K(x,f^{-1}A).
\] Probability distributions are row vectors and state functions are column vectors. Thus distributions evolve as \(\mu P\) and functions as \(Pv\).

We write \(H_t\) for a history, \(T\) for a time horizon, \(U_t\) for an environmental interface, \(V\) for a finite carrier universe, and \(\mathcal V\) for a domain used to assess persistence. The discount factor is \(\beta\in(0,1)\), the absorption rate is \(\lambda\), and the confidence failure probability is \(\alpha\). Object scale is denoted by \(s\) and the \nolinkurl{selection/evaluation} level by \(\ell\).

\section{1.2 Information, Distances, and Units}\label{information-distances-and-units}

Natural logarithms are used throughout unless stated otherwise; entropy and KL divergence are measured in nats. Examples explicitly marked in bits use base 2. Constants must change accordingly when the logarithm base changes. \[
D_{\rm TV}(P,Q)=\sup_A|P(A)-Q(A)|.
\] On a finite space, \(D_{\rm TV}=\frac12\sum_y|p_y-q_y|\). If \(P\ll Q\), define \[
D_{\rm KL}(P\|Q)=\int\log\frac{dP}{dQ}\,dP;
\] otherwise its value is \(+\infty\). We use \(0\log0=0\).

\hypertarget{p-met-01}{}\phantomsection\label{p-met-01}

\textbf{P-MET-01: Data processing for kernels and readouts.} For a measurable map \(f\) and a probability kernel \(K\), \[
D_{\rm TV}(f_\#P,f_\#Q)\le D_{\rm TV}(P,Q),\qquad
D_{\rm TV}(PK,QK)\le D_{\rm TV}(P,Q).
\] If \(f\) is a bimeasurable bijection, equality holds in the first inequality.

\textbf{Proof.} The first inequality restricts events to inverse images \(f^{-1}A\). For the second, express the difference as an integral against \(0\le K(x,A)\le1\) and use the bounded-function dual formulation of TV. For a bijection, apply the first inequality to \(f^{-1}\) to obtain the reverse inequality. \(\blacksquare\)

KL data processing is \textbf{K-INF-01}: applying a common stochastic channel cannot increase KL divergence. It follows from the conditional KL chain rule. For finite variables, this chain rule follows directly from the multiplicative factorization of a joint probability; the standard Borel version uses regular conditional kernels and the decomposition of relative entropy.

\hypertarget{p-met-02}{}\phantomsection\label{p-met-02}

\textbf{P-MET-02: The span bound.} For a bounded real function \(v\), write \(\operatorname{sp}(v)=\sup v-\inf v\). Then \[
|Pv-Qv|\le\operatorname{sp}(v)D_{\rm TV}(P,Q).
\] \textbf{Proof.} The integral difference of a constant is zero. Subtracting \((\sup v+\inf v)/2\) leaves a function whose supremum norm is half the span. Apply \(|\int g\,d(P-Q)|\le2\|g\|_\infty D_{\rm TV}(P,Q)\). \(\blacksquare\)

\textbf{K-INF-02: Pinsker's inequality.} \[
D_{\rm TV}(P,Q)\le\sqrt{\tfrac12D_{\rm KL}(P\|Q)}.
\] The constant shown corresponds to natural logarithms. The information chain rule and this inequality underlie Chapters 5, 6, and 22; they are not new UEOT postulates.

\section{1.3 Three Kinds of Error Must Remain Distinct}\label{three-kinds-of-error-must-remain-distinct}

\begin{itemize}
\tightlist
\item
  Pointwise or worst-case error: a bound valid for every declared reachable input.
\item
  Average error: a bound under a specified input distribution \(\mu\).
\item
  Estimation error: a bound holding with sampling probability at least \(1-\alpha\).
\end{itemize}

Average error is not worst-case error; error under the training distribution is not error under the deployment distribution; and a small error observed in one simulation is not a confidence upper bound. Chapter 16 states conditions under which these notions can be related legitimately.

\section{1.4 Rules for Measurable Quotients}\label{rules-for-measurable-quotients}

The ``set of equivalence classes under equality of kernels'' is not assumed to be a well-behaved measurable space. States are preferably defined as \nolinkurl{probability-measure-valued} random elements, or as sub-\(\sigma\)-algebras on a fixed underlying probability space. Conditions needed for sections, quotient manifolds, continuous selections, or weak limits are added explicitly. Pointwise results on finite spaces are kept distinct from results modulo null sets on general spaces.

\section{1.5 The Same Relation, the Same Metric, and the Same Evidence}\label{the-same-relation-the-same-metric-and-the-same-evidence}

A single zero-error relation may have many quantitative representations. For example, the discrete defects \(d(x,y)=\mathbf1_{x\ne y}\) and \(d'(x,y)=100d(x,y)\) have the same zero set, but the budgets \(d\le1\) and \(d'\le1\) are not equivalent. Real-valued quantities used in lossy quotients, developmental tubes, and control costs must therefore carry an explicit metric and scale.

Chapter 33 transfers budgets through calibrations \(a d\le d'\le b d\); retaining only the zero relation does not justify such numerical transfer. For probabilities, one must also distinguish the training measure \(\mu\), the conditional law of the future given a training record, and the occupation law \(\nu_n\) of the state--action pairs actually visited by that future process. These typically live on different spaces and cannot be substituted merely because all are called ``distributions.''

The Lean correspondences are \nolinkurl{MinCore.RelationDefectPresentation} and \nolinkurl{DefectCalibration}, together with \nolinkurl{UnifiedClosure.selectedTrainingDeployment} and \nolinkurl{coverageDeploymentRowLaw.} Full names and provenance appear in Appendix F.

\chapter{2. World Processes, Experimental Protocols, and Time}\label{ch-2}

\textbf{Chapter question.} In what order are actions, observations, and futures generated, and when do they belong to the same experiment?

\textbf{Reading route.} Construct the path law from the initial distribution and causal kernels, then declare the information available to each protocol. History augmentation makes a general process Markovian but does not provide finite memory. Deployment selected after training must retain its conditional order of generation.

\section{2.1 General Controlled Processes}\label{general-controlled-processes}

\textbf{D-WORLD.} At time \(t\), the history \[
H_t=(X_0,A_0,X_1,\ldots,A_{t-1},X_t)
\] takes values in the history space \(\mathcal H_t\). The environment, sensor records, internal structure, and other quantities may be components of \(X_t\). The next step is specified by a kernel \[
q_t(dx'\mid h,a)
\] and a policy is a causal kernel \(\pi_t(da\mid h)\). An initial distribution is specified first, followed by successive composition of the policy and state kernels.

\textbf{K-PROC-01: Existence of a path law.} For a countable time index, standard Borel spaces, and measurable probability kernels, the Ionescu--Tulcea extension theorem supplies a unique corresponding path probability law. This establishes the probabilistic consistency of the specified process model; it does not establish that the model describes a natural system.

\hypertarget{p-proc-01}{}\phantomsection\label{p-proc-01}

\textbf{P-PROC-01: Histories as states.} Every history-dependent process of the above form is a time-homogeneous Markov control process on the augmented state \(Z_t=(t,H_t)\).

\textbf{Proof.} Given \((t,h)\) and an action \(a\), sample \(x'\sim q_t(\cdot\mid h,a)\). The next augmented state is uniquely \((t+1,h,a,x')\). Its law does not depend on records preceding the augmented state. \(\blacksquare\)

This construction may grow without bound. The ability to make a process Markov does not mean that a useful finite memory has been found.

\section{2.2 Continuous Time}\label{continuous-time}

A continuous-time module uses one of the following explicitly selected models:

\begin{itemize}
\tightlist
\item
  a continuous semiflow \(\phi_t\);
\item
  a strong \nolinkurl{Markov/Feller} semigroup \(P_t\);
\item
  a finite-state generator \(L=(l_{xy})\);
\item
  a diffusion with appropriate existence, uniqueness, nonexplosion, and integrability assumptions, \[
  dX_t=b(X_t,a_t)\,dt+\sigma(X_t,a_t)\,dW_t.
  \]
\end{itemize}

The path space may be the continuous-path space \(C([0,T],\mathcal X)\) or an appropriate càdlàg path space. An arbitrary ``set of histories'' with nonmeasurable selections is not automatically a standard path model.

\section{2.3 Protocols and Interventions}\label{protocols-and-interventions}

\textbf{D-PROT.} A protocol specifies observation maps, implementable actions, readout order, time windows, stopping rules, a family of future tests, and implementation constraints. For a protocol \(i\) and a history \(h\), the induced record kernel is \(K_i(h,\cdot)\).

Unless stated otherwise, subsequent interventions take place after the history has already been formed, so different protocols share a historical reference distribution \(\mu\). If an experiment changes the mechanism generating histories, distinct environmental reference distributions must be introduced, a common structural kernel or dominating measure used, and the support of the comparison specified anew.

The observational conditional probability \(P(Y\mid A=a)\) and the interventional response \(P(Y\mid do(a))\) are different inputs. Connecting them requires a randomized experiment, a verifiable intervention model, or sufficient causal identification assumptions.

\section{2.4 Finite Protocols, Complete Futures, and Countability}\label{finite-protocols-complete-futures-and-countability}

Mathematically, a ``complete future'' is a declared future-record random element, not every fact about the universe. For an infinite future, a countable determining class of cylinder sets requires countable time and countably generated measurable structures.

The canonical prediction theorem in this core uses a finite or countable family of protocols. Uncountable action spaces may be used in control, but defining a canonical state simultaneously through all actions requires a common pointwise structural kernel with measurability, or a countable dense family together with continuity. Countable generation of the future space does not replace this safeguard.

\section{2.5 The Generative Order in Conditional Deployment}\label{the-generative-order-in-conditional-deployment}

Suppose a training record \(\omega\sim\mu\) is observed first and a model or policy \(s(\omega)\) is then selected. The declared conditional deployment kernel \(K_{s(\omega)}\) must generate the future \(h\). In the finite case, the joint mass is

\[J(\omega,h)=\mu(\omega)K_{s(\omega)}(h).\]

This is an actual construction. Its training marginal is \(\mu\), and its future marginal is \(\sum_\omega\mu(\omega)K_{s(\omega)}\). Selection may depend on training; the corresponding future-risk statement does not require selection to be independent of training. What must be verified is that, conditional on the training record, the future is indeed generated by this kernel.

Agreement of the future marginals for fixed candidates does not imply conditional agreement after data-dependent selection. Chapter 38 proves the tower identity for the actual joint law and then explains why an incorrect conditional law can invalidate the risk conclusion. In this way, ``learning → control → future'' is implemented in temporal order, without storing the desired risk equality as a certificate field.

\chapter{3. Factors, Covariance under Representation, and Object Candidates}\label{ch-3}

\textbf{Chapter question.} What remains the same predictive object after a change of descriptive coordinates?

\textbf{Reading route.} An embedded factor records the information retained from the same history; invertible recoding changes its expression. The proof preserves path laws by transporting dynamics, observations, and rewards together. Physical carriers, minimal coordinate counts, and object instances require separate judgments.

\section{3.1 Embedded Factors}\label{embedded-factors}

\textbf{D-FAC.} For fixed \((\mathcal H,\mathscr H,\mu)\) and a measurable statistic \(M=f(H)\), define \[
\mathscr A_f=\sigma(f)\subseteq\mathscr H.
\] Identify subalgebras modulo \(\mu\)-null sets on the same underlying space, with order \(\mathscr A_f\preceq\mathscr A_g\iff\mathscr A_f\subseteq\mathscr A_g\).

For standard-Borel-valued statistics, if \(f\) is measurable with respect to \(\sigma(g)\), Doob--Dynkin factorization gives the corresponding almost-sure version of \(f=\varphi\circ g\). The order therefore means that the latter statistic contains the information needed to recover the former.

An invertible recoding of a statistic does not change its embedded factor. The factors generated by two independent uniform bits are abstractly isomorphic but are different embedded subalgebras. This distinction retains object instances, not merely object types.

\section{3.2 Covariance under Representation}\label{covariance-under-representation}

\hypertarget{p-fac-01}{}\phantomsection\label{p-fac-01}

\textbf{P-FAC-01.} Let \(T:\mathcal X\to\mathcal X'\) be a bimeasurable bijection, and transport \[
P'(Tx,a,T B)=P(x,a,B),\qquad f'=f\circ T^{-1},
\] together with observations, interventions, and rewards. Then macroscopic path laws, predictive sufficiency, exact dynamical closure, and control values are preserved.

\textbf{Proof.} We have \(f'(Tx)=f(x)\). The one-step pushforwards coincide; successive conditioning and induction give equality of all finite path-record laws. Equality of path laws gives equality of the distributions of every declared future record. If rewards are present, their expectations agree term by term; taking the supremum over policies gives the same value. \(\blacksquare\)

This does not say that an arbitrary lossy sensor preserves an object. A lossy map satisfies only data processing.

\section{3.3 Factors, Coordinates, and Carriers}\label{factors-coordinates-and-carriers}

For example, \(Y_1=X_1\oplus X_2,\ Y_2=X_2\) is an invertible recoding. In the new coordinates, recovering \(X_1\) requires both \(Y_1,Y_2\). The minimum number of coordinates changes under recoding, while the predictive relation is correctly transported.

A counterexample in the other direction is provided by the same bit stored on two different pieces of hardware. The two may generate the same predictive factor while having different physical failure mechanisms. A predictive factor therefore does not uniquely determine a physical location. Realization is treated separately in Chapter 13.

\section{3.4 Factor Spectra and Scale}\label{factor-spectra-and-scale}

For a specified certificate, collect all embedded factors that pass it and order them by inclusion to obtain a factor spectrum. Independent processes \(A,B\) and the joint process \((A,B)\) may all be predictively closed. Unless a target, compression criterion, or composition condition is specified, dynamics does not supply a unique partition into objects.

Minimality of representation, minimality of carriers, physical inseparability, and irreducibility of a parent object are distinct questions. A parent object may contain genuine child objects; this does not conflict with the least sufficient state for a fixed prediction task.

\section{3.5 From Predictive Factors to an Organizational Interface}\label{from-predictive-factors-to-an-organizational-interface}

Preservation of a collection of responses by a factor does not imply preservation of every organizational predicate. The actual common-source core first registers observation, organization, \nolinkurl{currentProgram}, \nolinkurl{currentResource}, and safe as outputs to be retained, and then performs stable refinement for the same finite transition and reward structure. These attributes are recoverable in the quotient for this reason, not merely because the quotient is called an ``object.''

In \nolinkurl{OrganizationalCoreAssembly}, \nolinkurl{organizational_terminal_preserves_registered_interfaces} and \nolinkurl{organizational_terminal_is_coarsest} prove attribute preservation and coarseness relative to the specified initial distinctions, respectively. The physical meaning of the organizational labels still requires modeling, and coarseness is relative to the registered outputs, actions, rewards, and actual kernel. Chapter 35 connects this interface to survival and safe control on an actual source.

\chapter{4. Canonical Predictive Identity: Existence, Minimality, Refinement, and Updating}\label{ch-4}

\textbf{Chapter question.} What information must history retain for all the declared tests?

\textbf{Reading route.} Collect the conditional future laws of the protocols into a state. First show that every sufficient representation recovers it, then show that it is itself sufficient. Countability supplies a common probability-one set; recursive updating additionally requires the same causal source and closure under admissible continuation, not merely equal one-step marginals.

\section{\texorpdfstring{4.1 A \nolinkurl{Probability-Measure-Valued} Canonical State}{4.1 A  Canonical State}}\label{a-probability-measure-valued-canonical-state}

Let \(H\sim\mu\), let \(\mathcal I=\{i_1,i_2,\ldots\}\) be the protocol family, and suppose all future spaces are standard Borel; use a finite product for a finite protocol family. Fix measurable response kernels \(K_i(h,\cdot)\) and define \[
C_{\mathcal I}(h)=\bigl(K_{i_1}(h,\cdot),K_{i_2}(h,\cdot),\ldots\bigr).
\] Equip \(\mathcal P(\mathcal Y)\) with the evaluation \(\sigma\)-algebra. A countable product is again standard Borel. We call \(C_{\mathcal I}\) the canonical interventional predictive state. When the future record is the complete declared path, it is called FFIPS.

For an object candidate \(f\), write \(Y_f^+\) for its future and \(C^{f}_{\mathcal I}\) for its corresponding state. \textbf{An object need retain only the information required by its declared future and interfaces; it need not retain the future of the entire world.} Global FFIPS is a common comparison device and cannot be identified directly with an individual object.

\section{4.2 The Minimality Theorem}\label{the-minimality-theorem}

\hypertarget{p-pred-01}{}\phantomsection\label{p-pred-01}

\textbf{P-PRED-01: Canonical minimality for common protocols.} Let \(S=s(H)\) be a standard-Borel-valued statistic. If for every \(i\) there is a kernel \(L_i\) such that \[
K_i(H,\cdot)=L_i(S,\cdot)\quad\mu\text{-a.s.},
\] then there is a measurable \(\varphi\) such that \[
C_{\mathcal I}=\varphi(S)\quad\mu\text{-a.s.}
\] Moreover, \(C_{\mathcal I}\) is itself sufficient for every protocol. Thus \(\sigma(C_{\mathcal I})\) is the least embedded factor required by the common responses.

\textbf{Proof.} Every future space has a countable class of events determining probability measures. Intersect the probability-one sets for all protocols and their determining events. This countable intersection is a common probability-one set, on which \[
C_{\mathcal I}=(L_{i_1}(S,\cdot),L_{i_2}(S,\cdot),\ldots),
\] The right-hand side is a measurable function of \(S\). Conversely, each \(K_i\) is a coordinate of \(C_{\mathcal I}\). For a future event \(B\), the tower property gives \[
E[K_i(H,B)\mid C_{\mathcal I}]=K_i(H,B),
\] so the future kernel conditional on the canonical state equals the kernel conditional on the complete history. \(\blacksquare\)

Any other canonical minimal statistic satisfying the same requirements generates the same embedded factor. This conclusion is relative to the specified protocols, future, and reference distribution. It does not guarantee that the natural world admits a unique ultimate inventory of objects.

\section{4.3 The Obstruction for Uncountable Families}\label{the-obstruction-for-uncountable-families}

\textbf{N-PRED-01.} Let \(H\) be uniform on \([0,1]\), and let the actual future under every protocol be the constant zero. For each \(i\in[0,1]\), choose the following version of a conditional law: \[
K_i(h,\{1\})=\mathbf1_{\{h=i\}}.
\] Each version differs only on a null singleton, so a constant statistic is sufficient separately for each protocol. Yet the pointwise collection of all coordinates distinguishes every \(h\). Consequently, ``almost everywhere separately'' cannot be replaced without further assumptions by ``simultaneously on a common set.''

Possible extensions require a common structural kernel explicitly defined at every reachable point; or a countable dense protocol family with responses continuous in the protocol, where continuity and the chosen versions hold on a common probability-one set. Equality on the dense family then extends by continuity. If protocols have different historical distributions, common support must also be specified explicitly.

\section{4.4 Target Completeness and Refinement}\label{target-completeness-and-refinement}

\hypertarget{p-pred-02}{}\phantomsection\label{p-pred-02}

\textbf{P-PRED-02: Target pushforwards and protocol refinement.} If \(Y_1=T(Y_2)\), their canonical predictive kernels satisfy \[
C_1=T_\#C_2.
\] If \(\mathcal I_1\subseteq\mathcal I_2\), then \[
\sigma(C_{\mathcal I_1})\subseteq\sigma(C_{\mathcal I_2}).
\] For a countable increasing family, \[
\sigma(C_{\cup_n\mathcal I_n})=\bigvee_n\sigma(C_{\mathcal I_n}).
\]

\textbf{Proof.} A target transformation applies the same pushforward to each conditional kernel. The coordinates of the smaller protocol family are projections of those of the larger family. Every coordinate of the union family appears at some finite stage, so the two sides generate the same \(\sigma\)-algebra. \(\blacksquare\)

This also gives monotone refinement for finite future windows \(T_1<T_2\). Any finite plateau may still be broken by the next untested protocol: two histories can be constructed to agree on the first \(k\) tests and differ only on test \(k+1\).

\section{4.5 Recursive Updating of a Full-Future State}\label{recursive-updating-of-a-full-future-state}

\hypertarget{p-pred-03}{}\phantomsection\label{p-pred-03}

\textbf{P-PRED-03: Recursive updating for finite-alphabet processes.} Suppose all responses arise from one causal process. Assume the protocol family is closed under action prefixes and admissible continuations after positive-probability observations, and includes the corresponding joint future cylinder events. Then predictive equivalence is preserved by the next update, and there is a measurable version of \[
C_{t+1}=\mathcal U_t(C_t,a_t,o_{t+1})
\]

\textbf{Proof.} For the next observation \(o\) and any admissible continuation-record event \(B\), the subsequent conditional probability is \[
P(B\mid h,a,o)
=\frac{P(o,B\mid h,a)}{P(o\mid h,a)}.
\] The positive denominator and the numerator are protocol probabilities retained by the current canonical state. If \(h,h'\) have the same canonical state, all continuation kernels agree after updating. Finite alphabets and countably many kernel coordinates give measurability of the update. Zero-probability branches may be extended arbitrarily as mathematical kernels, but such extensions do not thereby receive empirical certification. \(\blacksquare\)

Equality of one-step marginals alone does not ensure recursive sufficiency. If the protocol does not allow subsequent adaptive operations, a state constructed for that protocol cannot be used directly for a larger class of control policies.

\section{4.6 The Formal Implementation Boundary of Recursive Sufficiency}\label{the-formal-implementation-boundary-of-recursive-sufficiency}

Sections 39.7 and 39.12 construct posterior-response recursion, an actual past--future joint law, and a minimal predictive consumer modulo null sets from a declared nonempty standard Borel latent state, the original P/O kernels, finite or countable protocols, and an actual complete-history program. The original joint law of the past and filter requires only countably generated observations. The jointly measurable ambient response update and the fresh/memory predictive kernels constructed from response arrays retain the assumptions of standard Borel observations and a countable action space with measurable singletons. When control memory is inherited, both response and memory are retained; the corresponding canonical-only conclusion applies only when a new future program starts from the registered empty memory. Models outside this scope remain subject to the conditions of this chapter and the protection for uncountable families in §4.3. These results do not infer an entire source automatically from finite data.

\nolinkurl{PredictionUpdate}, the Lean source for P-PRED-03, proves factorization of a Bayes update when the numerator and denominator already factor through the current state, and proves joint measurability under finite \nolinkurl{actions/observations} and the corresponding coordinate assumptions. It does not automatically produce sufficient coordinates from arbitrary data. Parameterized disintegration on general spaces is addressed by C1 in §34.6.

A finite controlled source also permits a weaker but useful construction: the support of the terminal microscopic state conditional on a coarse history. Chapter 37 proves from the actual historical conditional distribution that this support is precisely traceSupport, recursively updated using the recorded actions and the new observation. This support recursion expresses only which states remain possible; the full posterior weights still depend on probabilities. The distinction between positive-probability and zero-probability histories cannot be omitted.

\chapter{5. Predictive Information, Memory Complexity, and Compressibility}\label{ch-5}

\textbf{Chapter question.} What is lost when history is compressed, and when is a memory lower bound unavoidable?

\textbf{Reading route.} The information chain rule separates retained information from residual information, which is then converted into predictive distance. Zero-distortion and distinguishability bounds identify impossible compressions. They do not directly quantify computation time, energy use, or a unique information drive.

\section{5.1 The Information Retention Identity}\label{the-information-retention-identity}

Let \(M=f(H)\). Fix an experimental distribution for the moment, and assume the relevant mutual informations are finite.

\hypertarget{p-info-01}{}\phantomsection\label{p-info-01}

\textbf{P-INFO-01.} \[
I(H;Y)=I(M;Y)+I(H;Y\mid M).
\] Hence, if \(I(H;Y\mid M)\le\varepsilon\), \[
I(M;Y)\ge I(H;Y)-\varepsilon.
\] If \(M\) is discrete, \[
H(M)\ge I(H;Y)-\varepsilon.
\]

\textbf{Proof.} Since \(M\) is a function of \(H\), \(I(H,M;Y)=I(H;Y)\). Expand the joint mutual information in the other order, then use \(I(M;Y)\le H(M)\). \(\blacksquare\)

If the future \(Y=H\) copies a uniform variable with \(N\) values, an exactly sufficient state requires at least \(N\) values and entropy at least \(\log N\). No object-discovery algorithm can guarantee nontrivial compression of every system.

\section{5.2 From Information Residuals to Predictive Error}\label{from-information-residuals-to-predictive-error}

\hypertarget{p-info-02}{}\phantomsection\label{p-info-02}

\textbf{P-INFO-02.} For standard Borel variables \(H,Y,M,U\), suppose \(M,U\) are functions of the history and set \[
\varepsilon=I(H;Y\mid M,U).
\] Then \[
E D_{\rm TV}\bigl(P(Y\mid H),P(Y\mid M,U)\bigr)
\le\sqrt{\varepsilon/2}.
\]

\textbf{Proof.} Conditional mutual information is the average KL divergence between the two conditional kernels. Apply Pinsker's inequality conditionally and then Jensen's inequality for the square root. \(\blacksquare\)

\textbf{Corollary.} If two representations \(Z_1,Z_2\) have information residuals at most \(\varepsilon_1,\varepsilon_2\), respectively, then \[
E D_{\rm TV}(P(Y\mid Z_1),P(Y\mid Z_2))
\le\sqrt{\varepsilon_1/2}+\sqrt{\varepsilon_2/2}.
\] This establishes agreement of predictive laws, not necessarily identity of coordinate systems, physical carriers, or causal mechanisms.

\section{5.3 Predictive Rate--Distortion}\label{predictive-ratedistortion}

\textbf{D-RD.} Allow an encoder \(P(M\mid H,U)\) to use independent randomness, with decoder \(Q(\cdot\mid M,U)\). Define \[
R_{\rm obj}(\varepsilon)=
\inf I(H;M\mid U)
\] subject to \[
E D_{\rm TV}\bigl(P(Y\mid H,U),Q(\cdot\mid M,U)\bigr)
\le\varepsilon.
\] The encoder cannot read a future \(Y\) that has not yet occurred. This quantity measures the memory information required for a specified average predictive error.

\hypertarget{p-info-03}{}\phantomsection\label{p-info-03}

\textbf{P-INFO-03: Zero-distortion complexity.} If \(C=P(Y\in\cdot\mid H,U)\) is a discrete random element with \(H(C\mid U)<\infty\), then \[
R_{\rm obj}(0)=H(C\mid U).
\]

\textbf{Proof.} Zero average TV means the two kernels agree almost surely, so \(C\) is recoverable from \((M,U)\). Therefore \[
I(H;M\mid U)\ge I(C;M\mid U)=H(C\mid U).
\] Choose \(M=C\) and let the decoder output that kernel. This attains the lower bound. The proof also applies to randomized encoders. \(\blacksquare\)

As \(\varepsilon\) increases, the feasible encoder set expands, so \(R_{\rm obj}\) is nonincreasing. This is an information constraint; by itself it does not specify computation time, energy consumption, or biological memory capacity.

\section{5.4 Information Lower Bounds for Retaining Identity}\label{information-lower-bounds-for-retaining-identity}

\hypertarget{p-info-04}{}\phantomsection\label{p-info-04}

\textbf{P-INFO-04: A Fano-type identity bound.} Let \(J\) be uniform on \(K\ge2\) identities. Suppose a decoder based on the future record \(Y_T\) produces \(\widehat J\) with error probability \(e\). Then \[
I(J;Y_T)\ge\log K-h_2(e)-e\log(K-1).
\]

\textbf{Proof.} Let \(E=\mathbf1_{\{\widehat J\ne J\}}\) indicate an error. Given the record and the error indicator, the identity is known on a correct decoding, and at most \(K-1\) possibilities remain on an error. Thus \(H(J\mid Y_T)\le H(E)+e\log(K-1)\). Subtract this from \(\log K\). \(\blacksquare\)

The conditional binary version is as follows: if a future identity variable \(B\in\{0,1\}\) can be decoded from \((M,U)\) with error probability at most \(\epsilon\le1/2\), then \[
I(M;B\mid U)\ge H(B\mid U)-h_2(\epsilon).
\] This concerns information in an ensemble of distinguishable identities. It does not assert that an individual deterministic object must carry positive entropy.

\section{5.5 Omission and Packing Bounds}\label{omission-and-packing-bounds}

\hypertarget{p-info-05}{}\phantomsection\label{p-info-05}

\textbf{P-INFO-05.} Suppose two histories are identified by a representation \(R\), but their future kernels under some protocol have distance \(\delta>0\). Every decoder kernel depending only on \(R\) then has TV error at least \(\delta/2\) on at least one of the two histories.

\textbf{Proof.} The sum of the two errors is at least the distance between the two true kernels. \(\blacksquare\)

If \(N\) histories have pairwise response distances greater than \(2\varepsilon\), a deterministic representation with uniform error at most \(\varepsilon\) requires at least \(N\) values. Otherwise the pigeonhole principle contradicts the preceding triangle-inequality argument. If the input distribution is uniform on exactly these \(N\) histories, its representation entropy is at least \(\log N\).

\section{5.6 The Role of Information Obstructions in the Unified Theory}\label{the-role-of-information-obstructions-in-the-unified-theory}

In \nolinkurl{InformationPacking}, \nolinkurl{common_decoder_half_distance}, \nolinkurl{packing_forces_injective}, and \nolinkurl{packing_ncard_image} establish the precise obstructions in this chapter from the TV triangle relation for actual probability laws. P-INFO-05 is therefore a retained boundary in the compression ledger; it cannot be removed to claim that every endpoint is ``generated'' by the generators.

These lower bounds explain why some distinctions cannot be merged at a specified accuracy. They determine neither computation time nor energy consumption nor a unique information drive. After choosing TV, Bayes risk, or memory rate as \(\Phi\), one must still prove how that quantity enters the same policy evaluation and concrete measurement. Chapters 23, 33, and 34 treat gauges, quantitative calibration, and purpose faithfulness, respectively.

\chapter{6. Structured Openness, Environmental Interfaces, and Causal Boundaries}\label{ch-6}

\textbf{Chapter question.} Do the internal state and environmental interface together suffice to predict the object's future?

\textbf{Reading route.} Check the joint conditional kernel before studying separated internal and environmental representations. The XOR counterexample shows how separate tests can miss synergistic information. Causal shielding additionally requires intervention semantics; predictive sufficiency alone does not establish it.

\section{6.1 A Joint Closure Criterion}\label{a-joint-closure-criterion}

Decompose the history as \(H=(H^S,H^E)\), with internal state \(M=f(H^S)\) and interface \(U=g(H^E)\). For the candidate object's future \(Y_f^+\), define \[
\varepsilon_{\rm joint}
=I(H^S,H^E;Y_f^+\mid M,U).
\]

\hypertarget{p-int-01}{}\phantomsection\label{p-int-01}

\textbf{P-INT-01: A necessary and sufficient condition for structured sufficiency.} \[
Y_f^+\perp H\mid(M,U)
\iff C^f=\Psi(M,U)\quad\text{a.s.}
\] For an entire countable intervention family, this holds protocol by protocol when common versions are used.

\textbf{Proof.} The forward implication is the minimality assertion in P-PRED-01. For the reverse implication, use \(E[C^f(B)\mid M,U]=C^f(B)\) and the tower property to show that the future probabilities conditional on the complete history and on the representation coincide. \(\blacksquare\)

\textbf{N-INT-01: Small separate residuals do not suffice.} Let \(X,E\) be independent uniform bits, let \(Y=X\oplus E\), and let both \(M,U\) be constant. Then \[
I(X;Y)=I(E;Y)=0,\qquad I(X,E;Y)=\log2.
\] Separate tests of ``residual internal information'' and ``residual environmental information'' may therefore miss synergy.

The correct chain decomposition is \[
I(X,E;Y\mid M,U)
=I(X;Y\mid M,U)+I(E;Y\mid X,M,U).
\] The second term must additionally condition on \(X\). This edition uses the joint residual as the certification interface; two nonjoint indicators serve only as diagnostics.

\section{6.2 Least Internal--Environmental Factorization for Finite Structural Kernels}\label{least-internalenvironmental-factorization-for-finite-structural-kernels}

For all admissible \((x,e,i)\), let the finite response kernel \(K_i(x,e)\) be explicitly defined. Define \[
x\sim_Sx'\iff K_i(x,e)=K_i(x',e)\quad\forall e,i,
\] \[
e\sim_Ee'\iff K_i(x,e)=K_i(x,e')\quad\forall x,i.
\] Write \(M^*(x),U^*(e)\) for the corresponding equivalence classes.

\hypertarget{p-int-02}{}\phantomsection\label{p-int-02}

\textbf{P-INT-02.} There is a well-defined kernel satisfying \[
K_i(x,e)=Q_i(M^*(x),U^*(e)).
\] Every other exact separated representation \(K_i(x,e)=\widetilde Q_i(f(x),g(e))\) must refine \(M^*,U^*\) in the respective coordinates.

\textbf{Proof.} Replacing the internal point by an equivalent one, and then replacing the environmental point by an equivalent one, leaves the kernel unchanged. The quotient kernel is therefore well defined. If \(f(x)=f(x')\), then for every environment and protocol both kernels equal \(\widetilde Q_i(f(x),g(e))\), so \(x\sim_Sx'\). The environmental argument is identical. \(\blacksquare\)

Observational data alone do not supply a structural kernel at zero-probability combinations \((x,e)\). The pointwise full-product conclusion requires full support or an \nolinkurl{experimental/mechanistic} specification. For a general factorization modulo null sets, use P-INT-01.

\section{6.3 Causal Interfaces}\label{causal-interfaces}

Separate external factors into an interface \(U\) and remote variables \(E_{\rm far}\). \textbf{H-INT} asserts that, under the admissible mechanism-preserving interventions, \[
P(Y_f^+\mid M,U,do(E_{\rm far}=e))
\] does not depend on \(e\). This is an ideal form of a causal-shielding certificate. An empirical study must specify whether the intervention directly changes the internal mechanism or the interface, and whether any channels have been omitted.

Average and worst-case intervention residuals are, respectively, \[
\sup_i E_{\mu_i}D_{\rm TV}(K_i^{\rm full},K_i^{\rm obj}),
\qquad
\sup_i\operatorname*{ess\,sup}_{\mu_i}
D_{\rm TV}(K_i^{\rm full},K_i^{\rm obj}).
\] The latter is stronger. A single error symbol cannot be used interchangeably without stating which meaning is intended.

\section{6.4 Passive Observation Does Not Uniquely Determine Causal Direction}\label{passive-observation-does-not-uniquely-determine-causal-direction}

Take \(0<|a|<1\). In the first model, \(X\sim N(0,1),\ Y=aX+\epsilon\), where \(\epsilon\sim N(0,1-a^2)\) is independent. The second model reverses the generating order of \(X,Y\). The two joint Gaussian laws are identical, but \[
E_1[Y\mid do(X=x)]=ax,\qquad
E_2[Y\mid do(X=x)]=0.
\] No number of passive samples of this kind can distinguish them; selective interventions can. This is an identification limit for an object's causal boundary.

\section{6.5 Conditions for Uniqueness of a Minimal Markov Boundary}\label{conditions-for-uniqueness-of-a-minimal-markov-boundary}

Relative to a finite variable set \(W\), a subset \(B\subseteq W\) is a blanket for a target \(Y\) if \(Y\perp W\setminus B\mid B\). A minimal blanket is a boundary.

\hypertarget{p-int-03}{}\phantomsection\label{p-int-03}

\textbf{P-INT-03.} If conditional independence satisfies the semigraphoid axioms and intersection, the minimal boundary is unique.

\textbf{Proof.} Suppose \(B,C\) are both boundaries. Set \(A=B\setminus C\), \(D=C\setminus B\), \(F=B\cap C\), and \(R=W\setminus(B\cup C)\). The blanket properties give \(Y\perp D\mid A,F\) and \(Y\perp A\mid D,F\). Intersection yields \(Y\perp(A,D)\mid F\). Combining this with \(Y\perp R\mid A,D,F\) by contraction shows that \(F\) is also a blanket. Minimality forces \(F=B=C\). \(\blacksquare\)

A strictly positive joint distribution on a finite discrete full-product support ensures intersection: for fixed \(f\), a conditional kernel depending only on \(a\) and also only on \(d\) must, by full positive support, depend only on \(f\). Redundant copies \(B=A\) do not satisfy this condition and may yield several minimal boundaries. Uniqueness of an observational boundary still does not identify the set of causal parents.

\section{6.6 Separated Representation and Formation--Binding Are Not Interchangeable}\label{separated-representation-and-formationbinding-are-not-interchangeable}

The two input quotients in P-INT-02 establish only a well-defined two-sided factorization of a response kernel. Formation of a parent object additionally requires conditions on its future, interventions, carriers, and history. C4 separates \nolinkurl{FormationRelation} from \nolinkurl{BindingTransportCompatible}: formation error is amplified by an \(L\)-Lipschitz binding map and transport error is then added, giving the bound \(L\varepsilon_F+\varepsilon_B\).

Commutativity of formation does not imply commutativity of binding; finite counterexamples are present in the source. The claim ``the internal--external interface is closed, therefore a parent object has been generated'' must therefore supply these independent conditions. Section 34.4 gives the complete composition proposition and its proof.

\chapter{7. Dynamical Closure, General Processes, and Multiscale Consistency}\label{ch-7}

\textbf{Chapter question.} Can a compressed state evolve autonomously and remain compatible across scales?

\textbf{Reading route.} Exact results come from equality of projected rows within each observation fiber; approximate results accumulate through stepwise coupling under the same policy. Cross-scale proofs must insert the same underlying process. Finite-path error does not by itself yield infinite-time stability.

\section{7.1 Exact Markov Quotients}\label{exact-markov-quotients}

\hypertarget{p-dyn-01}{}\phantomsection\label{p-dyn-01}

\textbf{P-DYN-01: Strong lumpability.} Let \(X_t\) be a Markov chain, let \(f:\mathcal X\to\mathcal M\) be a measurable surjection, and suppose a measurable candidate macroscopic kernel \(\bar P\) has already been specified. The following are equivalent:

\begin{enumerate}
\def\labelenumi{\arabic{enumi}.}
\tightlist
\item
  For every microscopic initial distribution, \(f(X_t)\) is Markov with the same transition kernel \(\bar P\).
\item
  For every \(x,B\), \[
  P(x,f^{-1}B)=\bar P(f(x),B).
  \]
\end{enumerate}

\textbf{Proof.} Starting from 2, first condition on the complete microscopic history and then take conditional expectation given the macroscopic history. This gives the macroscopic Markov property. Starting from 1, take the first-step probability from microscopic initial state \(\delta_x\) to obtain 2. \(\blacksquare\)

On a finite state space, equality of projected rows within each fiber defines \(\bar P\) directly. On a general space, constancy on fibers alone still requires a measurable factorization condition. The controlled version must hold for every admissible action; closure under a single passive policy is insufficient.

\section{7.2 A Continuous-Time Generator Criterion}\label{a-continuous-time-generator-criterion}

\hypertarget{p-dyn-02}{}\phantomsection\label{p-dyn-02}

\textbf{P-DYN-02.} A finite CTMC generator \(L\) and a partition \(B_1,\ldots,B_m\) admit an exact partition quotient if and only if \[
\sum_{y\in B_j}l_{xy}
\] is constant within each \(B_i\) over all states \(x\). The common values form a macroscopic generator \(\bar L\).

\textbf{Proof.} If the semigroup quotient holds, differentiate \(e^{tL}\) at zero to obtain the condition. Conversely, let \(F\) be the block-indicator matrix. The condition is \(LF=F\bar L\); induction gives \(L^nF=F\bar L^n\), hence \(e^{tL}F=Fe^{t\bar L}\). \(\blacksquare\)

\section{7.3 Approximate Kernels and Path Bounds}\label{approximate-kernels-and-path-bounds}

Suppose the within-fiber diameter \[
\epsilon(f)=\sup_{f(x)=f(x'),a}
D_{\rm TV}(f_\#P_{x,a},f_\#P_{x',a})
\] is small and a measurable section \(s(m)\) exists. Choosing \(\bar P_{m,a}=f_\#P_{s(m),a}\) gives a uniform projection error at most \(\epsilon(f)\). A section always exists on finite spaces, but cannot be assumed without conditions on general spaces.

\hypertarget{p-dyn-03}{}\phantomsection\label{p-dyn-03}

\textbf{P-DYN-03: Path error under a common policy.} Suppose two finite-alphabet causal record processes have next-step kernel differences at most \(\epsilon_t\in[0,1]\) whenever their realized histories agree. If they use the same causal policy and have identical initial records, then \[
D_{\rm TV}(P_{0:T},\widehat P_{0:T})
\le1-\prod_{t=0}^{T-1}(1-\epsilon_t)
\le\sum_{t=0}^{T-1}\epsilon_t.
\]

\textbf{Proof.} As long as the histories have not diverged, actions can be sampled jointly and the next records maximally coupled, with conditional agreement probability at least \(1-\epsilon_t\). Induction gives the product lower bound for agreement of the entire path. The coupling inequality yields the conclusion. \(\blacksquare\)

If the initial record laws already differ by \(\epsilon_0^{\rm init}\), that error may be added. Bounds for time marginals follow by pushing forward the path bound. Approximate kernels constructed separately at arbitrary lags do not automatically form a semigroup.

\section{7.4 Cross-Scale Quotients of the Same Underlying Process}\label{cross-scale-quotients-of-the-same-underlying-process}

Let \(f_r=c_{sr}\circ f_s\), with errors \(\epsilon_r,\epsilon_s\) at the two scales relative to the same underlying kernel.

\hypertarget{p-dyn-04}{}\phantomsection\label{p-dyn-04}

\textbf{P-DYN-04.} On the reachable image \(f_s(\mathcal X)\), \[
\sup_{m,a}D_{\rm TV}
\bigl((c_{sr})_\#P_s(m,a),P_r(c_{sr}m,a)\bigr)
\le\epsilon_s+\epsilon_r.
\] In the exact case this is an intertwining relation.

\textbf{Proof.} For \(m=f_s(x)\), insert the pushforward of the common underlying process, \((f_r)_\#P(x,a)\). Data processing bounds the first term by \(\epsilon_s\), and the second is at most \(\epsilon_r\). \(\blacksquare\)

A supremum over the whole macroscopic space requires surjectivity or control of extensions at unreachable points. This conclusion does not follow automatically from legitimate coarse-graining of carrier sets.

\section{7.5 Multitime Interfaces for Non-Markov Processes}\label{multitime-interfaces-for-non-markov-processes}

Write \(\mathbf T[\pi]\) for the complete record law of a protocol tree \(\pi\). If a coarse protocol is implemented at the fine level through \(J\) and its readout transformed through \(C\), define \[
(\mathcal R\mathbf T)[\pi]=C_\#\mathbf T[J\pi].
\] Protocol lifting and readout pushforward are associative, so two successive interfaces can be composed. This compares only protocol trees that can actually be implemented; it does not require finite memory.

If the fine level has actions that cannot be expressed at the coarse level, a process quotient does not ensure equality of optimal values. Value preservation additionally requires policy completeness and reward correspondence; see Chapters 20 and 28.

\section{7.6 From Pointwise Error to Error Weighted by Actual Occupation}\label{from-pointwise-error-to-error-weighted-by-actual-occupation}

Uniform \(\epsilon_t\) bounds may be conservative. If the model error in row \((x,a)\) is \(e_t(x,a)\) and the actual prefix together with the common policy induces an occupation law \(\nu_t\), Chapter 36 proves, from actual PMF extension, that

\[D_{\rm TV}(\text{actual coarse path},\text{reference path})\le\sum_{t<T}E_{\nu_t}e_t.\]

If on a certified row set \(C_t\) one knows only that \(e_t\le\epsilon_t\in[0,1]\), then

\[E_{\nu_t}e_t\le\epsilon_t+(1-\epsilon_t)\nu_t(C_t^c).\]

This does not substitute an average error for a uniform error: the averaging law \(\nu_t\) is the actual visitation law of the process being compared, and its provenance is retained in the proof. If the training row distribution differs, P-STAT-09 or a new coverage result is still needed.

\chapter{8. Persistence: Invariant Structures, Viability, and Relative Identity}\label{ch-8}

\textbf{Chapter question.} Which distinct mathematical requirements can express persistence?

\textbf{Reading route.} Distinguish invariant limits of precompact trajectories, invariant limits of averaged time marginals, and viability kernels for almost-sure confinement. Finite models obtain safe actions by successive elimination. Partial observation, finite horizons, and pathwise empirical frequencies have their own additional conditions.

\section{8.1 Identity Branches and Persistence Domains}\label{identity-branches-and-persistence-domains}

A traceable historical branch need not be identified in advance with a fixed state set. To study how long it continues to satisfy specified structural requirements, choose explicitly an assessment domain \(\mathcal V\) and the exit time \[
\tau_{\mathcal V}=\inf\{t\ge0:X_t\notin\mathcal V\}.
\] Dynamic identity can also be represented by an augmented state \((X_t,\Theta_t)\) with a fixed domain in that space, or by a specified family of moving domains. Choosing a persistence domain is part of the model and task specification. Maximizing persistence probability is a further evaluative bridge.

\section{8.2 A Necessary Invariant Core for Deterministic Persistence}\label{a-necessary-invariant-core-for-deterministic-persistence}

\hypertarget{p-per-01}{}\phantomsection\label{p-per-01}

\textbf{P-PER-01.} Suppose an orbit of a continuous semiflow \(\phi_t\) eventually remains in a closed set \(\mathcal V\) and is precompact. Then \[
\omega(x)=\bigcap_{T>0}\overline{\{\phi_t(x):t\ge T\}}
\] is nonempty, compact, contained in \(\mathcal V\), and satisfies \(\phi_s(\omega(x))=\omega(x)\).

\textbf{Proof.} The orbit closure is compact, so the decreasing family of nonempty compact tail sets has nonempty intersection. Because the orbit remains in the closed domain, the intersection lies in that domain. If \(\phi_{t_n}(x)\to y\), continuity gives \(\phi_s(y)\in\omega(x)\). Conversely, take \(t_n>s\) and extract a convergent subsequence from the precompact sequence \(\phi_{t_n-s}(x)\). Its limit \(z\in\omega(x)\) satisfies \(\phi_s(z)=y\). \(\blacksquare\)

This is not a theorem of attraction from a neighborhood. Finite-time persistence does not imply a compact invariant core either: the translation \(x_t=x_0+t\) can remain in a sufficiently large interval throughout any prescribed finite window, while possessing no nonempty compact invariant set.

\section{8.3 Symmetry Reduction}\label{symmetry-reduction}

Suppose the action of a group \(G\) is equivariant with the dynamics, \[
\phi_t(gx)=g\phi_t(x),
\] and the orbit quotient is a well-behaved Hausdorff metrizable space. Then the quotient dynamics \(\bar\phi_t([x])=[\phi_t(x)]\) is well defined. If its orbit is precompact, P-PER-01 applies on the quotient. A moving form can thereby preserve relative identity.

Without equivariance, changing representatives may change the future on the quotient. A quotient cannot be taken merely because translation ``seems irrelevant.''

\section{8.4 Stochastic Persistence and an Occupation-Based Invariant Measure}\label{stochastic-persistence-and-an-occupation-based-invariant-measure}

\hypertarget{p-per-02}{}\phantomsection\label{p-per-02}

\textbf{P-PER-02.} Let \(P_t\) be a Feller semigroup on a Polish space that maps \(C_b\) into \(C_b\). If \[
\bar\mu_T=\frac1T\int_0^T\mu_0P_t\,dt
\] is tight over \(T\ge1\), there is a weakly convergent subsequence along some \(T_n\to\infty\), and every such limit \(\nu\) is invariant. If every time marginal is supported on a closed set \(\mathcal V\), then \(\nu(\mathcal V)=1\).

\textbf{Proof.} Prokhorov's theorem gives a weakly convergent subsequence. For \(f\in C_b\), \[
|(\bar\mu_TP_s-\bar\mu_T)f|
\le2s\|f\|_\infty/T.
\] The Feller property allows weak convergence to be applied to \(P_sf\), so \(\nu P_sf=\nu f\). The closed-support conclusion follows from Portmanteau. \(\blacksquare\)

These are averages of time-marginal distributions. The statement does not assert convergence of empirical frequencies along every sample path; that requires ergodic assumptions.

\section{8.5 A Finite-State Viability Kernel}\label{a-finite-state-viability-kernel}

For finite state and action spaces and an admissible persistence domain \(K_0=\mathcal V\), define recursively \[
K_{n+1}=\{x\in K_n:\exists a,\ P(K_n\mid x,a)=1\}.
\]

\hypertarget{p-per-03}{}\phantomsection\label{p-per-03}

\textbf{P-PER-03.} The recursion stabilizes after finitely many steps at \(K_\infty\). This is exactly the set of states from which some policy keeps the system in \(\mathcal V\) forever with probability one. A stationary deterministic policy preserving it exists.

\textbf{Proof.} The sets decrease and are finite. At stabilization, each state admits an action whose next state remains in the same set with probability one. Select one such action at each state; induction gives perpetual retention. Conversely, if a state permits perpetual retention, its first step must lead with probability one to states having the same property. Inductively, it cannot be removed at any deletion stage. \(\blacksquare\)

The conclusion concerns almost-sure persistence in a finite model. Long expected lifetime, high finite-horizon survival probability, and membership in this viability kernel are different requirements.

\section{8.6 Tangency Conditions for Continuous States}\label{tangency-conditions-for-continuous-states}

For a closed domain \(\mathcal V\subseteq\mathbb R^d\), define the contingent cone \[
T_{\mathcal V}(x)=
\left\{v:\liminf_{h\downarrow0}\frac{d(x+hv,\mathcal V)}h=0\right\}.
\]

\hypertarget{p-per-04}{}\phantomsection\label{p-per-04}

\textbf{P-PER-04: Tangency is necessary.} If a viable trajectory, differentiable at zero, starts at \(x\) and remains in \(\mathcal V\), then \(\dot x(0)\in T_{\mathcal V}(x)\).

\textbf{Proof.} Since \(x(h)=x+h\dot x(0)+o(h)\in\mathcal V\), the distance ratio is at most \(o(h)/h\). \(\blacksquare\)

\textbf{K-VIA-01.} For a differential inclusion satisfying the Marchaud conditions---upper semicontinuity, nonempty compact convex values, and linear growth---the corresponding viability theorem characterizes viability of a closed set by \(F(x)\cap T_{\mathcal V}(x)\ne\varnothing\). For a locally Lipschitz, nonexplosive single-valued ODE, Nagumo's positive-invariance criterion may be used. The original idea that constraints maintain identity at a boundary thus has a precise interface. Viability must remain distinct from retention under every control.

\section{8.7 Actual Paths of a Finite Source and Extension to Partial Observation}\label{actual-paths-of-a-finite-source-and-extension-to-partial-observation}

The current source \nolinkurl{ViabilitySource.p_per_03} for P-PER-03 constructs a genuine infinite-trajectory measure: for an initial PMF supported on the stable kernel, there is a stationary deterministic policy for which the event of remaining in that kernel at every time has probability one. A nonempty finite action space, the source PMF, and the safe set are explicit inputs; the stable kernel itself may be empty. ``Safe at every time'' has not been inserted as an assumption in a conclusion field.

Under partial observation, actions may depend only on the coarse history, not be selected separately for an unobserved microscopic state. Chapter 37 therefore requires one common certified action on the actual posterior support and performs backward recursion over all of its positive-probability successor supports, proving a necessary and sufficient condition for a finite-horizon policy to exist. This is a different theorem from the fully observed stationary infinite-horizon kernel in this chapter. One cannot directly take \(T\to\infty\) in the finite-support result.

\chapter{9. Recovery, Perturbations, and Canonical Recovery Potentials}\label{ch-9}

\textbf{Chapter question.} How can recovery after perturbation be proved, and how does recovery rate relate to return time?

\textbf{Reading route.} Expand the discrete or continuous drift inequality, then construct a recovery potential from the actual hitting time. Persistent noise generally leaves an error floor. A potential and an executable repair program are separate obligations; the latter must also be linked to the actual transition kernel.

\section{9.1 Discrete Recovery Tubes}\label{discrete-recovery-tubes}

\hypertarget{p-rec-01}{}\phantomsection\label{p-rec-01}

\textbf{P-REC-01.} Suppose \(R_n\ge0\), \(ER_0<\infty\), and \[
E[R_{n+1}\mid\mathcal F_n]\le\kappa R_n+\eta,
\quad0\le\kappa<1,\quad\eta\ge0.
\] Then \[
ER_n\le\kappa^nER_0+\eta\frac{1-\kappa^n}{1-\kappa},
\] and, for \(a>0\), \[
P(R_n\ge a)\le
\frac{\kappa^nER_0+\eta(1-\kappa^n)/(1-\kappa)}a.
\]

\textbf{Proof.} Taking total expectations gives a scalar recursion. Expand its geometric sum by induction. Markov's inequality gives the tail bound. \(\blacksquare\)

Persistent noise produces an error floor \(\eta/(1-\kappa)\). For example, \(x_{n+1}=(x_n+1)/2,\ x_0=0\) gives \(x_n=1-2^{-n}\), which does not tend to zero. Contraction of an individual recovery step relative to the just-perturbed point does not justify ignoring the next fresh disturbance.

\section{9.2 Continuous Stochastic Recovery}\label{continuous-stochastic-recovery}

\hypertarget{p-rec-02}{}\phantomsection\label{p-rec-02}

\textbf{P-REC-02.} Fix the initial state \(X_0=x\). Suppose \(W\ge c\,d(\cdot,\mathcal V)^2\) lies in an appropriate domain of the generator \(\mathcal L\), and assume the Dynkin formula, localization limits, and integrability conditions needed to make \(m(t)=EW(X_t)\) locally absolutely continuous. If \[
\mathcal LW\le-aW+b,\qquad a,c>0,\ b\ge0,
\] then \[
EW(X_t)\le e^{-at}W(x)+\frac ba(1-e^{-at}),
\] and division by \(c\) gives a mean-square distance bound.

\textbf{Proof.} Almost everywhere, \(m'(t)=E\mathcal LW(X_t)\le-am(t)+b\). Multiply by \(e^{at}\) and integrate. \(\blacksquare\)

A formal expression for \(\mathcal LW\) does not automatically justify Dynkin's formula. Explosion, a nonintegrable stochastic integral, or an incorrect generator domain must first be addressed in the model.

\section{9.3 A Canonical Potential from Recovery Time}\label{a-canonical-potential-from-recovery-time}

Let \(\tau_A=\inf\{n\ge0:X_n\in A\}\) and \(V_A(x)=E_x\tau_A\in[0,\infty]\), with value zero on \(A\).

\hypertarget{p-rec-03}{}\phantomsection\label{p-rec-03}

\textbf{P-REC-03.} For \(x\notin A\) with \(V_A(x)<\infty\), \[
PV_A(x)-V_A(x)=-1.
\]

\textbf{Proof.} First-step analysis and the Markov property give \(V_A(x)=1+E_xV_A(X_1)\). Finiteness also ensures integrability of the right-hand side. \(\blacksquare\)

\hypertarget{p-rec-04}{}\phantomsection\label{p-rec-04}

\textbf{P-REC-04: Drift bounds recovery time.} If a finite nonnegative function \(V\) satisfies, on \(A^c\), the inequality \(PV-V\le-c<0\), then \[
E_x\tau_A\le V(x)/c.
\]

\textbf{Proof.} Successive conditional expectations up to \(\tau_A\wedge N\) give \[
E V(X_{\tau_A\wedge N})+cE(\tau_A\wedge N)\le V(x).
\] Discard the nonnegative first term and apply monotone convergence. \(\blacksquare\)

In continuous time, the corresponding equation is \(\mathcal LV_A=-1\) under suitable extended-generator and stopped-Dynkin assumptions. Finite recovery does not require strict distance contraction: \(x\mapsto\max(x-1,0)\) always recovers on the nonnegative integers, but has Lipschitz constant 1.

\section{9.4 The Mechanism Profile of Ω}\label{the-mechanism-profile-of-ux3c9}

At least the following must be distinguished: recovery, invariance, periodicity, topological protection, quasi-stationarity, and viability under endogenous control. The equation \(x'=-x\) has a stable point on a contractible space, showing that nonzero first homology is not universally necessary. Neutral rotation on a circle persists without contraction, showing that strict contraction is not necessary either.

Organizational closure may be reported through a change in exit risk, a lifetime difference, or a hazard ratio after a process is disabled. Ratios require positive finite denominators, and hazards must actually exist; for an exit distribution with atoms, interval failure probabilities should be used. A directed cycle in a mechanism graph does not by itself ensure stability: gains around the cycle may instead produce divergence.

\section{9.5 Recovery Processes, Recovery Programs, and the Same Source Kernel}\label{recovery-processes-recovery-programs-and-the-same-source-kernel}

Bounds on recovery moments or return times do not automatically supply an executable repair program. \nolinkurl{RepairProgramSafeControlBridge} separately specifies a trusted repair substrate and a program encoding valid for the same source transition kernel. It reuses the program-recovery theorem to prove that the kernel implemented after recovery is exactly the already-certified safe source-action kernel.

One may therefore first prove the existence of a safe action through a survival fixed point, and then use the program bridge to prove that the recovered program implements that action. Program provenance, the interpreter, and \nolinkurl{encoding/decoding} correctness must still be specified; the contraction or drift inequalities in P-REC do not generate them. In the continuous branch, \nolinkurl{RecoveryDynkin} likewise explicitly consumes a process-specific \nolinkurl{DynkinExpectationCertificate} before deriving the consequences of the scalar differential inequality.

\chapter{10. Quasi-Stationarity, Spectra, and Persistence Windows}\label{ch-10}

\textbf{Chapter question.} Why can a system that may eventually fail exhibit a stable form conditional on survival?

\textbf{Reading route.} Retain the mass deficit of the substochastic killed kernel, then use either a conditional limit or positive eigenstructure. A left eigenvector gives the QSD, whereas the product of left and right eigenvectors gives the invariant weight of the Doob process. Mixing and survival bounds for the same initial law determine the useful time window.

\section{10.1 Killed Processes and QSDs}\label{killed-processes-and-qsds}

On the assessment domain \(\mathcal V\), define the subprobability semigroup \[
P_t^{\mathcal V}f(x)=E_x[f(X_t)\mathbf1_{\{t<\tau_{\mathcal V}\}}].
\] A probability measure \(q\) is a QSD if its distribution remains \(q\) when conditioned on survival.

\hypertarget{p-qsd-01}{}\phantomsection\label{p-qsd-01}

\textbf{P-QSD-01: A converse from conditional stabilization.} Suppose \[
\mu_t^c=\frac{\mu P_t^{\mathcal V}}{\mu P_t^{\mathcal V}\mathbf1}
\longrightarrow q
\quad\text{in TV},
\] with, for every finite \(s>0\), the positivity condition \(qP_s^{\mathcal V}\mathbf1>0\), and suppose its survival function is right-continuous at zero. Then \(q\) is a QSD, and for some \(\lambda\ge0\), \[
qP_s^{\mathcal V}=e^{-\lambda s}q.
\]

\textbf{Proof.} The semigroup property gives \[
\mu_{t+s}^c(f)=
\frac{\mu_t^c(P_s^{\mathcal V}f)}
{\mu_t^c(P_s^{\mathcal V}\mathbf1)}.
\] TV convergence permits passage to the limit for bounded functions; positivity of the denominator yields the QSD identity. Set \(c(s)=qP_s^{\mathcal V}\mathbf1\). The semigroup property gives \(c(s+t)=c(s)c(t)\). Positivity, right-continuity, and \(c(0)=1\) imply \(c(s)=e^{-\lambda s}\), and the subprobability property gives \(\lambda\ge0\). \(\blacksquare\)

Eventual almost-sure absorption implies \(\lambda>0\); the degenerate case without absorption permits zero.

\section{10.2 The Finite Perron Representation}\label{the-finite-perron-representation}

\textbf{K-PF-01.} For a primitive finite nonnegative matrix \(Q\), the Perron--Frobenius theorem supplies a unique positive dominant eigenvalue \(\rho>0\), positive left and right eigenvectors, and a dominant-mode asymptotic decomposition. Primitivity means that some matrix power is strictly positive entrywise.

When \(Q\) is a killed substochastic kernel, \(\rho\le1\). Normalize by \[
qQ=\rho q,\quad Qh=\rho h,\quad q\mathbf1=1,\quad qh=1.
\]

\hypertarget{p-qsd-02}{}\phantomsection\label{p-qsd-02}

\textbf{P-QSD-02.} The measure \(q\) is a QSD; the Doob transform \[
Q^\uparrow_{ij}=\frac{Q_{ij}h_j}{\rho h_i}
\] is a stochastic kernel, with invariant probability \(\varpi_i=q_i h_i\).

\textbf{Proof.} Since \(qQ^n=\rho^nq\), division by the survival mass \(\rho^n\) gives the QSD property. Row sums equal one by \(Qh=\rho h\). Moreover, \[
\sum_iq_i h_i Q^\uparrow_{ij}
=\frac{h_j}{\rho}(qQ)_j=q_jh_j.
\] The total mass is \(qh=1\). \(\blacksquare\)

Starting from \(\mu\), one has \(\mu Q^n\sim\rho^n(\mu h)q\). The measures \(q\) and \(\varpi\) describe, respectively, ``survival until the present time'' and the ``process conditioned on indefinite survival''; in general they differ. For a continuous-time Markov chain, use the killed generator \(K\), with \(qK=-\lambda q\) and \(Kh=-\lambda h\). Its transform is \[
K^\uparrow=\operatorname{diag}(h)^{-1}K\operatorname{diag}(h)+\lambda I.
\]

\section{10.3 Persistence Windows from the Same Initial Law}\label{persistence-windows-from-the-same-initial-law}

\hypertarget{p-qsd-03}{}\phantomsection\label{p-qsd-03}

\textbf{P-QSD-03.} For the same initial law \(\mu\), suppose that for \(t\ge t_0\), \[
\|\mu_t^c-q\|_{\rm TV}\le Ce^{-\gamma t},\qquad
P_\mu(\tau>t)\ge ce^{-\lambda t},
\] where \(C,c,\gamma>0,\lambda>0\). Given \(\epsilon>0\) and \(0<p\le c\), if the interval \[
\left[\max\{t_0,0,\log(C/\epsilon)/\gamma\},
\ \log(c/p)/\lambda\right]
\] is nonempty, every time in it simultaneously has conditional-stability error at most \(\epsilon\) and survival probability at least \(p\).

\textbf{Proof.} Solve the two exponential inequalities separately and intersect their solution sets. Equal endpoints of the closed interval still give one time; positive length requires a strict inequality. \(\blacksquare\)

When \(\lambda=0\), the survival lower bound gives no such finite upper endpoint. When \(C=0\), the conditional distribution is already exactly stable and the mixing lower endpoint is simply \(t_0\). A mixing bound for one initial law cannot be combined with a survival bound for another.

\section{10.4 A Spectral Version for Reversible Diffusions}\label{a-spectral-version-for-reversible-diffusions}

Suppose the killed generator \(L_D\), for a measure with \(m(D)<\infty\), is self-adjoint with compact resolvent on \(L^2(m)\) and has a positive first eigenfunction \(\phi_1\) with \[
-L_D\phi_n=\lambda_n\phi_n,\qquad
0<\lambda_1<\lambda_2\le\cdots.
\] Normalize \(\|\phi_1\|_2=1\). Let the initial probability density \(g\in L^2(m)\) satisfy \(g\ge0\) and \(\int g\,dm=1\), and write \(a_1=\langle g,\phi_1\rangle>0\).

\hypertarget{p-qsd-04}{}\phantomsection\label{p-qsd-04}

\textbf{P-QSD-04.} There are \(C_g,c_g>0,t_0\) such that \[
q(dx)=\frac{\phi_1(x)m(dx)}{\int\phi_1\,dm},\qquad
\|\mu_t^c-q\|_{\rm TV}\le C_ge^{-(\lambda_2-\lambda_1)t},
\] \[
P_g(\tau>t)\ge c_ge^{-\lambda_1t}\quad(t\ge t_0).
\]

\textbf{Proof.} The spectral expansion is \(P_t^Dg=e^{-\lambda_1t}a_1\phi_1+R_t\), with \(\|R_t\|_2\le e^{-\lambda_2t}\|g-a_1\phi_1\|_2\). Cauchy--Schwarz converts this to an \(L^1\) bound. The leading coefficient of the survival mass satisfies \(a_1\int\phi_1dm>0\); for sufficiently large \(t\), the remainder is less than half the leading term. Divide numerator and denominator by \(e^{-\lambda_1t}\) and compare the conditional density with \(q\) to obtain the stated rate. \(\blacksquare\)

\textbf{Current Lean provenance.} \nolinkurl{CompactSpectralSemigroupSource} derives the actual expansion and Parseval remainder bound from the generating relation between the same actual \(R=(I-L_D)^{-1}\) and \(S_t\), compact self-adjointness, and the original spectral assumptions. \nolinkurl{CompactKilledSpectralBinding.ofActualSource} uses these to construct the original SpectralData; \nolinkurl{actual_p_qsd_04} then consumes the conditional and survival laws of the same killed process. The original expansion and remainder are no longer conclusion fields supplied by the caller. Positivity of the leading coefficient and total leading-mode mass, together with injectivity of the resolvent, are also derived from the current inputs. Exact provenance and a common instance with a nonzero remainder appear in §§39.73--39.74.

General non-self-adjoint kernels may have nonorthogonal modes or Jordan terms. The same exact spectral-gap constant cannot be applied without controlling these contributions. General QSD existence and convergence results may be taken from the original research of Champagnat--Villemonais; this book uses them only as conditional modules. \href{https://arxiv.org/abs/1404.1349}{Original paper}

\section{10.5 Positive Eigenstructure Data and Spectral Existence as Separate Layers}\label{positive-eigenstructure-data-and-spectral-existence-as-separate-layers}

M-PE-01 reuses the normalization algebra of this chapter: given positive left and right eigendata, one can prove the QSD property, the stochasticity of the Doob kernel, stationary weights, and martingale properties of stochastic readouts. \nolinkurl{PositiveEigenstructure.qsdDoob_bundle} provides the declared algebra for given data. \nolinkurl{PrimitivePerronClosure} now constructs a complete KPF01Certificate from the original primitive-matrix hypothesis; \nolinkurl{CompactKilledSpectralBinding} constructs the spectral fields in this section from an actual generating relation. The responsibilities of these additional source proofs and of the shared algebraic interface are registered separately.

This explains why a ``unified positive eigenstructure'' is different from generating all long-term results with a single spectral theorem. The conditional limit in P-QSD-01, the same-initial-law window in P-QSD-03, the reversible spectral analysis in P-QSD-04, and general GOA stability each retain their own adapters. When support changes or the spectral gap is lost, those independent assumptions must be revisited.

\chapter{11. Historical Identity, Transport, and Developmental Laws}\label{ch-11}

\textbf{Chapter question.} How can identity continuation be tested through structural change, fission, and fusion?

\textbf{Reading route.} Directed historical edges record relations between instances, while transport maps compare their structures. The proofs accumulate error along the declared transport chain. Frozen endpoint kernels differ in type from actual one-step evolution, and neither resemblance nor common ancestry establishes identity of instances.

\section{11.1 Directed Continuation}\label{directed-continuation}

\textbf{D-ID.} An object occurrence \(O_t^\alpha\) is an instance carrying a time and a certificate. A continuation edge \[
\Gamma:O_t^\alpha\rightsquigarrow O_{t'}^\beta,\quad t<t'
\] records ancestry, structural transport, and a traceable mechanism. Composition in temporal order gives a directed historical path.

After a split \(A\to B,A\to C\), the simultaneously existing \(B,C\) are not the same instance merely because they share an ancestor. An undirected connected component records only lineage relatedness. Nor does fusion imply retrospectively that its predecessors were always one object.

\section{11.2 TCIC for Frozen Models}\label{tcic-for-frozen-models}

For structural models frozen at different times, use endokernels \(K_t:\mathcal M_t\to\mathcal P(\mathcal M_t)\) and transports \(\Gamma_{s,t}:\mathcal M_s\to\mathcal M_t\) satisfying \(\Gamma_{r,t}=\Gamma_{s,t}\circ\Gamma_{r,s}\). Exact commutation is \[
(\Gamma_{s,t})_\#K_s(m,\cdot)=K_t(\Gamma_{s,t}m,\cdot).
\]

\hypertarget{p-id-01}{}\phantomsection\label{p-id-01}

\textbf{P-ID-01: Accumulation of transport defects.} If adjacent endokernels commute with error at most \(\epsilon_t\), then \[
\sup_mD_{\rm TV}\bigl((\Gamma_{0,n})_\#K_0(m),
K_n(\Gamma_{0,n}m)\bigr)\le\sum_{t<n}\epsilon_t.
\]

\textbf{Proof.} At step \(n+1\), insert \((\Gamma_{n,n+1})_\#K_n(\Gamma_{0,n}m)\). Pushforward does not increase the previous error, and the new error is at most \(\epsilon_n\). Induction proves the claim. \(\blacksquare\)

These are endokernels of frozen models, not actual one-step temporal transitions. If \(P_t:\mathcal M_t\to\mathcal P(\mathcal M_{t+1})\), comparison of two actual histories instead has the form \[
(\Gamma_{t+1})_\#P_t(m,\cdot)
\approx\widehat P_t(\Gamma_tm,\cdot).
\] The next state must be transported by the map at the next time.

\section{11.3 Developmental Transport: DTRC}\label{developmental-transport-dtrc}

A structural descriptor \(\Theta_t\) may contain a predictive factor, transition kernel, interfaces, carrier families, and mechanism certificates. It belongs to a declared metric or pseudometric space; different components may retain vector-valued errors.

\textbf{D-DTRC.} A developmental model \[
\widehat\Theta_{t+1}=\Psi_t(\Theta_t,E_t,a_t)
\] must be constrained in advance by a mechanism, independent individuals, early data, or training environments. The test residual \[
r_t=d(\Theta_{t+1},\Psi_t(\Theta_t,E_t,a_t))
\] is evaluated with the model frozen. Predicted causal consequences must be checked, not merely the fit of the descriptors.

Any finite trajectory admits a zero-residual fit by a time-indexed lookup table. Zero training residual therefore does not establish a developmental mechanism.

\hypertarget{p-id-02}{}\phantomsection\label{p-id-02}

\textbf{P-ID-02: A developmental tube.} Under the same declared external inputs, let the reference trajectory satisfy \(\bar\Theta_{t+1}=\Psi_t(\bar\Theta_t)\). If the actual model residual is at most \(\epsilon_t\) and \(\Psi_t\) is \(L_t\)-Lipschitz, then \[
e_{t+1}\le L_te_t+\epsilon_t,
\] \[
e_n\le\Bigl(\prod_{j<n}L_j\Bigr)e_0+
\sum_{k<n}\epsilon_k\prod_{j=k+1}^{n-1}L_j,
\quad e_t=d(\Theta_t,\bar\Theta_t).
\]

\textbf{Proof.} Insert \(\Psi_t(\Theta_t)\), apply the triangle inequality and the Lipschitz assumption, and expand the resulting recursion. \(\blacksquare\)

If \(L_t\le L<1,\epsilon_t\le\epsilon\), this gives \(e_n\le L^ne_0+\epsilon(1-L^n)/(1-L)\). If the actual and reference inputs differ, an input-sensitivity error must also be included; these are then not the same trajectory model.

\section{11.4 Splitting, Fusion, and Fidelity}\label{splitting-fusion-and-fidelity}

A continuation record retains \[
\mathbf F_\Gamma=
(F_{\rm pred},F_{\rm causal},F_{\rm carrier},
F_{\rm control},F_{\rm heredity}),
\] where each component has its own comparison task, normalization, and confidence range. An unstudied component is left empty, not filled with one.

The branch profiles \((1,0)\) and \((0,1)\) have no natural unique percentage ranking. If the two axes have equal standing, interchange symmetry forces every symmetric scalar to give them the same score; breaking the tie requires an additional priority. The default is the componentwise partial order.

\section{11.5 Preserving the Same Object and Producing Offspring}\label{preserving-the-same-object-and-producing-offspring}

In \nolinkurl{LineageSemantics}, SameObject compares identity, whereas OffspringOf requires both a genuine parent--child relation and distinct identities. Accordingly, \nolinkurl{offspring_not_sameObject} and \nolinkurl{sameObjectRepair_not_offspring} hold. Structural similarity, common ancestry, and repair of the same object are different relations.

In P12, \nolinkurl{homeostaticMarginals_preserve_lifecycleParent} assumes that every organizational object on the carrier satisfies SameObject with the same parent, and derives identity preservation for the supported marginals at each finite time. This is a conditional identity bridge; closeness of marginal distributions does not automatically establish identity. The defect bounds in Chapter 33 must likewise not be interpreted directly as percentages of identity.

\chapter{12. Object Certificates and Causal Integrity Ω}\label{ch-12}

\textbf{Chapter question.} What independent claims are included when an object is said to be certified?

\textbf{Reading route.} List predictive, causal, historical, and realization obligations before defining failure under mechanism removal relative to a fixed certificate. A finite monotone failure family is reconstructed from its minimal destructive sets. When removal can restore function, retain the full response instead of imposing a monotone hypergraph.

\section{12.1 Layered Records}\label{layered-records}

\textbf{D-OBJ.} A local operational-object certificate contains at least:

\begin{enumerate}
\def\labelenumi{\arabic{enumi}.}
\tightlist
\item
  a candidate embedded factor, its declared own future, and its environmental interface;
\item
  predictive closure and any required dynamical or recursive closure;
\item
  causal-interface conditions under admissible interventions, together with their evidence;
\item
  transport rules covariant under representation;
\item
  protocols, support, time horizons, errors, and unresolved items.
\end{enumerate}

A historical object additionally includes directed continuation and TCIC/DTRC. Claims of ``local realizability,'' ``recovery,'' ``self-maintenance,'' ``heritability,'' or ``functional purpose'' require the corresponding modules separately.

\textbf{D-OBJ-STRONG.} A candidate that satisfies prespecified predictive, causal, historical, and nontriviality conditions is called a strong operational object in that research context. This defines an explicit certificate class and sufficient criteria for membership. It is not an unconditional necessary-and-sufficient theorem for all natural objects.

Nontriviality must at least specify which detectable distinctions are retained and the baseline against which compression or mechanistic explanation improves. The identity representation of the entire history does not automatically discover a macroscopic object; a wholly constant state does not automatically discover an informative identity either.

\section{12.2 Deriving Ω from a Certificate}\label{deriving-ux3c9-from-a-certificate}

First fix a baseline object certificate and an assessment horizon \(T\), then specify a finite mechanism set \(\mathcal M\). For a removal set \(D\subseteq\mathcal M\), define \[
\chi_{O,T}(D)=1
\] if and only if, after the specified removal operation, the branch still passes the \textbf{previously frozen} certificate. Direct side effects of removal belong to the intervention semantics and must be recorded as well.

If failure is monotone, \[
D\subseteq D',\ \chi(D)=0\Rightarrow\chi(D')=0,
\] define \[
\mathcal C_{\Omega,T}=\min_{\subseteq}\{D:\chi_{O,T}(D)=0\},
\qquad
\Omega_T(O)=(\mathcal M,\mathcal C_{\Omega,T}).
\]

\hypertarget{p-omg-01}{}\phantomsection\label{p-omg-01}

\textbf{P-OMG-01: Reconstruction from minimal damaging sets.} \[
\chi_{O,T}(D)=0
\iff \exists C\in\mathcal C_{\Omega,T},\ C\subseteq D.
\] \textbf{Proof.} For the forward implication, choose a minimal member among the finitely many failing subsets. The converse follows from monotonicity. \(\blacksquare\)

If removing an inhibitor restores function, failure is not monotone. Retain the full response table or a typed mechanism model; do not force the situation into a single hypergraph.

\section{12.3 Integrity Margins}\label{integrity-margins}

For a prespecified model distance \(d\) and a nonempty set \(\mathcal F_O\) of failure models, set \[
\Gamma_O(T)=\inf_{F\in\mathcal F_O}d(T,F).
\]

\hypertarget{p-omg-02}{}\phantomsection\label{p-omg-02}

\textbf{P-OMG-02.} \[
|\Gamma_O(T)-\Gamma_O(S)|\le d(T,S).
\] \textbf{Proof.} For each \(F\), \(d(T,F)\le d(T,S)+d(S,F)\). Take the infimum, then interchange \(S,T\). \(\blacksquare\)

If \(\widehat\Gamma_O>\eta\) and \(d(T,\widehat T)\le\eta\), the true margin is positive. A positive margin gives a controlled distance from that failure set. Merely lying outside a failure set need not give a positive margin if the failure set is not closed.

\section{12.4 Diagnostic Profiles and Original Parameters}\label{diagnostic-profiles-and-original-parameters}

The parameters \(L_c,L_o,\rho,\eta_{\rm pref}\) are retained as domain-specific diagnostic interfaces:

\begin{itemize}
\tightlist
\item
  \(L_c\): describe it through a vector of quantities such as the number of minimal realizations, the structure of damaging sets, and effective model complexity; avoid an arbitrary single total score.
\item
  \(L_o\): record resource flow, interface bandwidth, and dependence on external scaffolding separately, without conflating them with the rate of identity exit.
\item
  For a nonzero drift \(F\) and a set of optimal directions \(D^*\), one may use the signed alignment \[
  A(F,D^*)=\sup_{d\in D^*,\,d\ne0}
  \frac{\langle F,d\rangle}{\|F\|\|d\|}.
  \] Opposite directions give a negative value and orthogonal directions give zero; a sine alone would conflate agreement and opposition. If the drift is zero or no nonzero optimal direction exists, record the quantity as undefined or specify a task-dependent convention.
\item
  A drive-to-regulation capacity ratio requires compatible dimensions, a positive denominator, and independent measurement. It is not a universal psychological index.
\end{itemize}

\section{12.5 Operational-Object Roles and the Boundary of Autogenesis}\label{operational-object-roles-and-the-boundary-of-autogenesis}

\nolinkurl{OperationalObject} records predictive identity, formation provenance, distinguishability under interaction, compositional persistence, physical recoverability, and identity within semantic fibers separately. Passing one item does not fill a gap in another. In particular, statistical recovery of a parent candidate is an epistemic and identification input, not proof that the parent object has been physically constructed.

Chapter 34 expresses the endpoints of P0--P12 as a conditional object lifecycle. The final P12 audit is PARTIAL / EXPLICIT BOUNDARY: organizational construction, mean steady state, resource viability, and a historical identity chain have been obtained, but the inputs still include program provenance and a trusted implementation. CLOSED status of a work package cannot be rewritten as an assumption-free theorem of the origin of life or arbitrary object self-production.

\chapter{13. Realization of Predictive Identity, Redundancy, and Damage Duality}\label{ch-13}

\textbf{Chapter question.} Which readouts recover identity information, and how does redundancy determine access destruction?

\textbf{Reading route.} Sufficient readouts form an upward-closed family whose minimal members describe minimal realizations; hitting every realization produces the blocker. Keep the operation semantics explicit: access erasure changes readable information, whereas physical removal may also change power supply, dynamics, and repair capacity.

\section{13.1 Monotone Access Systems}\label{monotone-access-systems}

\textbf{D-CAR.} Let \(V\) be a finite carrier universe. For \(S\subseteq V\), let the readout \(R_S(H)\) satisfy restriction consistency: if \(S\subseteq T\), then \(R_S\) is a measurable function of \(R_T\).

For a fixed canonical identity \(C\), define \[
\mathcal A_0=\{S:\exists g_S,\ C=g_S(R_S)\},
\qquad
\mathcal M_0=\min_{\subseteq}\mathcal A_0.
\] The equalities use either a common pointwise interpretation or a common modulo-null interpretation under the same reference measure.

\hypertarget{p-car-01}{}\phantomsection\label{p-car-01}

\textbf{P-CAR-01.} The family \(\mathcal A_0\) is upward closed, and on a finite \(V\), every sufficient set contains a minimal sufficient set.

\textbf{Proof.} A larger readout recovers the smaller readout, which can then be followed by the original decoder. Among the finitely many sufficient subsets of a sufficient set, choose a minimal one. \(\blacksquare\)

The family \(\mathcal M_0\) is a clutter: none of its members strictly contains another. It may be empty if the entire device is insufficient, or equal to \(\{\varnothing\}\) if the target is constant. Infinite carrier universes may lack minimal members and are not covered by this theorem.

\section{13.2 Access Erasure and Blockers}\label{access-erasure-and-blockers}

Define \[
b_V(\mathcal M)=
\min_{\subseteq}\{D\subseteq V:\forall M\in\mathcal M,\ D\cap M\ne\varnothing\}.
\]

\hypertarget{p-car-02}{}\phantomsection\label{p-car-02}

\textbf{P-CAR-02.} Under pure access erasure, after deletion of \(D\), the identity \(C\) remains recoverable if and only if some \(M\in\mathcal M_0\) is disjoint from \(D\). Thus the minimal access-damaging sets form \(b_V(\mathcal M_0)\).

\textbf{Proof.} The remaining access set is \(V\setminus D\). It is sufficient if and only if it contains a minimal sufficient set. Negating this condition says exactly that every minimal sufficient set is hit. \(\blacksquare\)

\hypertarget{p-car-03}{}\phantomsection\label{p-car-03}

\textbf{P-CAR-03: Double blocker.} Let \(\mathcal M\) be a clutter on finite \(V\): no member properly contains another. \[
b_V(b_V(\mathcal M))=\mathcal M.
\]

\textbf{Proof.} For any \(S\subseteq V\), if it contains some \(M\in\mathcal M\), it hits every set that hits \(\mathcal M\), and hence hits \(b_V(\mathcal M)\). If it contains no \(M\), its complement \(V\setminus S\) meets every \(M\) and therefore contains a minimal hitting set \(D\) with \(S\cap D=\varnothing\). The sets hitting \(b_V(\mathcal M)\) are thus precisely the upward closure of \(\mathcal M\). Taking minimal members gives the result. \(\blacksquare\)

The degenerate conventions \(b(\varnothing)=\{\varnothing\}\) and \(b(\{\varnothing\})=\varnothing\) preserve the double duality. This finite combinatorial result is an application of standard blocker theory.

\section{13.3 When Access Damage Represents Actual Failure}\label{when-access-damage-represents-actual-failure}

\textbf{H-CAR-FRAME.} Interpreting an access result as physical failure over a time horizon requires at least:

\begin{enumerate}
\def\labelenumi{\arabic{enumi}.}
\tightlist
\item
  a one-to-one correspondence between deletion operations and carrier sets;
\item
  preservation of the remaining readouts and their dynamics as specified by the model;
\item
  equivalence between sufficiency of the remaining readout and the target certificate for that horizon;
\item
  inclusion of repair, compensation, power supply, control, and regeneration mechanisms in the model.
\end{enumerate}

A prohibition on repair cannot replace item 2. Consider two bit backups \(a,b\), where \(a\) also powers \(b\). Access damage requires erasing both, whereas physical deletion of \(a\) can cause both to fail subsequently.

There are therefore two related but distinct hypergraphs: the blocker of access realizations and \(\Omega_T\) for actual mechanism interventions. They can be identified only when the frame conditions hold.

\section{13.4 Defect Diameter and Decoding Radius}\label{defect-diameter-and-decoding-radius}

For a finite history and protocol table, define \[
e(S)=\max_{R_S(h)=R_S(h')}\max_i
D_{\rm TV}(p_{h,i},p_{h',i}).
\] Then \(e(S)=0\) exactly when the readout is sufficient, and adding carriers makes \(e(S)\) nonincreasing.

Define the optimal worst-case decoding radius by \[
r(S)=\inf_{\{q_{z,i}\}}\max_{h,i}
D_{\rm TV}(p_{h,i},q_{R_S(h),i}).
\]

\hypertarget{p-car-04}{}\phantomsection\label{p-car-04}

\textbf{P-CAR-04.} \[
\tfrac12e(S)\le r(S)\le e(S).
\] \textbf{Proof.} For every center, the distance between two true distributions in the same fiber is at most the sum of their errors relative to the center, giving the lower bound. Choosing the response of one representative history as the center in each readout fiber gives the upper bound. \(\blacksquare\)

This connects the response diameter at a common readout precisely with optimal decoding error. They cannot simply be assigned the same numerical value. Subsequent threshold recovery uses the explicitly defined \(e(S)\).

\section{13.5 From Preservation of Registered Interfaces to Physical Realization}\label{from-preservation-of-registered-interfaces-to-physical-realization}

The common-source core can place service or program attributes that must be retained among the outputs of a finite model and preserve them through the coarsest stable quotient. A carrier readout can then specify exactly which registered target it recovers. This strengthens the predictive-identity--carrier connection without removing H-CAR-FRAME.

For example, the same pair of redundant readouts may be implemented with independent power supplies or a single shared supply. Their access blockers agree while their actual damaging families differ. Even if a downstream policy has been formally proved to keep the source model safe, omission of a shared power supply remains a physical-adaptation issue. Numerical comparison of failure thresholds also requires the defect calibration of Chapter 33.

\chapter{14. Resolution, Coarse-Graining, and the Uncertainty Interval for Microscopic Realizations}\label{ch-14}

\textbf{Chapter question.} How much microscopic structure is determined by observations at a coarse resolution?

\textbf{Reading route.} Closure first gives the coarsening of minimal sufficient sets; Boolean deletion then gives blocker naturality. Two adjoint extensions bound all compatible microscopic realizations. Singleton active fibers characterize uniqueness, while the target, horizon, and underlying carrier remain fixed throughout the comparison.

\section{14.1 General Closure Systems}\label{general-closure-systems}

\textbf{D-RES.} On a finite carrier universe \(V\), a closure system \(\mathcal L\subseteq2^V\) contains \(V\) and is closed under intersections. Define \[
\operatorname{cl}_{\mathcal L}(S)=
\bigcap\{A\in\mathcal L:S\subseteq A\}.
\] A coarser scale has a smaller expressible family: \(\mathcal L_r\subseteq\mathcal L_s\).

For a fixed upward-closed property \(\mathcal A\), let \(\mathcal M_r=\min(\mathcal A\cap\mathcal L_r)\).

\hypertarget{p-res-01}{}\phantomsection\label{p-res-01}

\textbf{P-RES-01: Coarse-graining minimal sufficient sets.} \[
\mathcal M_r=
\min_{\subseteq}\{\operatorname{cl}_r(M):M\in\mathcal M_s\}.
\]

\textbf{Proof.} Every coarse minimal sufficient set \(A\) is also sufficient in the fine family, so it contains a fine minimal sufficient set \(M\). Its coarse closure still contains \(M\), belongs to the coarse family, remains sufficient, and lies within \(A\). Coarse minimality forces it to equal \(A\). Conversely, the coarse closure of every fine minimal sufficient set is sufficient, so it contains a coarse minimal sufficient set. The two families generate the same upward closure and therefore have the same minimal members. \(\blacksquare\)

\hypertarget{p-res-02}{}\phantomsection\label{p-res-02}

\textbf{P-RES-02: Composition law.} If \(\mathcal L_r\subseteq\mathcal L_s\subseteq\mathcal L_t\), then \[
\operatorname{cl}_r\operatorname{cl}_s=\operatorname{cl}_r.
\] On minimal-set families, define \(\mathcal R_{r\leftarrow s}(\mathcal M)=
\min\{\operatorname{cl}_r(M):M\in\mathcal M\}\). Then \[
\mathcal R_{r\leftarrow t}
=\mathcal R_{r\leftarrow s}\mathcal R_{s\leftarrow t}.
\]

\textbf{Proof.} Since \(S\subseteq\operatorname{cl}_sS\subseteq\operatorname{cl}_rS\), taking the coarse closure again gives the equality. Discarding nonminimal members containing other members at an intermediate step does not alter the upward closure generated by subsequent closures. \(\blacksquare\)

The property, target, time horizon, and underlying carrier universe must remain fixed. Changing them is not the same pure change of resolution.

\section{14.2 Boolean Deletion Semantics}\label{boolean-deletion-semantics}

A general closure system does not ensure blocker duality. For example, let \[
\mathcal L=\{\varnothing,\{1\},\{1,2\}\},\quad
\mathcal M=\{\{1,2\}\}.
\] If hitting sets are restricted to \(\mathcal L\), the first blocker is \(\{\{1\}\}\) and the second is again the same, rather than \(\mathcal M\).

Full access-deletion semantics therefore uses a finite atomic Boolean algebra. Coarsening is represented by a surjection \(\pi:V\twoheadrightarrow W\), where each coarse atom is an entire fiber. Define \[
\pi_*\mathcal M=\min\{\pi(M):M\in\mathcal M\}.
\]

\hypertarget{p-res-03}{}\phantomsection\label{p-res-03}

\textbf{P-RES-03: Naturality of blockers across scales.} \[
b_W(\pi_*\mathcal M)=\pi_*b_V(\mathcal M).
\]

\textbf{Proof.} A set \(T\subseteq W\) hits \(\pi_*\mathcal M\) if and only if the saturated inverse image \(\pi^{-1}T\) hits every \(M\). This holds if and only if \(\pi^{-1}T\) contains some fine minimal hitting set \(D\), equivalently \(\pi(D)\subseteq T\). Thus the coarse hitting-set family is precisely the upward closure of the images of fine minimal hitting sets. Take minimal members. \(\blacksquare\)

Coarse deletion here removes an entire fiber. It does not assert that its cost equals the cost of deleting a single fine atom.

\section{14.3 Three Adjoint Operations}\label{three-adjoint-operations}

For a fine upward-closed family \(\mathcal A_s\) and a coarse upward-closed family \(\mathcal A_r\), define \[
\operatorname{Res}_\pi(\mathcal A_s)
=\{B\subseteq W:\pi^{-1}B\in\mathcal A_s\},
\] \[
\operatorname{Ext}_\pi^-(\mathcal A_r)
=\{S\subseteq V:\exists B\in\mathcal A_r,\ \pi^{-1}B\subseteq S\},
\] \[
\operatorname{Ext}_\pi^+(\mathcal A_r)
=\{S\subseteq V:\pi(S)\in\mathcal A_r\}.
\]

\hypertarget{p-res-04}{}\phantomsection\label{p-res-04}

\textbf{P-RES-04: An adjoint triple.} \[
\operatorname{Ext}^-_\pi\dashv\operatorname{Res}_\pi
\dashv\operatorname{Ext}^+_\pi.
\] Explicitly, \[
\operatorname{Ext}^-_\pi(\mathcal A_r)\subseteq\mathcal A_s
\iff\mathcal A_r\subseteq\operatorname{Res}_\pi(\mathcal A_s),
\] \[
\operatorname{Res}_\pi(\mathcal A_s)\subseteq\mathcal A_r
\iff\mathcal A_s\subseteq\operatorname{Ext}^+_\pi(\mathcal A_r).
\]

\textbf{Proof.} Both sides of the first equivalence say that the full inverse image of every coarse sufficient set is fine sufficient. For the forward direction of the second, if \(S\in\mathcal A_s\), its saturated extension \(\pi^{-1}\pi(S)\) is also sufficient, so \(\pi(S)\in\mathcal A_r\). For the converse, take \(S=\pi^{-1}B\) and use the surjective identity \(\pi(\pi^{-1}B)=B\). \(\blacksquare\)

\section{14.4 The Exact Microscopic Uncertainty Interval}\label{the-exact-microscopic-uncertainty-interval}

For a coarse minimal family \(\mathcal D\), define \[
J^-_\pi(\mathcal D)=\{\pi^{-1}D:D\in\mathcal D\},
\] \[
J^+_\pi(\mathcal D)=
\min\{S\subseteq V:\pi(S)\in\uparrow\mathcal D\}.
\]

\hypertarget{p-res-05}{}\phantomsection\label{p-res-05}

\textbf{P-RES-05.} \[
\pi_*\mathcal C=\mathcal D
\iff
\uparrow J^-_\pi(\mathcal D)
\subseteq\uparrow\mathcal C
\subseteq\uparrow J^+_\pi(\mathcal D).
\]

\textbf{Proof.} If the coarse image is \(\mathcal D\), each \(D\) is the image of a fine minimal edge, so its full inverse image contains a fine sufficient edge. Every fine sufficient edge also has an image containing a coarse minimal edge. These give the two inclusions. Conversely, the left inclusion gives, for each \(D\), an \(F\in\mathcal C\) with \(F\subseteq\pi^{-1}D\). The right inclusion gives a \(D'\in\mathcal D\) with \(D'\subseteq\pi(F)\subseteq D\). No member of a minimal family strictly contains another, so \(D'=D=\pi(F)\). Moreover, every fine edge has an image containing a coarse minimal edge. Hence the coarse minimal image is exactly \(\mathcal D\). \(\blacksquare\)

\hypertarget{p-res-06}{}\phantomsection\label{p-res-06}

\textbf{P-RES-06: Object-relative uniqueness.} The two endpoints above coincide if and only if \[
|\pi^{-1}(w)|=1
\quad\forall w\in\bigcup\mathcal D.
\]

\textbf{Proof.} If all active fibers are singletons, a minimal section satisfying coarse sufficiency is the full inverse image, so the endpoints agree. If some active \(w\) has two atoms, take a coarse minimal edge containing \(w\), denoted by \(D\), and choose one atom from each fiber to form a section \(S\). This belongs to the upper endpoint but does not contain \(\pi^{-1}D\). Nor does it contain the full inverse image of any other coarse minimal edge, since such an edge cannot be properly contained in \(D\). Thus \(S\) does not belong to the lower endpoint. \(\blacksquare\)

Coarse-scale data generally determine a family of admissible realizations, not a unique microscopic boundary. Inactive fibers do not affect this object-specific conclusion.

\section{14.5 Resolution Types and Transport across Scales}\label{resolution-types-and-transport-across-scales}

The closure operations, Boolean deletion, and realization adjunctions in this chapter belong to H0-RES. They are related to the predictive quotients of H0-PQ, kernel commutation of H0-DYN, parent formation of H0-COMP, and historical transport of H0-HIST, but are not the same operator. Chapter 32 lists eight types and the permitted bridges.

\nolinkurl{ObjectScaleCalculus} composes scale maps and transports identity, purpose, GOD, GOA, repair, and steady-state properties when separate preservation certificates for those properties are supplied. The current \nolinkurl{WilsonianRGInterpretation} records only couplingFlow; it supplies no compatibility relation with an actual scale map or source dynamics. Its theorem that two constant flows differ cannot be called a nonuniqueness theorem for Wilson flows associated with a specified physical coarse-graining.

\chapter{15. Composition, Overlap, and the Formation of Parent Objects}\label{ch-15}

\textbf{Chapter question.} When do subsystems form a parent object with its own predictive and intervention structure?

\textbf{Reading route.} First control for spurious irreducibility caused by common external drivers, then check path nonfactorization, cut effects, and additional parent-level information separately. Lifting physical carriers to child coalitions requires a coverage relation. Formation, binding, and historical transport combine only when their respective conditions are supplied.

\section{15.1 Three Forms of Commonality}\label{three-forms-of-commonality}

A common external cause, shared predictive information, or a common objective among several systems does not alone constitute a parent object. A parent must have its own canonical future, dynamical interface, and historical certificate. Composition additionally asks whether it is reducible relative to a declared partition into subsystems.

For child paths \(\mathbf X_1,\ldots,\mathbf X_m\) and an external-interface record \(U\), define \[
I_\pi=
E_U D_{\rm KL}\left(P_{\mathbf X_{1:m}\mid U}
\middle\|\bigotimes_{B\in\pi}P_{\mathbf X_B\mid U}\right).
\]

\hypertarget{p-comp-01}{}\phantomsection\label{p-comp-01}

\textbf{P-COMP-01.} The equality \(I_\pi=0\) holds if and only if the path blocks of that partition are independent conditional on \(U\). For finitely many nontrivial partitions, \[
\Sigma_{\rm path}=\min_\pi I_\pi>0
\] if and only if no such path factorization exists.

\textbf{Proof.} Each conditional KL divergence is nonnegative, so its mean is zero if and only if the two conditional laws agree almost surely. The minimum of a finite family is positive if and only if every member is positive. \(\blacksquare\)

For a common-drive example, let \(X_1=U+\epsilon_1,\ X_2=U+\epsilon_2\) with conditionally independent noises. Raw correlation can be large while conditional integration is zero. If \(U\) is a consequence of internal outcomes, conditioning may instead introduce collider bias. The causal status of the external interface must be specified first.

\section{15.2 Interventional Irreducibility}\label{interventional-irreducibility}

For each nontrivial partition \(\pi\), specify an implementable cut \(\mathcal K_\pi\) that changes cross-block connections while preserving mechanisms within blocks. Compare the original record law \(P\) with the cut law \(P^\pi_{\rm cut}\).

One may use total variation, or \[
JS(P,Q)=\tfrac12D_{\rm KL}(P\|M)+\tfrac12D_{\rm KL}(Q\|M),
\quad M=(P+Q)/2.
\]

\hypertarget{p-comp-02}{}\phantomsection\label{p-comp-02}

\textbf{P-COMP-02.} One has \(0\le JS(P,Q)\le\log2\), with \(JS=0\iff P=Q\). Consequently, the minimum JS over a finite cut family is positive if and only if every cut changes the declared record law.

\textbf{Proof.} Nonnegativity follows from KL. The inequalities \(M\ge P/2,Q/2\) give the upper bound, while zero forces \(P=M=Q\). \(\blacksquare\)

Failure to detect a cut effect does not imply the complete absence of microscopic coupling. That interpretation requires additional observation completeness and cut-faithfulness: correctly cutting a genuinely relevant coupling must produce an effect in the admissible records.

\section{15.3 Composition Margins}\label{composition-margins}

Let \(d\) be a model distance. Assume each cut satisfies \(d(\mathcal K_\pi T,\mathcal K_\pi S)\le L_\pi d(T,S)\). Define \(\Gamma_{\rm comp}(T)=\min_\pi d(T,\mathcal K_\pi T)\).

\hypertarget{p-comp-03}{}\phantomsection\label{p-comp-03}

\textbf{P-COMP-03.} If \(L=\max_\pi L_\pi<\infty\), then \[
|\Gamma_{\rm comp}(T)-\Gamma_{\rm comp}(S)|
\le(1+L)d(T,S).
\]

\textbf{Proof.} The difference between each pair of distances is at most \(d(T,S)+d(\mathcal K_\pi T,\mathcal K_\pi S)\). Taking minima over the same finite index family preserves the bound. \(\blacksquare\)

For nonexpansive cuts, the constant is 2. This constant cannot be applied directly to a refitted-and-recut model unless nonexpansiveness has been proved.

\section{15.4 Additional Predictive Information at the Parent Level}\label{additional-predictive-information-at-the-parent-level}

Let \(C_1,\ldots,C_m\) be child-level canonical cores, \(C_P\) a parent-level core, and \(U\) an interface, with all relevant conditional entropies finite. Define \[
\Delta_P=H(C_P\mid C_1,\ldots,C_m,U).
\]

\hypertarget{p-comp-04}{}\phantomsection\label{p-comp-04}

\textbf{P-COMP-04.} If the parent representation \(M_P\) is exactly sufficient, then \[
H(M_P\mid C_1,\ldots,C_m,U)\ge\Delta_P.
\]

\textbf{Proof.} Minimality gives \(C_P=g(M_P,U)\). Conditional on the child cores and interface, the joint entropy \(H(M_P,C_P\mid C_{1:m},U)=H(M_P\mid C_{1:m},U)\) is at least \(H(C_P\mid C_{1:m},U)\). \(\blacksquare\)

A consistent process example preserves the parity \(S=X_1\oplus X_2\) of a bit pair forever. At each step, resample a uniform \(R\) and set \[
X'_1=R,\qquad X'_2=R\oplus S.
\] Given the complete history, the entire future of either individual child coordinate is an independent uniform sequence, so its individual predictive core is constant. The joint future, however, retains the parity \(S\) exactly. Under a uniform initial law, \(\Delta_P=\log2\). This demonstrates joint predictive information; it does not by itself prove that both coordinates are complete child objects.

\section{15.5 Coalitions of Child Objects as a Second Carrier Level}\label{coalitions-of-child-objects-as-a-second-carrier-level}

For a certified set of child objects \([m]\), define the upward-closed family of coalitions that recover the parental core and its minimal family \(\mathcal M_P^{\rm child}\). Under monotone erasure of logical interfaces, the blocker again gives the minimal damaging coalitions.

The intersection \(\bigcap_{M\in\mathcal M_P^{\rm child}}M\) consists of child objects required in every sufficient coalition. An empty intersection indicates redundant organizational possibilities, not the nonexistence of a parent object.

\section{15.6 Canonical Booleanization of Overlap Geometry}\label{canonical-booleanization-of-overlap-geometry}

For overlapping regions \(V_1,\ldots,V_m\subseteq V\), define the membership signature \[
s(x)=\{i:x\in V_i\},\quad A_J=\{x:s(x)=J\}.
\]

\hypertarget{p-comp-05}{}\phantomsection\label{p-comp-05}

\textbf{P-COMP-05.} The nonempty sets \(A_J\) are the atoms of the unique least Boolean algebra containing all \(V_i\).

\textbf{Proof.} Each point has a unique signature, so these sets partition \(V\). Moreover, \[
A_J=\bigcap_{i\in J}V_i\cap\bigcap_{i\notin J}V_i^c.
\] Every Boolean algebra containing the generating regions must contain these atoms, while each generating region is a union of the corresponding atoms. The resulting algebra is therefore least and unique. \(\blacksquare\)

A nerve recording only which intersections are nonempty does not generally recover complete membership signatures or physical deletion semantics.

For a child coalition \(S\), let \(V(S)=\bigcup_{i\in S}V_i\). If coalition sufficiency is defined exactly by sufficiency of its union region, then

\hypertarget{p-comp-06}{}\phantomsection\label{p-comp-06}

\textbf{P-COMP-06: Lifting physical carriers to coalitions.} \[
\mathcal M_P^{\rm child}
=\min_{\subseteq}
\bigcup_{M\in\mathcal M_P^{\rm phys}}
\min_{\subseteq}\{S:M\subseteq V(S)\}.
\]

\textbf{Proof.} The union region of a sufficient coalition contains a physical minimal carrier, so the coalition contains a minimal subcoalition covering that carrier. Conversely, every such cover is sufficient. \(\blacksquare\)

Logically disabling a child interface does not automatically delete a shared physical region. The two blockers have different underlying sets and operational meanings.

\section{15.7 Reconfiguration and Formation Events}\label{reconfiguration-and-formation-events}

If the regional structure changes, first pull back to a common underlying space through microscopic transport, then form the common Boolean refinement generated by both families of regions and compare carriers and mechanisms. Without nesting, a one-way resolution arrow must not be invented.

\begin{itemize}
\tightlist
\item
  Appearance of a parent object: a new parent certificate holds while the child branches continue.
\item
  Fusion: a new branch succeeds several predecessors, some of whose former branches no longer continue independently.
\item
  Splitting: one predecessor produces several simultaneously independent certified successors.
\item
  Birth/death: absence of a certified \nolinkurl{predecessor/successor} relative to the admissible continuation relation.
\item
  GOA bifurcation: a change in the long-term closed-loop structure, not identical to object birth or fusion.
\end{itemize}

\section{15.8 Composition Windows Preserving Individuality}\label{composition-windows-preserving-individuality}

\hypertarget{p-comp-07}{}\phantomsection\label{p-comp-07}

\textbf{P-COMP-07.} On a compact coupling-parameter interval \(K\), suppose the integration quantity \(D(k)\) is continuous and nondecreasing, and child fidelity \(F(k)\) is continuous and nonincreasing. Assume both threshold-feasible sets are nonempty, and set \[
k_D=\min\{k:D(k)\ge d_*\},\quad
k_F=\max\{k:F(k)\ge f_*\}.
\] Then the joint feasible set is \([k_D,k_F]\) if \(k_D\le k_F\), and is empty otherwise.

\textbf{Proof.} Continuity ensures attainment of the boundary values. Monotonicity makes the two feasible sets an upper and a lower interval, respectively. Their intersection is the stated interval. \(\blacksquare\)

Monotonicity is an assumption about the particular model. Natural systems may have several disjoint windows; stronger coupling cannot simply be assumed to improve integration.

\section{15.9 Actual Obligations for Composing Formation, Binding, and Transport}\label{actual-obligations-for-composing-formation-binding-and-transport}

Let \(f_0,f_1\) form candidates at the two ends, \(t_x\) transport underlying inputs, \(B_0,B_1\) implement binding, and \(t_p\) transport the parent level. Approximate commutation of formation and approximate commutation of binding are separate assumptions. If \(B_1\) is \(L\)-Lipschitz for the intermediate metric, then

\[d\bigl(B_1(f_1(t_xx)),t_p(B_0(f_0x))\bigr)\le L\epsilon_F+\epsilon_B.\]

For the proof, insert \(B_1(t_f(f_0x))\), apply Lipschitz continuity to the first part and the binding-transport bound to the second. C4's \nolinkurl{certified_formed_structural_continuation} additionally requires actual formation relations at both ends and budgets within the margin. It does not identify a structural distance bound with equality of parent-token identity. Complete types appear in §34.4.

\chapter{16. Certification from Finite Data: Responses, Carriers, and Protocols}\label{ch-16}

\textbf{Chapter question.} Which structures can finite samples certify, and which judgments must remain unresolved?

\textbf{Reading route.} Establish a simultaneous confidence event for fixed response rows, then transport distances, thresholds, and carrier conclusions within that event. Exact structural recovery relies on an independently supplied separation gap. Training distributions, deployment occupancy, and post-data selection retain their separate sources.

\section{16.1 Finite Response Tables and Simultaneous Confidence Bounds}\label{finite-response-tables-and-simultaneous-confidence-bounds}

Consider \(L\) history--protocol cells, with \(N\) independent future samples per cell and a common outcome alphabet of size \(K\). Independence is required only within each cell; the union bound across cells does not require independence. Write \(p_j,\widehat p_j\) for the true and empirical responses.

\hypertarget{p-stat-01}{}\phantomsection\label{p-stat-01}

\textbf{P-STAT-01: A simultaneous finite-alphabet TV bound.} \[
P\left(\max_{j\le L}D_{\rm TV}(p_j,\widehat p_j)>\eta\right)
\le L\,2^{K+1}e^{-2N\eta^2}.
\] Consequently, one may take \[
\eta_{N,K,L}(\alpha)=
\min\left\{1,\sqrt{\frac{(K+1)\log2+\log(L/\alpha)}{2N}}\right\}.
\]

\textbf{Proof.} For every subset \(A\) of outcomes, its empirical mass is a Bernoulli mean. Hoeffding's inequality bounds the two-sided error probability by \(2e^{-2N\eta^2}\). Take a union bound over at most \(2^K\) subsets and \(L\) cells. On a finite space, TV is exactly the maximum of these event differences. If the formula exceeds 1, the deterministic TV bound of 1 applies directly. \(\blacksquare\)

For Bernoulli responses, the sharper bound is available directly: \[
\eta_{N,L}(\alpha)=
\min\left\{1,\sqrt{\frac{\log(2L/\alpha)}{2N}}\right\}.
\] With different sample sizes across cells, allocate separate \(\alpha_j\) and then take the maximum valid cell radius. A continuous-valued conditional history cannot be treated as independently resampled at the same history merely because it has occurred once.

\section{16.2 Defect Error Does Not Add a Dependence on the Number of Carrier Candidates}\label{defect-error-does-not-add-a-dependence-on-the-number-of-carrier-candidates}

\hypertarget{p-stat-02}{}\phantomsection\label{p-stat-02}

\textbf{P-STAT-02.} On the simultaneous response-error event, \[
|\widehat e(S)-e(S)|\le2\eta
\quad\text{for every carrier }S.
\]

\textbf{Proof.} The estimation error of the TV distance between any response pair is at most the sum of its endpoint errors. Taking the maximum over the same readout fibers and protocol indices still gives a bound of \(2\eta\). The event already covers every underlying response, so no additional probability allocation for each \(S\) is needed. \(\blacksquare\)

This reduces the statistical union-bound burden, but does not remove the computational cost of enumerating \(2^{|V|}\) carriers. If readout maps are themselves estimated from data, the risk of incorrectly identifying their fibers must be accounted for separately.

\section{16.3 Approximate and Exact Families}\label{approximate-and-exact-families}

Define \[
\mathcal M_\delta=\min\{S:e(S)\le\delta\}.
\]

\hypertarget{p-stat-03}{}\phantomsection\label{p-stat-03}

\textbf{P-STAT-03: Recovery at a fixed tolerance.} If \[
\min_S|e(S)-\delta|>2\eta,
\] the empirical threshold \(\widehat e(S)\le\delta\) recovers \(\mathcal M_\delta\).

\textbf{Proof.} Every true defect is farther from the threshold than the possible estimation error. All classifications are unchanged; then take minimal members. \(\blacksquare\)

\hypertarget{p-stat-04}{}\phantomsection\label{p-stat-04}

\textbf{P-STAT-04: Recovery of the exact family.} Suppose a nonzero defect gap \[
\Delta=\min_{S:e(S)>0}e(S)>0
\] is supplied by a structural assumption or an independent proof, and \[
2\eta<\tau<\Delta-2\eta,
\] Then empirical thresholding recovers \(\mathcal M_0\).

\textbf{Proof.} A true zero defect has estimate at most \(2\eta\); a true nonzero defect has estimate at least \(\Delta-2\eta\). The two classes are separated. \(\blacksquare\)

If no nonzero candidate exists, the degenerate case in which all candidates are sufficient is handled directly. The gap condition \(\Delta>4\eta\) suffices for such a threshold to exist.

For a counterexample, take response parameters \(0.50,0.55\). The defect without a carrier is \(0.05\), while reading the history gives defect zero. With \(\delta=0.1\), even when \(\eta=0\), the minimal approximate family is \(\{\varnothing\}\) and differs from the exact minimal family.

\section{16.4 Three-Way Certification and Confidence Sets}\label{three-way-certification-and-confidence-sets}

For a defect confidence interval \[
[\ell_S,u_S]=
[\max(0,\widehat e(S)-2\eta),\min(1,\widehat e(S)+2\eta)],
\] when \(u_S\le\delta\), certify \(\delta\)-sufficiency; exclude it when \(\ell_S>\delta\), and otherwise leave the case unresolved.

Retain the minimal carrier families of all true defect tables compatible with the confidence intervals and monotonicity as a structural confidence set. A confidence interval containing zero does not prove that the true defect equals zero. Exact structural conclusions require a sampling bound \textbf{together with} a structural gap, not unconditional detection of equality to zero from finite noisy data.

\section{16.5 Recovery of Canonical Predictive Classes}\label{recovery-of-canonical-predictive-classes}

\hypertarget{p-stat-05}{}\phantomsection\label{p-stat-05}

\textbf{P-STAT-05.} For a finite history set, consider the distance \[
d(h,h')=\max_iD_{\rm TV}(p_{h,i},p_{h',i})
\] and suppose distinct true predictive classes are separated by at least \(\gamma>0\), with \(2\eta<\gamma/2\). Thresholding empirical distances at \(\gamma/2\) then recovers the exact classes.

\textbf{Proof.} Within-class empirical distances are at most \(2\eta\); between-class distances are at least \(\gamma-2\eta\). On this event the threshold adjacency relation is exactly the true equivalence relation. There is no need to infer classes through uncontrolled chaining of nearest neighbors. \(\blacksquare\)

With an unknown gap, the threshold relation may fail to be transitive. Report candidate partitions and unresolved pairs; connectivity in single-linkage clustering is not mathematical identity.

\section{16.6 Characteristic Embeddings for Continuous Futures}\label{characteristic-embeddings-for-continuous-futures}

Suppose a measurable positive-definite kernel \(k\) has a separable RKHS, satisfies \(k(y,y)\le1\), and has a Bochner-integrable feature map. Define \[
\mu_P=E_Pk(Y,\cdot),\qquad
\operatorname{MMD}_k(P,Q)=\|\mu_P-\mu_Q\|.
\] \textbf{K-KER-01.} On the declared distribution class, characteristic means \(\mu_P=\mu_Q\iff P=Q\). This provides a faithful representation of continuous future responses; it does not automatically guarantee learning of conditional laws for continuous histories. \href{https://www.jmlr.org/papers/volume13/gretton12a/gretton12a.pdf}{Original research}

\hypertarget{p-stat-06}{}\phantomsection\label{p-stat-06}

\textbf{P-STAT-06: Simultaneous embedding error.} For \(L\) fixed conditional responses, each with \(N\) independent samples, \[
\max_j\|\widehat\mu_j-\mu_j\|
\le r_N(\alpha)=
\frac{1+\sqrt{2\log(L/\alpha)}}{\sqrt N}
\] with probability at least \(1-\alpha\).

\textbf{Proof.} Cross inner products of independent centered features have mean zero, so \(E\|\widehat\mu-\mu\|^2\le1/N\). Jensen's inequality then bounds the mean norm by \(N^{-1/2}\). Replacing one sample changes the norm statistic by at most \(2/N\). McDiarmid's inequality bounds an upward deviation of \(t\) from its mean by \(e^{-Nt^2/2}\); take a union bound over the \(L\) responses. \(\blacksquare\)

If the interclass MMD gap is \(\gamma\) and \(r_N<\gamma/4\), thresholding at \(\gamma/2\) recovers the classes by the same proof as P-STAT-05.

\hypertarget{p-stat-07}{}\phantomsection\label{p-stat-07}

\textbf{P-STAT-07: Covariance under kernel transport.} For a bimeasurable bijection \(T\), define the transported kernel simultaneously by \(k_T(Tx,Ty)=k(x,y)\). Then \[
\operatorname{MMD}_{k_T}(T_\#P,T_\#Q)
=\operatorname{MMD}_k(P,Q).
\] \textbf{Proof.} Each of the three kernel expectations in squared MMD is unchanged by the variable transformation. \(\blacksquare\)

Small MMD does not generally imply small TV. For control-value bounds, one can instead use, directly for value functions of bounded RKHS norm, \(|Pv-Qv|\le\|v\|_{\mathcal H_k}\operatorname{MMD}(P,Q)\), or separately prove a conversion to TV. The characteristic property cannot substitute for an error-scale comparison.

\section{16.7 Protocol Plateaus and Repeated Testing}\label{protocol-plateaus-and-repeated-testing}

For a prespecified finite batch of additional protocols, suppose every response that should be distinguished is separated by at least \(\gamma\), and all embedding errors are below \(\gamma/4\). No detectable refinement within that batch is then missed. The error probability is controlled by the simultaneous bound.

This says nothing about untested protocols. Repeatedly inspecting data until a test passes requires independent holdout data, valid confidence sequences, or prespecified roundwise failure budgets with \(\sum_k\alpha_k\le\alpha\). A fixed \(95\%\) guarantee in every round does not automatically give a joint \(95\%\) guarantee.

\section{16.8 Dependent Trajectories, Discovery Consistency, and Model Selection}\label{dependent-trajectories-discovery-consistency-and-model-selection}

\hypertarget{p-stat-08}{}\phantomsection\label{p-stat-08}

\textbf{P-STAT-08: Discovery among finitely many candidates.} Consider \(m\) candidates with bounded losses \(\ell_f\in[0,1]\) and \(N\) independent validation samples. Set \[
u=\sqrt{\frac{\log(2m/\alpha)}{2N}}.
\] An empirical risk minimizer \(\widehat f\) satisfies \[
R(\widehat f)\le\min_fR(f)+2u
\] with probability at least \(1-\alpha\).

\textbf{Proof.} Hoeffding's inequality and a union bound give risk errors at most \(u\) for every candidate. Insert the empirical optimum and the empirical risk of a true optimal candidate. \(\blacksquare\)

For a stationary ergodic sequence, the ergodic theorem gives convergence of every sample average in a finite family and hence asymptotic consistency. Without additional mixing assumptions, it does not supply the finite-sample rate above.

For an absolutely regular (\(\beta\)-mixing) process, a coupling lemma applies with gaps of length \(b\) between \(m\) sample blocks: the TV distance between their joint law and that of independent copies with the same block marginals is at most \((m-1)\beta(b)\). This penalty is added to an independent-block concentration bound. Filtering the blocks by rare histories does not automatically preserve the same mixing assumptions or sample-size guarantees.

\section{16.9 Deployment Distributions and Coverage}\label{deployment-distributions-and-coverage}

\hypertarget{p-stat-09}{}\phantomsection\label{p-stat-09}

\textbf{P-STAT-09: Transport of average error.} For \(0\le e\le1\), if the deployment distribution satisfies \(\nu\ll\mu\) and \(d\nu/d\mu\le C\), then \[
E_\nu e\le C E_\mu e.
\] If \(D_{\rm TV}(\nu,\mu)\le\delta\), then \[
E_\nu e\le E_\mu e+\delta.
\]

\textbf{Proof.} Integrate the density bound for the first inequality. For the second, use P-MET-02 and \(\operatorname{sp}(e)\le1\). \(\blacksquare\)

States that are rare during training but frequently visited by a control policy may determine the entire failure risk. The effective support of object certification must therefore match the protocols, policies, and environments in actual use.

\section{16.10 The Good Event, Data-Dependent Selection, and Actual Deployment}\label{the-good-event-data-dependent-selection-and-actual-deployment}

Fix the registered row set before training. Let Good be the event that every true registered row is within its specified TV radius of the empirical row. The previous simultaneously calibrated Good subset is contained in this canonical event, so the latter has the same upper bound on bad-training mass. Selection of a representation, reference model, or policy may depend on the training record, provided that the needed fit, coverage, and numerical budgets hold on the same Good event.

Chapter 38 constructs the training--deployment joint law directly and proves that these Good-only assumptions suffice for the overall risk bound \(\alpha+(1-\alpha)\beta\). It does not require every extremely rare bad sample to satisfy a small budget. Sampling only some registered rows still does not automatically ensure coverage of deployment visits: Chapter 36 retains the actual \(\nu_n(C_n^c)\), and Chapter 37 identifies a class in which a zero coverage penalty follows from actual support viability.

\chapter{17. Identifiability of Inverse Objects, Dynamics, and Purposes}\label{ch-17}

\textbf{Chapter question.} When records are used to infer a model, is the identified object a parameter, a mechanism, or a purpose-equivalence class?

\textbf{Reading route.} Use experimental equivalence to identify what cannot be distinguished, then study separation through testing distance, Fisher null directions, and design matrices. Local full rank is not global uniqueness. A predictable-design condition cannot be replaced by a statement about the marginal noise law.

\section{17.1 Experimental Equivalence Classes and Optimal Testing}\label{experimental-equivalence-classes-and-optimal-testing}

If two models produce the same record law for every admissible experiment, no statistical rule based only on those records can distinguish them.

\hypertarget{p-inv-01}{}\phantomsection\label{p-inv-01}

\textbf{P-INV-01: A binary testing bound.} For two equally likely models with data laws \(P_0,P_1\), the least average error is \[
R^*=\tfrac12(1-D_{\rm TV}(P_0,P_1)).
\]

\textbf{Proof.} Let \(A\) be the region in which model one is selected. The error is \[
R(A)=\tfrac12[P_0(A)+P_1(A^c)]
=\tfrac12[1-(P_1(A)-P_0(A))].
\] The optimal event difference is TV. On a general space, comparing the densities relative to \(P_0+P_1\) gives an event attaining the optimum. \(\blacksquare\)

Accurate prediction, boundary identification, dynamics identification, and reward identification have different equivalence classes. Success in one does not substitute for success in all of them.

\section{17.2 Gauge Null Directions of Fisher Information}\label{gauge-null-directions-of-fisher-information}

A regular statistical model \(p_\theta\) has score \(s_\theta=\nabla\log p_\theta\) and \(\mathcal I=E_\theta s_\theta s_\theta^\top\).

\hypertarget{p-inv-02}{}\phantomsection\label{p-inv-02}

\textbf{P-INV-02.} If the likelihood is invariant along a smooth group orbit, every orbit tangent lies in \(\ker\mathcal I\).

\textbf{Proof.} For an orbit tangent \(v\), the directional derivative of the likelihood is zero. Thus \(v^\top s_\theta=0\) almost surely and consequently \(\mathcal I v=0\). \(\blacksquare\)

\hypertarget{p-inv-03}{}\phantomsection\label{p-inv-03}

\textbf{P-INV-03.} For conditionally independent experiments satisfying the zero-mean score condition, \[
\mathcal I_{\rm total}=\sum_e\mathcal I_e,\qquad
\ker\mathcal I_{\rm total}=\bigcap_e\ker\mathcal I_e.
\]

\textbf{Proof.} The joint score is a sum and the cross terms have zero expectation. A sum of positive-semidefinite quadratic forms is zero exactly when every term is zero. \(\blacksquare\)

If KL admits a second-order expansion in transverse directions and the Fisher quadratic form is bounded below by \(m\|h\|^2\), then, in a sufficiently small neighborhood, \[
D_{\rm KL}(P_\theta\|P_{\theta+h})\ge(m/4)\|h\|^2.
\] The proof bounds the \(o(\|h\|^2)\) remainder by \(m\|h\|^2/4\). This local second-order separation does not exclude global discrete equivalences; singular Fisher information also does not always mean complete nonidentifiability at higher orders.

\section{17.3 Identification of Developmental Laws by Design}\label{identification-of-developmental-laws-by-design}

Consider a prespecified linear parameter family \[
Y_t=\phi_t^\top\theta_*+\xi_t,\quad
G_N=\sum_{t=1}^N\phi_t\phi_t^\top.
\]

\hypertarget{p-inv-04}{}\phantomsection\label{p-inv-04}

\textbf{P-INV-04.} In the specified noiseless model with parameter space all of \(\mathbb R^d\), the design uniquely identifies the parameter if and only if \(G_N\succ0\).

\textbf{Proof.} Writing \(\Phi\) for the design matrix gives \(G_N=\Phi^\top\Phi\). Positive definiteness is equivalent to \(\ker\Phi=\{0\}\). If the matrix is singular, choose nonzero \(v\in\ker\Phi\); every \(\theta_*+cv\) produces the same responses. \(\blacksquare\)

When several environments share the same parameter, the null space of their combined Gram matrix is the intersection of the individual null spaces. If the environments actually have different parameters, pooling them does not create a common law.

\section{17.4 An OLS Bound for Predictable Designs}\label{an-ols-bound-for-predictable-designs}

\hypertarget{p-inv-05}{}\phantomsection\label{p-inv-05}

\textbf{P-INV-05.} Assume \(\phi_t\) is \(\mathcal F_{t-1}\)-measurable, \(\|\phi_t\|_\infty\le B\), and \[
E[e^{\lambda\xi_t}\mid\mathcal F_{t-1}]
\le e^{\lambda^2\sigma^2/2}
\quad\forall\lambda\in\mathbb R.
\] On samples satisfying \(G_N\succeq N\kappa I,\ \kappa>0\), except on an event of probability at most \(\alpha\), \[
\|\widehat\theta-\theta_*\|_2
\le\frac{\sigma B}{\kappa}
\sqrt{\frac{2d\log(2d/\alpha)}N}.
\]

\textbf{Proof.} For each coordinate \(Z_j=\sum_t\phi_{tj}\xi_t\), iteration of the conditional \nolinkurl{moment-generating-function} bound gives \(E e^{\lambda Z_j}\le e^{N\lambda^2\sigma^2B^2/2}\). Chernoff's inequality and a union bound over coordinates yield \(\|Z\|_2\le\sigma B\sqrt{2Nd\log(2d/\alpha)}\). Now use \(\widehat\theta-\theta_*=G_N^{-1}Z\). \(\blacksquare\)

If the Gram lower bound holds with probability \(1-\alpha_G\), the total coverage is at least \(1-\alpha-\alpha_G\). This is a conservative fixed-time bound; standard self-normalized methods provide sharper sequential versions. \href{https://arxiv.org/abs/1102.2670}{Original research}

For a counterexample, let \(\phi_t=\xi_t\) be independent uniform signs and take \(\theta_*=0,Y_t=\xi_t\). The noise is independent and sub-Gaussian and the Gram matrix is positive definite, but OLS is identically 1. Without the design--noise condition, the original vanishing-error bound fails.

If the actual model also has bias \(b_t\), the estimation-error bound must additionally include \[
\left\|G_N^{-1}\sum_t\phi_tb_t\right\|_2.
\] Measurement error, incorrect features, and nonstationary parameters cannot be placed in a ``mean-zero noise'' term without verification.

\section{17.5 Active Intervention Design}\label{active-intervention-design}

Pairs of candidate models form the distinctions to be resolved. An intervention \(i\) covers a pair when its true responses are distinguishable at the required confidence level and within the sample budget. Selecting a set of interventions that covers every target competitor pair gives a costed test-cover problem.

This is better suited to settings without known ground truth than testing only a boundary presumed correct. Greedy coverage is a computational method; in the worst case every intervention may be necessary. The response differences themselves are uncertain and must either be estimated using data separated from validation or handled by an explicit sequential confidence rule.

\section{17.6 Different Quotients for Identifying Objects, Programs, and Conditional Deployment}\label{different-quotients-for-identifying-objects-programs-and-conditional-deployment}

For a given encoding, \nolinkurl{ProgramPurposeInformationBoundary} proves that a complete left-inverse decoder exists if and only if the encoding is injective. Whether a physical seed can produce that encoding is a different question. If the same physical record corresponds to two distinct programs, both cannot be recovered without error at the same time.

To identify purpose from mechanism observations, C6 requires observationally equivalent mechanisms to induce the same policy ordering. Using a candidate after training instead requires its future conditional kernel, given the training record, to agree with the declared actual source. Marginal fit for each fixed candidate cannot remove this conditional common-source requirement. First determine which quotient the present evidence is intended to identify, and then design experiments that distinguish its classes.

\chapter{18. Feasible Futures, Value, and Purpose Bridges}\label{ch-18}

\textbf{Chapter question.} Why does the existence of an object not automatically determine what it should do?

\textbf{Reading route.} Specify implementable policies and a value bridge before comparing numerical objectives, policy orderings, and optimal sets. Reward shaping proves one class of purpose equivalence by telescoping. This equivalence neither uniquely identifies rewards nor derives normative value from an object certificate.

\section{18.1 Feasible paths and implementable policies}\label{feasible-paths-and-implementable-policies}

Given a world, an object branch, resource budgets, and action constraints, define the admissible policy class \(\mathfrak P_{\rm adm}\). Each policy induces a path law \(P^\pi\). One may impose a hard constraint \(P^\pi(\Gamma_{\rm adm})=1\), or a chance constraint \(P^\pi(\Gamma_{\rm adm})\ge p\). These define different problems.

Listing trajectories that satisfy state constraints does not ensure that every trajectory can be implemented by a causal controller. The set of feasible paths, the set of reachable path laws, and the class of implementable policies are recorded separately.

\section{18.2 Types of value}\label{types-of-value}

\textbf{D-BRIDGE.} A value bridge \[
B^{s,\ell}=
(\mathfrak P_{\rm adm},J,\text{horizon, budget, evaluator/level})
\] specifies how behavior at object scale \(s\) is ordered at evaluation level \(\ell\). Here \(J\) may be the expectation of a path functional, a risk functional, a priority ordering of constraints, or a vector of objectives. Bellman recursion applies directly only when the relevant dynamic consistency conditions hold.

For deterministic paths, a strictly increasing transformation preserves ordering. For random paths, however, transforming before taking expectations generally changes risk preferences. Positive affine transformations preserve all expected-utility comparisons; arbitrary increasing transformations do not.

\section{18.3 An object does not uniquely determine an objective}\label{an-object-does-not-uniquely-determine-an-objective}

\textbf{N-TEL-01.} Within the same dynamics and object certificate, if two actions \(a,b\) are feasible, the path objectives \[
J_a=\mathbf1_{\{A_0=a\}},\qquad J_b=\mathbf1_{\{A_0=b\}}
\] yield opposite optimal choices. Thus the existence of an object alone does not determine value.

Assigning the observed policy the unique highest score makes any behavior ``optimal'' after the fact and is unfalsifiable. A substantive interpretation of purpose must constrain the objective class in advance or identify it from an independent mechanism, then predict responses in environments and budgets not used in fitting.

\section{18.4 Value equivalence and reward shaping}\label{value-equivalence-and-reward-shaping}

Two evaluations are operationally purpose-equivalent if they induce the same set of optimal policies on a specified policy class and range of initial states. If both the dynamical realization and the optimal policy set are preserved, the corresponding GOD and well-defined GOA families are also preserved.

\hypertarget{p-tel-01}{}\phantomsection\label{p-tel-01}

\textbf{P-TEL-01: discounted potential equivalence.} Let \(\psi\) be a bounded potential, \(c\) a constant, and \(a>0\). Replace the one-step random reward by \[
r'(x,a_t,x')=a\,r(x,a_t,x')+\beta\psi(x')-\psi(x)+c.
\] Then, for every policy, \[
V'^\pi(x)=aV^\pi(x)-\psi(x)+\frac c{1-\beta},
\] and the optimal policies are unchanged.

\textbf{Proof.} Over the first \(T\) steps, the discounted potential terms telescope to \(-\psi(x_0)+\beta^T\psi(x_T)\). Boundedness makes the final term vanish, leaving a positive affine transformation. \(\blacksquare\)

For a finite horizon, the terminal reward must be adjusted accordingly; the boundary term cannot be omitted. This equivalence supplies a concrete direction of nonuniqueness in inverse reward identification.

\section{18.5 The purpose-realization certificate TVRC}\label{the-purpose-realization-certificate-tvrc}

Fix the bridge and mechanism parameters, then compare the following in held-out environments:

\begin{itemize}
\tightlist
\item
  the value regret of the observed policy;
\item
  the path laws generated by admissible optimal or near-optimal policy families;
\item
  policy and shadow-price responses to independent budget changes;
\item
  model predictions in anomalous and failure conditions.
\end{itemize}

The path residual should be defined by \[
\Delta_{\rm path}=
\inf_{\pi\in\mathfrak P^*_\epsilon}
D(P_{\rm obs},P^\pi),
\] rather than by selecting one optimal policy and requiring every actual behavior to match it. Equally optimal policies may produce different paths and long-term structures. This formulation prevents legitimate diversity among optimal policies from being misclassified as mechanism failure.

\section{18.6 Numerical objectives, orderings, and optimal sets}\label{numerical-objectives-orderings-and-optimal-sets}

An identical numerical objective gives an identical policy ordering, and an identical ordering gives an identical maximizing set; the converses generally fail. \nolinkurl{TeleologicalGauge} retains these different equivalence relations. Even two purposes with the same policy ordering need not belong to a specified algebraic dual-drive gauge orbit.

C6's \nolinkurl{mechanismObservation_identifies_teleologicalClass} identifies an ordering-equivalence class only under mechanism--purpose faithfulness. It supplies neither numerical rewards, Bellman realizability, nor survival feasibility. These are respectively the responsibilities of evaluation calibration, control theory, and viability theory. Section 33.5 gives the complete gauge characterization for the same numerical objective; Section 34.2 gives the boundary at the ordering level.

\chapter{19. Optimal Control and GOD}\label{ch-19}

\textbf{Chapter question.} Given an evaluation and admissible actions, how does an optimal control exist and dominate competing policies?

\textbf{Reading route.} In a finite discounted model, the contraction mapping theorem first yields a Bellman fixed point. Summation along arbitrary actual causal policies then verifies the upper bound and its attainment. GOD is therefore primarily an action correspondence. Continuous control additionally requires a selector, a well-defined process, and Itô conditions; safe optimality uses an explicit comparison class.

\section{19.1 Complete construction for finite discounted control}\label{complete-construction-for-finite-discounted-control}

Let the state space be finite, each state have a finite nonempty action set, the reward satisfy \(|r|\le R\), and the discount satisfy \(0<\beta<1\). Define \[
(Tv)(x)=\max_a[r(x,a)+\beta P_{x,a}v].
\]

\hypertarget{p-ctl-01}{}\phantomsection\label{p-ctl-01}

\textbf{P-CTL-01: Bellman existence and optimality.} The operator \(T\) has a unique fixed point \(V^*\), and value iteration converges from every bounded initial function. There is a stationary deterministic policy that is optimal at every state among all causal policies.

\textbf{Proof.} For any \(v,w\), the difference between two maxima is bounded by the maximum of the corresponding candidate differences. Hence \[
\|Tv-Tw\|_\infty\le\beta\|v-w\|_\infty.
\] The Banach fixed-point theorem on the complete finite-dimensional space gives the unique fixed point and convergence. Choose an action attaining the maximum at each state, defining \(\pi^*\). For any causal policy, \[
V^*(X_t)\ge r(X_t,A_t)+\beta E[V^*(X_{t+1})\mid H_t,A_t].
\] Discount, sum, and take expectations step by step. The return over the first \(T\) steps is at most \(V^*(x)-\beta^TEV^*(X_T)\). Boundedness makes the terminal term vanish. Under \(\pi^*\) every step is an equality, so the policy attains \(V^*\). \(\blacksquare\)

The value-iteration residual yields a stopping bound: \[
\|v-V^*\|_\infty\le\frac{\|Tv-v\|_\infty}{1-\beta}.
\] To prove it, insert \(Tv\) into the contraction estimate and rearrange.

\section{19.2 Finite horizons and survival evaluation}\label{finite-horizons-and-survival-evaluation}

For a terminal reward \(G\), set \[
V_T=G,\qquad
V_t(x)=\max_a[r_t(x,a)+E V_{t+1}(X')].
\] Backward induction gives the optimal value and a time-dependent policy.

If the sole objective is survival inside \(\mathcal V\) for \(H\) steps, use \[
V_0(x)=\mathbf1_{\mathcal V}(x),\quad
V_{h+1}(x)=\mathbf1_{\mathcal V}(x)\max_a\int V_h(y)P(dy\mid x,a).
\] The same induction proves that this is the maximal survival probability.

Survival-optimal actions generally depend on the horizon. After entering route A, suppose survival is \(0.9^H\). Route B has entry success probability \(0.8\) and subsequent per-step survival probability \(0.99\), giving \(H\)-step survival \(0.8(0.99)^{H-1}\). A is better over short horizons, B over long ones. The path model must include the route state; a comparison of memoryless labels is insufficient.

\section{19.3 Compact Feller control}\label{compact-feller-control}

\hypertarget{p-ctl-02}{}\phantomsection\label{p-ctl-02}

\textbf{K-CTL-01 / P-CTL-02.} Suppose the state and action spaces are compact metric spaces, the action set is fixed and nonempty, the reward is continuous, and the transition kernel is weakly continuous in state and action. Then \(T:C(\mathcal X)\to C(\mathcal X)\) is a \(\beta\)-contraction with a unique continuous fixed point and a measurable stationary optimal selector.

\textbf{Proof.} For bounded continuous \(v\), the integral term is continuous. The maximum theorem on the compact action space gives continuity of the maximum and a nonempty compact argmax. A measurable selection theorem gives a selector; contraction and optimality follow as above. \(\blacksquare\)

State-dependent action sets require a closed graph, compact values, and appropriate continuity assumptions. A measurable selector need not be continuous, so an optimal closed loop need not be Feller. Long-term occupancy results require separate verification.

\section{19.4 Continuous-time HJB verification}\label{continuous-time-hjb-verification}

Consider an implementable diffusion \(dX_t=b(X_t,a_t)dt+\sigma(X_t,a_t)dW_t\) with return \(E\int_0^\infty e^{-\rho t}r(X_t,a_t)dt\), where \(\rho>0\).

\hypertarget{p-ctl-03}{}\phantomsection\label{p-ctl-03}

\textbf{P-CTL-03.} Suppose there is a bounded \(C^2\) function \(V\) satisfying \[
\rho V(x)=\sup_a\{r(x,a)+\mathcal L^aV(x)\},
\] and a measurable selector attaining the supremum whose closed-loop SDE is well-defined. If all required localization, integrability, and stochastic-integral conditions hold, then \(V\) is the optimal value and the selector is optimal.

\textbf{Proof.} Apply Itô's formula to \(e^{-\rho t}V(X_t)\). The HJB equation gives \(r+\mathcal L^aV-\rho V\le0\) under every control. Integration and expectation yield \[
V(x)\ge E\left[\int_0^T e^{-\rho t}r\,dt+e^{-\rho T}V(X_T)\right].
\] Boundedness makes the terminal term vanish. The optimal selector makes the inequality an equality throughout. \(\blacksquare\)

If the value function has only viscosity regularity, the relevant standard control theory must be used; a classical gradient that does not exist cannot be inserted into the formula.

\section{19.5 The formal type of GOD}\label{the-formal-type-of-god}

Define \[
Q^*(x,a)=r(x,a)+\beta P_{x,a}V^*,\qquad
\mathcal A^*(x)=\arg\max_aQ^*(x,a).
\] \textbf{D-GOD.} GOD is primarily the optimal-action correspondence \(\mathcal A^*\). When a local drift is available, it may be mapped to \[
\mathcal D^*(x)=\{b(x,a):a\in\mathcal A^*(x)\}.
\] For a jump process one may use an optimal generator, and for a discrete system an optimal transition kernel; no Euclidean direction need be imposed.

The \(\epsilon\)-GOD contains actions satisfying \(Q^*(x,a)\ge V^*(x)-\epsilon\). If such actions are used at every step, total discounted regret is at most \(\epsilon/(1-\beta)\): sum the one-step Bellman defects along the actual trajectory with discounting. The one-step \(\epsilon\) here is not the same numerical tolerance as \(\epsilon\)-optimality of an entire policy.

\section{19.6 Formal HJB hypotheses and survival-constrained optimal control}\label{formal-hjb-hypotheses-and-survival-constrained-optimal-control}

The Lean implementation \nolinkurl{DiffusionHJBVerification.p_ctl_03} consumes a process-specific Model.ItoRun: discounted occupation measures, terminal probability laws, integrability of reward and drift, a finite-horizon Itô expectation identity, and a limit for discounted cumulative reward. It then proves that the HJB residual is nonpositive, the selector residual is zero, the bounded terminal term vanishes, and optimality follows. This certificate does not store the optimal-value conclusion, but it also does not currently construct a complete controlled stochastic-analysis substrate automatically from an arbitrary SDE. P-CTL-03 holds under the process conditions stated here; the source obligations of ItoRun cannot be removed.

Finite models admit another completed strengthening. A viability fixed point first produces action subtypes retaining the source kernel, and Bellman analysis is then performed on the constrained model with the same reward and transition. \nolinkurl{SafeOptimalControlFromViability.viable_restricted_optimal_dominates_all_causal_policies} compares all history-dependent randomized policies respecting this constraint. Its optimality is not claimed over the unrestricted policy class. Section 35.3 gives the complete proof.

\chapter{20. How Object Quotients Preserve Decisions and Value}\label{ch-20}

\textbf{Chapter question.} When does predictive compression preserve decisions, and what value loss follows from approximate compression?

\textbf{Reading route.} An exact quotient preserves actions, rewards, and transitions, so Bellman commutation gives equality of values. An approximate quotient first bounds the one-step residual and then propagates it using contraction of the Bellman operator. Refined occupancy bounds must use the respective visitation laws of the true optimal policy and the actual lifted policy.

This chapter uses finite models for which Chapter 19 already guarantees Bellman contraction, optimal values, and the required optimal policies, or models satisfying the corresponding compactness and measurable-selection conditions of that chapter. Rewards are bounded and the discount is unchanged. General noncompact models must establish these hypotheses separately; quotient identities alone do not establish the existence of an optimal policy.

\section{20.1 Exact control quotients}\label{exact-control-quotients}

\hypertarget{p-quo-01}{}\phantomsection\label{p-quo-01}

\textbf{P-QUO-01.} Suppose a surjection \(f:\mathcal X\to\mathcal M\) has the same action set throughout each fiber and satisfies \[
r(x,a)=\bar r(f(x),a),\qquad
f_\#P_{x,a}=\bar P_{f(x),a},
\] Then \[
V^*=\bar V^*\circ f,
\] and a lifted macro-optimal policy is micro-optimal.

\textbf{Proof.} For every macro function \(\bar v\), \(T(\bar v\circ f)=(\bar T\bar v)\circ f\). The macro fixed point therefore yields a micro Bellman fixed point; uniqueness gives equality of values. Action values agree term by term, so the argmax lifts. \(\blacksquare\)

Accurate prediction of an object's own future does not guarantee reward closure, and Markov closure under one policy does not guarantee closure under all actions. Both are necessary interface conditions when an object representation is used for control.

\section{20.2 Span-sensitive approximate quotient bounds}\label{span-sensitive-approximate-quotient-bounds}

Suppose reward and transition errors are bounded by \(\epsilon_r,\epsilon_p\), respectively, with the same action correspondence as above. Put \[
w=\bar V^*\circ f,\quad
\delta=\epsilon_r+\beta\epsilon_p\operatorname{sp}(\bar V^*),\quad
D=\delta/(1-\beta).
\]

\hypertarget{p-quo-02}{}\phantomsection\label{p-quo-02}

\textbf{P-QUO-02.} \[
\|V^*-w\|_\infty\le D,
\] and a macro-optimal policy \(\bar\pi^*\) has a lift \(\widehat\pi\) satisfying \[
0\le V^*-V^{\widehat\pi}\le2D.
\]

\textbf{Proof.} P-MET-02 bounds the error in each action's Bellman expression by \(\delta\), so the optimal Bellman residual of \(w\) is at most \(\delta\). The contraction residual bound proves the first inequality. Under the fixed lifted policy, the macro-optimal action attains the macro value, so the same residual bound also gives \(\|V^{\widehat\pi}-w\|\le D\). Adding the two bounds proves the result. \(\blacksquare\)

If \(|\bar r|\le R\), the bound may be further relaxed to \[
D\le\frac{\epsilon_r}{1-\beta}
+\frac{2\beta R\epsilon_p}{(1-\beta)^2}.
\] The actual span is preferable. This method was already present in mathematical version v2.1 and is not claimed as original to this edition.

\section{20.3 Action gaps and near-optimality error}\label{action-gaps-and-near-optimality-error}

\hypertarget{p-quo-03}{}\phantomsection\label{p-quo-03}

\textbf{P-QUO-03.} If \(\max_a|Q^*(x,a)-\widehat Q(x,a)|\le\eta\) and the estimated best action exceeds every other action by more than \(2\eta\), then it is also the unique true optimal action.

\textbf{Proof.} The true value of the estimated best action can decrease by at most \(\eta\), and a competitor's value can increase by at most \(\eta\). The difference remains strictly positive. \(\blacksquare\)

For an approximate quotient one may take \(\eta=D\): the action-value error is bounded by the local error \(\delta\) plus the future-value error \(\beta D\), which equals \(D\). An insufficient gap makes this certificate unresolved; it does not imply that the action necessarily changes.

If the macro policy has only a pointwise value error \(\epsilon_{\rm opt}\), its fixed-policy value-span bound must be reconsidered, or algorithmic error must be included in the stepwise action-defect accounting. It cannot still be treated as an exactly macro-optimal policy.

\section{20.4 An exact occupancy-weighted identity}\label{an-exact-occupancy-weighted-identity}

Fix a policy \(\pi\) and any bounded reference function \(w\). Define \[
b_\pi=r^\pi+\beta P^\pi w-w,\quad
d_{\mu}^\pi=(1-\beta)\sum_{t\ge0}\beta^t\mu(P^\pi)^t.
\]

\hypertarget{p-quo-04}{}\phantomsection\label{p-quo-04}

\textbf{P-QUO-04.} \[
V^\pi-w=(I-\beta P^\pi)^{-1}b_\pi,
\qquad
\mu(V^\pi-w)=\frac{E_{d_\mu^\pi}b_\pi}{1-\beta}.
\]

\textbf{Proof.} Subtract from the fixed-policy Bellman equation to obtain the linear equation. The Neumann series converges because \(\|\beta P^\pi\|_\infty<1\). Left multiplication by the initial distribution gives the second identity. \(\blacksquare\)

Thus actual decision error depends on the states visited by the policy, rather than only on worst-case error.

\section{20.5 From local model error to policy regret}\label{from-local-model-error-to-policy-regret}

Use the discounted state--action occupancy distribution \(d_\mu^\pi(x,a)\). Allow local error \[
\delta(x,a)=\epsilon_r(x,a)+
\beta\epsilon_p(x,a)\operatorname{sp}(\bar V^*).
\]

\hypertarget{p-quo-05}{}\phantomsection\label{p-quo-05}

\textbf{P-QUO-05: occupancy-weighted regret of a quotient policy.} \[
J(\pi^*;\mu)-J(\widehat\pi;\mu)
\le\frac{
E_{d_\mu^{\pi^*}}\delta+
E_{d_\mu^{\widehat\pi}}\delta}{1-\beta}.
\]

\textbf{Proof.} Take \(w=\bar V^*\circ f\) and define the one-step reference residual \[
b(x,a)=r(x,a)+\beta P_{x,a}w-w(x).
\] The macro Bellman advantage of every action is nonpositive, hence \(b(x,a)\le\delta(x,a)\). For a lifted macro-optimal action, the macro residual is zero, so \(b(x,\widehat\pi(x))\ge-\delta(x,\widehat\pi(x))\). Apply P-QUO-04 to the true optimal policy and the lifted policy, respectively, and subtract. \(\blacksquare\)

If the two occupancy distributions are dominated by a validation distribution \(\nu\) with density-ratio bounds \(C_*,C_{\widehat\pi}\), the upper bound becomes \[
\frac{(C_*+C_{\widehat\pi})E_\nu\delta}{1-\beta}.
\] This is a concrete connection between the coverage conditions of Chapter 16 and control quotients. When \(C_*\) is unknown, the expression cannot be reported as an actually computed guarantee.

\section{20.6 The order of composition for same-source quotients and constrained control}\label{the-order-of-composition-for-same-source-quotients-and-constrained-control}

The actual unified kernel first constructs a stable partition from the same finite model, registered labels, and reward, then obtains a control quotient by exact kernel factorization. If a survival constraint is added, the admissible action subtype must be proved to retain the original transition; it must not silently replace the model by an easier one.

This also explains why occupancy bounds must use actual policies: the estimated policy and the true optimal policy generally visit different states. The two occupancy terms in P-QUO-05 retain their respective sources. In the finite-risk construction, \(\nu_n\) is generated by the actual prefix and the same coarse-history policy. Error terms in different formulas must be aligned by state space, horizon, and policy type.

\chapter{21. GOA: Existence and Stability of Long-Term Structures}\label{ch-21}

\textbf{Chapter question.} What can the long-term behavior of an optimal policy be, and which stability claims need additional assumptions?

\textbf{Reading route.} Distinguish invariant laws, occupancy limits, recurrent classes, and QSDs before applying contraction or spectral margins. Section 21.7 places these objects in one finite source and compares unconditional paths, conditional survival, and constrained optimality while retaining the actions outside the safe set.

\section{21.1 Typed GOA families}\label{typed-goa-families}

GOD determines currently optimal actions, optimal policies generate paths, and GOA describes their long-term structures. The definitions depend on the model:

\begin{itemize}
\tightlist
\item
  deterministic: \(\omega\)-limit sets of optimal closed-loop trajectories satisfying precompactness conditions;
\item
  stochastic: weak limits of Cesàro averages of time marginals under optimal stationary policies;
\item
  multiple recurrent classes: initial-state-dependent class weights and invariant distributions within the classes;
\item
  metastability: QSDs and invariant occupancy of the Q-process, recorded separately;
\item
  nonstationary environments: an augmented environment state, or separately defined nonautonomous long-term structures.
\end{itemize}

The optimal-policy family may mean policies optimal simultaneously for all initial states or policies optimal for a specified initial distribution. One convention must be fixed. Under the latter, actions at states unreachable from that initial distribution may be arbitrary; global argmax conditions should not be required there.

\section{21.2 Occupancy limits of finite chains}\label{occupancy-limits-of-finite-chains}

\hypertarget{p-goa-01}{}\phantomsection\label{p-goa-01}

\textbf{P-GOA-01.} For a fixed finite stochastic kernel \(P\), the time averages \[
\bar\mu_N=N^{-1}\sum_{t=0}^{N-1}\mu_0P^t
\] have a convergent subsequence, and every limit \(\nu\) satisfies \(\nu P=\nu\).

\textbf{Proof.} The probability simplex is compact, and \[
\bar\mu_NP-\bar\mu_N=(\mu_0P^N-\mu_0)/N\to0.
\] Take the limit. \(\blacksquare\)

Finite chains in fact have a full Cesàro limit, obtained from the transient--recurrent decomposition and the Cesàro theorem in each irreducible class. General Feller processes use P-PER-02. A limit set need not be an attractor with an open basin of attraction; the word ``Attractor'' in GOA does not replace a proof of attraction.

\section{21.3 Dobrushin mixing margin}\label{dobrushin-mixing-margin}

For a finite stochastic kernel, define \[
\alpha(P)=\max_{x,x'}D_{\rm TV}(P_x,P_{x'}).
\] The standard finite-kernel contraction coefficient satisfies \(D_{\rm TV}(\mu P,\nu P)\le\alpha(P)D_{\rm TV}(\mu,\nu)\). This follows by decomposing the signed input difference into positive and negative parts of equal mass and applying convexity to mixtures of rows.

\hypertarget{p-goa-02}{}\phantomsection\label{p-goa-02}

\textbf{P-GOA-02.} If \(\alpha(P)<1\), the invariant probability is unique. If \(\mu P=\mu,\widehat\mu\widehat P=\widehat\mu\) and every row error is at most \(\epsilon\), then \[
D_{\rm TV}(\mu,\widehat\mu)\le\frac{\epsilon}{1-\alpha(P)}.
\]

\textbf{Proof.} Uniqueness follows immediately from contraction. Inserting \(\widehat\mu P\) gives \(d\le\alpha(P)d+\epsilon\); rearrange. \(\blacksquare\)

Since \(\alpha(\widehat P)\le\alpha(P)+2\epsilon\), the estimated kernel also has a unique invariant probability if the right-hand side is less than one. Without a mixing margin, nearby kernels may have entirely different long-term structures: the identity matrix and \(\bigl(\begin{smallmatrix}1-\epsilon&\epsilon\\0&1\end{smallmatrix}\bigr)\) have different numbers of recurrent classes.

\section{21.4 Stability of multiple structures with a preserved recurrent decomposition}\label{stability-of-multiple-structures-with-a-preserved-recurrent-decomposition}

Assume that the transient set and recurrent classes are unchanged by perturbation. Let \(Q\) be the transient transition block and \(R\) the matrix of direct transition probabilities into the classes, with \[
N=(I-Q)^{-1},\quad H=NR.
\] Use the maximum absolute row-sum matrix norm. Suppose \(\|\widehat Q-Q\|_\infty\le\epsilon_Q\), \(\|\widehat R-R\|_\infty\le\epsilon_R\), and \(\|N\|_\infty\epsilon_Q<1\).

\hypertarget{p-goa-03}{}\phantomsection\label{p-goa-03}

\textbf{P-GOA-03.} \[
\|\widehat H-H\|_\infty\le
B_H=
\frac{\|N\|_\infty^2\epsilon_Q+
\|N\|_\infty\epsilon_R}
{1-\|N\|_\infty\epsilon_Q}.
\] If the invariant distributions within the classes have TV errors \(\epsilon_j^{\rm stat}\), the Cesàro limits induced by the same initial law satisfy \[
D_{\rm TV}(\nu,\widehat\nu)
\le\tfrac12 B_H+\sum_j\widehat w_j\epsilon_j^{\rm stat}.
\]

\textbf{Proof.} The resolvent identity and Neumann series give \(\|\widehat N\|\le\|N\|/(1-\|N\|\epsilon_Q)\) and \(\|\widehat N-N\|\le\|N\|^2\epsilon_Q/(1-\|N\|\epsilon_Q)\). Use \(\|R\|_\infty\le1\) to bound \((\widehat N-N)R+\widehat N(\widehat R-R)\), proving the first inequality. The total \(L^1\) difference of initial-law-dependent absorption weights is at most \(B_H\). The triangle inequality for mixtures gives the second inequality. \(\blacksquare\)

The theorem does not apply when recurrent support changes.

\section{21.5 Spectral stability of symmetric killed kernels}\label{spectral-stability-of-symmetric-killed-kernels}

Let \(Q\) be symmetric, nonnegative, irreducible, and substochastic, with positive unit principal eigenvector \(u\) and spectral gap \(g=\rho-\lambda_2>0\). Let a symmetric perturbation \(E\) satisfy \(\|E\|_2\le\eta<g/2\), with \(\widehat Q=Q+E\) retaining the same properties.

\hypertarget{p-goa-04}{}\phantomsection\label{p-goa-04}

\textbf{P-GOA-04.} \[
|\widehat\rho-\rho|\le\eta,\qquad
D_{\rm TV}(u^2,\widehat u^2)\le\frac{2\sqrt2\,\eta}{g}.
\]

\textbf{Proof.} The Rayleigh extremal characterization bounds the eigenvalue error. Projecting \(\widehat u\) onto \(u^\perp\), the eigenvalue equation bounds the orthogonal component by \(\eta/(g-\eta)\le2\eta/g\). The angle between the two positive principal vectors is at most a right angle, so \(\|u-\widehat u\|_2\le\sqrt2\sin\theta\le2\sqrt2\eta/g\). Finally, \[
\tfrac12\sum_i|u_i^2-\widehat u_i^2|
\le\tfrac12\|u-\widehat u\|_2\|u+\widehat u\|_2
\le\|u-\widehat u\|_2.
\] The Q-process occupancy of a symmetric killed kernel is \(u_i^2\), which proves the claim. \(\blacksquare\)

This is a special spectral stability bound, not a universal bound for arbitrary non-self-adjoint QSDs.

\section{21.6 A finite risk budget does not unconditionally imply long-term attraction}\label{a-finite-risk-budget-does-not-unconditionally-imply-long-term-attraction}

Even with small one-step errors, \(\sum_{t<T}\epsilon_t\) may increase with \(T\) until the bound becomes the trivial probability bound one. Without contraction, spectral, or recurrence conditions, taking \(T\to\infty\) in a finite-path error bound does not establish GOA stability.

M-OI-01 unifies the argument ``a limit exists and observable residuals vanish, hence the limit is invariant.'' M-PE-01 unifies normalized consequences of given positive eigenstructure. Neither automatically supplies the mixing margin of P-GOA-02, the common recurrent decomposition of P-GOA-03, or the spectral gap of P-GOA-04. Chapter 32 therefore retains the respective adaptations within the GOA family.

\section{21.7 Stationarity, conditional survival, and value from one controlled source}\label{stationarity-conditional-survival-and-value-from-one-controlled-source}

This section instantiates the general mechanisms of Chapters 10, 19, and 21 in one actual finite controlled source and connects them to the object partition, support, and recovery of Section 35.7. The source kernel, exact control quotient, and viability fixed point already have Lean implementations. The parametric stationary laws, QSD, complete state--action conditional paths, common infinite experiment, and original unrestricted and constrained GOA are now connected by actual proofs in \nolinkurl{ConflictParametricSource}, Control, Stationary, QSD, Path, Infinite, and \nolinkurl{CommonConsumer.} The 8/5 model of Section 35.4 and the 9/64 registered experiment of Section 38.6 retain their own independent sources.

\subsection{21.7.1 The single controlled source and complete policy}\label{the-single-controlled-source-and-complete-policy}

Use the actual definitions \[
\begin{gathered}
X=\mathrm{Bool}\times\mathrm{Bool},\qquad A=\mathrm{Bool},\\
M=\mathrm{conflictRewardControl.micro},\\
P(x,a)=M.\mathrm{transitionPMF}(x,a).
\end{gathered}
\]

Write false and true as 0 and 1, respectively, and states as \(x=(b,u)\). The observation is \(q(b,u)=b\), the safe set is \(K=\{(1,0),(1,1)\}\), and the initial law is explicitly \(\nu_0=\delta_{(1,0)}\). Transition, reward, and discount all come from the same M: \[
P((b,u),a)\{(b',u')\}=\tfrac12{\bf1}_{b'=a},\qquad
r((b,u),a)=1-a,\qquad \gamma=\tfrac12.                 \tag{21.7-1}
\]

The PMF mass in (21.7-1) is obtained from M's real-valued transition through the existing \nolinkurl{transitionPMF} conversion theorem. \nolinkurl{ConflictParametricSource.transition_real} proves this source equality; an independently specified equal-valued table does not replace the proof. The internal bit is resampled by the transition and may retain its previous value.

For each \(\theta\in[0,1]\), define one complete causal policy \(\pi_\theta\): at every time and every visible history its action PMF is \[
\pi_\theta(0)=\theta,\qquad \pi_\theta(1)=1-\theta.     \tag{21.7-2}
\]

This definition includes zero-probability histories. The actual finite-path law is constructed by the causal product/bind recursion of the production path API. Thus (21.7-2) specifies the action distribution \textbf{conditional on every history}, rather than merely an unconditional action frequency; independence is not inferred from a marginal frequency.

Put \[
p_\theta(0)=\theta,\quad p_\theta(1)=1-\theta,\qquad
\mu_\theta(b,u)=p_\theta(b)/2.
\]

\subsection{21.7.2 The unconditional closed loop: one-step convergence to a unique invariant law for every parameter}\label{the-unconditional-closed-loop-one-step-convergence-to-a-unique-invariant-law-for-every-parameter}

\hypertarget{m-21-1}{}\phantomsection\label{m-21-1}

\textbf{M-21.1.} For every \(\theta\in[0,1]\), the actual closed-loop transition induced by (21.7-1) and (21.7-2) is \[
L_\theta(x,y)
=\sum_{a\in A}\pi_\theta(a)P(x,a)\{y\}
=\mu_\theta(y).                                      \tag{21.7-3}
\] Every probability row \(\eta\) satisfies \(\eta L_\theta=\mu_\theta\). Consequently, \(\mu_\theta\) is the unique invariant probability row, every initial distribution equals it after one step, and the Dobrushin coefficient is zero.

\textbf{Proof.} For fixed \(y=(b',u')\), only the term \(a=b'\) can be nonzero, giving \(p_\theta(b')/2\). Summing over \(x\) gives \[
(\eta L_\theta)(y)=\sum_x\eta(x)\mu_\theta(y)=\mu_\theta(y).
\] Taking \(\eta=\mu_\theta\) proves invariance. If \(\eta L_\theta=\eta\), the same equality gives \(\eta=\mu_\theta\). All rows are equal, so all pairwise row TV distances vanish. \(\blacksquare\)

\hypertarget{m-21-2}{}\phantomsection\label{m-21-2}

\textbf{M-21.2.} Let the actual marginal at step t be \(\nu_t=\nu_0L_\theta^t\). For every positive integer \(N\), \[
\bar\nu_N=\frac1N\sum_{t=0}^{N-1}\nu_t
=\frac1N\nu_0+\frac{N-1}{N}\mu_\theta.                \tag{21.7-4}
\] Therefore \[
D_{\rm TV}(\bar\nu_N,\mu_\theta)
=\frac1N D_{\rm TV}(\nu_0,\mu_\theta)
=\frac{1+\theta}{2N}\longrightarrow0,                \tag{21.7-5}
\] \[
\bar\nu_NL_\theta-\bar\nu_N
=\frac{\mu_\theta-\nu_0}{N}.                         \tag{21.7-6}
\]

\textbf{Proof.} M-21.1 gives, for \(t\ge1\), the equality \(\nu_t=\mu_\theta\). Substitution into the average gives (21.7-4). TV is positively homogeneous in the same signed difference, and \(D_{\rm TV}(\delta_x,\mu)=1-\mu(x)\), yielding (21.7-5). Applying one transition gives (21.7-6). \(\blacksquare\)

Equation (21.7-4) averages probability laws in the Cesàro sense; it does not yet assert a theorem about empirical frequencies along individual sample trajectories. The production cesaroRow uses parameter \(n\) for \(N=n+1\), so the index shift must be handled in adaptation.

In this example, explicit calculation supplies M-OI's vanishing-residual and limit-existence conditions; subsequence extraction from the finite simplex can be reused. General P-GOA-01 itself still asserts only subsequence existence and invariance of each subsequential limit. The stronger full-sequence convergence and uniqueness here come from (21.7-3), not from a stronger reading of the general theorem.

\subsection{21.7.3 The killed kernel comes from the same actual process}\label{the-killed-kernel-comes-from-the-same-actual-process}

Define the first exit from the safe set by \[
\tau=\inf\{t\ge1:X_t\notin K\},
\] and put \(S_T=\{X_1,\ldots,X_T\in K\}\), with \(S_0\) the entire sample space.

Identify K with Bool through \(u\mapsto(1,u)\). Restrict the same closed loop \(L_\theta\) to K; do not specify an independent matrix: \[
Q_\theta(u,u')=L_\theta((1,u),(1,u'))=\frac{1-\theta}{2}. \tag{21.7-7}
\]

Its row sum is \(1-\theta\), and the missing mass \(\theta\) is exactly the original process's probability of exiting K in that step. Deleting the states outside K and renormalizing produces a different kernel, not the original killed kernel (21.7-7).

Set \[
H(u,u')=\tfrac12,\quad \rho=1-\theta,\quad
q_*(u)=\tfrac12,\quad h_*(u)=1.
\] Then \(Q_\theta=\rho H\) and \(H^2=H\).

\hypertarget{m-21-3}{}\phantomsection\label{m-21-3}

\textbf{M-21.3.} For \(0\le\theta<1\), these data satisfy every requirement of the existing BiEigenData and KilledQSDGOA: \[
Q_\theta\ge0,\quad Q_\theta{\bf1}=\rho{\bf1},\quad
0<\rho\le1,\quad q_*Q_\theta=\rho q_*,
\quad \sum_uq_*(u)=\sum_uq_*(u)h_*(u)=1.              \tag{21.7-8}
\] The Doob transform and invariant weight are \[
Q_\theta^{h_*}(u,u')
=\frac{Q_\theta(u,u')h_*(u')}{\rho h_*(u)}
=\tfrac12=H(u,u'),\qquad q_*h_*=q_*.                 \tag{21.7-9}
\]

\textbf{Proof.} Nonnegativity follows from \(\theta\in[0,1)\). Summing two terms makes \(Q_\theta{\bf1}\) equal to \(1-\theta\). For every \(u'\), \[
(q_*Q_\theta)(u')=2\cdot\tfrac12\cdot\tfrac{\rho}{2}
=\rho/2=\rho q_*(u').
\] Both normalizations follow by direct summation. Since \(\rho>0\), the division in (21.7-9) is legitimate. \(\blacksquare\)

There is no need to assume an unknown positive eigenvector via Perron--Frobenius existence: this finite instance computes the positive eigenstructure explicitly and then reuses the generator's consequences.

\hypertarget{m-21-4}{}\phantomsection\label{m-21-4}

\textbf{M-21.4.} For every initial probability row \(\eta\) within K and every \(T\ge1\), \[
\eta Q_\theta^T=\rho^Tq_*,\qquad
\sum_{u'}(\eta Q_\theta^T)(u')=\rho^T.                \tag{21.7-10}
\] When \(\theta<1\), the terminal internal-bit distribution conditional on survival to T is exactly \(q_*\).

\textbf{Proof.} From \(H^2=H\) we have \(Q_\theta^T=\rho^TH\), and every probability row satisfies \(\eta H=q_*\). Summing gives the mass in (21.7-10); normalize when \(\rho^T>0\). \(\blacksquare\)

The QSD initial law in \nolinkurl{survival_mass_pow} cannot directly replace \(\nu_0=\delta_{(1,0)}\): that general theorem starts from \(q_*\). Equation (21.7-10) supplies the necessary adaptation to arbitrary initial laws within K. At \(T=0\), the actual conditional terminal internal-bit law remains \(\delta_0\), not \(q_*\). The all-steps QSD assertion starting from \(q_*\) includes time zero and has different quantifiers.

\subsection{21.7.4 Actual finite risk and conditional sample spaces}\label{actual-finite-risk-and-conditional-sample-spaces}

\hypertarget{m-21-5}{}\phantomsection\label{m-21-5}

\textbf{M-21.5.} For \(\theta\in[0,1]\) and \(T\ge0\), the finite path of the same \((P,\pi_\theta,\nu_0)\) satisfies \[
\Pr_{\pi_\theta}(S_T)=(1-\theta)^T,\qquad
\Pr_{\pi_\theta}(\tau\le T)=1-(1-\theta)^T.           \tag{21.7-11}
\] For \(T\ge1\), however, the terminal failure probability is \[
\Pr_{\pi_\theta}(X_T\notin K)=\theta.                \tag{21.7-12}
\]

\textbf{Proof.} After any surviving history, the next step remains in K exactly when action 1 is selected. Its conditional probability is always \(1-\theta\) by (21.7-2). Thus \(s_{T+1}=(1-\theta)s_T\) with \(s_0=1\), and induction proves (21.7-11). This agrees with the actual killed-kernel mass (21.7-10). Equation (21.7-12) follows from the actual marginal \(\nu_T=\mu_\theta\) in M-21.1. \(\blacksquare\)

For example, at \(\theta=1/8,T=2\), \[
\Pr(\tau\le2)=15/64,\qquad \Pr(X_2\notin K)=1/8.
\] The difference \(7/64\) is exactly the probability of exiting K and then returning. This return within the two operating steps is recovery under the original physical transition, not execution of a separate program-recovery protocol.

If the certified-row predicate is \(C(y,a)\iff a=1\), actual uncertified-row occupancy at each decision time is \(\theta\). It vanishes only for \(\theta=0\); the same \(\pi_\theta\) cannot borrow the zero occupancy of \(\pi_0\). If instead every row is certified and the exact reference kernel is used, occupancy can be zero while the reference model's actual physical failure risk remains (21.7-11). These are different certificate designs, and their favorable terms cannot be selectively combined.

\subsubsection{21.7.4.1 An exact identity for conditioning the whole path}\label{an-exact-identity-for-conditioning-the-whole-path}

For \(\theta<1\) and any fixed T, the full finite state--action path law conditional on \(S_T\) equals the path law of \(\pi_0\) over the same horizon.

Indeed, every surviving path has action 1 throughout. Its \(\pi_\theta\) mass is therefore the mass of the same path under \(\pi_0\), multiplied by \((1-\theta)^T\). Nonsurviving paths have zero conditional mass. Divide by (21.7-11). This explains how conditioning produces the fair internal kernel (21.7-9), while also showing that the two unconditional deployment laws differ. This simple constant-hazard instance does not establish a general identity between a Q-process and finite-window conditioning.

\subsubsection{21.7.4.2 Interpretation on an infinite path space}\label{interpretation-on-an-infinite-path-space}

The finite identity (21.7-11) already implies \[
0<\theta\le1\quad\Longrightarrow\quad
\lim_{T\to\infty}\Pr_{\pi_\theta}(\tau\le T)=1.         \tag{21.7-13}
\]

To interpret (21.7-13) as almost-sure eventual failure in the same infinite experiment, \nolinkurl{ConflictParametricInfinite} first computes the PMF of the complete one-step record \((A_n,X_{n+1})\) from the original source and proves that, in this example, it is independent of the current state. It then constructs the countable product of that record law. Fixing \(X_0=(1,0)\) and recursively forming complete prefixes, \nolinkurl{finite_record_cylinder} identifies every finite cylinder law term by term with the original production causal-path PMF; \nolinkurl{record_prefix_consistent} proves prefix consistency. Thus the source and policy are indeed the same. Continuity from above for the events \(S_T\) gives \[
\Pr(\tau=\infty)=\lim_T(1-\theta)^T=0
\] for \(0<\theta\le1\). At \(\theta=0\), this probability equals one.

\nolinkurl{persistent_survival_null}, \nolinkurl{actual_eventual_failure_ae}, and \nolinkurl{zero_theta_survives_ae} prove these assertions in the same infinite experiment. For positive parameters the conclusion is almost everywhere, not pointwise: null infinite-safe sample paths are not deleted to manufacture pointwise failure. The general assertion ``for each finite T there exists a feasible policy'' still cannot be changed to ``there exists one policy that survives forever.''

\subsection{21.7.5 Optimality, viability, and invariant laws on the full space}\label{optimality-viability-and-invariant-laws-on-the-full-space}

\hypertarget{m-21-6}{}\phantomsection\label{m-21-6}

\textbf{M-21.6.} For every \(\theta\in[0,1]\), \[
J_M(\pi_\theta;\nu_0)
=\mathbb E\!\left[\sum_{n=0}^\infty2^{-n}(1-A_n)\right]
=2\theta.                                          \tag{21.7-14}
\] The original model M has infinite discounted optimal value 2 at every state, and its unique Bellman-greedy action is 0. Its closed loop is \(L_1\), whose unique invariant law \(\mu_1\) is supported entirely outside K.

\textbf{Proof.} One-step rewards are nonnegative and bounded. The expected finite sum is \(\theta\sum_{n<T}2^{-n}\), and monotone convergence gives (21.7-14). Every causal policy has reward at most 1 at each step, hence total value at most 2; always choosing 0 attains it. For the constant function 2, the two Bellman action returns are 2 and 1, so action 0 is uniquely greedy. \(\blacksquare\)

Every \(L_\theta\) therefore has a unique invariant law reached in one step, but this property alone does not identify the policy optimizing M's specified reward. Nor does optimality ensure preservation of the registered safety condition.

\subsubsection{21.7.5.1 The actual action domain outside K must be retained}\label{the-actual-action-domain-outside-k-must-be-retained}

The current production \nolinkurl{viableRestrictedModel} uses the restriction \[
A_K(x)=\{a:x\in K\Rightarrow \operatorname{supp}P(x,a)\subseteq K\}.
\] By (21.7-1), \[
A_K(x)=\{1\}\ \ (x\in K),\qquad A_K(x)=\{0,1\}\ \ (x\notin K). \tag{21.7-15}
\]

Let R be this actual constrained model, retaining M's transition, reward, and discount. \textbf{Equation (21.7-15) must not be rewritten as allowing only action 1 at every x.}

\hypertarget{m-21-7}{}\phantomsection\label{m-21-7}

\textbf{M-21.7.} The optimal value and unique greedy policy of R are \[
V_R(b,u)=2(1-b),\qquad g_R(b,u)=b.                   \tag{21.7-16}
\] Its full closed-loop matrix is \[
D((b,u),(b',u'))=\tfrac12{\bf1}_{b'=b},              \tag{21.7-17}
\] and its invariant probabilities are exactly \[
\{\lambda\mu_0+(1-\lambda)\mu_1:0\le\lambda\le1\}.   \tag{21.7-18}
\] From \(\nu_0=\delta_{(1,0)}\), its marginals after one step and its long-term limit are \(\mu_0\), and its entire actual path law agrees with that of \(\pi_0\).

\textbf{Proof.} Inside K, the only legal action is 1; the process remains in K and all rewards are zero. Outside K, always choosing 0 retains the process outside K and gives reward 1 at each step, attaining the global upper bound 2. Direct comparison with the candidate value gives Bellman returns 2 and 0 for actions 0 and 1 outside K; inside K only action 1 is available. This proves (21.7-16) and (21.7-17).

For any probability row \(\eta\), write its visible-bit mass as \(\eta_b=\sum_u\eta(b,u)\). One application of D gives \[
(\eta D)(b,u)=\eta_b/2.
\] Invariance therefore requires precisely a fair internal-bit distribution within each visible fiber, with arbitrary total weights on the two fibers. This gives (21.7-18). Starting inside K, all actual successors remain in K and the greedy policy always selects 1, so causal recursion agrees step by step with \(\pi_0\). \(\blacksquare\)

Thus R's full \nolinkurl{GreedyInvariantGOASet} is \textbf{not a singleton}. Its sole member supported in K is \(\mu_0\), and the CesàroLimitGOA from the specified initial state is also the singleton \(\mu_0\). The full matrix D has Dobrushin coefficient 1 because rows in the two visible fibers have disjoint supports. A theorem requiring strict contraction on the full space cannot establish global uniqueness for D.

Restricting the state type separately to K gives kernel H, which does have the unique invariant law \(q_*\). This is a restriction with an explicit state-inclusion map, not a silent replacement of R's original state type. The full D also differs from \(L_0\): they agree only inside K, while their actions and successors differ outside K.

\pagebreak

\subsection{21.7.6 Endpoints and invalid inferences}\label{endpoints-and-invalid-inferences}

\begin{longtable}[]{@{}
  >{\raggedright\arraybackslash}p{(\columnwidth - 6\tabcolsep) * \real{0.2500}}
  >{\raggedright\arraybackslash}p{(\columnwidth - 6\tabcolsep) * \real{0.2500}}
  >{\raggedright\arraybackslash}p{(\columnwidth - 6\tabcolsep) * \real{0.2500}}
  >{\raggedright\arraybackslash}p{(\columnwidth - 6\tabcolsep) * \real{0.2500}}@{}}
\toprule\noalign{}
\begin{minipage}[b]{\linewidth}\raggedright
Case
\end{minipage} & \begin{minipage}[b]{\linewidth}\raggedright
\(\theta=0\)
\end{minipage} & \begin{minipage}[b]{\linewidth}\raggedright
\(0<\theta<1\)
\end{minipage} & \begin{minipage}[b]{\linewidth}\raggedright
\(\theta=1\)
\end{minipage} \\
\midrule\noalign{}
\endhead
\bottomrule\noalign{}
\endlastfoot
Unique invariant law of \(L_\theta\) & \(\mu_0\), inside K & \(\mu_\theta\), with mass in both fibers & \(\mu_1\), outside K \\
Probability of any failure by positive T & 0 & \(1-(1-\theta)^T\) & 1 \\
Terminal failure probability at positive T & 0 & \(\theta\) & 1 \\
Killed row sum \(\rho\) & 1 & Strictly between 0 and 1 & 0 \\
BiEigenData / positive-survival QSD & Applicable; no actual killing mass & Applicable & Inapplicable: scale\_pos fails \\
Doob kernel & H & H & Zero normalization denominator; not definable as the Doob kernel of this theorem \\
\(J_M(\pi_\theta)\) & 0 & \(2\theta\) & 2 \\
\end{longtable}

At T=0, survival probability is 1 and probability of previous failure is 0 for every parameter. Natural-number exponentiation uses \(0^0=1\). The actual initial internal-bit law is \(\delta_0\), so the assertion that the conditional terminal internal bit has law \(q_*\) strictly retains T≥1.

This section refutes the following unconditional inferences: a unique stationary law implies zero failure; a QSD supported in K implies safety of unconditional deployment; a unique limit from a safe initial state implies a unique GOA on the full state space; reward optimality implies viability; the same state type or invariant law implies the same actual policy or path source. It does not refute existing theorems with actual optimality, invariant-domain, initial-state, contraction, or conditioning hypotheses.

\subsection{21.7.7 Lean correspondence and indispensable source bridges}\label{lean-correspondence-and-indispensable-source-bridges}

\nolinkurl{InvariantLawGOA}, \nolinkurl{CesaroLimitGOA}, and KilledQSDGOA describe long-term objects of a given kernel and contain no optimality field. \nolinkurl{GreedyInvariantGOASet} explicitly uses the Bellman-greedy closed loop of the supplied control model. Each name in this section is used with that distinction: the existence of a stable distribution does not confer reward optimality or normative value.

The existing \href{../formalization/ueot-core/UEOT/V3/FiniteCesaroInvariant.lean}{\nolinkurl{FiniteCesaroInvariant}} and \href{../formalization/ueot-core/UEOT/V3/Compression/OccupationLimitInvariance.lean}{\nolinkurl{OccupationLimitInvariance}} can consume the stochastic matrices, initial simplex points, and Cesàro limits of this section. BiEigenData, \nolinkurl{qsdDoob_bundle}, and the killed wrapper in \href{../formalization/ueot-core/UEOT/V3/Compression/PositiveEigenstructure.lean}{\nolinkurl{PositiveEigenstructure}} can consume the explicitly computed positive eigenstructure. The general interfaces were already proved. Aligning Lθ with actual policy paths, Qθ with the actual killed source, and R's greedy policy with its actual state-dependent action domain was the formal obligation originally registered here. It was subsequently completed by the \nolinkurl{ConflictParametric} series and their common consumers; Section 39.71 gives the exact scope. The source requirements are retained here, but are no longer listed as currently unresolved.

\href{../formalization/ueot-core/UEOT/V3/Compression/TheoryCompletion/SameObjectViableGOA.lean}{\nolinkurl{SameObjectViableGOA}} proves that when an actual same-object greedy policy preserves a nonempty persistence domain, an invariant law supported in that domain exists. This interface additionally needs a complete Objecthood parent, a control specification, and their source equalities. This section does not construct those fields merely from the four-state type. Its conclusion does not require a unique invariant law on the entire state space. Stronger global uniqueness and mixing endpoints separately require actual contraction or isolation conditions and cannot be applied to D, which has two closed classes.

The common-source instance thus connects M-QD's exact observation quotient, finite-path propagation, M-OI's unconditional occupancy, and M-PE's conditional survival structure. It changes neither the four generating mechanisms nor the classification of the 106 P-IDs, and does not complete the separate conditions for physical program sources, parent formation, lineage, or empirical dual-drive mechanisms.

\chapter{22. Path Information Cost, Reachable Action, and Variational Realization}\label{ch-22}

\textbf{Chapter question.} How much path information cost is necessary to change the probability of a future event?

\textbf{Reading route.} Projection to an event indicator gives a necessary lower bound; event tilting shows when it is attained among all absolutely continuous path laws. An actual control problem optimizes over its implementable subclass of path laws. Projection of a diffusion may also discard noise information, so control energy need not equal state-path KL divergence.

\section{22.1 An information lower bound for survival events}\label{an-information-lower-bound-for-survival-events}

Let \(P^0\) be a baseline path law, \(Q\ll P^0\) a reachable path law, and \(A\) an event. Write \(q=P^0(A),p=Q(A)\).

\hypertarget{p-kl-01}{}\phantomsection\label{p-kl-01}

\textbf{P-KL-01.} \[
D_{\rm KL}(Q\|P^0)\ge
d_{\rm Bern}(p\|q).
\] \textbf{Proof.} Apply KL data processing to the measurable path map \(\gamma\mapsto\mathbf1_A\). \(\blacksquare\)

If \(q=0\), absolute continuity forces \(p=0\); a positive probability requires a change of reachable support or infinite KL in this comparison. If \(q=1\), absolute continuity likewise forces \(p=1\).

\section{22.2 Exact event I-projection}\label{exact-event-i-projection}

\hypertarget{p-kl-02}{}\phantomsection\label{p-kl-02}

\textbf{P-KL-02.} Given \(0<q<1\) and target \(0\le p\le1\), among all absolutely continuous path laws, \[
\inf_{Q:Q(A)\ge p}D_{\rm KL}(Q\|P^0)
=
\begin{cases}
0,&p\le q,\\
d_{\rm Bern}(p\|q),&p>q.
\end{cases}
\] For \(p>q\), an attaining law is \[
\frac{dQ^*}{dP^0}
=\frac pq\mathbf1_A+
\frac{1-p}{1-q}\mathbf1_{A^c}.
\]

\textbf{Proof.} When \(p\le q\), the baseline itself is feasible. When \(p>q\), Bernoulli KL is strictly increasing in the event probability \(r\ge p\), since its derivative is \(\log[r(1-q)/(q(1-r))]>0\). The displayed density changes only the masses of the two events and preserves their conditional path laws, attaining the lower bound. At \(p=1\), interpret the boundary using \(0\log0=0\). \(\blacksquare\)

This is the minimal information change among all path laws. A constrained feedback controller may not implement \(Q^*\), so the actual minimal control cost can only be greater. Existence of an exponential tilt does not itself show that a natural system performs this optimization.

\section{22.3 KL divergence of discrete Markov paths}\label{kl-divergence-of-discrete-markov-paths}

With the same initial law and appropriate rowwise absolute continuity of kernels \(P_t,Q_t\), \[
D_{\rm KL}(P_{0:T}\|Q_{0:T})
=\sum_{t=0}^{T-1}E_P
D_{\rm KL}(P_t(\cdot\mid H_t)\|Q_t(\cdot\mid H_t)).
\]

\hypertarget{p-kl-03}{}\phantomsection\label{p-kl-03}

\textbf{P-KL-03, proof.} The path density ratio is the product of successive conditional density ratios. Taking logarithms gives a sum, followed by successive conditional expectations. Distinct initial laws contribute an additional initial KL term. \(\blacksquare\)

This chain rule also applies to history-dependent kernels and provides an executable error account for general processes.

\section{22.4 An explicit CTMC formula}\label{an-explicit-ctmc-formula}

Consider finite nonexplosive CTMCs with the same initial law. Suppose the positive support of the controlled rates \(l^u_{ij}\) is contained in that of the baseline rates \(l^0_{ij}\). Write \(\lambda_i^u=\sum_{j\ne i}l^u_{ij}\). Then

\hypertarget{p-kl-04}{}\phantomsection\label{p-kl-04}

\textbf{P-KL-04.} \[
D_{\rm KL}(P^u_{[0,T]}\|P^0_{[0,T]})
=
E_u\int_0^T
\sum_{j\ne X_t}
\left[
l^u_{X_tj}\log\frac{l^u_{X_tj}}{l^0_{X_tj}}
-l^u_{X_tj}+l^0_{X_tj}
\right]dt.
\]

\textbf{Proof.} The path log-likelihood ratio is the sum of log rate ratios at jumps, minus the time integral of the difference of exit rates. Take expectation under the controlled law and use the compensator of the jump counts to replace the jump sum by its rate integral. \(\blacksquare\)

Zero controlled-rate terms use \(0\log0=0\). Time- or history-dependent rates require the corresponding predictability and compensator integrability conditions.

\section{22.5 Diffusions: noise space and state-path space must be distinguished}\label{diffusions-noise-space-and-state-path-space-must-be-distinguished}

On the reference probability space, let \[
dX_t=b^0(t,X_t)dt+\sigma(t,X_t)dW_t^0.
\] Let \(u_t\) be progressively measurable, and suppose the exponential \[
Z_T=\exp\left(\int_0^T u_t\,dW_t^0-
\tfrac12\int_0^T\|u_t\|^2dt\right)
\] is a true martingale with \(EZ_T=1\), for example under an applicable Novikov condition. After the change of measure \(d\mathbf Q=Z_Td\mathbf P^0\), the process \(W^u=W^0-\int u\,dt\) is Brownian motion and the drift of \(X\) becomes \(b^0+\sigma u\).

\hypertarget{p-kl-05}{}\phantomsection\label{p-kl-05}

\textbf{P-KL-05: the observability-corrected control-energy relation.} If \(E_{\mathbf Q}\int_0^T\|u_t\|^2dt<\infty\), then on the full probability space, \[
D_{\rm KL}(\mathbf Q\|\mathbf P^0)
=\tfrac12E_{\mathbf Q}\int_0^T\|u_t\|^2dt.
\] For the path laws observing only \(X\), in general one has only \[
D_{\rm KL}(Q_X\|P_X^0)
\le\tfrac12E_{\mathbf Q}\int_0^T\|u_t\|^2dt.
\]

\textbf{Proof.} Substitute \(dW^0=dW^u+u\,dt\) into \(\log Z_T\). Under \(\mathbf Q\), the stochastic integral has mean zero by square integrability, proving the first identity. Projection onto the path of \(X\) and data processing give the second inequality. \(\blacksquare\)

If \(Z_T\) is measurable with respect to the path of \(X\), no information is lost under projection and the second inequality is an equality. A common sufficient setting has a square invertible diffusion matrix, driving coefficients and controls determined by observable state history, and well-defined SDEs and changes of measure, so that Brownian increments and the density ratio can be reconstructed from the state path.

\textbf{N-KL-01.} Take two-dimensional noise, \(\sigma=(1,0)\), and \(u=(0,1)\). Changing the second noise coordinate leaves the state-path law of \(X=W_1\) unchanged, hence state KL is zero while control energy is \(T/2\). This rules out unconditional equality of state-path KL and total control energy for arbitrary degenerate diffusions.

For precise Girsanov and entropy conditions, see the original mathematical literature. Here the form used explicitly retains true-martingale and observability conditions. \href{https://leonard.perso.math.cnrs.fr/papers/2012-Leonard-Girsanov\%20theory\%20under\%20a\%20finite\%20entropy\%20condition.pdf}{Léonard's paper}

Combining with P-KL-01, even the upper-bound form gives \[
\tfrac12E_{\mathbf Q}\int\|u\|^2dt
\ge D_{\rm KL}(Q_X\|P_X^0)
\ge d_{\rm Bern}(Q_X(A)\|P_X^0(A)).
\] The necessary control-energy lower bound for persistence events is therefore retained. It is not automatically a thermodynamic work bound: energy units, temperature, and the relation to dissipation require an independent physical realization.

\section{22.6 Information-cost conclusions and continuous-process adaptation}\label{information-cost-conclusions-and-continuous-process-adaptation}

Probability constructions for finite paths and finite CTMCs, relative-entropy chains, and actual common-source kernels can be expanded within their respective modules. The continuous-diffusion branch explicitly retains \nolinkurl{TerminalGirsanovData}; recovery moments retain \nolinkurl{DynkinExpectationCertificate}, and HJB retains ItoRun. These interfaces specify measure and expectation identities to be supplied by valid process theorems. They do not automatically complete all stochastic calculus.

\nolinkurl{StochasticSubstrate} connects these interfaces to P11 without turning control energy into unconditional thermodynamic work. If the density is not measurable from the observed path, Section 22.5 uses only the data-processing upper bound. The unit calibration of Chapter 33 must also be supplied separately.

Sections 39.9--39.15 provide actual all-time Gaussian and continuous Brownian sources, conditional moments for the natural filtration, and a genuine predictable L² terminal integral. Sections 39.17 and 39.19 further construct the actual continuous integral process for the same input, deriving a square-compensated martingale from actual time energy and conditional second-order relations; Brownian partition sums of squares also have a grid-error bound from the original source. Exponential densities on the natural filtration are normalized and connected to actual RN derivatives. Section 39.22 now derives actual quadratic variation and actual weighted first- and second-order stochastic sums along deterministic vanishing-mesh partitions for the same full predictable scalar input. Sections 39.26--39.28 derive same-input scalar Itô and Dynkin formulas from actual Taylor sums, the full predictable vector domain and conditional processes, the original OU conditional kernel, and actual \nolinkurl{TerminalGirsanovData} with the original KL consumer for constant controls. Section 39.30 constructs the full changed-measure path and original KL consumer for an actual finite feedback table; Section 39.31 adds a continuous version for the same full vector input. Global first- and second-derivative bounds and other conditions remain explicit. Section 39.34 further derives the full actual changed-measure path for the same bounded predictable input, the integral under the constructed \(Q\), and the original terminal KL result. Section 39.35 gives deterministic-partition QV for the same fixed full joint input. Section 39.39 constructs the original ItoRun from the same actual measure and connects the original control consumer for a nonconstant evaluation contract; Sections 39.42--39.43 connect actual finite-state calculus and a genuine martingale on the Q closure. General controlled SDEs, uncovered control domains, and HJB models retain their source obligations.

\chapter{23. Π/Φ, Constraint Geometry, and a Rigorous Dual-Drive Formulation}\label{ch-23}

\textbf{Chapter question.} Can resources and information serve as common drives, and how are their values and dimension identified?

\textbf{Reading route.} Distinguish an algebraic gauge preserving the objective, variational duality for fixed features, and the rank condition for a common two-dimensional representation. These are separate tasks. The three-response counterexample from the same regenerative source shows that organization and growth do not require every response to admit a smooth two-dimensional representation.

\section{23.1 A common interface for resources and information}\label{a-common-interface-for-resources-and-information}

The original \(\Pi\) denotes resources, returns, or functional efficiency; \(\Phi\) denotes predictive deficiency, decision-relevant uncertainty, or information cost. One may use \[
J=\Pi-\lambda_\Phi\Phi-C.
\] This is a usable template; the bridge determines units, signs, measurements, and comparisons. If \(\Phi\) denotes information gain rather than an information deficit, its sign should be changed accordingly.

Resource and information constraints can jointly determine feasible policies. This does not automatically establish that every system has exactly two effective response dimensions.

\section{23.2 Gauge freedom}\label{gauge-freedom}

\hypertarget{p-ddh-01}{}\phantomsection\label{p-ddh-01}

\textbf{P-DDH-01.} \[
\Pi'=\Pi+\lambda\chi,\qquad
\Phi'=\Phi+\chi
\] preserves \(\Pi'-\lambda\Phi'=\Pi-\lambda\Phi\).

\textbf{Proof.} Substitute and cancel. \(\blacksquare\)

Behavior therefore identifies at most the objective-equivalence class constrained by that behavior. Independent resource measurements, prediction experiments, or selective interventions anchor the semantic axes. Two arbitrarily rotated latent variables cannot automatically be called the physical and informational axes.

\section{23.3 Finite exponential families and constraint duality}\label{finite-exponential-families-and-constraint-duality}

Let \(\Gamma\) be a finite path set, \(p_0\) a strictly positive baseline, and \(F:\Gamma\to\mathbb R^k\) a prespecified feature map. Define \[
\psi(\lambda)=\log\sum_\gamma p_0(\gamma)e^{\lambda^\top F(\gamma)},
\quad
q_\lambda(\gamma)=p_0(\gamma)e^{\lambda^\top F(\gamma)-\psi(\lambda)}.
\]

\hypertarget{p-ddh-02}{}\phantomsection\label{p-ddh-02}

\textbf{P-DDH-02.} \[
\nabla\psi=E_{q_\lambda}F,\qquad
\nabla^2\psi=\operatorname{Cov}_{q_\lambda}(F).
\]

\textbf{Proof.} Differentiate the finite sum term by term. The first derivative is a normalized weighted mean; the second is the second moment minus the outer product of the mean. \(\blacksquare\)

\hypertarget{p-ddh-03}{}\phantomsection\label{p-ddh-03}

\textbf{P-DDH-03: variational minimality under moment constraints.} If the target moment is \(b=E_{q_\lambda}F\), then \(q_\lambda\), among all \(q\) satisfying \(E_qF=b\), is the unique minimizer of \(D_{\rm KL}(q\|p_0)\).

\textbf{Proof.} \[
D(q\|p_0)=D(q\|q_\lambda)+\lambda^\top E_qF-\psi(\lambda).
\] For fixed moments, the last two terms are constant and the first is nonnegative, vanishing only at \(q=q_\lambda\). \(\blacksquare\)

Boundary moments need not be attained by finite \(\lambda\); parameter limits or reduced support may be necessary. The construction is testable only when features and budgets are independently specified. Choosing \(F=\log(q/p_0)\) after the fact would express every observed law as a tilt.

The covariance and variational duality of finite exponential families are standard theory. \href{https://www.emerald.com/ftmal/article/1/1-2/1/1332382/Graphical-Models-Exponential-Families-and}{Wainwright--Jordan}

\section{23.4 A common two-dimensional bottleneck}\label{a-common-two-dimensional-bottleneck}

\textbf{H-DDH.} In fixed budget coordinates \(B\) and across all environments under comparison, there is a common differentiable map \(z(B)\in\mathbb R^2\) such that \[
m_e(B)=g_e(z(B)).
\]

\hypertarget{p-ddh-04}{}\phantomsection\label{p-ddh-04}

\textbf{P-DDH-04.} At the same budget point \(B\), the vertical stack of the environmental response Jacobians has rank at most 2.

\textbf{Proof.} The chain rule gives \(J_e=Dg_e\,Dz\). Stacking leaves the same \(2\times k\) matrix \(Dz\) as a right factor. \(\blacksquare\)

The matrix \(Dz(B)\) may rotate between budget points; stacking across points requires an additional common linear-subspace condition. Rank two in each environment also does not imply common rank two: stacking \(\operatorname{diag}(1,1,0)\) and \(\operatorname{diag}(0,1,1)\) produces rank three.

A single scalar objective \[
J(a,B)=-\tfrac12\|a-B\|^2
\] has optimal response \(a^*=B\) in arbitrarily high dimension. The number of objectives, number of departments, latent dimension, and response rank must therefore remain distinct.

\subsection{A three-response boundary in the same regenerative source}\label{a-three-response-boundary-in-the-same-regenerative-source}

This dimensional restriction can also be tested inside the actual source of the core overview. Fix \(\delta,\gamma\) and vary \((s,b,q)\) over the admissible open set. Start from healthy \(\theta_0\) and run its actual care program. Use three event probabilities from the public record: continuation of the original identity, a strict birth, and a strict birth whose newborn is damaged. The 36 primitive atoms and independent output damage give

\[F(s,b,q)=(s,b,bq),\qquad
DF=\begin{pmatrix}1&0&0\\0&1&0\\0&q&b\end{pmatrix},\qquad
\det DF=b>0.\]

The three readouts concern actual retained, newborn, and newborn H/D. The public receipt's \nolinkurl{PublicReceipt.newborn} retains the output observation, and \nolinkurl{diagnostic_observedBody} ensures that this observation reveals actual H/D; no additional sensor permission is needed. The third quantity is a joint event probability, not the conditional damage rate \(q\). With \(\delta=1/8,\gamma=1/4\), the point \((s,b,q)=(3/4,3/8,1/3)\) lies in the interior of the positive class, and all original conditions remain valid in a neighborhood.

If a common \(C^1\) encoder \(h\) into \(\mathbb R^2\) and a common \(C^1\) decoder \(g\) existed on such a three-dimensional neighborhood with \(F=g\circ h\), the chain rule would give \(\operatorname{rank}DF\le2\), contradicting the display. Thus these three public responses cannot be generated completely on that neighborhood by the same smooth two-dimensional latent variable. The proof uses only the original finite event law and the chain rule. This source calculation was supplied in the strengthening as a written corollary with exact verification; it was not added as a corresponding Lean endpoint.

This is a local dimension bound for source-parameter responses. It does not make the parameters independently manipulable budgets in reality, and it does not exclude an explicitly restricted two-dimensional parameter family, a Π/Φ model with additional conditions, or a nonsmooth encoding. The general conditional C5 theorem and the positive organization, recovery, and growth results for this source remain valid. What must be withdrawn is the unconditional extrapolation that all these responses fit a common smooth two-dimensional space.

\section{23.5 Anchors and spectral certification}\label{anchors-and-spectral-certification}

If a rank-two response matrix factors as \(J=AB^\top\), every \(R\in GL(2)\) gives the same matrix: \[
J=(AR)(BR^{-\top})^\top.
\] Two independent, predefined, oriented intervention axes reduce the allowed basis changes to positive rescaling along each axis. Specifying units may further fix the scales.

\hypertarget{p-ddh-05}{}\phantomsection\label{p-ddh-05}

\textbf{P-DDH-05: an effective-dimension threshold.} If \(\|\widehat S-S\|_2\le\eta\), every singular-value error is at most \(\eta\). If \[
\sigma_3(S)+\eta<\tau<\sigma_2(S)-\eta,
\] exactly two estimated singular values exceed \(\tau\).

\textbf{Proof.} The variational characterization of singular values gives the Lipschitz perturbation bound; insert the threshold inequalities. \(\blacksquare\)

This certifies effective dimension at a specified threshold. It is recovery of an exactly rank-two model only with the additional structural assumption \(\sigma_3=0\). Failure of a finite-sample third singular value to differ significantly from zero does not prove that its true value is zero.

\section{23.6 Using Π/Φ in applications}\label{using-ux3c0ux3c6-in-applications}

The resource side may represent task returns, material or energy budgets, time, control amplitude, or computational cost. The information side may represent predictive KL, required memory rate, Bayes risk, or policy-relevant uncertainty. State how each quantity is measured, whether it can be independently intervened upon, and how it determines behavior.

If the value of information is already fully represented in belief-space control, an information-gain reward should not be added mechanically unless it expresses a new preference or cost. Chapter 27 makes this distinction precise.

\section{23.7 Complete gauge characterization is separate from rank testing}\label{complete-gauge-characterization-is-separate-from-rank-testing}

For fixed \(\lambda\), P-DDH-01 has a converse: if \(\Pi'-\lambda\Phi'=\Pi-\lambda\Phi\), take \(\chi=\Phi'-\Phi\) to obtain \(\Pi'=\Pi+\lambda\chi\). \nolinkurl{DualDriveGaugeCompleteness.same_objective_iff_gauge_related} therefore completely characterizes the gauge orbit of the same numerical objective. Chapter 33 gives the precise conditions for identification by anchors.

This algebraic iff does not establish a common two-dimensional response bottleneck in natural systems. C5's common-rank necessity, singular-value errors, and three-valued conclusions serve another purpose: compatibility means the test has not excluded the representation, not that the latent variables have been uniquely identified. Two independent physical axes still require independent measurements and interventions.

\chapter{24. Multilevel Value, Integrability, and Coordination Mechanisms}\label{ch-24}

\textbf{Chapter question.} When are evaluations at different levels compatible, and when do local actions increase parent-level value?

\textbf{Reading route.} Retain the partial order of multiple objectives before checking whether a direction field has a global potential. The chain rule gives instantaneous value change and a coordination threshold. These local sufficient conditions guarantee neither global optimality nor the normative justification of a coordination institution.

\section{24.1 Multiple objectives and common optimality}\label{multiple-objectives-and-common-optimality}

For evaluations \(J_1,\ldots,J_m\) at different levels, retain the Pareto partial order by default. If \[
\bigcap_{\ell=1}^m\mathfrak P_\ell^*\ne\varnothing,
\] a jointly optimal policy exists. An empty intersection means no exactly common optimum exists; negotiation, constraints, or additional weights are required.

For two policies with value vectors \((1,0),(0,1)\), no symmetric scalarization can choose uniquely between them if neither level has priority and exchange symmetry is required. Additional weights belong to the value bridge and do not follow from objecthood.

\section{24.2 When a direction field has a global potential}\label{when-a-direction-field-has-a-global-potential}

On a connected smooth manifold, let a \(C^1\) one-form be denoted by \(\omega\).

\hypertarget{p-ali-01}{}\phantomsection\label{p-ali-01}

\textbf{P-ALI-01: global integrability criterion.} \[
\omega=dV
\iff
\oint_C\omega=0
\quad\text{for every piecewise smooth closed curve }C.
\]

\textbf{Proof.} The forward implication is the fundamental theorem of calculus. Conversely, fix a base point \(x_0\) and define \(V(x)=\int_{\gamma_{x_0,x}}\omega\). Vanishing integrals around closed curves ensure path independence. Differentiating the endpoint in local coordinates gives \(dV=\omega\). \(\blacksquare\)

Local closedness \(d\omega=0\) does not guarantee a global potential on a space with topological holes. A specified metric converts a drift into a one-form; the metric is part of the model, with no privileged Euclidean form exempt from coordinate transport.

\textbf{K-GEO-01.} On a compact oriented Riemannian manifold without boundary, appropriately smooth one-forms have a Hodge decomposition \(\omega=dV+\delta B+h\). Boundary conditions must be specified when a boundary is present. The decomposition separates gradient, circulating, and harmonic components. Nongradient motion does not imply the absence of long-term control value.

\section{24.3 Child-level updates and parent-level value}\label{child-level-updates-and-parent-level-value}

On a real Hilbert direct-sum policy space, let the parent value \(J_P\) and child values \(J_i\) be differentiable. Put \[
g_i=\nabla_iJ_P,\quad
e_i=\nabla_iJ_i-g_i,\quad
\dot a_i=\alpha_i\nabla_iJ_i,\quad\alpha_i>0.
\]

\hypertarget{p-ali-02}{}\phantomsection\label{p-ali-02}

\textbf{P-ALI-02.} \[
\frac{dJ_P}{dt}
=\sum_i\alpha_i(\|g_i\|^2+\langle g_i,e_i\rangle)
\ge\sum_i\alpha_i\|g_i\|(\|g_i\|-\|e_i\|).
\]

\textbf{Proof.} Apply the chain rule, substitute the update, and use Cauchy--Schwarz. \(\blacksquare\)

If every nonzero parent-gradient component satisfies \(\|e_i\|<\|g_i\|\) and at least one is nonzero, parent value strictly increases. This instantaneous sufficient condition does not guarantee convergence to a global optimum.

\section{24.4 A coordination threshold}\label{a-coordination-threshold}

Replace the joint child field \(G\) by \(G_\eta=(1-\eta)G+\eta g\), where \(g=\nabla J_P\) and \(0\le\eta\le1\).

\hypertarget{p-ali-03}{}\phantomsection\label{p-ali-03}

\textbf{P-ALI-03.} If \(\langle g,G\rangle<0\), the instantaneous derivative of parent value is positive exactly when \[
\eta>
\eta_*=\frac{-\langle g,G\rangle}
{\|g\|^2-\langle g,G\rangle}.
\]

\textbf{Proof.} Solve the linear inequality obtained from \(\dot J_P=(1-\eta)\langle g,G\rangle+\eta\|g\|^2\). \(\blacksquare\)

This gives an interpretable mechanism-design template. Whether taxation, incentives, permissions, communication, or feedback control implements such a field mixture requires a separate model. The formula itself does not establish the normative legitimacy of an institution.

\section{24.5 Bridges between evaluation levels and other hierarchies}\label{bridges-between-evaluation-levels-and-other-hierarchies}

H0-VAL differs from predictive quotients, physical resolution, and levels of selection. The parent-value gradient, child update field, and weights in this chapter must be compared in the same policy space. A parent--child object map alone does not generate an inner product, a value, or \(\eta_*\).

In particular, local language about a ``driving direction'' cannot obscure the global period obstruction of P-ALI-01. A mathematical drift may contain gradient, circulating, and harmonic components. Interpreting which component relates to purpose also requires an explicit evaluation mechanism. Chapter 34 retains this faithfulness interface.

\chapter{25. Reproduction, Heredity, and Positive-Operator Evolution}\label{ch-25}

\textbf{Chapter question.} How do actual parent--offspring mechanisms generate mean growth, hereditary transmission, and selection decompositions?

\textbf{Reading route.} Fix individuals, types, generations, and sampling units before defining the same mean operator from offspring counts. Price's identity is an algebraic decomposition of that operator, and Perron theory describes a primitive mean model. Extinction, martingale limits, and mean frequencies of actual random populations remain separate issues.

\section{25.1 From objects to an offspring mechanism}\label{from-objects-to-an-offspring-mechanism}

Reproduction is not necessary for objecthood. Before this chapter applies, parents and candidate offspring should already have the object certificates of Chapter 12 and the historical relations of Chapter 11; detached material is not automatically offspring. Fix a finite type set \(\{1,\ldots,d\}\), a generation length, and a counting unit. An individual of type \(i\) produces a random offspring-count vector \(N^{(i)}\in\mathbb N^d\). Define \[
M_{ij}=E N_j^{(i)}<\infty,\qquad b_i=\sum_jM_{ij}.
\] Here \(M\) is the mean reproduction operator and \(b_i\) the mean total offspring count. If \(b_i>0\), put \(K_{ij}=M_{ij}/b_i\), so that \[
M=D_bK.
\] This uniquely separates reproductive intensity from type transmission \textbf{sampled in proportion to offspring intensity}. When \(b_i=0\), that row has no offspring and observations do not define \(K_{i\cdot}\). An arbitrary probability row leaves \(M\) unchanged but cannot be claimed to have been identified.

Observe that \[
K_{ij}=\frac{EN_j^{(i)}}{E\sum_kN_k^{(i)}}
\] is generally different from \(E[N_j/\sum_kN_k\mid\sum_kN_k>0]\), obtained by first sampling a nonempty family uniformly and then sampling a child uniformly within that family. Transmission probabilities and hereditary information require a fixed sampling unit.

\section{25.2 Mean-frequency dynamics and Price's identity}\label{mean-frequency-dynamics-and-prices-identity}

In the deterministic mean-intensity model, let type frequencies be a row vector \(p\) with \(\bar b=pb>0\). Then \[
p'=\frac{pM}{\bar b}.
\] This is not an exact formula for \(E[Z_{n+1}/|Z_{n+1}|]\) in a general finite random population. The expectation of a random ratio and the ratio of expectations must remain distinct.

\hypertarget{p-evo-01}{}\phantomsection\label{p-evo-01}

\textbf{P-EVO-01: exact decomposition into selection and transmission.} For numerical parent-type traits \(z_i\) and offspring-type traits \(z'_j\), put \((Kz')_i=\sum_jK_{ij}z'_j\). Then \[
p'z'-pz=
\frac{\operatorname{Cov}_p(b,z)}{\bar b}
+\frac{\sum_i p_ib_i[(Kz')_i-z_i]}{\bar b}.
\]

\textbf{Proof.} We have \(p'z'=\sum_i p_ib_i(Kz')_i/\bar b\). Add and subtract \(\sum_i p_ib_iz_i/\bar b\); the first remaining term is \((E_p[bz]-E_p[b]E_p[z])/\bar b\). Zero-reproduction rows are multiplied by zero, so the unidentified \(K\) on those rows is irrelevant. \(\blacksquare\)

This is an algebraic identity, not an independent causal model of selection. A causal interpretation additionally requires \(b,K\) to have mechanistic meaning under the relevant interventions.

\section{25.3 Hereditary information and heritable predictive kernels}\label{hereditary-information-and-heritable-predictive-kernels}

Growth alone does not establish heredity. If all rows \(K_{i\cdot}\) coincide, offspring type is independent of parent type, however large \(b_i\) may be. Under a specified parent--offspring joint distribution one may report \(I(F;F')\), but it depends on the type partition.

\hypertarget{p-evo-02}{}\phantomsection\label{p-evo-02}

\textbf{P-EVO-02: irrelevant labels inflate hereditary information.} If a finite label \(T\) is independent of \((F,F')\) and copied without error, then \[
I((F,T);(F',T))=I(F;F')+H(T).
\]

\textbf{Proof.} Expand by the chain rule. Independence removes the cross terms, while the shared label contributes \(I(T;T)=H(T)\). \(\blacksquare\)

The core therefore prioritizes a canonical object tied to offspring behavior. For a prespecified lineage-sampling mechanism, observation target, and countable intervention family, map parent histories to the family of conditional laws of all declared offspring futures. By P-PRED-01 this \textbf{lineage predictive state} is the minimal sufficient representation of that prediction problem. It excludes labels irrelevant to the specified future, without declaring them meaningless for all other tasks.

Lineage sampling may select one random lineage under a fixed inheritance rule or the full descendant tree. These are different future spaces; their entropies and growth rates cannot be interchanged.

\section{25.4 Perron growth and composition}\label{perron-growth-and-composition}

\textbf{K-PF-01 (reused and expanded): the finite primitive nonnegative matrix theorem.} Fix a nonnegative matrix M on a nonempty finite type set. If all entries of \(M^k\) are strictly positive for some positive integer k, there exist \(R>0\), a positive column vector \(r\), and a positive row vector \(l\) such that \[
Mr=Rr,\quad lM=Rl,\quad l\mathbf1=1,\quad lr=1.
\] The principal eigenvalue \(R\) is simple, all other eigenvalues have modulus strictly less than \(R\), and \(R^{-n}M^n\to rl\). This chapter invokes Perron--Frobenius only under these explicit finite-primitivity assumptions.

\hypertarget{p-evo-03}{}\phantomsection\label{p-evo-03}

\textbf{P-EVO-03: mean growth and reproductive value.} For a nonzero nonnegative initial count row \(z_0\), \[
R^{-n}z_0M^n\to(z_0r)l,\qquad
\frac{z_0M^n}{z_0M^n\mathbf1}\to l.
\] If \(w>0\) satisfies, for every nonnegative \(z\), the equality \((zM)w=\rho zw\), then \(\rho=R\) and \(w\) is a positive multiple of \(r\).

\textbf{Proof.} Left-multiply the matrix limit in K-PF-01 by the initial count row for the first assertion; the limiting denominator is \(z_0r>0\). For the second, use coordinate unit vectors to obtain \(Mw=\rho w\), then invoke Perron uniqueness of a positive eigenvector. \(\blacksquare\)

\textbf{Current Lean source.} \nolinkurl{PrimitivePerronClosure.certificateOfPrimitive} takes only the same M and \nolinkurl{Matrix.IsPrimitive} M as external inputs. Positive right data, left data and dual normalization, the full power limit, strict complex spectral dominance, and uniqueness of the positive eigendirection are derived from the general source. \nolinkurl{actual_root_simple} and \nolinkurl{PrimitivePerronSimplicity} also prove algebraic multiplicity one of the actual principal root of the original real matrix and its complexification. \nolinkurl{actual_p_evo_03} directly invokes the original growth, composition, and evaluation consumers, without replacing the general quantifiers by a numerical finite instance. Section 39.72 supplies the decomposition and an independent example.

The right vector \(r\) gives relative reproductive value for a linear count evaluation; the left vector \(l\) gives long-term mean type composition. ``Canonical'' here is relative to the specified positive linear reproduction mechanism, not a universally unique value for all objects.

\section{25.5 Random populations and the boundary of mean models}\label{random-populations-and-the-boundary-of-mean-models}

Suppose offspring of different individuals are independent conditional on type, reproduction laws are constant across generations, and \(Z_n\) is the actual population count. Then \[
E[Z_{n+1}\mid\mathcal F_n]=Z_nM.
\]

\hypertarget{p-evo-04}{}\phantomsection\label{p-evo-04}

\textbf{P-EVO-04: reproductive-value martingale.} With finite initial counts, \[
W_n=R^{-n}Z_nr
\] is a nonnegative integrable martingale, and \(EW_n=Z_0r\).

\textbf{Proof.} Its conditional expectation is \(R^{-(n+1)}Z_nMr=R^{-n}Z_nr\). Nonnegativity and finite first moments give integrability. \(\blacksquare\)

\textbf{K-MAR-01.} A nonnegative martingale converges almost surely to a finite \(W_\infty\). This alone gives neither \(EW_\infty=EW_0\), \(L^1\) convergence, nor \(W_\infty>0\) on nonextinction. Those require additional uniform integrability and offspring moment conditions. This edition does not include the corresponding deeper branching-process theorems as unconditional core results. Random extinction, finite-population drift, and deterministic mean growth must be reported separately.

\section{25.6 Bridging actual lineage through the same reproduction operator}\label{bridging-actual-lineage-through-the-same-reproduction-operator}

\nolinkurl{LineageEvolution}'s \nolinkurl{lineage_price_bridge} consumes the same transmission kernel as the Price algebra. Its \nolinkurl{lineage_perron_bridge} consumes the same matrix and KPF01Certificate as P-EVO-03 and requires every positive reproductive edge to have actual offspring semantics. Thus the spectral conclusion cannot come from an unrelated matrix. Primitivity alone does not supply the offspring-identity condition. Mean-operator indices are types, not permanent individuals: identifying them would make a positive diagonal entry require an individual to be its own distinct offspring. Actual identities across generations use the rectangular birth and clock sources proved in Chapter 39; the normalization obstruction for a strictly ordered finite permanent lineage remains in force.

Repair of the same identity and production of an offspring with a distinct identity are formalized as mutually exclusive roles. Spectral existence, growth, and random-population results retain this chapter's conditions. The conclusions do not state that every object spontaneously forms a lineage satisfying Perron assumptions.

\chapter{26. Transitions in Individuality and Endogenous Purpose Bridges}\label{ch-26}

\textbf{Chapter question.} How can parent-object formation, reproductive advantage, and endogenous purpose be connected over time?

\textbf{Reading route.} Separate object scale from level of selection, then specify the bridge from behavior to reproduction. An explicit frequency formula proves concentration under fixed fitness; frequency-dependent cycles mark its boundary. An endogenous mechanism must still account for its program, resources, and external scaffolding.

\section{26.1 Object scale differs from level of selection}\label{object-scale-differs-from-level-of-selection}

Object scale \(s\) describes candidate parent--child decompositions and observation interfaces. Selection level \(\ell\) specifies the units counted, inherited, and differentially reproduced. They may coincide, but need not. A cooperative group may form a parent object over a finite horizon without group-level reproduction; a replicator may borrow another object's maintenance mechanism.

The core uses four components of an \textbf{operational ETI certificate}, without claiming a logically necessary and sufficient definition of every natural transition in individuality:

\begin{enumerate}
\def\labelenumi{\arabic{enumi}.}
\tightlist
\item
  The parent level has predictive, persistence, and cut-irreducibility certificates.
\item
  Children have valid realization and resolution interfaces into the parent level.
\item
  Parent-level units have an identifiable reproduction and transmission mechanism.
\item
  Within-group conflict, child fidelity, and parent reproductive advantage lie in a prespecified admissible region.
\end{enumerate}

The first two alone establish parent-object formation, not an evolutionary transition in individuality. The third alone does not establish an irreducible parent object.

\section{26.2 Mechanism internalization and withdrawal of scaffolding}\label{mechanism-internalization-and-withdrawal-of-scaffolding}

Partition maintenance mechanisms into candidate internal mechanisms \(I\) and external scaffolding \(S\). Under the same resource budget and comparable environmental protocols, perform a declared withdrawal or replacement intervention on the scaffolding and recompute object certificates and persistence.

Failure after removing \(S\) establishes dependence \textbf{under that intervention and replacement class}; it does not rule out every other external realization. Persistence after withdrawal shows that this scaffold is unnecessary for that certificate, but does not automatically prove complete causal closure. Report which causal and resource edges the intervention changes, so that loss of power is not confused with organizational disintegration.

\section{26.3 How selection induces a control objective}\label{how-selection-induces-a-control-objective}

Let a heritable controller type \(\theta\), in environment \(e\), implement policy \(\pi_\theta\), generating path law \(Q_\theta^e\). Specify a reproduction mechanism \[
R_e(\theta)=\mathcal R_e(Q_\theta^e)>0.
\] The behavior-to-reproduction map \(\mathcal R_e\) is the \textbf{purpose bridge}. It may depend on survival, output, or another explicit mechanism; an object certificate does not uniquely specify it.

\hypertarget{p-brg-01}{}\phantomsection\label{p-brg-01}

\textbf{P-BRG-01: concentration under fixed fitness without mutation.} For finitely many types with fixed \(R_i>0\), \[
p_i(n+1)=\frac{p_i(n)R_i}{\sum_jp_j(n)R_j}
\quad\Longrightarrow\quad
p_i(n)=\frac{p_i(0)R_i^n}{\sum_jp_j(0)R_j^n}.
\] Let \(R_*=\max_{i:p_i(0)>0}R_i\). The total mass of types below \(R_*\) decays exponentially, while the maximizers retain their initial relative proportions.

\textbf{Proof.} Induction gives the explicit formula. Divide numerator and denominator by \(R_*^n\). Suboptimal terms decay as \((R_i/R_*)^n\), and optimal terms remain unchanged. \(\blacksquare\)

Under these conditions, evolution therefore induces behavioral selection for maximal \(R\). It does not create controllers absent from the initial support. With a mutation kernel \(K\), long-term composition is governed by the Perron structure of \(D_RK\), rather than generally concentrating on a single maximal-\(R\) type.

\section{26.4 No unconditional principle that evolution is optimization}\label{no-unconditional-principle-that-evolution-is-optimization}

\textbf{N-BRG-01: frequency-dependent cycling.} On the interior of the simplex, take an antisymmetric rock--paper--scissors matrix \(A\) with zero row and column sums, and let \[
\dot p_i=p_i(Ap)_i.
\] Since \(p^\top Ap=0\), this is replicator dynamics. Moreover, \[
\frac{d}{dt}\log(p_1p_2p_3)=\mathbf1^\top Ap=0.
\] A noncentral interior trajectory cannot converge to the boundary or the unique interior equilibrium: its product is a fixed positive number below the central value. Thus not every selection trajectory converges to an optimal point. Antisymmetry also prevents a general interpretation of this dynamics as the flow of one strictly increasing potential.

If two environments reverse the ordering of \(R_1,R_2\), their fixed-fitness orderings do not jointly specify a unique optimal type. Generalization across environments must specify an environment distribution, a worst-case objective, or an adaptation mechanism.

\hypertarget{p-brg-02}{}\phantomsection\label{p-brg-02}

\textbf{P-BRG-02: selection cannot distinguish behaviorally equivalent types.} In the same environment, suppose reproduction depends entirely on \(Q_\theta^e\), with no type-dependent mutation or other reproduction channel. If \(Q_\theta^e=Q_{\theta'}^e\), then \(R_e(\theta)=R_e(\theta')\), and the types cannot be differentially selected by that reproduction bridge alone.

\textbf{Proof.} Substitute equal arguments into the same \(\mathcal R_e\). \(\blacksquare\)

This is an interface-consistency result: attaching different semantic labels to controllers does not automatically change the declared selection dynamics. Different mutation neighborhoods or material costs change the model and require renewed examination.

\section{26.5 The closed-loop architecture of endogenous purpose}\label{the-closed-loop-architecture-of-endogenous-purpose}

The full loop is a temporal recursion: \[
\theta_n \to \pi_{\theta_n}\to Q_{\theta_n}^{e_n}
\to\text{reproduction/learning feedback}\to\theta_{n+1}.
\] Every arrow requires an independent mechanism. Explaining survival or growth in a particular environment can strengthen the original paper's account of function and purpose, but it does not deduce a unique moral value from existence. Cross-level value conflicts retain vector form or an explicit coordination mechanism, as in Chapter 24.

\section{26.6 Source obligations in a conditional constructive closure}\label{source-obligations-in-a-conditional-constructive-closure}

P7/P12 start from an \nolinkurl{OntogeneticSeed} containing both physical carriers and program provenance. A trusted substrate constructs the organization and enters recovery and maintenance processes. Existing no-go results show that identical physical components with different program sources cannot both be recovered correctly from the physical components alone. This differs from the question of which evaluation selects offspring; the latter requires a reproduction--purpose mechanism bridge.

Endogeneity can therefore mean that mechanisms become maintained or transmitted by the current system over an explicit temporal process. It cannot mean that the program, interpreter, objective, and environmental laws used in the proof never required sources. Section 34.5 lists the connected lifecycle and its retained premises in full.

\chapter{27. Reflexive Structure, Belief Learning, and Agency}\label{ch-27}

\textbf{Chapter question.} How can a system that learns or modifies itself retain a valid state and decision description?

\textbf{Reading route.} Include structure in the state, then construct Bayesian updating within a specified model class. Free information that may be ignored has nonnegative decision value, but information gain need not improve task return. Objective error, resource costs, and permission constraints must be explicitly accounted for in the actual decision problem.

\section{27.1 Including modifiable structure in the state}\label{including-modifiable-structure-in-the-state}

Let the physical state be \(X_t\) and controller or model structure be \(S_t\). If \[
P(X_{t+1},S_{t+1}\mid H_t,a_t)
=K(X_t,S_t,a_t;\cdot),
\] then \(Z_t=(X_t,S_t)\) is a controlled Markov state. Self-modification does not mathematically preclude a state-space description, although the resulting space may be large or admit no finite compression.

\hypertarget{p-ref-01}{}\phantomsection\label{p-ref-01}

\textbf{P-REF-01: reflexive state augmentation.} Under the displayed condition and measurable-kernel assumptions, every causal policy determines a unique path law of \(Z\) from its initial law and \(K\). A deterministic modification \(S_{t+1}=F(S_t,X_t,a_t,\xi_t)\) is a special case.

\textbf{Proof.} Conditional independence gives the one-step kernel factorization; apply the kernel extension of Chapter 2. If compressed to \(\phi(Z_t)\), the controlled closure conditions of Chapter 7 must be established again. Reflexivity does not grant automatic closure. \(\blacksquare\)

\section{27.2 The Bayes interface for finite hidden-state models}\label{the-bayes-interface-for-finite-hidden-state-models}

Let both the static unknown parameter \(\theta\in\Theta\) and hidden state \(x\in\mathcal X\) be finite. The known model family has transition \(P_\theta(x'|x,a)\) and observation kernel \(O_\theta(y|x',a)\). The current belief is \(b_t(\theta,x)=P(\theta,X_t=x\mid\text{observed history})\). Then \[
b_{t+1}(\theta,x')=
\frac{O_\theta(y_{t+1}|x',a_t)
\sum_xP_\theta(x'|x,a_t)b_t(\theta,x)}
{\sum_{\vartheta,z}O_\vartheta(y_{t+1}|z,a_t)
\sum_xP_\vartheta(z|x,a_t)b_t(\vartheta,x)}.
\] An observation with zero denominator is impossible in the current model. The interface should return a model conflict; it cannot insert an arbitrary posterior and continue to claim certification.

\hypertarget{p-ref-02}{}\phantomsection\label{p-ref-02}

\textbf{P-REF-02: belief sufficiency.} In this model, the conditional laws of the next observation and updated belief depend only on \(b_t,a_t\). A bounded discounted task can therefore be rewritten as a belief-state control problem.

\textbf{Proof.} Total probability over hidden state and parameter gives the predictive observation law in the denominator; Bayes' formula gives the update. The next belief is a deterministic function of the current belief, action, and observation, so its kernel depends only on \((b_t,a_t)\). The conditional expectation of a bounded reward is handled in the same way. \(\blacksquare\)

This is sufficiency within the model. If the actual mechanism lies outside the model family, parameters drift, or a modification action changes \(P_\theta,O_\theta\), the corresponding mechanism state must be augmented or handled through the error interfaces of Chapter 16.

\section{27.3 Value of information is not information gain}\label{value-of-information-is-not-information-gain}

\hypertarget{p-ref-03}{}\phantomsection\label{p-ref-03}

\textbf{P-REF-03: free information has nonnegative decision value.} In a one-shot decision problem, suppose a signal \(Y\) does not change the world distribution or action set, costs nothing, and may be ignored. For finitely many actions and integrable utility, \[
E_Y\max_a E[u(a,\theta)\mid Y]
\ge \max_a E[u(a,\theta)].
\]

\textbf{Proof.} Fix any action \(a\). For each \(Y\), the maximum is at least the conditional utility of that action. Take expectations, then maximize over the fixed actions. \(\blacksquare\)

If the signal reveals an independent coin irrelevant to every action's payoff, mutual information can be \(\log2\) while decision value is zero. Information with costs, physical effects, or interference with action may have negative net value. Thus \(\Phi\) cannot unconditionally mean maximizing whatever information is available.

\section{27.4 Objective revision and misspecification risk}\label{objective-revision-and-misspecification-risk}

\hypertarget{p-ref-04}{}\phantomsection\label{p-ref-04}

\textbf{P-REF-04: regret under uniform objective error.} Suppose, on the same policy set, \[
\sup_\pi|J_g(\pi)-J_{\hat g}(\pi)|
\le Ld(g,\hat g),
\] and \(\pi_g,\pi_{\hat g}\) are respectively optimal. Then \[
J_g(\pi_g)-J_g(\pi_{\hat g})\le2L d(g,\hat g).
\]

\textbf{Proof.} Add and subtract \(J_{\hat g}(\pi_g),J_{\hat g}(\pi_{\hat g})\). The middle term is nonpositive by optimality; each remaining term is at most \(Ld\). \(\blacksquare\)

The bound must hold uniformly over all relevant policies. It neither converts normative disagreement into an already measured distance nor guarantees small costs for irreversible actions. If the reward error is at most \(\epsilon\) at every step, discounted return allows \(Ld=\epsilon/(1-\beta)\).

\section{27.5 Operational components of agency}\label{operational-components-of-agency}

For a fixed task and initial conditions, measure separately:

\begin{itemize}
\tightlist
\item
  the implementable policy set and its effect on future distributions;
\item
  the ability to identify the environment and use feedback;
\item
  the ability to maintain one's own realization, models, and permission boundaries;
\item
  the ability to extend or revise control structures;
\item
  the ability to recover, roll back, and stop under cost and safety constraints.
\end{itemize}

\hypertarget{p-ref-05}{}\phantomsection\label{p-ref-05}

\textbf{P-REF-05: monotonicity under feasible-choice expansion.} Under the same evaluation \(J\), if \(\Pi_1\subseteq\Pi_2\), then \(\sup_{\Pi_2}J\ge\sup_{\Pi_1}J\).

\textbf{Proof.} Take the supremum over the larger set. \(\blacksquare\)

If a new capability necessarily consumes resources or removes old policies, the inclusion may fail and the theorem does not apply. Expanding permissions is not a free increase in freedom.

For maintenance capability, one may compare matched interventions by \[
\Delta_{\rm self}(T)
=P^{\rm active}(\tau>T)-P^{\rm disabled}(\tau>T).
\] A causal interpretation requires comparable resources and external conditions apart from the mechanism under test. This quantity differs from the maintenance information cost of Chapter 22.

\section{27.6 Three interfaces of meaning}\label{three-interfaces-of-meaning}

Functional meaning is a mechanism's counterfactual contribution to object certification, persistence, or task value. Intentional or representational meaning concerns protocol-defined relations of internal representations to environmental variables and action consequences. Normative meaning concerns which outcomes merit pursuit and requires an additional value bridge and intersubjective constraints.

The core can formalize some observable relations in the first two layers and represent constraints in the third. It cannot derive experience, consciousness, or a unique ultimate value from predictive compression, reproductive advantage, or a gradient formula. Explicit boundaries preserve the verifiable content of the unified architecture.

\section{27.7 General belief recursion and implementable policy existence}\label{general-belief-recursion-and-implementable-policy-existence}

Under a standard Borel latent space and explicit countable-generation conditions, \nolinkurl{C1ParametricBelief} constructs a jointly measurable parameterized posterior and aligns it with an existing actual joint law. Existence of conditional distributions for fixed parameters does not suffice for a simultaneously measurable choice across all beliefs and actions. C1 specifically fills that interface.

Policy existence for finite-support control additionally requires a viability recursion. Verifying actual support for a fixed policy proves the safety of that policy. Treating that support as an unchanged input while choosing a different policy may change future reachability. The backward recursion and forward actual-support proof of Chapter 37 avoid this circularity without assuming the support of an unknown optimal policy.

\chapter{28. Complete Interfaces, Algorithms, and Composition Contracts}\label{ch-28}

\textbf{Chapter question.} How can proved interfaces be passed to downstream models without losing their assumptions?

\textbf{Reading route.} The contract table specifies inputs and outputs, and process interfaces compose by path pushforward. Successive refinement proves termination and coarseness of finite stable partitioning. Control lifting, exact support decisions, and physical interpretations each require their own additional proof.

\section{28.1 Table of module contracts}\label{table-of-module-contracts}

An interface is both a mathematical type and a specification of the information an implementation must supply. A downstream module may not silently assume a missing premise.

\begin{longtable}[]{@{}
  >{\raggedright\arraybackslash}p{(\columnwidth - 6\tabcolsep) * \real{0.2500}}
  >{\raggedright\arraybackslash}p{(\columnwidth - 6\tabcolsep) * \real{0.2500}}
  >{\raggedright\arraybackslash}p{(\columnwidth - 6\tabcolsep) * \real{0.2500}}
  >{\raggedright\arraybackslash}p{(\columnwidth - 6\tabcolsep) * \real{0.2500}}@{}}
\toprule\noalign{}
\begin{minipage}[b]{\linewidth}\raggedright
Module
\end{minipage} & \begin{minipage}[b]{\linewidth}\raggedright
Required inputs
\end{minipage} & \begin{minipage}[b]{\linewidth}\raggedright
Outputs and certificates
\end{minipage} & \begin{minipage}[b]{\linewidth}\raggedright
Principal failure states
\end{minipage} \\
\midrule\noalign{}
\endhead
\bottomrule\noalign{}
\endlastfoot
Process & \nolinkurl{State/history} spaces, initial law, time, measurable transition and observation kernels & Path laws under protocols & No measurable kernel; inconsistent support \\
Protocol & Intervention order, available information, action constraints, environment version & Legitimate causal protocol families & Future information or unimplementable actions \\
Factor & Candidate readout, underlying world, predictive target, horizon & Factor and its relative predictive state & Comparing labels only; changing the future target \\
Interface & \nolinkurl{Internal/external} variables, intervention semantics, conditional kernels & Joint residual and separated-closure certificates & Missing synergistic information or support coverage \\
Dynamics & Action consistency, quotient map, pushforward kernel & Exact or approximate closure & History dependence omitted from the state \\
Persistence & Evaluation domain, recovery set, initial law, exit mechanism & \nolinkurl{Invariance/recovery/survival/QSD} profile & Failure of compactness, absorption, or spectral conditions \\
Continuation & Historical edges, developmental transport, tolerated errors & Directed identity paths and fidelity profiles & Treating resemblance as lineage; uncalibrated transport \\
Realization & Physical readout sets, access and destruction operations & Carrier families, access blockers, mechanism-failure families & Confusing readout deletion with physical removal \\
Resolution & Domains at two scales, maps, closures, visible blocks & Transport rules and IR--UV intervals & Nonmeasurable quotients or failure of deletion naturality \\
Composition & Candidate parent future, cuts, conditional environment & Parent irreducibility and child fidelity & Common drivers; cuts changing the target \\
Statistics & Sampling design, confidence level, coverage, model class & Error events or unresolved structural families & Omitted data dependence; deployment outside distribution \\
Control & Feasible policies, \nolinkurl{rewards/risks}, horizon, discount & Value, GOD, policy error & Undeclared purpose bridge; misuse of average bounds \\
GOA & Specified optimal policy family and closed-loop model & Long-term \nolinkurl{sets/distributions} and stability & Near-bifurcation behavior, no mixing, nonunique policies \\
Evolution & Individual units, reproduction law, sampling, mutation & Mean operator, hereditary kernel, growth structure & Confusing mean models with random ratios \\
Learning & Model family, observation kernel, structural updates, costs & Belief updates and feasible-policy expansion & Zero likelihood, misspecification, unauthorized modification \\
\end{longtable}

The modules allow different mathematical implementations. Finite tables, stochastic differential equations, path kernels, and executable simulators are possible backends. A shared interface does not require a shared physical ontology.

\section{28.2 An object-certificate data structure}\label{an-object-certificate-data-structure}

The following schema illustrates a specification; it does not assert that general automatic proof software already exists. Numerical entries must carry units, metrics, quantifiers, and applicability domains.

\begin{verbatim}
{
  "schema": "UEOT-Core/ObjectCertificate/3.0",
  "object_id": "example-object",
  "world_version": "specified-model-v1",
  "factor": {
    "readout": "f",
    "target": "factor-relative-future",
    "horizon": 8,
    "protocol_family": "declared-family"
  },
  "claims": [
    {
      "id": "predictive-closure",
      "status": "certified-conditional",
      "metric": "TV",
      "quantifier": "uniform-on-declared-support",
      "bound": 0.01,
      "failure_probability": 0.05,
      "assumptions": ["finite-state", "independent-repeated-sampling"],
      "evidence": ["model-id", "data-id", "proof-id"]
    }
  ],
  "continuation": {"incoming_edges": [], "transport_model": "declared"},
  "realization": {"universe": ["r1", "r2", "r3"], "operation": "access"},
  "persistence": {"type": "finite-horizon-survival", "initial_law": "mu"},
  "purpose_bridge": null,
  "unresolved": ["physical-deletion-equivalence"]
}
\end{verbatim}

Recommended states are defined, proved-conditional, certified-conditional, estimated, hypothesized, refuted, unresolved, and out-of-scope. A zero numerical estimate cannot automatically become proved-conditional. A null purpose bridge is not filled by default with survival maximization.

Every certificate should retain the model version, data scope, protocol support, proof identifier, error sources, and expiry or drift-check conditions. Instance identifiers are tracking tools, not mathematical proofs of identity.

\section{28.3 Composability of path interfaces}\label{composability-of-path-interfaces}

A process interface from model \(A\) to model \(B\) consists of:

\begin{itemize}
\tightlist
\item
  a protocol lift \(J_{AB}:\mathcal U_B\to\mathcal U_A\);
\item
  a measurable path readout \(C_{AB}:\Omega_A\to\Omega_B\);
\item
  naturality under the specified initial law and every declared protocol: \[
  (C_{AB})_\#P_A^{J_{AB}u}=P_B^u.
  \]
\end{itemize}

\hypertarget{p-api-01}{}\phantomsection\label{p-api-01}

\textbf{P-API-01: composition of exact and approximate process interfaces.} If \(A\to B\) and \(B\to C\) have these types, then \[
J_{AC}=J_{AB}\circ J_{BC},\qquad
C_{AC}=C_{BC}\circ C_{AB}.
\] Exact interfaces compose exactly. If their TV defects over all corresponding protocols are bounded by \(\epsilon_{AB},\epsilon_{BC}\), the composite defect is at most \(\min(1,\epsilon_{AB}+\epsilon_{BC})\).

\textbf{Proof.} In the exact case, substitute naturality twice. In the approximate case, insert \((C_{BC})_\#P_B^{J_{BC}u}\) between the composite pushforward and the target, then use the triangle inequality and P-MET-01. \(\blacksquare\)

The identity interface uses identity maps; associativity follows from function composition. Time coarsening, readout resolution, and protocol restriction can fit this form, provided their path and protocol maps are explicit. An informal analogy between domains is not a process interface.

\section{28.4 Control interfaces require additional structure}\label{control-interfaces-require-additional-structure}

To lift a process interface to a control interface, one must also supply an action correspondence, implementable policy lifting, corresponding initial conditions, and transport of returns and constraints. Naturality of observed paths alone does not imply that every fine-level policy has a coarse representation or that optimal values coincide.

Chapter 20 gives the standard exact-quotient contract: feasible actions, rewards, and quotient transitions agree within fibers. If action maps are used, their surjectivity onto feasible actions must be checked state by state, and policy lifts must be measurable. An approximate interface should output \[
(\epsilon_r,\epsilon_p,\beta,\operatorname{span}\bar V^*,D)
\] and the domain of validity. An MMD interface must pass its kernel and function class; its numerical value cannot be inserted into a TV field.

\section{28.5 The coarsest stable partition algorithm for a known finite model}\label{the-coarsest-stable-partition-algorithm-for-a-known-finite-model}

Let \(X\) be a finite state set, \(A\) a finite action set common to all states, \(P_a\) known transition matrices, \(r(x,a)\) the reward, and \(o(x)\) the output label to preserve. The initial partition \(\mathcal P_0\) groups states with identical labels and identical rewards for all actions. At each iteration, split every current block by the signature \[
\operatorname{sig}_n(x)
=\left(\text{current block}(x),
\left(\sum_{y\in B}P_a(x,y)\right)_{a\in A,B\in\mathcal P_n}\right)
\] until the partition stops changing.

\hypertarget{p-alg-01}{}\phantomsection\label{p-alg-01}

\textbf{P-ALG-01: correctness and coarseness of finite stable partitioning.} The algorithm terminates in finitely many steps. Its terminal partition is the coarsest controlled-lumpable partition refining \(\mathcal P_0\). Its quotient therefore preserves the original outputs, rewards, and one-step distributions for all actions, and by induction the corresponding output laws at every finite horizon.

\textbf{Proof.} Iteration only splits blocks, never merges them. Every strict change increases the number of blocks, so there are at most \(|X|-|\mathcal P_0|\) strict changes. At termination, states in one block have the same transition probability to every quotient block under every action, establishing closure. Let \(\mathcal Q\) be any stable partition refining \(\mathcal P_0\), and assume inductively that it refines \(\mathcal P_n\). Each \(\mathcal P_n\) block is a union of \(\mathcal Q\) blocks; states in a common \(\mathcal Q\) block have equal transition probabilities to these unions. The next iteration therefore cannot split a \(\mathcal Q\) block. Thus \(\mathcal Q\) refines the terminal partition, proving coarseness. Finite-horizon preservation follows by kernel recursion. \(\blacksquare\)

\begin{verbatim}
partition := group_by(output, complete_reward_vector)
repeat:
    signature[x] := (block_id[x],
                     transition_mass_to_each_block_for_each_action)
    refined := split_each_block_by(signature)
    if refined == partition: return quotient(partition)
    partition := refined
\end{verbatim}

This is the coarsest stable quotient on the specified Markov state domain, not unconditionally the FFIPS of the full history space; sufficiency of the state for history remains to be proved. Floating-point tolerance does not have the transitivity of exact equality. An approximate algorithm must output defect and threshold certificates, rather than identifying rounded numbers as strictly equal.

\section{28.6 Composite certification and failure propagation}\label{composite-certification-and-failure-propagation}

When a premise in a proof-dependency graph fails, the affected downstream conclusion becomes unresolved or invalidated; its numerical error is not set to zero. For example:

\begin{itemize}
\tightlist
\item
  If deployment support exceeds sampling support, withdraw the uniform error certificate while retaining algebraic proofs within the model.
\item
  If hidden states invalidate quotient closure, finite one-step readout prediction may remain valid, while control and full-future identity require review.
\item
  If destruction changes the actual resource budget, the access blocker remains a combinatorial object, but the original interpretation of physical necessity fails.
\item
  If the objective changes, object identity may remain unchanged, but GOD and GOA must be recomputed for the new task.
\item
  If a mixing margin disappears, finite-horizon error bounds may remain valid, while long-term stationary-error bounds cannot be reused.
\end{itemize}

\section{28.7 Minimal interfaces for domain extensions}\label{minimal-interfaces-for-domain-extensions}

A domain module must provide variables and units, state or path mechanisms, experimental or intervention operations, readouts and noise, implementable controls, object candidates, domain conservation laws or constraints, and an error model. It may invoke core theorems. A proposed new general theorem must, in turn, remove domain-specific premises and supply a fresh proof and counterexample analysis.

Successful cross-domain unification first unifies structures and interfaces. It does not set every domain constant to one number or replace domain equations with shared terminology.

\section{28.8 The implemented finite-support decision interface}\label{the-implemented-finite-support-decision-interface}

The new finite decision procedure takes finite states and actions, an observation map, an initial PMF, time-indexed source-support relations, a certification predicate, and a horizon. Mathematical correctness directly concerns actual PMF support. An executable version may use finite relation tables, provided they are proved to coincide exactly with source support.

A True result means precisely that all nonempty initial observation supports are viable and hence that an actual policy with zero uncertified occupancy exists. False negates only that property in the same model, horizon, and certification conditions. Floating-point thresholds or absence of an edge from finite samples do not certify exact support. Chapter 37 gives the complete recursion, correctness proof, algorithm, and actual dead-end example.

\chapter{29. Application Method: From Models to Testable Certificates}\label{ch-29}

\textbf{Chapter question.} How does an actual research project move from a mechanism model to a falsifiable certificate?

\textbf{Reading route.} Each application declares its objects, inputs, operations, and budget before specifying a construction, baseline comparisons, and failure criteria. This chapter explains how to use the preceding theorems. Synthetic models, within-model proofs, and empirical support from real systems are reported separately.

\section{29.1 Redundant equipment and fault-tolerant maintenance}\label{redundant-equipment-and-fault-tolerant-maintenance}

\textbf{Inputs.} Component states \(X_t\), a load interface \(U_t\), maintenance actions \(a_t\), failure and repair transitions, and a service output \(f(X_t)\) to preserve. Declare whether the evaluation concerns service-predictive identity, a material individual, or a physical maintenance mechanism.

\textbf{Construction.} Form the predictive partition for service futures. Enumerate accessible readouts \(R_S\) and compute the defect \(e(S)\) of Chapter 13 to obtain minimal carriers of service information. Separately, use actual component-failure interventions to construct the physical mechanism-failure family. The two admit a combined interpretation only when the physical realization conditions hold.

\textbf{Decision.} A reward may be effective service minus maintenance cost, with hard safety constraints. Compute GOD for the quotient MDP and check the value error and action gap of Chapter 20. If the closed loop has a mixing margin, report a service-state GOA. With permanent failure, use survival or QSD analysis; the ordinary stationary law at an absorbing state is not a healthy steady state.

\textbf{Testable outputs.} Minimal readout families, physical cut sets, maintenance policies, service-interruption probabilities, and error bounds. Compare with the full-state model, single-component heuristics, and a no-maintenance baseline under the same budget.

\textbf{Failure criteria.} Load-induced transition drift, repair actions affecting unobserved hidden modes, or treating sensor removal as removal of a power-supply component each invalidate the corresponding interface.

\section{29.2 Steady-state maintenance, stochastic dynamics, and development}\label{steady-state-maintenance-stochastic-dynamics-and-development}

\textbf{Inputs.} A continuous-state diffusion \[
dX_t=b(X_t,u_t)dt+\sigma(X_t,u_t)dW_t,
\] nonexplosion conditions, noise units, a candidate persistence domain \(\mathcal V\), and a recovery set \(A\).

\textbf{Construction.} Seek a nonnegative potential satisfying \(\mathcal LV\le-aV+b\) to obtain a moment recovery bound. For return to \(A\), use a negative drift \(\mathcal LV\le-c\) outside the set, or solve \(\mathcal LE=-1\) under suitable boundary conditions. For exit questions, analyze the killed semigroup, conditional survival distribution, and mixing time; no return to one point along every trajectory is required.

\textbf{Developmental interface.} Include the structural stage \(s\) in the state and define stage-to-stage transport \(\Gamma_{s\to s'}\). Test predictive and interface fidelity after transport under subsequent environmental protocols. Geometric change alone does not prove interrupted identity, while uncalibrated transport does not prove continuation.

\textbf{Testable outputs.} Recovery speed, noise floor, mean return-time bounds, survival curves, stage-transport errors, and identity branches. Fix recovery sets and exit thresholds before testing.

\textbf{Failure criteria.} A potential decreasing only along finitely many training trajectories, unestablished nonexplosion, or calling the width of a noisy equilibrium an exact invariant point does not establish a persistence theorem.

\section{29.3 Multi-agent cooperation and parent-object formation}\label{multi-agent-cooperation-and-parent-object-formation}

\textbf{Inputs.} Child states and policies, a shared environmental interface, communication and resource coupling, and candidate parent futures. Fix the conditional environment before designing cuts, so that a common external driver does not create spurious irreducibility.

\textbf{Construction.} Chapter 15 supplies distances between complete parent-level path laws and cut laws, together with irreducibility margins. Chapter 24 evaluates angles and errors between child update fields and the parent-value gradient. When scanning coupling strength, measure both parent integrity and child fidelity to locate possible coupling windows.

\textbf{Decision.} A coordination mechanism must change actual update fields or incentives, not merely the objective written by an analyst. For multiple objectives, retain the Pareto set or prespecify a weight-selection rule.

\textbf{Testable outputs.} Parent-object certificates, valid cut families, child-loss profiles, coordination thresholds, and stable intervals across environments. Communication frequency or correlation alone is not a sufficient measure of parent-object strength.

\section{29.4 Memory compression and closed-loop decisions in learning systems}\label{memory-compression-and-closed-loop-decisions-in-learning-systems}

\textbf{Inputs.} Observation histories, action logs, prediction targets, deployment environment families, admissible memory states, and update rules.

\textbf{Construction.} Define required memory by predictive sufficiency rather than reconstruction of all history. Information bounds give lower limits on exact or approximate memory cost; Chapter 16 tests predictive classes in finite contexts. Verify action closure for the compressed state as well: good prediction does not guarantee correct control.

\textbf{Decision.} Construct a controller through belief updating or an explicitly learnable state. Propagate errors to policy regret through TV, span, or occupancy-weighted interfaces. Evaluate information acquisition by its actual effect on task decisions; producing novel labels is not itself improvement.

\textbf{Testable outputs.} Compression--prediction-loss--control-regret curves, deployment coverage certificates, failure detection, and rollback conditions. Retain an uncompressed model and simple history windows as baselines, using the same data and computational budgets.

\textbf{Safety boundary.} Tool permissions, resource budgets, irreversible actions, and shutdown mechanisms constrain feasible policies. A model may not delete them in the name of expanding choices. Mathematical capability expansion does not constitute external authorization.

\section{29.5 Evolutionary individuality and hereditary mechanisms}\label{evolutionary-individuality-and-hereditary-mechanisms}

\textbf{Inputs.} Parent--offspring correspondences, counting units, generation lengths, environmental interventions, and child- and parent-level readouts. Record zero-offspring families so reproductive intensity and conditional transmission can be estimated separately.

\textbf{Construction.} Certify objects and lineage before estimating \(M\), \(b\), and \(K\). Use states relevant to heritable futures and remove labels irrelevant to the current prediction target. Compute left and right Perron structures for mean growth; report extinction and variance separately for actual finite populations.

\textbf{Transition tests.} Apply matched interventions to external scaffolding, regulation of child-level conflict, and parent-level reproduction. Assess the four-component certificate of Chapter 26. A change in growth rate does not replace the full ETI evidence.

\textbf{Testable outputs.} Transmission identifiability, parent irreducibility, reproductive value, degree of mechanism internalization, and the conclusions that still depend on sampling or environmental assumptions.

\section{29.6 Mechanism identification and scientific model comparison}\label{mechanism-identification-and-scientific-model-comparison}

\textbf{Inputs.} Candidate model families, implementable experiments, observable variables, and a budget. First quotient by observational equivalence, then ask whether new experiments can separate the remaining classes.

\textbf{Construction.} In regular parametric models, inspect intersections of Fisher null spaces across protocols. In linear problems, inspect the accumulated Gram matrix. Local full rank is not global identifiability, and singular Fisher information does not invariably exclude higher-order identification.

\textbf{Experiment selection.} For a finite candidate class, choose protocols covering model pairs that remain unseparated. Continuous parameters require additional continuity, compactness, or local design theory. Protocols must be implementable, and separation strength must enter sample complexity.

\textbf{Testable outputs.} Identifiable quotients, inseparable mechanism classes, target model pairs for the next experiment, finite-sample intervals, and extrapolation limits. The empirical content of unification lies in these falsifiable interfaces, not in an ability to rename objects after every outcome.

\section{29.7 Delivery from a proved model to an experimental conclusion}\label{delivery-from-a-proved-model-to-an-experimental-conclusion}

An application report should specify the fixed source or model family, actual sampling units, available information, selection rule, order of deployment generation, common failure event, and error budget. A bare ``passed'' label does not reveal whether the conclusion concerns prediction, access, actual physical maintenance, or task risk.

The three-valued methods of C2/C5 allow certification, rejection, or an unresolved result. C7 distinguishes failures of protocol, model, interpretation, and core mathematics. The scientific-closure status still retains unverified real-world support. Research on actual systems should use falsifiable model contracts and competing explanations; correctness of a synthetic model is not a universal conclusion about nature.

\chapter{30. Complete Examples and Independent Executable Checks}\label{ch-30}

\textbf{Chapter question.} Can the abstract interfaces be computed step by step in finite models while exposing their limits?

\textbf{Reading route.} Compute the predictive quotient, redundant carriers, control value, and stationary law from one four-state kernel, then compare other mechanisms through continuous and reproductive examples. Numerical outputs retain their historical identity. Finite enumeration checks implementation consistency but replaces neither general proof nor experiments on natural systems.

Every example in this chapter is an explicitly specified mathematical model or synthetic experiment, not empirical evidence from a natural system. Appendix D embeds the complete code and output of the historical v3.0 run. Numerical checks examine implementation consistency and counterexamples; they do not replace the preceding general proofs.

\section{30.1 A finite model connecting prediction, carriers, control, and GOA}\label{a-finite-model-connecting-prediction-carriers-control-and-goa}

The microstate is \(x=(s,z)\in\{0,1\}^2\), and the action is \(a\in\{0,1\}\). At the next step, \(z'\) is an independent fair bit, with \[
P(s'=1\mid s,a)=
\begin{pmatrix}
0.15&0.85\\
0.90&0.85
\end{pmatrix}_{s,a},
\qquad r(s,z,a)=s-0.12a,\qquad\beta=0.9.
\] Thus the microscopic kernel is \[
P((s',z')\mid(s,z),a)=\tfrac12
p_{sa}^{s'}(1-p_{sa})^{1-s'}.
\]

The prediction target is the future of \(s\), and protocols allow both actions and their causal compositions. The variable \(z\) contributes nothing to this future, while the responses of \(s=0,1\) under action zero differ by \(0.75\), so the two classes cannot merge. Finite stable partitioning gives \(\{\{(0,0),(0,1)\},\{(1,0),(1,1)\}\}\), and \(f(s,z)=s\) is an exact control quotient.

With three readouts \(r_1=s,r_2=s,r_3=z\), \[
\mathcal M_0=\{\{r_1\},\{r_2\}\},\qquad
b(\mathcal M_0)=\{\{r_1,r_2\}\}.
\] A set containing either readout of \(s\) has zero defect; every other set has defect \(0.75\). Both redundant readouts must therefore be removed to destroy access sufficiency. This does not establish that removing two physical components is the sole way to destroy the equipment.

\section{30.2 Independent sampling and exact-family recovery}\label{independent-sampling-and-exact-family-recovery}

In each of the eight \((s,z,a)\) cells, independently generate \(N=20000\) Bernoulli responses, for 160000 synthetic samples in total, using random seed 30905 and \(\alpha=0.05\). The Bernoulli form of P-STAT-01 gives \[
\eta=\sqrt{\frac{\log(2\cdot8/0.05)}{40000}}
=0.0120086646.
\] The maximum empirical error recorded in v3.0 was \(0.00545<\eta\). With defect threshold \(\tau=0.2\), \[
2\eta<0.2<0.75-2\eta,
\] so the exact-structure gap condition holds and the correct carrier family and blocker are recovered.

The positive gap away from zero defect is proved from the specified model, not inferred as an exact equality from data alone. In an unknown continuous model, this confidence radius alone would not establish exact carrier recovery.

\section{30.3 Optimal value and certifiable actions}\label{optimal-value-and-certifiable-actions}

The true optimal policy and value are \[
\pi^*(0)=1,\quad\pi^*(1)=0,\qquad
V^*=(7.7717277487,\;8.9445026178).
\] For each \(s\), construct an estimated quotient using its \(z=0\) representative cell. This gives \[
\widehat V=(7.7698153658,\;8.9411000664),
\qquad \widehat\pi=(1,0).
\]

\begin{longtable}[]{@{}
  >{\raggedright\arraybackslash}p{(\columnwidth - 4\tabcolsep) * \real{0.3000}}
  >{\raggedleft\arraybackslash}p{(\columnwidth - 4\tabcolsep) * \real{0.4000}}
  >{\raggedright\arraybackslash}p{(\columnwidth - 4\tabcolsep) * \real{0.3000}}@{}}
\toprule\noalign{}
\begin{minipage}[b]{\linewidth}\raggedright
Quantity
\end{minipage} & \begin{minipage}[b]{\linewidth}\raggedleft
Historical v3.0 result
\end{minipage} & \begin{minipage}[b]{\linewidth}\raggedright
Interpretation
\end{minipage} \\
\midrule\noalign{}
\endhead
\bottomrule\noalign{}
\endlastfoot
Actual maximum value error & 0.00340255 & Directly computable only with the true model known \\
Span bound \(D\) using actual model defect & 0.05745151 & Deterministic conditional bound \\
Span bound \(D_{\rm conf}\) using the simultaneous confidence radius & 0.12659009 & Bound on a response event of probability at least 95\% \\
True regret of the lifted policy & 0 & The true policy happens to agree in this example \\
Estimated action gap at state zero & 0.62344368 & Exceeds \(2D_{\rm conf}\); certifiable \\
Estimated action gap at state one & 0.17797859 & Below \(2D_{\rm conf}\); unresolved by this certificate \\
\end{longtable}

Every 95\% statement is relative to the declared repeated-sampling experiment. Using the data-dependent span of \(\widehat V\) does not require another union bound: the response-error event is a uniform kernel bound, and on that event the span inequality holds for every corresponding bounded function.

\section{30.4 Long-term structure and finite paths}\label{long-term-structure-and-finite-paths}

The true closed-loop kernel of the lifted policy is \[
P^\pi=\begin{pmatrix}.15&.85\\.10&.90\end{pmatrix}.
\] Its Dobrushin coefficient is \(0.05\), and its invariant distribution is \[
\mu=(0.1052631579,\;0.8947368421).
\] The estimated closed-loop invariant distribution is \((0.1055867977,0.8944132023)\). Its TV error is approximately \(0.00032364\), below the confidence transport bound \[
\eta/(1-0.05)=0.01264070.
\] This long-term conclusion uses the mixing condition of the true closed loop. A generally unknown model requires separate certification of the margin.

Starting from \(s_0=0\), enumeration of eight-step paths gives TV distance \(0.00206291\) between true and estimated path laws, below the bound based on the uniform actual model defect \(0.00545\), namely \(1-(1-0.00545)^8=0.04277733\). A uniform defect over all actions is conservative but type-consistent and valid.

\section{30.5 Continuous recovery: the Ornstein--Uhlenbeck model}\label{continuous-recovery-the-ornsteinuhlenbeck-model}

Let \[
dX_t=-\kappa X_tdt+\sigma dW_t,\quad\kappa>0.
\] An integrating factor gives \[
X_t=e^{-\kappa t}X_0+
\sigma\int_0^te^{-\kappa(t-u)}dW_u.
\] For \(X_0=x_0\), Itô isometry yields \[
EX_t^2=e^{-2\kappa t}x_0^2+
\frac{\sigma^2}{2\kappa}(1-e^{-2\kappa t}).
\] This agrees with the recovery bound from \(\mathcal L(x^2)=-2\kappa x^2+\sigma^2\). With initial law \(N(0,\sigma^2/(2\kappa))\), the Gaussian convolution preserves that law, producing an invariant probability. Two solutions sharing the noise differ by \(e^{-\kappa t}(x_0-y_0)\), giving a coupling contraction.

There is strict recovery and a noise floor, but the state does not remain fixed at zero. Specifying exit from an interval requires additional killed-process analysis; an ordinary invariant distribution does not imply permanent confinement to any bounded interval.

\section{30.6 Conditional survival differs from conditioning on indefinite continuation}\label{conditional-survival-differs-from-conditioning-on-indefinite-continuation}

Take the killed CTMC generator \[
K=\begin{pmatrix}-2&1\\2&-4\end{pmatrix},
\] with additional absorption rates 1 and 2 in its two states. Its eigenvalues are \(-3\pm\sqrt3\), so \[
\lambda=3-\sqrt3=1.2679491924,
\quad q=(0.7320508076,0.2679491924).
\] Compute a positive right eigenvector \(h\) and normalize by \(qh=1\). Then \[
\varpi=q\odot h=(0.7886751346,0.2113248654).
\] Here \(q\) is the QSD conditional on being alive now, while \(\varpi\) is the invariant law of the Doob process conditioned for indefinite continuation. The code separately verifies \(qK=-\lambda q\), zero row sums for the transformed generator, and \(\varpi K^\uparrow=0\).

For Brownian motion absorbed at the endpoints of \((0,L)\) with generator \(\frac12\partial_{xx}\), direct solution of the Dirichlet eigenvalue equation gives \[
\phi_n(x)=\sin(n\pi x/L),\quad
\lambda_n=\frac{n^2\pi^2}{2L^2},
\] \[
q(x)=\frac{\pi}{2L}\sin(\pi x/L),\quad
\varpi(x)=\frac2L\sin^2(\pi x/L).
\] The spectral mixing gap is \(\lambda_2-\lambda_1=3\lambda_1\). This fixed ratio alone does not guarantee a wide common window for arbitrary error and survival requirements; initial-law coefficients and thresholds must still enter P-QSD-03.

\section{30.7 Three counterexamples to spurious unification}\label{three-counterexamples-to-spurious-unification}

\textbf{Omitted synergy.} For independent fair \(X,E\) and \(Y=X\oplus E\), \(I(Y;X)=I(Y;E)=0\), whereas \(I(Y;X,E)=\log2\). Separate tests of internal and environmental marginal residuals cannot replace joint closure.

\textbf{Unobservable control cost.} With two-dimensional Brownian noise, observe \(X=W^1\) and add unit drift only to the second noise coordinate. Over unit time the full noise-path KL is \(1/2\), while observed state-path KL is zero. Equality between state-path KL and total control energy therefore requires observability; the data-processing inequality remains valid.

\textbf{Two dimensions in each environment need not mean two shared dimensions.} In common budget coordinates \(B\in\mathbb R^3\), let the two environmental response Jacobians be \(\operatorname{diag}(1,1,0)\) and \(\operatorname{diag}(0,1,1)\). Each has rank two, but their vertical stack has rank three, so at that point they cannot both factor through the same smooth two-dimensional intermediate variable. A common Π/Φ dual-drive model must test this cross-environment consistency.

\section{30.8 Numerical examples of reproduction and coordination}\label{numerical-examples-of-reproduction-and-coordination}

Let \[
M=\begin{pmatrix}1.8&0.2\\0.4&0.8\end{pmatrix},\quad
p=(0.4,0.6),\quad z=(0,1).
\] Then \(b=(2,1.2)\), \(\bar b=1.52\), and \[
p'=(0.6315789474,0.3684210526).
\] The mean-trait change \(-0.2315789474\) decomposes into selection \(-0.1263157895\) and transmission \(-0.1052631579\). The principal growth factor is approximately \(1.8744562647\), not simply the largest row sum 2.

For hierarchical coordination, take parent gradient \(g=(1,0)\) and joint child field \(G=(-1,1)\). Then \(\eta_*=1/2\). Mixing above this threshold increases parent value instantaneously, without proving long-term optimality or improvement in every child value.

\section{30.9 Limits of the v3.0 numerical checks}\label{limits-of-the-v3.0-numerical-checks}

The original v3.0 text records 2371 passing assertions. These include blocker involution for all 20 antichains on three vertices and, for all six surjections from three atoms to two, blocker naturality, left and right adjunctions, exact realization intervals, and active-fiber uniqueness. Additional checks concern carrier recovery, exact quotients, approximate value bounds, action certification, stable partitioning, path TV, stationary perturbations, QSD/Doob distinctions, Price's identity, reward shaping, and counterexamples. Of the assertions, 2316 enumerate finite-resolution instances; they are not 2316 distinct general theorems.

These checks prove neither every conclusion for arbitrary large systems nor verification by a formal proof assistant such as Lean or Coq. General conclusions depend on the written proofs and explicit standard mathematical hypotheses. Synthetic data do not establish the universality of UEOT in nature.

\section{30.10 Formal finite witnesses added in the maintenance edition}\label{formal-finite-witnesses-added-in-the-maintenance-edition}

The v3.0 numerical record is retained in full. The maintenance edition adds actual Lean witnesses: the strictly weaker posterior-support model and the Bool model with a safe policy but a covered dead end in Chapter 37; and the experiment in Chapter 38 with 256 samples from each of two actual registered source rows, selecting between two distinct deployment laws using the first sample.

In the latter model, selection probabilities are \(7/8,1/8\), the exact actual joint failure probability is \(9/64\), and Good-only calibration gives the general guarantee \(17/32<1\). The rare all-true record has positive probability but fails the small-budget condition, showing that requiring every record to pass the numerical check is too strong. Exact values, the general bound, and Lean sources are listed separately in Chapter 38 and must not be combined with the old floating-point output as one experiment.

\chapter{31. Unified Assembly: What the Core Can Jointly Establish}\label{ch-31}

\textbf{Chapter question.} What can separately correct modules jointly guarantee under common conditions?

\textbf{Reading route.} Specify the shared event and each module's hypotheses before invoking predictive, carrier, control, path, and long-term results. A union bound combines failure probabilities. It does not create coverage of untested protocols, physical realization, or additional long-term stability conditions.

\section{31.1 An assembly theorem for the finite operational core}\label{an-assembly-theorem-for-the-finite-operational-core}

Combining modules into a testable conclusion serves the project's aim of unification more directly than listing similar mathematical tools separately.

\hypertarget{p-core-01}{}\phantomsection\label{p-core-01}

\textbf{P-CORE-01: finite object--realization--control certification assembly.} Suppose:

\begin{enumerate}
\def\labelenumi{\arabic{enumi}.}
\tightlist
\item
  A finite controlled process has correctly specified transitions, rewards, action semantics, and observation targets. In the true model, the candidate \(f\) satisfies controlled and reward closure, or has declared uniform defects.
\item
  A finite history--protocol response table has a fixed future target. Independent repeated sampling within each cell gives a simultaneous response-error event \(E_{\rm resp}\) of probability at least \(1-\alpha_{\rm resp}\).
\item
  The true predictive-class gap satisfies \(\gamma>4\eta\), and the positive access-carrier defect gap satisfies \(\Delta>4\eta\). Thresholds obey Chapter 16's conditions.
\item
  Readout fibers, the finite carrier domain, and access operations are fixed. Physical-destruction interpretations additionally satisfy Chapter 13's realization conditions.
\item
  A bounded discounted control problem and policy lifting are valid. On an event \(E_{\rm ctl}\), reward and transition errors are at most \(\epsilon_r,\epsilon_p\), respectively; the event has failure probability at most \(\alpha_{\rm ctl}\).
\item
  Persistent historical identity additionally has valid historical edges, transport, and object certificates. Long-term stationary-error claims additionally have appropriate mixing or recurrent-structure assumptions. Failure probabilities of these auxiliary certification events sum to at most \(\alpha_{\rm aux}\).
\end{enumerate}

Then, with probability at least \[
1-\alpha_{\rm resp}-\alpha_{\rm ctl}-\alpha_{\rm aux}
\] (with no repeated charge when the same event is reused), all the following hold:

\begin{itemize}
\tightlist
\item
  The predictive classes of the tested response problem, the exact minimal access-carrier family, and its blocker are correctly recovered.
\item
  Approximate quotient value error is at most \[
  D=\frac{\epsilon_r+\beta\epsilon_p\operatorname{sp}(\bar V^*)}{1-\beta},
  \] the lifted policy's regret is at most \(2D\), and an estimated best action with gap greater than \(2D\) is certifiable.
\item
  Under the same policy, the \(T\)-step path error is bounded by the stepwise coupling composition bound of Chapter 7.
\item
  Additional long-term hypotheses give the corresponding GOA stability bound; on a certified transport chain, historical fidelity errors accumulate as in Chapter 11.
\item
  If the estimated distance margin from a fixed failure set exceeds the object-representation error, that representation perturbation has not crossed into the failure set.
\end{itemize}

\textbf{Proof.} The union bound gives the common-event probability without independence between module events. On that event, P-STAT-05 recovers predictive classes and P-STAT-04 recovers the carrier family. The finite blocker is a deterministic function, so it is recovered as well. P-QUO-02 and P-QUO-03 give value, regret, and action certification. P-DYN-03 gives the finite-path bound for the same implemented policy. Long-term conclusions use their declared P-GOA conditions rather than an unsupported infinite-horizon limit of a finite bound. Historical transport follows the corresponding edgewise P-ID recursion. Finally, the distance to a fixed nonempty failure set is 1-Lipschitz, giving robustness. \(\blacksquare\)

The theorem does not identify every untested protocol from a finite response table or recover every physical mechanism from access readouts. It assembles existing certificates under common premises rather than adding an unproved universal criterion for object existence.

\section{31.2 Core closure and explicit open interfaces}\label{core-closure-and-explicit-open-interfaces}

At this point, objects, persistence, composition, purpose, learning, and evolution share a language of processes, protocols, factors, realization, and certification. The core's unification consists in the following:

\begin{itemize}
\tightlist
\item
  Objects are jointly certified by candidate futures, \nolinkurl{internal/environmental} interfaces, and historical mechanisms, rather than prespecified solely by names or material collections.
\item
  Explicit quotient conditions connect macroscopic compression to microscopic dynamics, finite statistics, and control value.
\item
  Operational semantics connect the combinatorial duality of carriers and destruction to physical realization, while distinguishing access from physical destruction.
\item
  Survival, recovery, conditional stability, and optimal long-term behavior have different but compatible mathematical types.
\item
  Heredity and learning can change future control structures, while purpose and normative bridges remain explicit.
\end{itemize}

The following remain research questions, not implicit theorems:

\begin{enumerate}
\def\labelenumi{\arabic{enumi}.}
\tightlist
\item
  Measurable existence conditions for predictive quotients and minimal realizations in infinite, continuous, noncompact environments.
\item
  Effective complexity of jointly discovering objects and protocols under restricted experimentation, strongly dependent sampling, and structural drift.
\item
  Provable algorithms recovering large carrier hypergraphs and physical mechanism cut sets without enumerating all subsets.
\item
  Transferable identity transport and GOA stability in nonstationary, multiscale, and near-bifurcation settings.
\item
  Independent experimental support for a common low-dimensional Π/Φ mechanism, and criteria for adding drive dimensions.
\item
  Identification of purpose bridges from specific reproduction, learning, and coordination mechanisms rather than arbitrary post hoc fitting of objectives to behavior.
\item
  Further machine proofs for general-process adaptations still open, and preregistered real-system tests matching the hypotheses of completed finite common-source connections.
\end{enumerate}

A sustainable route toward a unified theory of truth is for the common architecture to support increasingly explicit, provable, and falsifiable connections, not to write unestablished conclusions as settled axioms.

\section{31.3 Using this document}\label{using-this-document}

For a new task, locate the corresponding input contract in Chapter 28, then read the definitions, theorems, and counterexamples on which it depends. Preserve quantifiers, support, protocols, and error types when deriving a new result. Domain modules must not retroactively alter the premises of the general core. A genuine counterexample updates the relevant certificate and proof dependencies; it is not avoided by renaming the object.

This document is the main general mathematical version. Subsequent revisions should maintain it as a continuous text while preserving version and verification records. Readers need not locate an earlier edition to complete a central proof, and other domain branches are not prerequisites for using it.

\section{31.4 From conditional assembly to an actual common-source closure}\label{from-conditional-assembly-to-an-actual-common-source-closure}

Chapters 35--38 realize the assembly principles in the same finite source. Stable refinement preserves registered interfaces; the source fixed point supplies safe actions; constrained Bellman analysis compares all admissible causal policies; and a recovery program returns the same safe source kernel. Under partial observation, actual historical posterior support equals recursive support. Backward viability exactly characterizes existence of a policy with zero uncertified occupancy, and the finite relation procedure is proved to reconnect to actual support. The empirical model and selected policy on the statistical Good event then obtain an overall risk guarantee through the actual training--deployment joint law.

Each arrow has explicit assumptions and theorem sources. None removes reward provenance, program provenance, the physical realization framework, or calibration of the external conditional deployment law. The P0--P12/C1--C7 boundaries of Chapter 34 therefore coexist with this stronger positive chain. Progress in unification consists in removing unnecessary premises, deriving genuinely derivable conditions, and retaining indispensable ones, rather than renaming all premises as an object axiom.

Sections 21.7 and 35.7 further connect the exact observation quotient, actual policy, occupancy limit, and killed conditional structure of one source. They give a positive assembly while distinguishing its questions: one policy may have a unique stationary distribution while accumulating probability of previous failure; a safe constrained optimum may persist from a specified initial state while having multiple invariant laws over the entire state space. Sections 36.8 and 38.10 respectively test whether finite representations retain all escape mass and whether statistical selection still consumes the same actual sampling and deployment. These connections make MINCORE's quantitative obligations, H0's hierarchical distinctions, and specific theorems constrain one another through actual sources. They do not remove the conditions still required for purpose, formation, and empirical validation.

\chapter{32. Typed Unification: Eight Kinds of Hierarchy and Four Generating Mechanisms}\label{ch-32}

\textbf{Chapter question.} How can proof mechanisms be reused while retaining the types of different hierarchical relations?

\textbf{Reading route.} Eight relations first distinguish their objects, and four generators then extract reusable descent, error-chain, limit-residual, and positive-eigenstructure arguments. The decomposition of 106 source results describes reuse along registered routes; it is not an unconditional derivation of the entire theory from four axioms.

\section{32.1 The Mathematical Form of the Unification Problem}\label{the-mathematical-form-of-the-unification-problem}

The results of the preceding 31 chapters show that many problems share probability kernels, factors, and mechanisms for transporting error. Unification nevertheless requires preserving the types of the data. For example, a predictive quotient \(q:X\to Y\) is a state map; carrier coarse-graining \(\pi:V\to W\) acts on physical or access atoms; a genealogical relation acts on object instances with identities; and a value ordering acts on policies. Such maps may all be described as moving to a higher level, but they cannot therefore be composed directly in a single formula.

This chapter organizes them through typed relations. An interface must specify at least its domain, codomain, preserved structure, and comparison criterion. Composition is meaningful only when the intermediate objects agree or a proved adapter connects them. This also explains why formal review must check source identity in addition to types: two kernels of the same type may describe different processes.

\section{32.2 Eight Kinds of Hierarchical Relation}\label{eight-kinds-of-hierarchical-relation}

The eight H0 relations are listed below. Their distinctions concern mathematical content; they do not divide the theory into eight unrelated disciplines.

\begin{longtable}[]{@{}
  >{\raggedright\arraybackslash}p{(\columnwidth - 4\tabcolsep) * \real{0.3333}}
  >{\raggedright\arraybackslash}p{(\columnwidth - 4\tabcolsep) * \real{0.3333}}
  >{\raggedright\arraybackslash}p{(\columnwidth - 4\tabcolsep) * \real{0.3333}}@{}}
\toprule\noalign{}
\begin{minipage}[b]{\linewidth}\raggedright
Type
\end{minipage} & \begin{minipage}[b]{\linewidth}\raggedright
Principal mathematical objects
\end{minipage} & \begin{minipage}[b]{\linewidth}\raggedright
Typical conclusions available
\end{minipage} \\
\midrule\noalign{}
\endhead
\bottomrule\noalign{}
\endlastfoot
PQ: prediction and quotients & Response families, fiber compatibility, measurable statistics, controlled kernels & Minimal sufficient representations, kernel descent, control quotients \\
DYN: dynamics across scales & A common source kernel, commuting diagrams, semigroups, one-step errors & Cross-scale consistency, finite-path error bounds \\
RES: resolution and realization & Closure systems, atom maps, upward-closed families, adjunctions & Coarse-graining of minimal realizations, intervals of realization uncertainty \\
COMP: composition and parent objects & Joint child-path laws, cuts, parent predictive kernels, covers & Conditional irreducibility, additional parent-level information, composition windows \\
HIST: historical continuation & Directed instances, kernel transport, developmental models & Accumulation of frozen-kernel defects and developmental errors \\
VAL: value and evaluation & Preferences, differential forms, policy gradients, weights & Integrability, value preservation, coordination thresholds \\
SEL: selection and reproduction & Parent--offspring relations, mean reproduction operators, transmission kernels & Price decompositions, Perron compositions, selection indistinguishability \\
GOA: long-term structures & Invariant laws, QSDs, recurrent classes, spectral gaps, occupation limits & Conditional existence, uniqueness, mixing, and stability \\
\end{longtable}

For example, one exact controlled quotient can connect PQ and DYN. If rewards and actions are also preserved, it can connect VAL; if actual carrier readouts are supplied as well, it can connect RES. Without these additional data, the definition of a quotient alone implies neither parent formation nor the existence of a long-term invariant law.

Representative Lean sources include \nolinkurl{QuotientDescent}, \nolinkurl{DynamicsKernel}, PathError, \nolinkurl{ClosureResolution}, \nolinkurl{CompositionPathSource}, \nolinkurl{TransportDefect}, \nolinkurl{AlignmentGlobalExactness}, \nolinkurl{LineageEvolution}, and \nolinkurl{FiniteCesaroInvariant.} \nolinkurl{H0_HIERARCHY_INVENTORY} is an inventory of these relations. It does not define a single operator that replaces all eight kinds.

\section{32.3 M-QD: Compatibility and Unique Descent}\label{m-qd-compatibility-and-unique-descent}

Let \(q:X\twoheadrightarrow Y\) be surjective and let \(g:X\to Z\). Define fiber compatibility by

\[q(x)=q(x')\Longrightarrow g(x)=g(x').\]

\hypertarget{m-32-1}{}\phantomsection\label{m-32-1}

\textbf{M-32.1: Unique descent at the set level.} Fiber compatibility holds if and only if there is a unique \(\bar g:Y\to Z\) such that \(g=\bar g\circ q\).

\textbf{Proof.} Under compatibility, choose for each \(y\) an \(x_y\) with \(q(x_y)=y\), and set \(\bar g(y)=g(x_y)\). Compatibility makes the value independent of the choice and gives \(\bar g(q(x))=g(x)\). If \(h\circ q=g\), substitution of \(x_y\) yields \(h(y)=g(x_y)=\bar g(y)\). Conversely, factorization through the same \(\bar g\) immediately gives fiber compatibility. \(\blacksquare\)

\textbf{Lean.} \nolinkurl{UEOT.V3.Compression.QuotientDescent.fiberCompatible_iff_existsUnique_descend.}

This theorem is a set-theoretic statement. If \(X,Y,Z\) carry measurable structures, measurability of \(\bar g\) still has to be proved. Chapter 4 uses \nolinkurl{probability-measure-valued} coordinates and a countable protocol family to obtain measurable factorization on one common probability-one set; \nolinkurl{countableAEFamily_descend} reuses that mechanism. The singleton-null-set counterexample in Section 4.3 continues to apply to unrestricted uncountable versions.

The same mechanism applies to two-sided separated representations. If \(K:X\times E\to Z\) is invariant under both an internal equivalence relation and an environmental equivalence relation, first replacing \(x\) and then replacing \(e\) produces a well-defined \(\bar K\). Every other exact separated representation must refine response equivalence. The Lean declarations \nolinkurl{existsUnique_twoSidedDescend} and \nolinkurl{exactSeparatedFactorization_refines_responseEq} implement this argument.

In the counted compression, M-QD-01 generates P-PRED-01 and P-INT-02. The remaining structures needed for measurable Markov kernels, control-value preservation, algorithm termination, and conditional disintegration remain in their respective retained adapters.

\section{32.4 M-TC: Reusing Error Bounds Along One Transport Chain}\label{m-tc-reusing-error-bounds-along-one-transport-chain}

Suppose that the comparison quantities \(e_n\ge0\), transport constants \(L_n\ge0\), and local defects \(\epsilon_n\ge0\) satisfy

\[e_{n+1}\le L_ne_n+\epsilon_n.\]

\hypertarget{m-32-2}{}\phantomsection\label{m-32-2}

\textbf{M-32.2: Weighted error chain.} For every \(N\),

\[e_N\le \left(\prod_{j<N}L_j\right)e_0+
\sum_{k<N}\epsilon_k\prod_{k<j<N}L_j.\]

An empty product is one.

\textbf{Proof.} The case \(N=0\) is immediate. Multiply the bound at \(N\) by the nonnegative number \(L_N\), add \(\epsilon_N\), and write this last term as the \(k=N\) term with an empty product. This gives the bound at \(N+1\). \(\blacksquare\)

\textbf{Lean.} \nolinkurl{UEOT.V3.Compression.TransportCertificate.weighted_chain_bound}; the developmental source wrapper is \nolinkurl{development_pipeline_via_weighted_chain.}

Two-stage triangle bounds, composition of exact commuting diagrams, and chains with \(L_n=1\) are special cases of this mechanism. Probability paths also admit a multiplicative version: if each stage preserves a fraction at least \(1-\epsilon_n\) of the common mass that has not diverged, then the common mass over the entire path is at least \(\prod_n(1-\epsilon_n)\). This gives the path-TV bound in P-DYN-03. The actual PMF extensions and common-mass calculation in PathError prove that these fractions arise from the actual kernels, rather than merely satisfying an abstract recurrence.

The five counted endpoints of M-TC-01 are P-API-01, P-ID-01, P-DYN-04, P-DYN-03, and P-ID-02. They share a calculation mechanism; each source retains its probability, measurability, reward, and local-defect conditions.

\section{32.5 M-OI: Residuals That Vanish in the Limit}\label{m-oi-residuals-that-vanish-in-the-limit}

Suppose that \(\mu_k\) converges weakly to \(\mu\), that \(P\) preserves the class of bounded continuous test functions in use, and that for every such \(f\),

\[\mu_k(Pf)-\mu_k(f)\longrightarrow0.\]

\hypertarget{m-32-3}{}\phantomsection\label{m-32-3}

\textbf{M-32.3: Observable residuals imply invariance.} If the test class determines probability measures, then \(\mu P=\mu\).

\textbf{Proof.} Weak convergence applied separately to \(f\) and \(Pf\) gives \(\mu(Pf)-\mu(f)=0\). Since the test class determines measures, the two measures agree. \(\blacksquare\)

\textbf{Lean.} \nolinkurl{UEOT.V3.Compression.OccupationLimitInvariance.invariant_of_continuous_observable_residual.}

For finite Cesàro occupation measures, the residual is \((\mu_0P^N-\mu_0)/N\); for continuous Feller occupation measures, it is controlled by two temporal boundary terms of length \(s\). These yield P-GOA-01 and P-PER-02, respectively. Tightness, weakly convergent subsequences, and preservation of closed support each require their own proofs. The argument does not establish existence of a stationary law without a limit or the requisite compactness.

\section{32.6 M-PE: Normalized Positive Eigenstructure}\label{m-pe-normalized-positive-eigenstructure}

Let \(M\) be a finite nonnegative matrix, let \(\rho>0\), and let \(q\) and \(h\) be componentwise positive row and column vectors satisfying

\[qM=\rho q,\qquad Mh=\rho h,\qquad q\mathbf1=1,\quad qh=1.\]

\hypertarget{m-32-4}{}\phantomsection\label{m-32-4}

\textbf{M-32.4: QSD/Doob structures and stochastic readouts.} Set

\[M^h_{ij}=\frac{M_{ij}h_j}{\rho h_i},\qquad \varpi_i=q_ih_i.\]

Then \(M^h\) is stochastic and \(\varpi\) is its invariant probability vector. If \(M\) is substochastic, \(qM^n=\rho^nq\) gives the corresponding QSD relation. If a componentwise nonnegative, random row-vector process \(Z_n\) is adapted and integrable and satisfies

\[E[Z_{n+1}\mid\mathcal F_n]=Z_nM,\]

then \(W_n=\rho^{-n}Z_nh\) is a nonnegative integrable martingale.

\textbf{Proof.} The row sum is \((Mh)_i/(\rho h_i)=1\), and stationarity follows from

\[\sum_iq_ih_iM^h_{ij}=\frac{h_j}{\rho}(qM)_j=q_jh_j.\]

The total mass is \(qh=1\). Repeated application of the left eigenvector equation gives the QSD relation. Linearity of conditional expectation and the right eigenvector equation give the martingale relation. \(\blacksquare\)

\textbf{Lean.} \nolinkurl{UEOT.V3.Compression.PositiveEigenstructure.qsdDoob_bundle} and \nolinkurl{stochasticRightEigen_bundle}; the counted endpoints are P-QSD-02 and P-EVO-04.

This proposition consumes eigenstructure data. It does not prove that an arbitrary \(M\) admits those data. Perron existence, uniqueness, strict spectral dominance, and the limit of normalized powers remain the responsibility of the applicable K-PF-01 theorem or KPF01Certificate. Positive eigenstructure must also be distinguished from the deeper moment conditions required for stochastic branching processes.

\section{32.7 The Actual Compression Decomposition of the 106 Source Results}\label{the-actual-compression-decomposition-of-the-106-source-results}

The frozen compression ledger records

\[106=11\text{ generated endpoints}+89\text{ domain adapters}+6\text{ retained boundaries}.\]

The 11 generated endpoints are P-PRED-01 and P-INT-02; P-API-01, P-ID-01, P-DYN-04, P-DYN-03, and P-ID-02; P-PER-02 and P-GOA-01; and P-QSD-02 and P-EVO-04.

The six retained boundaries are P-INFO-05, P-INV-01, P-INV-02, P-DDH-01, P-EVO-02, and P-BRG-02. They limit unconditional extensions of compression, identification, gauge claims, or selection. These boundaries are part of the unified theory; they are not unresolved remainders to be removed.

Minimality of the minimal core is relative to the registered theorem surface and permitted reuse routes: after removal of a generator, a specified witness can no longer be derived by the remaining registered routes. The ablation rules prohibit constructing a replacement generator ad hoc from lower-level libraries during the test. This is architectural nonredundancy, not absolute logical independence or a claim that four axioms uniquely generate all 106 propositions.

The earlier count of 633 modules and the later count of 650 are counts of source snapshots. Module numbers do not strengthen a theorem logically. A canonical reader-facing entry point can be simplified, but removing retained adapters cannot legitimately remove their hypotheses.

\chapter{33. Relational Defects, Quantitative Calibration, and Gauge Covariance}\label{ch-33}

\textbf{Chapter question.} Why can two representations with the same zero-discrepancy relation still not share numerical budgets?

\textbf{Reading route.} A positive-rescaling counterexample first shows quantitative nonuniqueness. Two-sided calibration then transports zero relations, errors, and limits. Exact scaling supports a stronger covariance of error chains. The complete gauge of one numerical objective still differs from equality of ordering and does not establish a two-dimensional bottleneck in natural responses.

\section{33.1 Zero Relations and Quantitative Defects}\label{zero-relations-and-quantitative-defects}

Given a set \(X\) with a relation \(R\), a defect presentation is a nonnegative function \(d:X\times X\to\mathbb R_{\ge0}\) satisfying

\[d(x,y)=0\iff R(x,y).\]

Only the zero set is prescribed. Symmetry, the triangle inequality, and dynamical contraction require separate proofs. \nolinkurl{RelationDefectPresentation} is precisely this minimal interface.

\hypertarget{m-33-1}{}\phantomsection\label{m-33-1}

\textbf{M-33.1: A relation does not uniquely determine a quantitative defect.} If \(d(x_0,y_0)>0\), then for every \(c>0\), the defect \(d_c=cd\) has the same zero relation, whereas, whenever \(c\ne1\), \(d_c(x_0,y_0)\ne d(x_0,y_0)\). Consequently, no rule that reads only the relation can recover the original defect values for every such presentation.

\textbf{Proof.} Multiplication by a positive number preserves and reflects zero. If a rule recovered both \(d\) and \(2d\), the common relation would require the same output, but at \((x_0,y_0)\) the required values would differ. This is a contradiction. \(\blacksquare\)

\textbf{Lean.} \nolinkurl{MinCore.RelationDefectPresentation.rescale_preserves_exact_relation} and \nolinkurl{rescale_changes_nonzero_defect}; \nolinkurl{MinCore.no_relation_only_quantitative_recovery.}

For another example, take three points \(0,1,2\), put zero on the diagonal, and set \(d(0,1)=d(1,2)=1\) and \(d(0,2)=3\), symmetrically. The zero relation is still equality, but the triangle inequality fails. This is \nolinkurl{relationDefectPresentation_does_not_force_triangle.} Calling a quantitative index a structural distance does not supply the missing proof.

\section{33.2 Two-Sided Calibration}\label{two-sided-calibration}

Let \(d_P,d_Q\) be defects on the same pairs of inputs. For \(a>0,b\ge0\), the calibration condition is

\[a d_P(x,y)\le d_Q(x,y)\le b d_P(x,y)\quad\forall x,y.\]

\hypertarget{m-33-2}{}\phantomsection\label{m-33-2}

\textbf{M-33.2: Transport of zero relations, budgets, and convergence to zero.} Under the calibration condition:

\begin{enumerate}
\def\labelenumi{\arabic{enumi}.}
\tightlist
\item
  \(d_P(x,y)=0\iff d_Q(x,y)=0\);
\item
  \(d_P\le\epsilon\) implies \(d_Q\le b\epsilon\);
\item
  \(d_Q\le\delta\) implies \(d_P\le\delta/a\);
\item
  for every sequence, or the corresponding filter limit, \(d_P\to0\iff d_Q\to0\).
\end{enumerate}

\textbf{Proof.} The first claim follows by nonnegative squeezing and \(a>0\). The second and third use the upper bound and division by a positive number, respectively. For the fourth, apply squeezing to \(0\le d_Q\le bd_P\) and \(0\le d_P\le d_Q/a\). \(\blacksquare\)

\textbf{Lean.} \nolinkurl{MinCore.DefectCalibration.zero_defect_iff}, rel\_iff, \nolinkurl{forward_budget}, \nolinkurl{reverse_budget}, and \nolinkurl{tendsto_zero_iff.}

Equality of the zero relations need not be stored as an additional calibration field: it follows from the inequalities. Conversely, equality of zero relations alone supplies no calibration constants. For instance, the defects \(d\) and \(cd\) allow the scale factor \(c\) to be arbitrarily large.

\section{33.3 Quasi-Triangle Inequalities and Cross-Stage Contraction}\label{quasi-triangle-inequalities-and-cross-stage-contraction}

\hypertarget{m-33-3}{}\phantomsection\label{m-33-3}

\textbf{M-33.3: The triangle constant after calibration.} If \(d_P\) satisfies the triangle inequality, then

\[d_Q(x,z)\le\frac ba\bigl(d_Q(x,y)+d_Q(y,z)\bigr).\]

\textbf{Proof.} Apply the target upper calibration, the source triangle inequality, and the source lower calibration twice:

\[d_Q(x,z)\le b d_P(x,z)
\le b[d_P(x,y)+d_P(y,z)]
\le\frac ba[d_Q(x,y)+d_Q(y,z)].\quad\square\]

This does not assert that \(b/a\) is always optimal, or that the target retains a triangle inequality with the original constant one.

For a map \(F:X\to Y\), suppose that the source metrics satisfy

\[d_{P,Y}(Fx,Fx')\le Ld_{P,X}(x,x'),\qquad L\ge0,\]

and that the input lower calibration constant is \(a_X>0\) while the output upper calibration constant is \(b_Y\ge0\). The same calculation yields

\[d_{Q,Y}(Fx,Fx')\le\frac{b_YL}{a_X}d_{Q,X}(x,x').\]

\textbf{Lean.} \nolinkurl{MinCore.DefectCalibration.triangle_quasi} and \nolinkurl{MinCore.calibrated_cross_stage_contraction.}

\section{33.4 Covariance of Error Chains Under Exact Rescaling}\label{covariance-of-error-chains-under-exact-rescaling}

General two-sided calibration gives bounds. If, instead, each stage has an exact scale change \(d'_n=c_nd_n\), with \(c_n>0\), its coefficients can be transported exactly.

\hypertarget{m-33-4}{}\phantomsection\label{m-33-4}

\textbf{M-33.4: Stagewise scales and error chains.} From \(e_{n+1}\le L_ne_n+\epsilon_n\) and \(e'_n=c_ne_n\), one obtains

\[e'_{n+1}\le L'_ne'_n+\epsilon'_n,
\quad L'_n=\frac{c_{n+1}}{c_n}L_n,
\quad\epsilon'_n=c_{n+1}\epsilon_n.\]

The expansion in M-32.2 becomes the original expansion multiplied throughout by \(c_N\).

\textbf{Proof.} Multiply the original one-step bound by \(c_{n+1}\) and substitute \(e_n=e'_n/c_n\). Each product of ratios \(c_{j+1}/c_j\) telescopes, leaving the correct terminal-stage scale in both the initial term and every error term. \(\blacksquare\)

\textbf{Lean.} \nolinkurl{MinCore.rescale_contraction_transformed}, \nolinkurl{rescale_local_budget_transformed}, and \nolinkurl{rescaled_weighted_chain_bound.} Exact positive rescaling at one stage also preserves the triangle constant; see \nolinkurl{rescale_triangle_preserved.}

Changing units while retaining the old error thresholds and Lipschitz constants compares different quantities. The covariance theorem resolves this mathematical issue. Whether a particular \(c_n\) arises from actual physical renormalization requires a separate model-compatibility relation.

\section{33.5 The Complete Gauge of a Fixed Dual-Drive Objective}\label{the-complete-gauge-of-a-fixed-dual-drive-objective}

On a common policy set \(\mathfrak P\), let \(\Pi,\Phi\) be real-valued functions. Fix \(\lambda\in\mathbb R\) and put

\[J(\pi)=\Pi(\pi)-\lambda\Phi(\pi).\]

\hypertarget{m-33-5}{}\phantomsection\label{m-33-5}

\textbf{M-33.5: Complete gauge characterization.} For another pair \(\Pi',\Phi'\),

\[J'=J\iff\exists\chi:\mathfrak P\to\mathbb R,
\quad\Pi'=\Pi+\lambda\chi,\quad\Phi'=\Phi+\chi.\]

\textbf{Proof.} Substitution proves the implication from right to left by cancellation. Conversely, take \(\chi=\Phi'-\Phi\); rearranging \(J'=J\) gives \(\Pi'-\Pi=\lambda\chi\). \(\blacksquare\)

\textbf{Lean.} \nolinkurl{UEOT.V3.Compression.TheoryCompletion.UnifiedClosure.same_objective_iff_gauge_related.}

If \(\Phi'=\Phi\), equality of objectives immediately gives \(\Pi'=\Pi\). If \(\Pi'=\Pi\) and \(\lambda\ne0\), it gives \(\Phi'=\Phi\). When \(\lambda=0\), \(\Phi\) does not enter the objective and cannot be identified from it. The corresponding declarations are \nolinkurl{pi_identified_of_phi_anchor}, \nolinkurl{phi_identified_of_pi_anchor}, and \nolinkurl{unanchored_drive_is_not_uniquely_identified.}

Equality of numerical objectives is stronger than equality of policy orderings. The values \((0,1,2)\) and \((0,1,4)\) on three policies induce the same ordering but are not numerically equal. They therefore cannot be related by an algebraic gauge that preserves \(J\) exactly. Teleological equivalence in \nolinkurl{TeleologicalGauge} must retain this distinction from an algebraic orbit.

\section{33.6 From Formal Calibration to Identifiable Mechanisms}\label{from-formal-calibration-to-identifiable-mechanisms}

The complete gauge theorem determines neither the physical dimensions of \(\Pi,\Phi\) nor whether responses across environments require only two dimensions. Resource measurements, information tasks, intervention axes, a common budget, and the error model remain mechanism inputs. The rank necessity condition in Chapter 23 is a separate falsifiable proposition: a common-response rank demonstrably greater than two at the same budget point rules out the stated common two-dimensional factorization. Compatibility with rank two does not uniquely identify two latent variables.

The chapter thus provides two complementary results: it identifies quantities that cannot be recovered without calibration, and establishes which budgets, limits, and gauges can be transported under explicit calibration. Both constrain the quantitative expression of a unified theory.

\chapter{34. Complete Interfaces for Operational Objects, Purpose, and Conditional Lifecycles}\label{ch-34}

\textbf{Chapter question.} What source obligations are carried by the terms object, purpose, and lifecycle?

\textbf{Reading route.} Six object roles and the responsibilities of P0--P12 provide a reading map, followed by a checkable formation--binding--transport proof. The program nonidentifiability counterexample identifies inputs that cannot be deleted. Completion of a work package is distinct from empirical support for a real mechanism.

\section{34.1 Six Roles of an Object}\label{six-roles-of-an-object}

Given a candidate domain, experimental protocol, and process, at least six judgments must be distinguished.

\begin{longtable}[]{@{}
  >{\raggedright\arraybackslash}p{(\columnwidth - 4\tabcolsep) * \real{0.3333}}
  >{\raggedright\arraybackslash}p{(\columnwidth - 4\tabcolsep) * \real{0.3333}}
  >{\raggedright\arraybackslash}p{(\columnwidth - 4\tabcolsep) * \real{0.3333}}@{}}
\toprule\noalign{}
\begin{minipage}[b]{\linewidth}\raggedright
Role
\end{minipage} & \begin{minipage}[b]{\linewidth}\raggedright
Mathematical question
\end{minipage} & \begin{minipage}[b]{\linewidth}\raggedright
Evidential obligation
\end{minipage} \\
\midrule\noalign{}
\endhead
\bottomrule\noalign{}
\endlastfoot
Predictive identity & Does the representation preserve the specified future responses? & Fiber compatibility and measurable factorization \\
Formation provenance & Does the candidate arise from the given formation process? & Membership in the actual formation relation or generated image \\
Interaction-based discrimination & Do registered probes reveal its differences from other candidates? & Pairwise distinguishability or a corresponding bound \\
Constitutive persistence & Does the dynamics preserve the constitutive conditions? & The actual source and a persistence certificate \\
Recoverable realization & Is there an effective repair path or program after perturbation? & A recovery policy, carrier, and trusted realization \\
Semantic identity & Which instances are interpreted as the same object? & A specified identity relation and preservation bridge \\
\end{longtable}

\nolinkurl{PredictiveIdentity}, \nolinkurl{FormedProvenance}, \nolinkurl{InteractionDelimitedAt}, \nolinkurl{ConstitutivelyPersists}, \nolinkurl{PhysicallyRecoverableOn}, and \nolinkurl{SemanticFiberIdentity} in \nolinkurl{OperationalObject} represent these roles. They do not share a universal field that automatically implies the other five.

Statistical identification and physical formation are especially distinct. An algorithm may identify the class of an existing parent object with high probability without constructing that object in the world. Conversely, a generative mechanism may reliably produce a structure without making it uniquely identifiable under the current observation protocol.

\section{34.2 Necessary Conditions for Purpose Contracts and a Mechanism Bridge}\label{necessary-conditions-for-purpose-contracts-and-a-mechanism-bridge}

Let a policy set \(\mathfrak P\) carry a comparison relation \(\preceq\). A complete real-valued representation \(J\) would satisfy

\[\pi\preceq\pi'\iff J(\pi)\le J(\pi'),\]

The totality and transitivity of the real order are therefore necessary conditions on \(\preceq\).

\hypertarget{m-34-1}{}\phantomsection\label{m-34-1}

\textbf{M-34.1: An obstruction to numerical representation of an incomparable contract.} If two distinct policies are each comparable only with themselves, no such complete real-valued representation exists.

\textbf{Proof.} Of any two real numbers, one is at most the other. A complete representation would therefore force at least one cross-comparison to hold, contrary to the contract. \(\blacksquare\)

\textbf{Lean.} \nolinkurl{TheoryCompletion.incomparableBoolContract_has_no_numerical_representation}, together with the corresponding theorem excluding a valid dual-drive representative.

Now let \(M\) be a set of mechanisms, \(o:M\to O\) an observation map, and \(J_m\) the mechanism-indexed objective. Mechanism--purpose faithfulness requires

\[o(m)=o(m')\Longrightarrow J_m,J_{m'}\text{ induce the same order on the same policy domain}.\]

Under this hypothesis, mechanisms with the same observation have the same teleological equivalence class. The proof applies faithfulness and passes to equivalence classes; the exact endpoint is \nolinkurl{ScientificClosure.mechanismObservation_identifies_teleologicalClass.}

For a counterexample showing the need for this bridge, take a constant observation \(o(m)=*\) and two mechanisms that assign opposite objectives to Bool policies. Their observations agree while their purposes oppose one another. Thus \nolinkurl{equalMechanismObservation_does_not_determine_purpose} rules out identification without faithfulness. This counterexample requires a purpose interpretation to have independent content; it does not deny that a specific mechanism can support one.

\section{34.3 The Responsibilities of P0--P12}\label{the-responsibilities-of-p0p12}

\begin{longtable}[]{@{}
  >{\raggedright\arraybackslash}p{(\columnwidth - 4\tabcolsep) * \real{0.3333}}
  >{\raggedright\arraybackslash}p{(\columnwidth - 4\tabcolsep) * \real{0.3333}}
  >{\raggedright\arraybackslash}p{(\columnwidth - 4\tabcolsep) * \real{0.3333}}@{}}
\toprule\noalign{}
\begin{minipage}[b]{\linewidth}\raggedright
Package
\end{minipage} & \begin{minipage}[b]{\linewidth}\raggedright
Main questions organized
\end{minipage} & \begin{minipage}[b]{\linewidth}\raggedright
Boundaries that must be retained
\end{minipage} \\
\midrule\noalign{}
\endhead
\bottomrule\noalign{}
\endlastfoot
P0 & Semantic distinctions among operational objecthood, purpose, GOD, GOA, and scale & Definitions and synonyms are not external facts \\
P1 & Purpose contracts, numerical representatives, gauges, and orderings & Representation existence, faithfulness, and gauge identification are distinct \\
P2 & Agency, control, GOD, and GOA on the same object & Actual controlled sources, purpose realization, and viability hypotheses \\
P3/P4 & Identification from finite data and of changing object classes & A common Good event, separation margins, protocols, and sample sources \\
P5 & Recurrence and mean homeostasis under organizational faults & Repair drift, fault envelopes, and the same mixture kernel \\
P6 & Resource accumulation and maintenance budgets & Initial reserves, supply, costs, and the expectation semantics used \\
P7 & Assembly of a repair organization from a seed & Program provenance, trusted interpretation, \nolinkurl{coding/decoding}, and the physical source \\
P8 & Lineage, reproduction, and Price/Perron bridges & The same operator, actual offspring edges, and spectral conditions \\
P9 & Composition of scales and transport of properties & Separate preservation of six properties; no physical RG compatibility law is defined \\
P10 & General observation, kernel, and belief interfaces & The relevant measurable and topological regularity \\
P11 & Required stochastic-analysis interfaces & Process obligations for ItoRun, Dynkin, and Girsanov \\
P12 & Connection of supplied bridges into a conditional lifecycle & PARTIAL / EXPLICIT BOUNDARY; it does not synthesize all sources \\
\end{longtable}

CLOSED means that a package has completed its agreed deliverables, which may include conditional propositions and impossibility results. It does not mean that the obligations in the right-hand column have been automatically discharged.

\section{34.4 Quantitative Composition of Formation, Binding, and Transport}\label{quantitative-composition-of-formation-binding-and-transport}

Let \(X_0,X_1\) be underlying input domains, \(F_0,F_1\) formation-output domains, and \(P_0,P_1\) parent-realization domains. Suppose that

\[f_i:X_i\to F_i,\quad B_i:F_i\to P_i,\]

with transport maps \(t_X:X_0\to X_1\), \(t_F:F_0\to F_1\), and \(t_P:P_0\to P_1\). Equip \(F_1,P_1\) with the respective distances and assume

\[d_F(f_1(t_Xx),t_F(f_0x))\le\epsilon_F,\]

\[d_P(B_1(t_Fu),t_P(B_0u))\le\epsilon_B,\]

\[d_P(B_1u,B_1v)\le Ld_F(u,v),\qquad L\ge0.\]

\hypertarget{m-34-2}{}\phantomsection\label{m-34-2}

\textbf{M-34.2: The FBT bound for realized structures.}

\[d_P(B_1(f_1(t_Xx)),t_P(B_0(f_0x)))\le L\epsilon_F+\epsilon_B.\]

If both endpoint outputs also satisfy the actual formation relation and the right-hand side is within a prespecified structural tolerance, one obtains the corresponding continuation certificate for the formed structures.

\textbf{Proof.} Insert \(B_1(t_F(f_0x))\) between the endpoints. Bound the first segment by the Lipschitz property of \(B_1\) and the formation error, and the second by the binding-transport error. The formation relations and the tolerance remain separate assumptions; the triangle inequality does not create them. \(\blacksquare\)

\textbf{Lean.} \nolinkurl{ScientificClosure.approximate_fbt_realized_bound}, \nolinkurl{certified_structural_continuation}, and \nolinkurl{certified_formed_structural_continuation.}

Without the binding-transport condition, formation maps may commute exactly while parent outputs change arbitrarily, invalidating the conclusion. The declaration \nolinkurl{formationTransport_does_not_imply_bindingTransport} gives an explicit Bool counterexample. A structural-distance bound also does not imply equality of parent-identity tokens.

\section{34.5 The Actual P12 Endpoint and the Program-Provenance Counterexample}\label{the-actual-p12-endpoint-and-the-program-provenance-counterexample}

The stronger P12 endpoint \nolinkurl{p12_terminal_conditional_finite_lifecycle} consumes aligned data: a finite physical PMF source \(P\), a viability fixed point \(K\), a physical repair certificate, a trusted repair substrate, finite program-representation coding and decoding, a seed with both a physical component and program provenance, and a proof that the actual program preserves the policy after performing repair. The fault kernel must lie in the stated joint fault envelope, with fault proportion \(\epsilon<1\).

Let \(S\) be the organizational maintenance system constructed from these data, and let \(\mu_0\) be the Dirac initial law of the actually assembled organization. Write \(D_n\) for the probability that its \(n\)th actual marginal lies in the set of illegitimate organizations. Let \(\kappa>0\) be the effective repair drift and \(\lambda\ge0\) the mean fault-injection burden. In the source, these are determined by the same fault proportion and the same system; they are not two freely chosen parameters.

The mean-homeostasis conclusion states that, for every \(\eta>0\), all sufficiently large \(N\) satisfy

\[\frac1N\sum_{n<N}D_n\le\frac\lambda\kappa+\eta.\]

It does not say that every path eventually remains error-free. If the initial potential is \(V_0\), the variable cost is \(c_v\), and the fixed maintenance cost is \(c_m\), the resource assumptions have the form

\[c_v\frac{V_0}{\kappa}\le R_0,\qquad
c_m+c_v\frac\lambda\kappa\le s.\]

They imply that the expected costs computed from the actual organizational marginals satisfy every cumulative resource budget:

\[\sum_{n<N}E[c(Z_n)]\le R_0+Ns.\]

The guarantee concerns the expected-cost ledger in the definition; it does not imply absence of bankruptcy on every random path.

The same lifecycle endpoint also includes convergence to zero of the changing-class discoverer's error probability for the same parent object; support of the assembled initial law in the organizational carrier; preservation of the same parent at every supported marginal through an explicit SameObject bridge; distinctness of the actual parent and the given child; the Price decomposition for the same reproduction kernel; and the relevant goal property under a supplied scale-preservation bridge. Its proof invokes the actual assembly, mean-homeostasis, expected-resource, identity, lineage, and scale endpoints and then combines their conclusions. It does not return these conclusions as fields of an arbitrary record.

Here ``finite'' refers first to the contract of finite states, actions, and program representations. The conclusions include marginals at all natural-number times and asymptotic averages. They \textbf{do not prove that the organization must die in finite time}, and they do not establish indefinite persistence under arbitrary conditions. Finite lifetime, absorption, and resource exhaustion require further hypotheses.

This conditional endpoint is not a completed realization by an actual append-history family of individuals. If a normalized reproduction kernel on a finite permanent universe of individuals is also required to increase actual generation number strictly at every transition, the combined requirements are impossible. Section 39.2 retains the strict counterexample; Sections 39.20--39.21 use the correct replacement: a rectangular channel from parent types to offspring events, with positive actual birth mass. Generation increase cannot be folded into a finite cycle to apply the endpoint.

\hypertarget{m-34-3}{}\phantomsection\label{m-34-3}

\textbf{M-34.3: Program provenance cannot be recovered from indistinguishable physical components.} If two valid seeds have the same physical component but require different program outputs, no deterministic synthesizer reading only the physical component can output the correct program for both.

\textbf{Proof.} A function has only one value at a given input, contradicting the two distinct required outputs. \(\blacksquare\)

\textbf{Lean.} \nolinkurl{TheoryCompletion.no_physicalOnly_synthesizer_for_samePhysical_distinctSeeds}; \nolinkurl{UnifiedClosure.same_physical_seed_program_reconstruction_no_go.}

For a supplied encoding \(e:\mathrm{Program}\to\mathrm{Information}\), if the program domain has a fallback value, existence of a total left-inverse decoder is equivalent to injectivity of \(e\). A left inverse implies injectivity; conversely, an injective encoding permits the unique preimage on its image and the fallback off its image. The declaration \nolinkurl{perfect_program_decoder_iff_injective_information} proves both directions. Physical existence of the encoding remains a separate question.

\section{34.6 The Scientific Interfaces C1--C7}\label{the-scientific-interfaces-c1c7}

\textbf{C1: Parameterized belief.} On a standard Borel latent space, under the relevant countable-generation conditions, a jointly measurable posterior kernel is constructed from the supplied belief and action--observation joint kernel. Conditional-probability factorization ensures that each fixed-parameter posterior section agrees almost everywhere with the existing posterior; pushforward under the observation distribution then produces the belief transition. The declarations \nolinkurl{c1_parameterized_belief_recursion} and \nolinkurl{c1_countablyGenerated_belief_recursion} connect these constructions to the same actual joint law. Null observations still permit version freedom, and the construction does not extend unconditionally to arbitrary measurable spaces.

\textbf{C2: Three-valued decisions with drift.} If the error radius between a true quantity \(d\) and its estimate \(\widehat d\), together with the drift allowance, totals \(r\), then the true quantity lies in \([\widehat d-r,\widehat d+r]\). Certification is sound when the upper endpoint is at most the threshold, rejection is sound when the lower endpoint exceeds it, and the remaining case is undecided. The proof uses interval containment alone. \nolinkurl{IntervalCertificate.decideAt_certified_sound}, \nolinkurl{decideAt_rejected_sound}, and \nolinkurl{decideAt_ambiguous_iff} retain all three outputs.

\textbf{C3: Common-channel calibration.} For one and the same Markov channel, data processing shows that the output comparison does not increase the TV distance used, giving calibration constant one for the specified parent-binding construction. The declaration \nolinkurl{commonMarkov_calibrates_parentBinding} does not cover a setting in which the channel changes with every sample without an error allowance.

\textbf{C4: Conditional FBT.} M-34.2 supplies this interface. It does not automatically generate a parent-formation relation or binding compatibility.

\textbf{C5: A necessary rank condition for two common drives.} If every environment \(m_e(B)=g_e(z(B))\) shares the same two-dimensional \(z\), the stacked Jacobian at a common budget point factors through the same right factor \(Dz\), so its rank is at most two. The declaration \nolinkurl{commonTwoDrive_rankNecessary} and the singular-value perturbation results support rejection, compatibility, or an undecided outcome. Compatibility does not establish unique identification.

\textbf{C6: Purpose faithfulness.} The mechanism bridge following M-34.1 supplies this interface. Identification of an ordering class still leaves control realization and viability conditions to be verified independently.

\textbf{C7: Evidence and kinds of failure.} A single observational failure may refute a sampling protocol, a source model, a physical interpretation, or the application hypotheses of a formal proposition. It cannot be called a refutation of the mathematical core without distinguishing these possibilities. Conversely, successful compilation of the core cannot protect an incorrect world model. The historical evidence has a limitation: retention of every interrupted attempt cannot be verified. A new persistent-record mechanism does not retrospectively repair that past. External empirical support in the current scientific-interface record remains UNVERIFIED.

\section{34.7 Scale Properties and the Uncompleted RG Interface}\label{scale-properties-and-the-uncompleted-rg-interface}

\nolinkurl{ObjectScaleMap} is composable. If a predicate is preserved by each of two maps, successive application of the two preservation proofs establishes preservation by their composite. The endpoint \nolinkurl{p9_terminal_objectScale_semantic_transport} separately consumes six certificates: identity, purpose, GOD, GOA, repair, and homeostasis.

The current \nolinkurl{wilsonianCouplingFlow_nontrivial_of_twoCouplings} states only that a coupling domain with two distinct points admits two distinct constant-flow records. Such a record has no compatibility equation with an actual quotient map, dynamics, or defect calibration. The book therefore does not promote it to the claim that a given scale map admits two valid Wilsonian flows. Such a new claim would require a compatibility definition followed by a proof of existence or nonuniqueness of compatible flows.

\chapter{35. Prediction, Carriers, Viability, Repair, and Optimal Control on One Finite Source}\label{ch-35}

\textbf{Chapter question.} Can prediction, physical safety, repair, and control genuinely consume the same finite source?

\textbf{Reading route.} Fix the kernel, reward, and registered readout before constructing a stable quotient, a safety fixed point, and a constrained Bellman problem. A source-kernel equality links the repair program to safety. Safety, optimality, and identity preservation retain their own hypotheses; none automatically grants the others.

\section{35.1 A Common Mathematical Basis}\label{a-common-mathematical-basis}

Fix a finite state space \(X\), a finite nonempty action space \(A\), one source kernel \(P(x,a)\in\mathrm{PMF}(X)\), a reward \(r(x,a)\), a discount \(0<\beta<1\), and a registered output \(o(x)\). Organizational, program, resource, and safety labels to be preserved may be included as components of \(o\). Every subsequent construction must preserve the same \(P\) that it claims to use; another kernel of the same type cannot replace it.

Stable refinement as in Chapter 28 gives the coarsest stable partition \(\sim\), with quotient map \(q:X\to X/{\sim}\). The initial partition preserves \(o\) and the complete vector of action rewards. Stability makes \(q_\#P(x,a)\) depend only on \(q(x),a\). The exact quotient kernel and reward follow, and P-QUO-01 gives value equality and lifting of a macro-optimal policy.

\textbf{Lean.} \nolinkurl{exactControlFromStablePartition} and \nolinkurl{organizationalTerminalControl} in \nolinkurl{OrganizationalCoreAssembly}; \nolinkurl{registered_organizational_attributes_preserved} and \nolinkurl{unified_finite_core_control_and_viability} in \nolinkurl{UnifiedCoreClosureCertificate.}

\section{35.2 Deriving a Safe Region and Actual Paths from the Same Kernel}\label{deriving-a-safe-region-and-actual-paths-from-the-same-kernel}

Given a safe set \(V\subseteq X\), apply the deletion recursion in Chapter 8 to obtain its greatest fixed point \(K\). Its action condition is

\[x\in K\Longrightarrow\exists a:\operatorname{supp}P(x,a)\subseteq K.\]

\hypertarget{m-35-1}{}\phantomsection\label{m-35-1}

\textbf{M-35.1: Safety at all times on the same source.} For every \(\mu_0\) supported on \(K\), there is a deterministic stationary policy \(\pi\) such that the infinite trajectory constructed from the original \(P\) and \(\mu_0\) satisfies

\[\Pr_{P,\pi,\mu_0}(\forall n, X_n\in K)=1.\]

\textbf{Proof.} At each point choose an action satisfying the fixed-point condition. If the \(n\)th marginal is supported on \(K\), the original \(P\) under these actions is also supported on \(K\), so the same holds for the \(n+1\) marginal. The marginals of the actual path extension agree with this recursion. Intersecting the probability-one events over all \(n\) gives the claim. \(\blacksquare\)

\textbf{Lean.} \nolinkurl{ViabilitySource.p_per_03}; the unified same-source endpoint is \nolinkurl{UnifiedClosure.single_source_terminal_object_viability_all_times.}

The set \(K\) may be empty. In that case no probability initial law is supported on \(K\), and the proposition does not guarantee viability of a real initial state. In the earlier model of Section 30.1, if the service bit must always equal one, every action has positive probability of leaving that state. Its corresponding almost-sure viability kernel is therefore empty, despite high average service quality and the existence of an ordinary invariant law.

\section{35.3 Constrained Optimal Control from the Safe Fixed Point}\label{constrained-optimal-control-from-the-safe-fixed-point}

Define the admissible actions at each state by

\[A_K(x)=\{a\in A: x\in K\Rightarrow\operatorname{supp}P(x,a)\subseteq K\}.\]

On \(K\), nonemptiness follows from the fixed-point property. Outside \(K\), every action satisfies the implication. The constrained model retains the original reward and kernel and restricts only the action domain.

\hypertarget{m-35-2}{}\phantomsection\label{m-35-2}

\textbf{M-35.2: Optimality over all causal policies subject to the safety constraint.} The constrained Bellman operator

\[(T_Kv)(x)=\max_{a\in A_K(x)}[r(x,a)+\beta P(x,a)v]\]

has a unique fixed point \(V_K^*\) and a deterministic stationary policy attaining it. Every history-dependent randomized policy that selects only actions in \(A_K(X_n)\) at each step has discounted return at most \(V_K^*\). If its initial law is supported on \(K\), every such policy is almost surely safe on the same source.

\textbf{Proof.} The action sets are finite and nonempty, so the Bellman contraction in Chapter 19 gives the fixed point and an attaining policy. For any constrained policy, the Bellman upper bound holds for every supported action. Conditional expectation with respect to its actual history, discounted summation, and passage to the limit in the bounded terminal term yield domination of all causal policies. Safety follows from the support induction in M-35.1; the competing policies need not be stationary. \(\blacksquare\)

\textbf{Lean.} \nolinkurl{UnifiedClosure.exists_viable_source_action}, \nolinkurl{viable_restricted_greedy_preserves_source_kernel}, and \nolinkurl{viable_restricted_optimal_dominates_all_causal_policies.}

The optimal policy and the safety constraint here use the same kernel. The conclusion makes no claim of optimality over unconstrained policies.

\section{35.4 A Finite Example Separating Safety from Unconstrained Optimality}\label{a-finite-example-separating-safety-from-unconstrained-optimality}

Take \(x=(s,z)\in\{0,1\}^2\), with \(z'\) independently fair at each step. The actions are waiting, \(a=0\), and repair, \(a=1\). Repair always sets \(s'=1\). Waiting retains zero when \(s=0\); when \(s=1\), it retains one with probability \(3/4\) and changes to zero with probability \(1/4\). Let

\[r(s,z,a)=s-a,\qquad\beta=1/2.\]

The map \(q(s,z)=s\) preserves every action kernel and reward. For the safe set \(V=\{s=1\}\), the fixed point is \(K=V\), and repair is the only admissible action in \(K\). The constrained safe value is zero because the reward at each step is \(1-1=0\).

The unconstrained values are \(V(0)=0,V(1)=8/5\). At state zero, repair has action value \(-1+(1/2)(8/5)=-1/5\), while waiting has value zero. At state one, waiting has value \(1+(1/2)(3/4)(8/5)=8/5\), and repair has value \(4/5\). Thus the unconstrained optimal policy waits, incurring positive failure risk in the healthy state.

This example is formalized by the separate \nolinkurl{WaitingRepairControl} source. actual\_pmf, \nolinkurl{actual_reward}, and \nolinkurl{actual_discount} bind the source just described; \nolinkurl{optimal_value}, \nolinkurl{dominates_all_causal}, and \nolinkurl{actual_greedy_attains} provide the original infinite causal-control consumers. \nolinkurl{restricted_optimal_zero}, \nolinkurl{restricted_all_causal}, and \nolinkurl{actual_safe_attainment} provide the actual constrained model.

The declaration \nolinkurl{actual_macro_optimal} also consumes the original exact control quotient. This process differs from the \nolinkurl{zero/two-value} model of Section 21.7. It demonstrates that unconstrained optimality and safety-constrained optimality can differ; registered outputs and exact quotients cannot remove that difference in value constraints.

\section{35.5 A Recovery Program Returns the Same Safe Kernel}\label{a-recovery-program-returns-the-same-safe-kernel}

Let \(T\) be a trusted program-execution substrate, and suppose that program \(p\) is valid for one of the safe action selectors in the preceding section. The repair process jointly updates the physical state and program representation. Under the proved \nolinkurl{coding/decoding} and recovery conditions, its projected next-physical-state kernel is exactly \(P(x,a_p(x))\).

\hypertarget{m-35-3}{}\phantomsection\label{m-35-3}

\textbf{M-35.3: A source-kernel identity bridge for program recovery.} If the recovery program's execution kernel equals the original safe source action kernel, then execution from a state in \(K\) remains supported on \(K\).

\textbf{Proof.} The recovery theorem first gives equality of kernels. Transport the proved inclusion \(\operatorname{supp}P(x,a_p(x))\subseteq K\) along that equality. Program recovery and physical safety follow from their own hypotheses and are then connected by the common kernel. \(\blacksquare\)

\textbf{Lean.} \nolinkurl{UnifiedClosure.reconstructed_program_source_kernel_agrees_with_safe_action}, \nolinkurl{reconstructed_program_preserves_source_viability}, and \nolinkurl{reconstructed_program_joint_and_source_safety.}

This bridge neither synthesizes \(p\) from unknown physical material nor asserts that the recovered program attains the constrained optimal value in M-35.2. To obtain a safety-constrained optimal program, one must also prove that it implements a constrained Bellman selector. Safety alone does not confer optimality.

\section{35.6 The Unification Obtained and the Inputs Retained}\label{the-unification-obtained-and-the-inputs-retained}

One finite source can connect registered outputs, a stable \nolinkurl{predictive/control} quotient, access targets and carriers, a viability fixed point, an actual infinite safe path law, constrained optimal control, and source-kernel identity for a recovery program. Each step has a separate proof; the reader need not identify several different models as a single physical world.

Retained inputs include the reward and safety task, the physical meaning of registered organizational attributes, the actual source model, initial viability, and the program substrate. They are consistent with the provenance impossibility results in Chapter 34. The next three chapters ask whether policies can actually be constructed under partial observation and how using a data-selected model after finite training yields a genuine risk guarantee.

\section{35.7 Object Partitions, Actions, Supports, and Recovery in One Instance}\label{object-partitions-actions-supports-and-recovery-in-one-instance}

The four-state source fixed in Section 21.7 comes from the existing \nolinkurl{conflictRewardControl.micro.} Its state is \((b,u)\in\{0,1\}^2\); action \(a\) makes the next visible bit equal to \(a\), while the internal bit is resampled fairly. The reward is \(1-a\) and the discount is \(1/2\). This section consistently uses that complete M with initial law \(\delta_{(1,0)}\), rather than the different transition or reward in Section 35.4.

For the registered readout \(q(b,u)=b\), every action satisfies \(q_\#P(x,a)=\delta_a\), and rewards do not distinguish states in the same visible fiber. The existing \nolinkurl{survivalRead_strongLumpable} and \nolinkurl{conflictRewardControl} therefore supply an exact control quotient with nonzero reward. The actual macro type is q's \nolinkurl{ReachableState} subtype; its correspondence with Bool requires an explicit map. If the internal bit is also registered as a program attribute to be preserved, the original two-class partition no longer suffices. The existing \nolinkurl{hiddenProgramProtectedModel} retains this distinction.

Although the earlier P-ALG witness \nolinkurl{registeredSurvivalModel} uses the same physical kernel, its reward field is zero and cannot replace the complete M here. The existing \nolinkurl{ConflictOrganizationalSource.registeredModel} explicitly uses M.reward and M.transition. The declarations \nolinkurl{registered_reward_actual} and \nolinkurl{registered_transition_actual} bind it to the original source; \nolinkurl{initial_stable}, \nolinkurl{terminal_eq_initial}, and \nolinkurl{terminal_relation} prove directly that its original P-ALG terminal partition is exactly the kernel of q. \nolinkurl{actual_macro_read} gives the explicit map from the actual \nolinkurl{ReachableState} subtype to Bool. The \nolinkurl{protectedModel} adds the protected internal bit on the same nonzero-reward source, while \nolinkurl{protected_terminal_relation} and \nolinkurl{protected_hidden_cannot_merge} retain the corresponding finer partition.

The actual supports are

\[
\operatorname{supp}P(x,1)=K,\qquad
\operatorname{supp}P(x,0)=K^c,\qquad K=\{b=1\}.
\]

Inside K, action 1 is precisely the action that preserves the domain. Under this chapter's constrained-action definition, both actions remain admissible outside K. M-21.7 gives the constrained optimal value state by state: zero on K and two outside it. The full greedy selector preserves the current visible bit and has two closed classes and all their mixture invariant laws. Only from the specified safe initial law does it have the same actual path law as the policy constantly choosing 1. Persistence and optimality thus have an actual common source, together with indispensable initial-law and comparison-class restrictions.

\nolinkurl{ConflictOrganizationalSource.Posterior.actual_posterior_support} proves the posterior support of this source using the original \nolinkurl{fiberTerminalPosterior} and exhibits a useful scope boundary. At time zero the only positive-probability observation is 1, with actual support the singleton \((1,0)\). At every positive time and positive-probability actual history, the internal bit has been resampled fairly on the last step, so the terminal support equals the current visible fiber of that history. This instance exhibits no memory gain from a history that strictly narrows the current fiber. The separate strict-posterior and dead-end witnesses in Chapter 37 remain essential.

For the same policy that chooses action 0 with probability \(\theta\) at every history, if certification requires \(a=1\), the actual uncertified occupancy is \(\theta\) at every step. Its probability of at least one failure in T steps is \(1-(1-\theta)^T\). If all actual rows are certified and the exact actual reference kernel is chosen, uncertified occupancy can be zero while reference risk remains that positive quantity. A different certificate may certify only action 1 and use a zero-failure reference path, but must pay the occupancy budget of the policy itself. These parameters belong to their respective complete certificates; zero reference risk from one cannot be combined with zero occupancy from another.

The recovery program can also be connected to the same P. An explicitly sourced recovery step may first reconstruct the program while temporarily leaving the physical state unchanged. The next organizational step executes recovered action 1 and then enters K through P. The program representation is additional content in the joint recovery state; it is not the internal bit u resampled at each step in this example. Nor are the two organizational steps two unchanged physical P steps. The declarations \nolinkurl{first_tick_restores}, \nolinkurl{first_tick_physical_unchanged}, \nolinkurl{two_tick_physical_source}, and \nolinkurl{two_tick_recovered} in \nolinkurl{ConflictOrganizationalSource} construct and consume the original \nolinkurl{JointRepairKernel} and \nolinkurl{RecoveredJointTarget.} For every ordinary controller and every single damaged copy, the first tick restores the actual program without changing the physical state; the second executes the recovered program and enters K through the original P. \nolinkurl{damaged_not_safe_after_restore} retains the counterexample in which a damaged physical initial state is still unsafe after the first tick. Recovery reopens the possibility of safe future operation; it does not erase the earlier interruption from the actual history.

This common instance connects object descriptions, feasible actions, constrained value, supports, probability risks, and the source conditions for recovery. It does not replace the physical framework for carrier damage, parent formation, lineage, or external experiments. The precise scope of each mathematical connection gives the philosophical discussion of objects, persistence, and meaning a definite content.

\textbf{Sources.} \href{../formalization/ueot-core/docs/compression/theory_completion/unified_closure/interface_research_v1/UNIFIED_FINITE_SOURCE_EXAMPLE.md}{Unified finite-source study and arrow-by-arrow table}, \href{../formalization/ueot-core/UEOT/V3/Compression/TheoryCompletion/UnifiedClosure/SafeRewardConflictWitness.lean}{\nolinkurl{SafeRewardConflictWitness}}, and \href{../formalization/ueot-core/UEOT/V3/Compression/TheoryCompletion/UnifiedClosure/OrganizationalSufficiencyWitness.lean}{\nolinkurl{OrganizationalSufficiencyWitness}}. The nonzero-reward partition, posterior, and two-clock recovery added here, together with the complete parameter family in Section 21.7, are documented in the \href{../formalization/ueot-core/docs/compression/theory_completion/unified_closure/conflict_family_book_source_v1/RESULTS_AND_OPEN_OBLIGATIONS.md}{named-source package for the original book}. Read the exact quantifiers and endpoints together with Section 21.7 and Appendix F.

\chapter{36. Actual Finite Paths, Posterior Decomposition, and Deployment Coverage}\label{ch-36}

\textbf{Chapter question.} How does rowwise model error enter the risk of an entire deployment path under the actual policy?

\textbf{Reading route.} Construct the posterior and occupancy from the actual microscopic prefix. A TV identity retaining the prefix then yields stepwise accumulation. Uncertified rows are charged according to their own actual occupancy. With a finite observation image, escape mass must be retained rather than deleted and renormalized.

\section{36.1 First Fix the Generating Law of the Experiment}\label{first-fix-the-generating-law-of-the-experiment}

Fix a finite nonempty state space \(X\), a finite nonempty action space \(A\), an observation map \(q:X\to Y\), an initial probability \(p_0\), and a time-dependent Markov source

\[P_n(x,a)\in\mathrm{PMF}(X).\]

An observation history contains more than the current observation. Set

\[H_0^Y=Y,\qquad H_{n+1}^Y=H_n^Y\times(A\times Y),\]

and define microscopic histories \(H_n^X\) recursively in the same way. The map \(Q_n:H_n^X\to H_n^Y\) retains each action and replaces each state by its \(q\)-observation. Here the capital \(Q_n\) denotes the history projection; the reference dynamics below is denoted separately by \(\widehat P_n\).

The policy \(\pi_n(g)\in\mathrm{PMF}(A)\) reads the full coarse history \(g\). Given a microscopic prefix \(h\), first use \(\pi_n(Q_nh)\) to draw an action \(a\) and then draw the next state according to \(P_n(\operatorname{last}h,a)\). The initial law \(p_0\) and this recursion generate the actual microscopic path law \(p_n\), whose actual coarse path law is \(\bar p_n=(Q_n)_\#p_n\).

Every subsequent proof uses this generating relation. Training data, a reference model, or another PMF cannot be renamed \(p_n\) merely because their state types agree. Probability calculations use finitely supported PMFs. When measures are involved, finite spaces carry the discrete measurable structure, and a general observation space retains the assumption of measurable singletons. These facts make maps from the finite source and actual finitely supported paths measurable; they do not assert automatic measurability of arbitrary infinite-history policies. PMF extension is the substantive construction here, not merely an existence assertion about an abstract joint law.

\textbf{Lean.} causalLaw and pmfExtend in PathError; \nolinkurl{riskMicroControlKernel} and \nolinkurl{riskHistoryMap} in \nolinkurl{UnifiedClosure}; \nolinkurl{coverageDeploymentRowLaw} in \nolinkurl{DeploymentCoverage.}

\section{36.2 The Actual Posterior on a Common History Fiber}\label{the-actual-posterior-on-a-common-history-fiber}

When \(\bar p_n(g)>0\), define the terminal-state posterior by

\[
\rho_n^g(x)=
\frac{\displaystyle\sum_{\substack{h\in H_n^X\\Q_nh=g\\\operatorname{last}h=x}}p_n(h)}
{\bar p_n(g)}.
\]

It is supported on the current observation fiber \(q^{-1}(\operatorname{last}g)\), but may be much smaller than that fiber because earlier observations and actions exclude some states. Given a coarse history and an action, the actual next-observation kernel is

\[
K_n(g,a)(y)=\sum_x\rho_n^g(x)\,[q_\#P_n(x,a)](y).
\]

Combining the action with this posterior uses an explicit hypothesis: all microscopic prefixes in the same coarse-history fiber use \(\pi_n(g)\). Thus, for positive-mass \(g\),

\[
\Pr(X_n=x,A_n=a\mid G_n=g)=\rho_n^g(x)\pi_n(g)(a).
\]

This is not an independence assertion for arbitrary training--deployment couplings. It follows from the action-generation order specified here. If the policy also reads hidden microscopic information absent from \(g\), this factorization generally fails.

\hypertarget{m-36-1}{}\phantomsection\label{m-36-1}

\textbf{M-36.1: Finite-fiber decomposition and zero-mass branches.} At every positive-mass coarse history, \(\rho_n^g\) is a probability distribution and satisfies

\[
\bar p_n(g)\rho_n^g(x)
=\sum_{\substack{Q_nh=g\\\operatorname{last}h=x}}p_n(h).
\]

Assigning any probability completion to a zero-mass coarse history preserves the weighted equality. The kernel \(K_n\) constructed from these posteriors agrees with the next coarse-history extension of the same \(p_n\).

\textbf{Proof.} At positive mass, cancel the denominator. Summing the numerator over terminal states gives \(\bar p_n(g)\). At zero mass, the nonnegative terms in the fiber have total sum zero, so each vanishes; any completion multiplied by zero also vanishes. Finally, rearrange the finite sums over \(g,a,x,x'\) and factor out the common \(\pi_n(g)(a)\) on each fiber. This gives the projected next-step path law. \(\blacksquare\)

\textbf{Lean.} \nolinkurl{riskFiberPosterior}, \nolinkurl{risk_map_mul_fiberPosterior}, \nolinkurl{risk_conditionedRow_mass}, \nolinkurl{risk_map_extend_apply}, and \nolinkurl{risk_map_extend_eq} in \nolinkurl{FinitePMFProjection.} The terminal-state posterior is the pushforward of the existing history-fiber posterior under last. The new support bridge reuses that decomposition instead of assuming another posterior that matches the actual law.

A posterior at a zero-mass history is only a choice of version. Inclusion of a state in such a completion cannot be used to require certification of that state in the actual experiment.

\section{36.3 Exact TV Decomposition When the Prefix Is Retained}\label{exact-tv-decomposition-when-the-prefix-is-retained}

For probability distributions on finite sets, use the convention

\[
\operatorname{TV}(p,r)=\frac12\sum_z|p(z)-r(z)|
=\sup_E|p(E)-r(E)|.
\]

For a common prefix law \(p\) and two successor-kernel families \(K,L\), write the extended joint law as

\[
(p\ltimes K)(h,z)=p(h)K(h)(z).
\]

\hypertarget{m-36-2}{}\phantomsection\label{m-36-2}

\textbf{M-36.2: The exact row-error formula for a common prefix.}

\[
\operatorname{TV}(p\ltimes K,p\ltimes L)
=\sum_h p(h)\operatorname{TV}(K(h),L(h)).
\]

For a common successor kernel,

\[
\operatorname{TV}(p\ltimes K,r\ltimes K)\le\operatorname{TV}(p,r).
\]

\textbf{Proof.} For the first identity, factor out the nonnegative \(p(h)\) termwise and interchange the finite sums. The second follows from Markov data processing. Retaining the entire prefix in fact gives equality, but the upper bound suffices for recursion. If the prefix is discarded and only the mixed terminal states are compared, the first identity generally becomes only ``\(\le\)'', since signed errors from different prefixes may cancel. \(\blacksquare\)

\textbf{Lean.} \nolinkurl{coverage_tv_extend_same_initial} and \nolinkurl{coverage_tv_extend_same_kernel}; the corresponding mixture bound reuses \nolinkurl{risk_tv_bind_same_weights.} The first result fixes the weights of subsequent errors to the actual prefixes. An arbitrary training-row distribution cannot be substituted on the right-hand side of the equality.

\section{36.4 From Row Errors to Actual Path Error}\label{from-row-errors-to-actual-path-error}

Let \(\widehat P_n(y,a)\) be a reference kernel with initial observation law \(q_\#p_0\), and generate actions using the same \(\pi_n(g)\). Write \(\widehat p_n\) for the reference path law and define

\[
d_n(x,a)=\operatorname{TV}\bigl(q_\#P_n(x,a),\widehat P_n(qx,a)\bigr).
\]

The state--action occupancy law in actual deployment is

\[
\nu_n(x,a)=
\sum_{\operatorname{last}h=x}p_n(h)\pi_n(Q_nh)(a).
\]

It is the pushforward of the actual prefix--action joint PMF under \((h,a)\mapsto(\operatorname{last}h,a)\), and therefore has total mass one.

\hypertarget{m-36-3}{}\phantomsection\label{m-36-3}

\textbf{M-36.3: A path bound weighted by actual occupancy under the same policy.}

\[
\operatorname{TV}(\bar p_T,\widehat p_T)
\le
\min\left\{1,\ \sum_{n<T}\sum_{x,a}\nu_n(x,a)d_n(x,a)\right\}.
\]

\textbf{Proof.} Put \(\delta_n=\operatorname{TV}(\bar p_n,\widehat p_n)\). Insert an intermediate law using the actual coarse prefix \(\bar p_n\) and the reference successor kernel. The triangle inequality splits the next-step error into a change of kernel at a common prefix and a change of prefix under a common kernel. The second term is at most \(\delta_n\). By M-36.2, the common-history action factor, and convexity of TV, the first is at most

\[
\sum_g\bar p_n(g)\sum_a\pi_n(g)(a)\sum_x\rho_n^g(x)d_n(x,a).
\]

Use M-36.1 to rewrite this as \(\sum_{x,a}\nu_n(x,a)d_n(x,a)\). The initial error is zero. Induction and summation, together with the bound of one on probability TV, finish the proof. \(\blacksquare\)

\textbf{Lean.} \nolinkurl{coverage_projected_extension_error}, \nolinkurl{coverage_projected_causalLaw_error}, \nolinkurl{coverage_rowLaw_expectation}, \nolinkurl{coverage_same_policy_row_tv}, and \nolinkurl{coverage_same_policy_path_error.} They use the actual prefix of the original causalLaw. Equality between the occupancy used and actual occupancy is not supplied as an additional conclusion-valued input field.

\section{36.5 Certifying Only Some Deployment Rows}\label{certifying-only-some-deployment-rows}

Let \(C_n(x,a)\) specify the certified source rows. Suppose that \(0\le\varepsilon_n\le1\) and that \(d_n(x,a)\le\varepsilon_n\) on those rows. On the remaining rows use only the general probability-distance bound \(d_n\le1\).

\hypertarget{m-36-4}{}\phantomsection\label{m-36-4}

\textbf{M-36.4: The exact linear charge for uncertified occupancy.}

\[
\sum_{x,a}\nu_n(x,a)d_n(x,a)
\le
\varepsilon_n+(1-\varepsilon_n)\nu_n(C_n^c).
\]

If the actual uncertified mass also satisfies \(\nu_n(C_n^c)\le\kappa_n\), the bound may be replaced by

\[
\varepsilon_n+(1-\varepsilon_n)\kappa_n
\]

For a microscopic failure event \(E\subseteq H_T^X\) and a coarse event \(F\subseteq H_T^Y\), if \(E\subseteq Q_T^{-1}F\), one obtains

\[
\Pr_{p_T}(E)\le
\min\left\{1,\ \Pr_{\widehat p_T}(F)
+\sum_{n<T}\bigl[\varepsilon_n+(1-\varepsilon_n)\kappa_n\bigr]\right\}.
\]

\textbf{Proof.} Split the finite sum over \(C_n\) and \(C_n^c\). The first part is at most \(\varepsilon_n\nu_n(C_n)\), and the second is at most \(\nu_n(C_n^c)\); their sum gives the first inequality. Substitution of the occupancy upper bound uses \(1-\varepsilon_n\ge0\). Event containment, the event bound for TV, and M-36.3 then give the risk inequality. \(\blacksquare\)

\textbf{Lean.} \nolinkurl{coverage_uncovered_indicator_expectation}, \nolinkurl{coverage_expected_error_le}, \nolinkurl{coverage_same_policy_certified_path_error}, and \nolinkurl{coverage_same_policy_microscopic_failure_risk.} The budget-monotonicity interface \nolinkurl{ObservableActionCoverageContinuationCheck.deployment_coverage_budget_le_of_actual_occupancy_cap} replaces an upper bound on the actual \(\nu\) by a supplied \(\kappa\). This proved monotone substitution does not manufacture missing coverage evidence.

If training observes only rows in \(C_n\) without controlling the deployment probability of entering \(C_n^c\), the risk budget still contains unknown actual occupancy. Increasing the sample count alone cannot eliminate that term.

\section{36.6 How Action Certification Bounds Actual Coverage}\label{how-action-certification-bounds-actual-coverage}

Assign an admissible action set \(\Gamma_n(g)\) to every actual positive-mass coarse history. Assume

\[
a\in\Gamma_n(g),\quad \rho_n^g(x)>0
\quad\Longrightarrow\quad C_n(x,a),
\]

and suppose that at those histories the policy satisfies

\[
\pi_n(g)(\Gamma_n(g)^c)\le\kappa_n.
\]

\hypertarget{m-36-5}{}\phantomsection\label{m-36-5}

\textbf{M-36.5: A support-certified action budget implies an actual coverage budget.}

\[
\nu_n(C_n^c)\le\kappa_n.
\]

\textbf{Proof.} By M-36.1, write \(\nu_n(C_n^c)\) as a finite triple sum of coarse-history, action, and posterior-state masses. At a positive history, every action in \(\Gamma_n(g)\) is certified for all states of positive posterior mass, so only actions outside that set can contribute. The conditional contribution at each positive history is at most its action leakage mass \(\kappa_n\). Zero-mass coarse histories, actions, and posterior states contribute nothing. Sum over the probability weights \(\bar p_n(g)\). \(\blacksquare\)

\textbf{Lean.} \nolinkurl{supported_history_actions_bound_actual_coverage} and \nolinkurl{posterior_supported_actions_bound_actual_coverage} in \nolinkurl{PosteriorSupportCoverageCheck.} The passed minimal check directly consumes the existing \nolinkurl{riskFiberPosterior} without reconstructing the conditional-probability library.

A stronger sufficient condition uses the entire current observation fiber:

\[
\Gamma_n(y)\subseteq\bigcap_{q(x)=y}\{a:C_n(x,a)\}.
\]

This condition is independent of the history posterior and is convenient as a fixed action mask. It may, however, exclude feasible history-dependent policies; Section 37.6 gives a strict witness.

\section{36.7 The Training Distribution Cannot Replace the Deployment Posterior}\label{the-training-distribution-cannot-replace-the-deployment-posterior}

Here \(\nu_n\) lives on \(X\times A\), whereas the training-experiment law \(\mu\) usually lives on a sample-record space \(\Omega\). They do not even have the same type. Even if the space \(X\times A\) is supplied with a separate training-row distribution \(\lambda_n\), an additional relation is needed to substitute one expectation for the other.

For example, for \(0\le d\le1\), if \(\operatorname{TV}(\nu_n,\lambda_n)\le\delta_n\), then

\[
\mathbb E_{\nu_n}d\le\mathbb E_{\lambda_n}d+\delta_n.
\]

If \(\nu_n\ll\lambda_n\) and its density is at most \(M_n\), then

\[
\mathbb E_{\nu_n}d\le M_n\mathbb E_{\lambda_n}d.
\]

These are applications, on the appropriate common space, of the distribution-shift bounds already given in Chapter 16. TV calibration of a source row against its empirical row in the response space does not imply that the distribution of rows visited in deployment is close to the training design. The two distances concern different spaces and serve different mathematical purposes.

\section{36.8 The Finite Actual Observation Image and Escape Mass That Must Be Retained}\label{the-finite-actual-observation-image-and-escape-mass-that-must-be-retained}

A readout \(q:X\to Y\) of finite X has a finite full image \(S=q(X)\) even when Y is infinite. A reference kernel may send mass outside S, so those points cannot simply be deleted and the remainder renormalized. This section gives a written extension preserving the original reference process, actual policy, and occupancy. The Lean status of the complete path adapter is recorded separately from the more elementary finite-encoding lemma.

\hypertarget{m-36-6}{}\phantomsection\label{m-36-6}

\textbf{M-36.6: Exact TV encoding with finite support on one side.} Suppose that Y has measurable singletons, S is finite, and the probability measures \(\mu,\nu\) satisfy \(\mu(S^c)=0\). Retain each point of S separately and send all exterior points to one additional symbol \(\bot\), obtaining a measurable map \(c:Y\to S\sqcup\{\bot\}\). Then

\[
\operatorname{TV}(\mu,\nu)=\operatorname{TV}(c_\#\mu,c_\#\nu).
\]

\textbf{Proof.} Measurability follows from the measurable singleton fibers of finite S and the measurable complement. Data processing bounds the right-hand side by the left. For the reverse inequality, take an arbitrary measurable A and let B contain exactly the retained points corresponding to \(A\cap S\). Then

\[
\mu(A)-\nu(A)\le\mu(A\cap S)-\nu(A\cap S)
=(c_\#\mu)(B)-(c_\#\nu)(B).
\]

Bound the right-hand side by encoded TV and apply the same inequality to \(A^c\). Both measures have total mass one, so this also bounds \(\nu(A)-\mu(A)\). Take the supremum over A. \(\blacksquare\)

For \(\mu=\delta_0,\nu=(\delta_0+\delta_1)/2,S=\{0\}\), the escape mass is \(1/2\), and encoded TV remains \(1/2\). Deleting 1 and renormalizing would incorrectly give zero. If both measures can have mass outside S, for example \(\mu=\delta_1,\nu=\delta_2\), merging exterior points reduces TV from one to zero. The one-sided concentration condition is essential.

An infinite ambient space \(Y=\mathbb N\) also distinguishes source error from differences between reference laws. Take \(S=\{0,2\}\) and

\[
\mu=\tfrac13\delta_0+\tfrac23\delta_2,\qquad
\nu_e=\tfrac12\delta_0+\tfrac14\delta_2+\tfrac14\delta_e,
\quad e\in\{5,7\}.
\]

Both references have escape mass \(1/4\), but \(\operatorname{TV}(\mu,\nu_e)=5/12\): the half-\(L^1\) difference on their finite atoms is \(\tfrac12(1/6+5/12+1/4)\). Meanwhile \(\operatorname{TV}(\nu_5,\nu_7)=1/4\), whereas the two encoded references agree and have TV zero. Source--reference distance therefore includes more than escape probability. Exact comparison against a fixed finitely supported source does not make the encoding a global isometry between all reference laws. \nolinkurl{FiniteObservationEscapeBoundary} formally proves this combined boundary.

\subsection{36.8.1 From Finite Encoding to the Same Actual Path}\label{from-finite-encoding-to-the-same-actual-path}

In this subsection, the reference dynamics kernel is denoted by \(Q_n\) (the \(\widehat P_n\) of the preceding sections). It is distinct from the history projection \(Q_n:H_n^X\to H_n^Y\) in Section 36.1; this subsection no longer uses the latter notation for that projection.

Take \(Z=S\sqcup\{\bot\}\). Choose \(x_\star\) from the nonempty support of the initial PMF and set \(y_\star=q(x_\star)\). The decoder d restores retained points and sends \(\bot\) to \(y_\star\). Decode every observation of a full Z-history while retaining its actions, then apply the original full policy π to obtain \(\bar\pi\). This defines the policy on impossible and escaped histories as well, without creating an action distribution from nothing.

For retained points \(s\in S\), set \(\bar Q_n(s,a)=c_\#Q_n(s,a)\); at \(\bot\), set \(\bar Q_n(\bot,a)=\delta_\bot\). These are normalized actual PMF rows. Put \(\bar q=c\circ q\). Mapping an actual microscopic history through \(\bar q\) and decoding it returns its original q-history, so \(\bar\pi\) assigns exactly the same action PMF there as π. The original actual microscopic path law and state--action occupancy are preserved exactly.

The original reference process may leave S and later return, whereas the finite reference is absorbed at \(\bot\). Thus they are not the same process under coordinatewise encoding; their full histories must be compared.

Let \(L_n\) be the finite set of full n-step observation--action histories whose observations all lie in S. Let \(i_n,j_n\) be the natural action-preserving injections into the original Y-history and Z-history spaces. Equality of the initial laws and the actual multiplicative recursion for path masses give, for every \(\ell\in L_n\),

\[
\bar\nu_n(j_n\ell)=\nu_n(i_n\ell).
\]

The inductive step uses equality of prefix masses, equality of policies on retained histories, and equality of reference-row masses at retained successors. Full-path matching is derived rather than imposed as a separate assumption.

Now encode both full history spaces into the same \(L_n\sqcup\{\bot\}\), matching each retained point to its corresponding \(\ell\) and sending every nonretained history to \(\bot\). The two reference laws agree at all retained points, and normalization gives equality of the remaining mass. The actual observation path is concentrated on retained histories, so applying M-36.6 on each side gives exact equality between the original path TV and its finite-reference comparison. A history that escaped and later returned is still nonretained; the terminal observation alone does not suffice.

\hypertarget{m-36-7}{}\phantomsection\label{m-36-7}

\textbf{M-36.7: A finite-observation-image path bound preserving the original source.} Under the preceding conditions, retain the original source P, full policy π, initial law, and reference Q. Writing \(w_n\) for the actual state--action occupancy, one has

\[
\operatorname{TV}(\mu_T,\nu_T)
\le \sum_{n<T}\sum_{x,a}w_n(x,a)
\operatorname{TV}(q_\#P_n(x,a),Q_n(q(x),a)).
\]

\textbf{Proof.} Every actual source-observation row is concentrated on S, so M-36.6 preserves its TV distance from the corresponding reference row. Apply the actual path bound on finite Z to \(\bar q,\bar Q,\bar\pi\) and rewrite using the established equalities of path TV, row TV, and actual occupancy. \(\blacksquare\)

With error \(\varepsilon_n\) on certified rows, the same actual-occupancy budget in M-36.4 remains valid. Replacing actual uncertified occupancy by \(\kappa_n\) still requires \(\varepsilon_n\le1\). For a \textbf{measurable} event F in the original Y-history space and an envelope of the same microscopic E, event transport gives

\[
\Pr_{\rm actual}(E)\le
\min\{1,\nu_T(F)+B_T\}.
\]

Here \(\nu_T(F)\) is the reference event probability under the original Q and π. Measurability of F replaces the measurability that was automatic for finite Y.

The basic encoding statement M-36.6 applies to general probability measures, including nonatomic references. At this chapter's path level, both source and reference remain PMFs. Infinite Y permits reference PMFs with countably infinite support; it does not mean that a path interface for arbitrary nonatomic continuous kernels has been constructed.

The first version uses the full \(q(X)\), retaining all source states and globally certified rows. Using a smaller image reachable under a particular policy would additionally require an actual-support induction, a revised certification domain, and proof that the old and new occupancies agree. Calling the omitted states unreachable cannot replace these source proofs.

\textbf{Lean scope.} The production interfaces in this edition include exact encoding in \nolinkurl{FiniteObservationEscapeCheck}, the same-policy full-path identity \nolinkurl{PathObservationEscapeCheck.escape_path_tv}, and the original-space path, certified-occupancy, and microscopic-event risk endpoints in \nolinkurl{ArbitraryObservationRiskCheck.} \nolinkurl{FiniteObservationEscapeBoundary} also proves the positive-escape numerical boundary above. The original finite-Y interfaces are retained; the new endpoints extend their scope through the source identities proved in this section. Encoding uses noncomputable definitions and does not add an already executed general algorithm. See \href{../formalization/ueot-core/docs/compression/theory_completion/unified_closure/interface_research_v1/FINITE_OBSERVATION_ESCAPE.md}{\nolinkurl{FINITE_OBSERVATION_ESCAPE}} for the written proof source, and \href{../formalization/ueot-core/docs/compression/theory_completion/unified_closure/REGISTERED_INTERFACES_CLOSURE_V1.md}{interface closure} and \href{LEAN_API_SCOPE.md}{\nolinkurl{LEAN_API_SCOPE}} for the current declarations and acceptance scope.

\chapter{37. Finite-Support Viability Under Partial Observation: An Equivalence Theorem and an Executable Check}\label{ch-37}

\textbf{Chapter question.} When can a controller reading only coarse history use certified rows throughout a finite horizon?

\textbf{Reading route.} First identify actual posterior support on histories of positive probability with the combinatorial recursion. Backward viability then selects one common action and covers every nonempty successor. Necessity and sufficiency use the same source. Certification of the current action does not imply that its prefix can be safely extended.

\section{37.1 The Support Recursion and Its Necessary Boundaries}\label{the-support-recursion-and-its-necessary-boundaries}

Continue to fix \(p_0,q,P\) from Chapter 36. Certification concerns the finite horizon \(0\le n<T\) and the state--action rows \(C_n(x,a)\) used at each step. No additional terminal safe-state condition is imposed by default. If the task requires a terminal set, the terminal condition of the recursion must be changed; the present theorem does not supply one automatically.

For a set of possible terminal states \(S\subseteq X\), a recorded action \(a\), and a next observation \(y\), define

\[
\operatorname{Succ}_n(S,a,y)=
\{x':q(x')=y\ \land\
\exists x\in S,\ P_n(x,a)(x')>0\}.
\]

The initial observation support is

\[
S_0(y)=\{x:p_0(x)>0,\ q(x)=y\}.
\]

Recursively, for full coarse histories, set

\[
S_{n+1}(g,(a,y))=\operatorname{Succ}_n(S_n(g),a,y).
\]

This trace support is determined by \(p_0,q,P\) and the recorded actions; policy probabilities are not part of its definition. It includes the complete observation constraints and cannot be replaced by an arbitrary representative or the whole terminal-observation fiber.

\textbf{Lean.} \nolinkurl{initialSupport}, \nolinkurl{successorSupport}, and traceSupport in \nolinkurl{SupportViabilityClosure.} The mathematical support bridge does not require \(Y\) to be enumerated as a finite set. Since \(X\) is finite, observations with actual mass always lie in the finite image \(q(X)\).

\section{37.2 Equivalence Between the Support Recursion and Actual Probability}\label{equivalence-between-the-support-recursion-and-actual-probability}

Define microLaw and coarseLaw as in Chapter 36. Every result below concerns an arbitrary fixed full policy \(\pi\).

\hypertarget{m-37-1}{}\phantomsection\label{m-37-1}

\textbf{M-37.1: The actual support bridge at positive histories.}

\begin{enumerate}
\def\labelenumi{\arabic{enumi}.}
\item
  An initial coarse observation has positive mass if and only if \(S_0(y)\) is nonempty.
\item
  If a coarse history \(g\) has positive mass, then

  \[
  S_n(g)=\operatorname{supp}\rho_n^g.
  \]
\item
  For the next coarse history \((g,(a,y))\),

  \[
  \bar p_{n+1}(g,(a,y))>0
  \iff
  \bar p_n(g)>0\ \land\ \pi_n(g)(a)>0\
  \land\operatorname{Succ}_n(S_n(g),a,y)\ne\varnothing.
  \]
\item
  At such a positive successor history, the actual terminal-state posterior support is exactly \(\operatorname{Succ}_n(S_n(g),a,y)\).
\end{enumerate}

\textbf{Proof.} A finite sum of nonnegative terms is positive if and only if at least one term is positive. A product of two finite nonnegative probability factors is positive if and only if both factors are positive. The microscopic next-step support therefore decomposes into the previous microscopic support, policy action support, and source-row support. Induction on coarse histories shows that \(x\in S_n(g)\) holds exactly when an actual positive-mass microscopic history projects to \(g\) and has terminal state \(x\).

The reverse construction in this induction particularly uses the same coarse-history policy. Replacing a microscopic prefix by any other prefix in its fiber leaves the action with positive mass under the same \(\pi_n(g)\), so it can be recombined with a source successor in the support. Finally, the existing support theorem for the fiber posterior converts these realization sets into posterior supports. \(\blacksquare\)

\textbf{Lean.} \nolinkurl{mem_microLaw_succ_support_iff}, \nolinkurl{mem_coarseLaw_zero_support_iff}, \nolinkurl{mem_traceSupport_iff_realized}, \nolinkurl{mem_coarseLaw_succ_support_iff}, \nolinkurl{traceSupport_eq_fiberTerminalPosterior_support}, and \nolinkurl{successorSupport_eq_fiberTerminalPosterior_support.}

The theorem identifies supports only. It makes no claim that different policies have identical complete posterior probabilities at a common reachable coarse history. At a zero-mass history, trace support remains combinatorially defined, but actual posterior completions may differ; claims 2 and 4 do not extend to those histories.

\hypertarget{m-37-2}{}\phantomsection\label{m-37-2}

\textbf{M-37.2: A supportwise characterization of zero uncertified occupancy.}

\[
\nu_n(C_n^c)=0
\iff
\forall g\in\operatorname{supp}\bar p_n\
\forall a\in\operatorname{supp}\pi_n(g)\
\forall x\in S_n(g),\ C_n(x,a).
\]

\textbf{Proof.} First, the pushforward definition from the actual prefix--action PMF gives

\[
(x,a)\in\operatorname{supp}\nu_n
\iff
\exists g\in\operatorname{supp}\bar p_n,\
a\in\operatorname{supp}\pi_n(g),\ x\in S_n(g).
\]

The finite nonnegative sum \(\nu_n(C_n^c)\) is zero if and only if every positive-mass state--action pair is certified. Substitute the support characterization. \(\blacksquare\)

\textbf{Lean.} \nolinkurl{mem_deploymentRowLaw_support_iff}, \nolinkurl{zero_uncertified_mass_iff_support}, and \nolinkurl{zero_uncovered_iff_support_certified.} Correctness of the actual support is proved internally, rather than supplied by the user as a matching hypothesis of the final theorem.

\section{37.3 Backward Support-Viability Recursion}\label{backward-support-viability-recursion}

Let \(V(n,k,S)\) express support viability starting at time \(n\), with \(k\) action steps remaining and current possible terminal-state set \(S\):

\[
V(n,0,S):=\mathrm{True},
\]

\[
\begin{aligned}
V(n,k+1,S):\iff \exists a\in A:\;&
\bigl[\forall x\in S,\ C_n(x,a)\bigr]\\
&\land
\bigl[\forall y,\ \operatorname{Succ}_n(S,a,y)\ne\varnothing
\\
&\qquad \Rightarrow V(n+1,k,\operatorname{Succ}_n(S,a,y))\bigr].
\end{aligned}
\]

There must be \textbf{one and the same action} available to all currently possible states. The successor condition must cover \textbf{every observation branch with positive-probability possibility}. Replacing the first condition by \(\forall x\exists a\) would permit access to an unobserved hidden state. Replacing the second by existence of one good successor would admit other positive-mass failure branches.

\textbf{Lean.} viableWith and Viable. Structural recursion uses the remaining step count \(k\), while the source and certification may vary with absolute time \(n\).

To obtain a full policy defined everywhere, supply an explicit base policy \(\pi_{\rm base}\). Define \(\pi^V\) as follows: if \(n<T\) and \(V(n,T-n,S_n(g))\), choose one common action from the existential clause and put a Dirac distribution on it; otherwise use \(\pi_{\rm base}(n,g)\).

This is a mathematical selection construction. It does not assert that every zero-mass history uses the fallback: a zero-mass history satisfying \(V\) may also use a recursively selected action. Properties at all actual positive histories are unaffected by those choices. A nonempty action type or an actual PMF fallback is an input needed for a full policy; an empty action space cannot be concealed by an instruction to choose an arbitrary action.

\section{37.4 Necessary and Sufficient Conditions for Policy Existence}\label{necessary-and-sufficient-conditions-for-policy-existence}

\hypertarget{m-37-3}{}\phantomsection\label{m-37-3}

\textbf{M-37.3: Finite-support viability is equivalent to an actual zero-occupancy policy.} Given a full \(\pi_{\rm base}\), the following statements are equivalent:

\[
\forall y,\ S_0(y)\ne\varnothing\Rightarrow V(0,T,S_0(y));
\]

\[
\exists\pi,\ \forall n<T,\ \nu_n^\pi(C_n^c)=0.
\]

On the right, \(\nu_n^\pi\) is the deployment occupancy law actually generated by the same \(p_0,q,P\) and the quantified policy \(\pi\). A reference law is not fixed first while the controller is changed arbitrarily afterward.

\textbf{Proof of sufficiency.} Use \(\pi^V\) from the preceding section. Positive initial coarse histories satisfy the hypothesis. If an actual positive history at time \(n\) satisfies the appropriate \(V\), the sole action in the policy support is the common certified action chosen by the recursion. By M-37.1, every actual positive successor corresponds to a nonempty Succ; the universal successor clause therefore preserves \(V\) there. Forward induction covers all actual positive histories, and M-37.2 gives \(\nu_n(C_n^c)=0\) at each step.

\textbf{Proof of necessity.} Take a policy as on the right. At any actual positive history, its action PMF has at least one positive-mass action. M-37.2 certifies that action jointly for all states in \(S_n(g)\). For every nonempty Succ, M-37.1 gives positive actual mass to the corresponding coarse successor, so the induction hypothesis applies to the remaining steps. This yields \(V(n,T-n,S_n(g))\); apply it to the initial supports.

Both directions use support bridges for the same actual law. Neither circularly assumes that the policy to be constructed has the posterior of a reference policy. \(\blacksquare\)

\textbf{Lean.} \nolinkurl{viablePolicy_support_step}, \nolinkurl{viablePolicy_supported_histories_viable}, \nolinkurl{viablePolicy_zero_uncovered}, \nolinkurl{supported_viable_of_zero_uncovered}, and \nolinkurl{initial_viability_iff_exists_zero_uncovered_policy.}

For \(T=0\), the left terminal condition is true; on the right, the explicit base policy suffices and there are no times requiring certification. The theorem also shows that, for this objective of almost-sure stepwise certification, success of a randomized history policy implies success of a deterministic history policy selected by the support recursion. It implies neither infinite-horizon viability, reward optimality, learning optimality, nor robustness to approximate support tests.

\section{37.5 Executable Decision on Finite Relation Tables}\label{executable-decision-on-finite-relation-tables}

The mathematical PMF source \(P_n\) may be defined noncomputably. Finiteness of \(X\) alone does not make its support test executable. The execution layer therefore supplies an explicit Boolean relation table \(R_n(x,a,x')\) and proves

\[
R_n(x,a,x')\iff P_n(x,a)(x')>0.
\]

This is an exact representation theorem. Failure to observe \(x'\) in a finite sample does not establish this equivalence.

For a finite set \(S\), enumerate successors using \(R\), group them by observation equality, and then test recursively over actions whether one common certified action has every nonempty successor pass. Even if \(Y\) is not enumerated as a finite type, only \(q(X)\) must be traversed.

\hypertarget{m-37-4}{}\phantomsection\label{m-37-4}

\textbf{M-37.4: Correctness of the executable check.} If the relation table represents the support of the actual \(P\) exactly, the finite-relation check returns true if and only if \(V(n,k,S)\) holds. Hence it returns true for all nonempty initial supports if and only if the actual zero-occupancy policy in M-37.3 exists.

\textbf{Proof.} Induct on \(k\). Exactness of the relation table identifies every constructed finite successor set with Succ. Finite action enumeration represents the existential quantifier, and enumeration of the finite observation image represents all nonempty successors. Rewriting the Boolean and logical connectives at each level proves equivalence with \(V\); then apply M-37.3. \(\blacksquare\)

\textbf{Lean.} \nolinkurl{finiteViabilityCheck} and \nolinkurl{finiteViabilityCheck_correct.} The actual executable layer consists of \nolinkurl{finiteRelationSuccessor}, \nolinkurl{finiteRelationViabilityCheck}, and \nolinkurl{finiteRelationViabilityCheck_correct.} It first proves equality of its output with the former check and then reuses that check's correctness, without rebuilding the policy-existence proof.

The following pseudocode gives the complete control structure. It requires exactSupport to be a proved exact relation table, an enumerable image of observe, and decidable certified.

\begin{verbatim}
ViabilityCheck(n, k, S):
    if k == 0:
        return true
    for a in Actions:
        if not all(certified(n, x, a) for x in S):
            continue
        passes = true
        for y in image(observe, States):
            Next = {z in States |
                    observe(z) == y and
                    exists x in S: exactSupport(n, x, a, z)}
            if Next is nonempty and not ViabilityCheck(n+1, k-1, Next):
                passes = false
                break
        if passes:
            return true
    return false
\end{verbatim}

Results may be cached by \((n,k,S)\); there are at most \(2^{|X|}\) state supports. This exponential support space is a substantive cost of partial observation. The current Lean correctness theorem claims neither a new polynomial-complexity result nor measured performance on large instances merely from a symbolic algorithm.

\section{37.6 A Strict Improvement over the Entire Current Observation Fiber}\label{a-strict-improvement-over-the-entire-current-observation-fiber}

Take four microscopic states \((b,s)\in\{0,1\}^2\): \(b=0\) is the initial phase, \(b=1\) the subsequent phase, and \(s\) a hidden side. The initial law assigns each of \((0,0),(0,1)\) mass \(1/2\). Both source actions move an initial-phase state to the subsequent phase on the same side. In the subsequent phase, the matching action \(a=s\) retains the state, while a nonmatching action returns it to the initial phase on the same side. There are three observation values:

\[
q(0,0)=\mathrm{some}(0),\qquad
q(0,1)=\mathrm{some}(1),\qquad
q(1,0)=q(1,1)=\mathrm{none}.
\]

The first action is fixed at zero. The second remembers the initial observation and selects side \(s\). All actions are certified in the initial phase, but only \(a=s\) is certified in the subsequent phase.

\hypertarget{m-37-5}{}\phantomsection\label{m-37-5}

\textbf{M-37.5: Actual posterior support is strictly less restrictive than the whole current fiber.} The two actual positive-mass coarse histories

\[
g_L=(\mathrm{some}(0),(0,\mathrm{none})),\qquad
g_R=(\mathrm{some}(1),(0,\mathrm{none}))
\]

have mass \(1/2\) each. Their actual terminal-state posterior supports are \(\{(1,0)\}\) and \(\{(1,1)\}\), respectively, and the memory policy satisfies \(\nu_1(C_1^c)=0\). Yet there is no common certified action for the entire current observation fiber \(q^{-1}(\mathrm{none})\).

\textbf{Proof.} Both initial branches have mass \(1/2\), and the first action and transition are deterministic, producing exactly the two positive histories above. The first can arise only from the left state, and the second only from the right. The existing fiber-posterior support theorem gives the respective singleton supports. On each support the policy chooses its matching side, so the zero-leakage case of M-36.5 gives zero actual uncovered mass. The whole current fiber contains both states; one action cannot be both zero and one. \(\blacksquare\)

Furthermore, \(g_\varnothing=(\mathrm{none},(0,\mathrm{none}))\) has an impossible initial observation and hence mass zero. Even assigning an empty admissible action set to that branch, with conditional action leakage one, does not affect the actual \(\nu\). This also checks the zero-mass boundary.

\textbf{Lean.} \nolinkurl{two_positive_coarse_histories}, \nolinkurl{posterior_supports_are_singletons}, \nolinkurl{both_terminal_states_reachable}, \nolinkurl{actual_uncovered_probability_zero}, \nolinkurl{whole_current_fiber_has_no_common_action}, and \nolinkurl{null_history_boundary} in \nolinkurl{PosteriorSupportStrictWitness.} The current public endpoint formalizes singleton posterior support. The book does not report a separate independently audited declaration that the posterior PMF equals a Dirac mass, nor a hand-calculated memoryless-policy risk consequence as an endpoint compiled in this check.

\section{37.7 A Viable Problem Can Still Contain a Locally Safe Dead End}\label{a-viable-problem-can-still-contain-a-locally-safe-dead-end}

Take \(X=A=\{0,1\}\) and a constant singleton observation. The initial state is fair. The first actual transition is \(X_1=A_0\); later transitions retain the state. Certification is defined by

\[
C_n(x,a)\iff(n=0)\ \lor\ (x=0).
\]

The horizon is \(T=2\). Every action is certified at the first step; at the second, actions are certified only at state zero.

\hypertarget{m-37-6}{}\phantomsection\label{m-37-6}

\textbf{M-37.6: Existence of a full viable policy does not make every certified prefix extendible.} The full policy always choosing zero has zero uncertified occupancy at both steps. Choosing action one at the first step is also a certified prefix, but it actually reaches support \(\{1\}\), from which no action is certified at the second step.

\textbf{Proof.} First action zero gives \(X_1=0\), so every second-step action is certified. First action one gives \(X_1=1\), and \(C_1(1,a)\) is false for every action. First-step certification itself is independent of the action. The bad prefix is therefore certified but has no safe extension. \(\blacksquare\)

\textbf{Lean.} \nolinkurl{a_complete_zero_uncovered_policy_exists}, \nolinkurl{deadend_policy_second_step_not_covered}, and \nolinkurl{viable_source_with_a_covered_deadend} in \nolinkurl{SafePrefixDeadend.} The passed relation check actually evaluated:

\[
\begin{array}{c|c}
\text{Input}&\text{Output}\\ \hline
(n=0,k=2,S=\{0,1\})&\mathrm{true}\\
(n=1,k=1,S=\{1\})&\mathrm{false}\\
(n=1,k=1,S=\{0\})&\mathrm{true}
\end{array}
\]

The counterexample distinguishes two quantifier patterns. Requiring every already covered partial policy to admit a next-step extension is a stronger sufficient contract. Backward \(V\) instead excludes dead-end actions at the current step and gives a necessary and sufficient condition. It cannot be paraphrased as saying that checking only the current action's safety is sufficient.

\chapter{38. From Registered Sampling to Actual Joint Risk After Data-Dependent Selection}\label{ch-38}

\textbf{Chapter question.} How can a model and policy selected after training still receive a guarantee on actual deployment failure?

\textbf{Reading route.} Construct the training--selected-deployment joint law, then split it over one Good event and its complement. Rowwise calibration, the selected policy's own occupancy, and a common failure event are connected on Good. Bad samples are charged by their actual probability mass; they need not all satisfy a small budget.

\section{38.1 The Joint Experiment of Training Records and Future Paths}\label{the-joint-experiment-of-training-records-and-future-paths}

Let the training space \(\Omega\) and the future-path space \(H\) used in this chapter be finite, with training probability \(\mu\in\mathrm{PMF}(\Omega)\). Candidates are indexed by \(c\in\mathcal C\), each specifying an actual future-deployment PMF \(K_c\in\mathrm{PMF}(H)\). The selector \(s:\Omega\to\mathcal C\) may use the entire training record.

Define the training--future joint law directly by

\[
J(\omega,h)=\mu(\omega)K_{s(\omega)}(h).
\]

This means that the training record is actually generated first, followed by a future path drawn from the selected conditional deployment kernel. Training and the selected model may be dependent; the selector need not be independent of training. Finiteness of \(\Omega\) makes the selected finite future law measurable. The candidate index itself need not be finite.

\hypertarget{m-38-1}{}\phantomsection\label{m-38-1}

\textbf{M-38.1: Marginals and the tower formula for the actual joint law.} The law \(J\) is a probability distribution with training marginal \(\mu\) and future marginal \(\sum_\omega\mu(\omega)K_{s(\omega)}\). For a common future event \(E\subseteq H\),

\[
J(\Omega\times E)
=\sum_\omega\mu(\omega)r(\omega),
\qquad r(\omega)=K_{s(\omega)}(E).
\]

\textbf{Proof.} Summation over \(h\) uses total mass one of every \(K_c\). Summation over \(\omega\) gives the future marginal. Two finite sums of the event indicator give the tower formula. If training is originally given as a probability measure on a finite space, \nolinkurl{Measure.toPMF} and its inverse toMeasure identity ensure that the same training law is used. \(\blacksquare\)

\textbf{Lean.} \nolinkurl{selectedTrainingDeployment_mass}, \nolinkurl{selectedTrainingDeployment_training_marginal}, \nolinkurl{selectedTrainingDeployment_path_marginal}, \nolinkurl{selectedTrainingDeployment_failure_tower}, and \nolinkurl{selectedTrainingDeployment_failure_integral.} The measurable interfaces are \nolinkurl{finite_training_selection_measurable} and \nolinkurl{selectedDeployment_measure_measurable.}

Fixing the common event \(E\) matters. Different macro representations may use different \(F_c\), but each must cover the same microscopic failure \(E\). Model selection cannot simultaneously redefine failure.

\section{38.2 Checking the Budget Only on the Calibration Good Event}\label{checking-the-budget-only-on-the-calibration-good-event}

Let \(G\subseteq\Omega\) be a good event. Suppose that on \(G\) the actual conditional risk satisfies \(r(\omega)\le b(\omega)\).

\hypertarget{m-38-2}{}\phantomsection\label{m-38-2}

\textbf{M-38.2: A joint bound with a data-dependent budget.}

\[
J(\Omega\times E)
\le
\mu(G^c)+\sum_{\omega\in G}\mu(\omega)b(\omega).
\]

No numerical budget is required of \(b\) outside the good event, and no separate assumption makes \(b\) nonnegative everywhere on \(\Omega\). The proof uses it only on \(G\).

\textbf{Proof.} Split the sum in M-38.1 over \(G\) and \(G^c\). Substitute the conditional bound on the good event and use risk at most one on the bad event. \(\blacksquare\)

If on \(G\) one has the uniform bound \(r\le\beta\), with \(0\le\beta\le1\), and if \(\mu(G^c)\le\alpha\), then:

\hypertarget{m-38-3}{}\phantomsection\label{m-38-3}

\textbf{M-38.3: The exact mixture form for a fixed good-event budget.}

\[
J(\Omega\times E)
\le\min\{1,\alpha+(1-\alpha)\beta\}.
\]

\textbf{Proof.} Write \(q=\mu(G^c)\) for the actual bad mass. M-38.2 gives \(q+(1-q)\beta=\beta+(1-\beta)q\). Since \(1-\beta\ge0\), substitute \(q\le\alpha\) and take the minimum with the probability bound one. If \(\alpha>1\), the right-hand side is clipped to one: the statement remains valid but is uninformative. \(\blacksquare\)

\textbf{Lean.} \nolinkurl{selectedTrainingDeployment_data_dependent_bound}, \nolinkurl{selectedTrainingDeployment_good_bad_bound}, and \nolinkurl{selectedTrainingDeployment_failure_le.} This probability algebra is the standard mixture bound. The substantive new connection is that \(G\) supplied by the existing statistical certificate and the constructed actual \(J\) share the source, training law, selected policy, and failure event.

\section{38.3 Fixed Registered Rows and the Canonical Calibration Event}\label{fixed-registered-rows-and-the-canonical-calibration-event}

Fix a registered source-row set \(\mathcal R\) before training. Each row \(j\) has an actual finite response law \(p_j\), from which \(N_j\ge1\) independent identically distributed responses are sampled, giving empirical law \(\widehat p_j(\omega)\). Independence between different source rows is not required by the finite union bound itself, though a product-sampling model may additionally provide it.

For radii \(\rho_j\ge0\), the canonical good event is

\[
G=\{\omega:\forall j\in\mathcal R,\
\operatorname{TV}(p_j,\widehat p_j(\omega))\le\rho_j\}.
\]

For response-space cardinality \(m\), a finite event union bound and Hoeffding's inequality give

\[
\mu(G^c)\le
\alpha:=\sum_{j\in\mathcal R}
2\cdot2^m\exp(-2N_j\rho_j^2).
\]

The upper bound may exceed one, so numerical informativeness must be checked in applications. Enumeration of all response subsets gives a direct bound matching the current formal certificate; no optimal concentration constant is claimed.

If an earlier statistical certificate first supplies a good event \(G_0\) satisfying these rowwise properties, then \(G_0\subseteq G\), so \(\mu(G^c)\le\mu(G_0^c)\). No additional failure probability is charged.

\textbf{Lean.} The original equal-sample-count version uses \nolinkurl{measure_finite_row_calibration_bad_le}, \nolinkurl{coverage_registered_source_calibration}, and the canonical Good of \nolinkurl{RegisteredGoodJointCheck.} Section 38.10 connects rowwise budgets to the dependent row-indexed sample types in \nolinkurl{VariableRegisteredSamples} and \nolinkurl{VariableRegisteredJointCheck.} See Appendix F.3.1 for the exact interface scope of the arbitrary-Y extension.

The definition of canonical \(G\) contains the unknown true source \(p_j\). It is therefore an actual calibration event used in the analysis. Writing it as a finite conjunction does not make its truth directly observable after training. An executable selection rule must be defined on all records; the theorem uses calibration only on records that do belong to \(G\).

\section{38.4 From Rowwise Fit to Selected Conditional Risk}\label{from-rowwise-fit-to-selected-conditional-risk}

For each training record one may select a quotient map \(q_\omega\), an initial law, a reference kernel \(\widehat P^\omega\), and a causal policy \(\pi^\omega\). The microscopic source \(P\) retains its identity, and the registered rows remain those fixed before training.

Assume only on \(G\) that:

\begin{enumerate}
\def\labelenumi{\arabic{enumi}.}
\tightlist
\item
  On registered source rows, the empirical microscopic response pushed through \(q_\omega\) has TV fitting error at most \(\eta_n(\omega)\) relative to the selected reference row.
\item
  The selected actual policy's uncertified occupancy satisfies \(\nu_n^\omega(C_n^c)\le\kappa_n(\omega)\).
\item
  The macro event \(F_\omega\) covers the common microscopic failure \(E\), and reference risk plus path budget is at most some \(\beta\).
\end{enumerate}

Comparison of a single true row with its selected reference row requires only

\[
\varepsilon_n(\omega)=\min\{1,\rho_n+\eta_n(\omega)\}.
\]

The first segment calibrates the actual microscopic row against its empirical row, with no increase of TV under Markov pushforward. The second fits the empirical pushforward to the reference row. The resulting allowance is \(\rho+\eta\); there is no justification for adding a second \(\rho\). A different comparison chain through representative rows must be charged according to its actual steps.

By M-36.4, the actual conditional risk on \(G\) satisfies

\[
r(\omega)\le b(\omega):=
\min\left\{1,\
\Pr_{\widehat p_T^\omega}(F_\omega)+
\sum_{n<T}\bigl[
\varepsilon_n(\omega)+
(1-\varepsilon_n(\omega))\kappa_n(\omega)
\bigr]\right\}.
\]

If \(b(\omega)\le\beta\) for all \(\omega\in G\), M-38.3 gives the total risk bound. The selector may change the model after training. The same \(G\) already calibrates all registered rows simultaneously, so selection of one candidate does not require another union bound over candidates.

\textbf{Lean.} \nolinkurl{coverage_empirical_registered_microscopic_failure_risk} and \nolinkurl{registered_selected_joint_failure_of_cap_on_good.} This endpoint requires hfit, hcap, and hchecked only on Good. The data-dependent version directly consumes M-38.2 and need not replace all training records by one common \(\beta\).

The reference event probability and initial PMF remain inputs of the selected model. Without a separate computational or interval-certification procedure, the entire mathematical expression cannot be described as a quantity computed solely from observable data.

\section{38.5 Eliminating the Coverage Input Through Support Viability}\label{eliminating-the-coverage-input-through-support-viability}

\hypertarget{m-38-4}{}\phantomsection\label{m-38-4}

\textbf{M-38.4: Support viability in the statistical joint-risk bound.} Suppose that, on the same \(G\), every selected candidate's actual \(p_0^\omega,q_\omega,P,C\) has all nonempty initial observation supports satisfying the backward viability condition of M-37.3. Use the corresponding viablePolicy and retain a full fallback policy as input. No separate \(\kappa\) or upper bound on actual \(\nu\) is then needed: the theorem internally derives

\[
\forall n<T,\quad \nu_n^\omega(C_n^c)=0.
\]

If the empirical fit and reference budget of that same selected policy satisfy, on \(G\),

\[
\min\{1,\Pr_{\widehat p_T^\omega}(F_\omega)
+\sum_{n<T}\varepsilon_n(\omega)\}\le\beta,
\]

then the actual joint risk still satisfies \(\min\{1,\alpha+(1-\alpha)\beta\}\).

\textbf{Proof.} Apply \nolinkurl{viablePolicy_zero_uncovered} record by record to obtain zero coverage deficit for its actual source kernel and policy. Substitute \(\kappa=0\) into the preceding section's Good-only endpoint. The good event remains the same \(G\), so \(\alpha\) is charged only once. \(\blacksquare\)

\textbf{Lean.} \nolinkurl{viableFamilyPolicy} and \nolinkurl{viable_registered_selected_joint_failure_on_good} in the original \nolinkurl{ViableRegisteredJointCheck} are established production interfaces. In Section 38.10, \nolinkurl{ViableVariableRegisteredJointCheck} reuses the same \nolinkurl{viableFamilyPolicy} to connect rowwise samples and the selected Good-event envelope, internally deriving that policy's own zero occupancy on the same Good event. Its exact version scope is given there. The sampling example in Section 38.6 instead proves actual zero uncovered occupancy directly; it is not presented as a specialized instance of this selected viability policy.

\section{38.6 A Complete Witness for Data Selection, Actual Deployment, and Joint Risk}\label{a-complete-witness-for-data-selection-actual-deployment-and-joint-risk}

Take \(X=A=\{0,1\}\), \(q=\mathrm{id}\), and horizon \(T=1\), with initial state zero. Next state one is failure. The source rows for the two actions satisfy

\[
p_0(1)=\frac18,\qquad p_1(1)=\frac14.
\]

Here the subscript in \(p_a\) denotes the action, not the initial law. The source rows are independent of the current state. Registered certification covers only the two action rows at time zero and current state zero.

Training draws \(N=256\) independent responses from each registered row and forms the full training PMF as the product of the two row experiments. A sample record therefore has \(512\) Boolean coordinates. Define the selection rule

\[
A(\omega)=\text{the first response of the action-0 source row}.
\]

The selector uses actual training data and is not claimed to optimize a learning objective. The reference model uses the empirical law from all \(256\) responses in each action row, whereas the deployed action depends only on the first coordinate above.

\subsection{38.6.1 Selection Changes the Actual Deployment Law}\label{selection-changes-the-actual-deployment-law}

The first coordinate has actual marginal \(p_0\), so

\[
\Pr(A=0)=\frac78,\qquad \Pr(A=1)=\frac18.
\]

Both selection probabilities are strictly positive, and the respective actual future failure probabilities are \(1/8\) and \(1/4\). The selected deployment PMFs are therefore distinct. The candidate index may contain the full record because the reference model reads all empirical data; actual deployment still depends on \(A(\omega)\).

\subsection{38.6.2 Registered Coverage Comes from the Same Actual Experiment}\label{registered-coverage-comes-from-the-same-actual-experiment}

There is only one action step, the initial state is certainly zero, and both action rows are registered. Hence, for every training record,

\[
\nu_0^\omega(C_0^c)=0.
\]

This follows directly from the actual initial PMF and policy extension, not from a guess about a training-row distribution. The empirical reference is exactly the corresponding empirical row, so the fitting remainder is \(\eta=0\).

\subsection{38.6.3 The Good Event and Its Budget}\label{the-good-event-and-its-budget}

Take \(\rho=1/8\). On the canonical good event, the two empirical failure probabilities are bounded by

\[
\widehat p_0(1)\le\frac18+\frac18=\frac14,\qquad
\widehat p_1(1)\le\frac14+\frac18=\frac38.
\]

Whichever action the data select, reference failure is at most \(3/8\). Adding path error \(1/8\) gives at most

\[
\beta=\frac12.
\]

Two registered rows, four subsets of the Boolean response space, and the two-sided Hoeffding factor give

\[
\alpha
=2\cdot4\cdot2\exp\!\left(-2\cdot256\cdot(1/8)^2\right)
=16e^{-8}\le\frac1{16}.
\]

An existing numerical exponential bound proves the final step; a floating-point approximation does not replace the formal inequality. Thus

\[
\Pr_J(E)\le\frac12+8e^{-8}\le\frac{17}{32}<1.
\]

\subsection{38.6.4 Exact Actual Risk}\label{exact-actual-risk}

M-38.1 and the actual selection marginal give

\[
\Pr_J(E)=
\frac78\cdot\frac18+\frac18\cdot\frac14
=\frac9{64}.
\]

The constructed law is therefore an actual joint law of training and the selected future, not a constant-policy risk renamed adaptive selection. The value \(17/32\) is a conservative certificate from the general calibration method; \(9/64\) is the exact value in this fully specified finite model. Neither should be mislabeled as an optimal or sharp version of the other. If the two source rows are completely known, the direct conditional-risk bound \(1/4\) already gives a better general upper bound. The purpose of the statistical certificate is to verify its data-calibration chain, not to outperform exact calculation in a known model.

\subsection{38.6.5 Why a Small Budget Cannot Be Required for Every Record}\label{why-a-small-budget-cannot-be-required-for-every-record}

Consider the record in which all \(512\) training coordinates are one, denoted by \(\omega_{\rm all}\). Its actual probability is strictly positive; in the product model it equals

\[
\left(\frac18\right)^{256}\cdot
\left(\frac14\right)^{256}>0.
\]

At this record, action one is selected and both empirical rows are pure-failure distributions. The clipped \nolinkurl{reference-risk-plus-calibration} budget is therefore one, violating \(\beta=1/2\).

Thus the earlier requirement that every training record pass a small-budget check fails for this actual sampling experiment. The Good-only endpoint uses the budget only on the calibration event and charges bad records through their actual bad mass, covering precisely this case. The all-one record is indeed outside canonical \(G\).

\textbf{Lean.} \nolinkurl{selectedAction_law}, \nolinkurl{selectedAction_pmf}, \nolinkurl{selection_probability}, \nolinkurl{actual_deployments_distinct}, \nolinkurl{two_positive_deployment_choices}, \nolinkurl{checkedBudget_on_good}, \nolinkurl{confidenceBudget_eq}, \nolinkurl{confidenceBudget_small}, \nolinkurl{registered_selected_joint_bound}, \nolinkurl{actual_joint_failure_exact}, \nolinkurl{allTrue_checkedBudget}, \nolinkurl{all_record_budget_false}, \nolinkurl{allTrue_positive_probability}, and \nolinkurl{registered_selection_joint_witness} in \nolinkurl{RegisteredSelectionWitness.} Product sampling and the actual coordinate marginals and independence reuse \nolinkurl{StatisticalRiskWitness}; sample provenance from the target source is not assumed anew.

\section{38.7 Sharpness of the Mixture Bound and Necessary Conditions}\label{sharpness-of-the-mixture-bound-and-necessary-conditions}

\textbf{A sharp two-branch example.} The training bad event has mass \(1/4\) and the good event mass \(3/4\). Future failure probability is \(1/2\) on the good branch and one on the bad branch. The selector chooses distinct actual deployment laws on the two positive-mass branches, so

\[
\Pr_J(E)=\frac14+\frac34\cdot\frac12=\frac58.
\]

This attains \(\alpha+(1-\alpha)\beta\) in M-38.3. The Lean declarations \nolinkurl{sharp_joint_failure_probability}, \nolinkurl{sharp_mixture_is_attained}, and \nolinkurl{selected_deployments_are_nonconstant} provide the example. Thus the bad-event charge cannot be removed using only the general information \(\alpha,\beta\).

\textbf{A stronger conclusion for a fixed policy.} If the actual future law is completely independent of the training record and \(G\) has positive mass, then \(G\) is nonempty. A single good record gives the risk bound for that fixed law itself, so the same constant risk need not incur another charge \(\alpha\). The earlier \nolinkurl{StatisticalRiskWitness} certificate for fixed deployment risk \(143/256\) remains \(143/256\) under the new joint law; the weaker \(685/1024\) is not presented as the principal improvement.

\textbf{The conditional deployment law must be actual.} If a certificate uses a pure-safe future kernel while external deployment is changed to a pure-failure law, the former cannot constrain the latter. \nolinkurl{wrong_conditional_law_counterexample} supplies an actual construction of this necessary boundary. Knowledge of a training marginal and some unconditional candidate marginals also does not unconditionally identify the conditional kernel after data selection. The positive theorem concerns the explicitly constructed experiment. Identifying an external system with that generating experiment still requires model evidence.

\section{38.8 Actual Formal Checks and the Status of the Book's Proofs}\label{actual-formal-checks-and-the-status-of-the-books-proofs}

The checks of the added mathematics come from actual Lean compilation and printed axiom reports for fixed source files. The table records the checking history of early mathematical drafts for Chapters 37--38, preserving the successes, failures, and integration status of those batches. The scope adopted in current production modules is recorded separately in Section 38.10, Appendix F, and \nolinkurl{LEAN_API_SCOPE.} Historical draft receipts are not receipts for the present publication build, and descriptions of integration that was pending at the time should not be read as the current source status.

\begin{longtable}[]{@{}
  >{\raggedright\arraybackslash}p{(\columnwidth - 4\tabcolsep) * \real{0.3333}}
  >{\raggedright\arraybackslash}p{(\columnwidth - 4\tabcolsep) * \real{0.3333}}
  >{\raggedright\arraybackslash}p{(\columnwidth - 4\tabcolsep) * \real{0.3333}}@{}}
\toprule\noalign{}
\begin{minipage}[b]{\linewidth}\raggedright
Check
\end{minipage} & \begin{minipage}[b]{\linewidth}\raggedright
Actual result
\end{minipage} & \begin{minipage}[b]{\linewidth}\raggedright
Conclusions supported
\end{minipage} \\
\midrule\noalign{}
\endhead
\bottomrule\noalign{}
\endlastfoot
check-03: support bridges and policy equivalence & Full exit 0; 12 printed items used only the standard three axioms & The actual PMF connection and equivalence in M-37.1--M-37.3 \\
check-04: registered Good endpoint & This item exited 0; the first finite-check draft in the same batch failed & The Good-only endpoint passed; the batch retains its historical failed status \\
check-06: relation check and thin A→B bridge & Both exited 0; the check actually returned true, false, true & M-37.4, the dead-end witness, and M-38.4 \\
check-08: third draft of the complete registered-selection witness & Full exit 0; 12 items used standard axioms; no warnings & The actual sampling, selection, joint risk, and all-one-record boundary in Section 38.6 \\
Previous round: second draft of the strict-posterior-support witness & This item exited 0; the first draft of a separate extension failed in that round & M-37.5; the failed history of the mixed batch is not rewritten \\
\end{longtable}

The standard axiom set is a subset of propext, Classical.choice, and Quot.sound; these printed results contain no sorryAx. They establish kernel checking of the corresponding declarations under the listed Lean environment and hypotheses. They establish neither empirical confirmation of the physical model nor a constructive executable program for every declaration. Exact source paths, fully qualified principal theorem names, compiled-source hashes, and relationships to historical inputs are given in Appendix F and the accompanying source index.

\section{38.9 The Scope of the Completed Mathematical Chain}\label{the-scope-of-the-completed-mathematical-chain}

The book has now connected the following steps individually: define actual finite paths from a supplied actual source and registered information; derive conditional posteriors and state--action occupancies from those paths; decide finite viability under partial observation by support recursion and construct a policy; combine its actual coverage with registered empirical calibration; construct the genuine joint law of training and the data-selected conditional deployment kernel; and finally bound the probability of a common microscopic failure event.

This chain still starts from explicit mathematical inputs: the actual source, observation protocol, exact support or an effective exact representation, training-generation protocol, failure event, and task-certification conditions. Application in an open world requires checking those inputs. Finite-support viability does not automatically learn unknown transition probabilities; finite-horizon zero uncovered occupancy is not an infinite-lifetime theorem; and a general risk certificate does not uniquely identify the source of purpose. Recording these boundaries in the theorems gives subsequent empirical and philosophical work precise interfaces to test.

\section{38.10 Actual Rowwise Sampling and the Selected Envelope on One Good Event}\label{actual-rowwise-sampling-and-the-selected-envelope-on-one-good-event}

This section connects the mathematical scope of the preceding three sections to the new actual types. \nolinkurl{VariableRegisteredSamples}, \nolinkurl{VariableRegisteredJoint}, and \nolinkurl{ViableVariableRegisteredJoint} are formalized in the production modules of this edition. Their joint endpoints allow arbitrary Y and explicitly require measurability of the selected event on the same Good event. Together with the four modules in Section 36.8, they constitute the probability-interface extension described here. The original common-Fin-N interface, existing witnesses, and historical compilation records for the first finite-Y version remain available.

\subsection{38.10.1 Realizing Unequal Sample Counts by an Actual Coordinate Experiment}\label{realizing-unequal-sample-counts-by-an-actual-coordinate-experiment}

Fix a finite registered row set J, a finite response space X, actual rows \(p_j\in\mathrm{PMF}(X)\), a sample allocation \(N:J\to\mathbb N\), and radii \(\rho:J\to\mathbb R\), with \(N_j>0,\rho_j\ge0\) for every actual row. These inputs are fixed before observing data. Define the coordinate type and product experiment by

\[
I_N=\{(j,k):j\in J,\ 0\le k<N_j\},\qquad
\Omega_N=X^{I_N},\qquad
\mu_N=\bigotimes_{(j,k)\in I_N}p_j.
\]

The sample \(S_{j,k}(\omega)=\omega(j,k)\) has actual marginal exactly \(p_j\), and samples within each row are independent. The actual empirical law is thus built from that row's \(N_j\) coordinates; a common-length array is not padded with repetitions that have no sampling interpretation.

\hypertarget{m-38-5}{}\phantomsection\label{m-38-5}

\textbf{M-38.5: Actual registered calibration with rowwise sample counts.} The experiment satisfies

\[
\mu_N(G^c)\le
\alpha_N:=\sum_{j\in J}2\cdot2^{|X|}
\exp(-2N_j\rho_j^2),
\quad
G=\{\omega:\forall j,\operatorname{TV}(p_j,\widehat p_j(\omega))\le\rho_j\}.
\]

The same conclusion holds for any actual training measure and sample family whose samples are measurable, independent within each row, and have the corresponding actual row marginals. Independence between different registered rows is not required.

\textbf{Proof.} For each j, apply single-row calibration with its actual \(N_j\), intersect the resulting good events, and apply a finite union bound to the complements. That intersection is contained in canonical G, so the bad mass of canonical G is bounded by the same budget. Coordinate-projection theorems for the product experiment separately provide measurability, correct marginals, and the required independence. \(\blacksquare\)

\textbf{Lean.} \nolinkurl{VariableRegisteredSamples.productRowLaw}, \nolinkurl{productRowSample_law}, \nolinkurl{productRowSample_independent}, \nolinkurl{registered_row_calibration}, and \nolinkurl{product_experiment_calibration.} The general combination of rowwise statements does not require a finite training space; the subsequent joint-deployment endpoint still uses one. If J is empty, the coordinate experiment remains well-defined, the budget is zero, and all rowwise conditions are empty conjunctions. This does not automatically give deployment nonempty viability or zero physical-failure risk.

\subsection{38.10.2 Using an Envelope Only for the Candidate Actually Selected on Good}\label{using-an-envelope-only-for-the-candidate-actually-selected-on-good}

Fix the production registration type \(J=\mathrm{CoverageRegisteredRow}(C,T)\), where each j specifies an actual time \(n_j<T\), state, action, and certification proof. The actual source P, certification condition C, horizon T, N, error radii, and microscopic failure event E are all fixed before data. A candidate c may contain an observation map, initial law, reference kernel, causal policy, and coarse event F. The choice \(c=s(\omega)\) may depend on the full training record.

For every \(n<T\), the current joint endpoint uses a nonnegative time-indexed radius \(\rho_n\). Its budget is therefore the special case of M-38.5 with \(\rho_j=\rho_{n_j}\). The more general rowwise radii of the Samples layer cannot be reported as part of the current Joint endpoint's full signature.

On the same canonical G, require the selected candidate's empirical fit, actual-occupancy bound, and numerical budget. The general-Y route of Section 36.8 additionally requires measurability of the coarse event F selected by that record. The envelope condition is only

\[
\forall\omega\in G,\qquad
E\subseteq H_{q_{s(\omega)}}^{-1}(F_{s(\omega)}).
\]

This is not an envelope condition on every candidate c, nor on selected candidates at bad training records. The failure event E remains the same fixed microscopic event; its meaning cannot be redefined after seeing the record.

\hypertarget{m-38-6}{}\phantomsection\label{m-38-6}

\textbf{M-38.6: Selected-Good joint risk with unequal sample counts.} If, on the common G above, the selected policy's reference event risk plus actual path budget is at most \(\beta\in[0,1]\) after clipping, then the actual training--selection--deployment joint law satisfies

\[
J(E)\le \min\{1,\alpha_N+(1-\alpha_N)\beta\}.
\]

\textbf{Proof.} M-38.5 gives \(\mu(G^c)\le\alpha_N\). At each \(\omega\in G\), derive the source-row error locally from that record's actual row calibration and empirical fit. Its selected policy's own actual occupancy and envelope then bound actual conditional deployment failure by β. Finally, split the same joint law over G and its complement. Apply the tower identity M-38.1 and the mixture bound M-38.3. There is no second envelope-good event to be charged. \(\blacksquare\)

\textbf{Lean.} \nolinkurl{VariableRegisteredJointCheck.RegisteredCalibrationGood}, \nolinkurl{registeredConfidenceBudget}, and \nolinkurl{registered_selected_joint_failure_of_cap_on_good.} The sample type is \(\forall j,\mathrm{Fin}(N_j)\to\Omega\to X\); hfit, hcap, hF, henvelope, and hchecked are all restricted to the same \nolinkurl{RegisteredCalibrationGood.} This interface retains finite Ω, X, and A; Y has a measurable structure with measurable singletons, and hF explicitly gives measurability of the selected coarse event on Good. The candidate type need not be finite.

The confidence budget consequently has a specific source: fixed registered rows, actual \(N_j\), the source response space X, and the radii. The original rows are first calibrated on X; a data-selected q then uses that same calibration by TV data processing. A large candidate type, or allowing infinite Y through the mathematical route in Section 36.8, does not itself create another selection-failure budget. This does not make every selected model correct without conditions: empirical fit, the selected event envelope, measurability of the original event, and the numerical risk check retain their distinct obligations and must meet on the same G.

\subsection{38.10.3 Eliminating Occupancy Internally for the Same Support-Viability Policy}\label{eliminating-occupancy-internally-for-the-same-support-viability-policy}

\hypertarget{m-38-7}{}\phantomsection\label{m-38-7}

\textbf{M-38.7: A viable policy in the rowwise-Good joint-risk bound.} Suppose that every nonempty initial observation support of the candidate selected in M-38.6 satisfies Chapter 37's Viable condition on the same G, and construct \nolinkurl{viableFamilyPolicy} using an explicit fallback. That policy's own uncertified occupancy is then zero at every step on G. Its joint-risk bound needs only the reference event risk plus

\[
\sum_{n<T}\min\{1,\rho_n+\eta_n(\omega)\},
\]

and satisfies the same mixture bound.

\textbf{Proof.} Fix a good record and apply the actual support-viability theorem to obtain zero occupancy of \nolinkurl{viableFamilyPolicy} at every \(n<T\). Substitute \(\kappa=0\) into M-38.6. The reference path, actual deployment kernel, and internal coverage proof all use that same policy. \(\blacksquare\)

\textbf{Lean.} \nolinkurl{ViableVariableRegisteredJointCheck.viable_registered_selected_joint_failure_on_good.} hinit, hfit, hF, henvelope, and hchecked share \nolinkurl{VariableRegisteredJointCheck.RegisteredCalibrationGood}; the caller no longer supplies hcap. This interface has the same observation-space and training-space scope as the preceding subsection. Eliminating the occupancy input neither proves zero physical failure nor creates a full policy for an empty action type; fallback remains an explicit input.

These connections do not cover correlated single-trajectory sampling, data-dependent changes to N after training, additional unregistered rows, or arbitrary nonatomic reference kernels. They connect the specified actual sampling and deployment conditions more tightly, removing a duplicated group of premises from the common-certificate assembly in Section 31.1 while retaining the substantive research questions.

\chapter{39. Common-source constructions for execution, faults, reproduction, and purpose}\label{ch-39}

\textbf{Chapter question.} Which program, path, birth, payment, and posterior inputs can be derived from one actual construction?

\textbf{Reading route.} This chapter records successive source constructions and is longer than the preceding chapters. Identify the source, clock, available information, and comparison policy class retained by each section before tracing the existing consumer it supplies. Similar terminology does not permit different source packages to be combined into a stronger model.

\textbf{Proof order and historical wording.} The sections retain the order in which sources were constructed and connected to existing consumers. Wording about completion, such as ``this round,'' ``subsequently,'' or ``still required,'' must therefore first be located at the source-construction stage recorded in the section. A stage obligation is resolved only when a later section explicitly supplies its source, theorem, and adopted scope. Model assumptions, quantifier restrictions, and uncovered cases do not disappear merely because the text is later revised. The table provides thematic entry points; the indicated sections retain the full conditions.

\begin{longtable}[]{@{}
  >{\raggedright\arraybackslash}p{(\columnwidth - 4\tabcolsep) * \real{0.3333}}
  >{\raggedright\arraybackslash}p{(\columnwidth - 4\tabcolsep) * \real{0.3333}}
  >{\raggedright\arraybackslash}p{(\columnwidth - 4\tabcolsep) * \real{0.3333}}@{}}
\toprule\noalign{}
\begin{minipage}[b]{\linewidth}\raggedright
Reading question
\end{minipage} & \begin{minipage}[b]{\linewidth}\raggedright
Main entry points
\end{minipage} & \begin{minipage}[b]{\linewidth}\raggedright
Connection to earlier chapters
\end{minipage} \\
\midrule\noalign{}
\endhead
\bottomrule\noalign{}
\endlastfoot
How an actual program generates faults, births, costs, and optimality & 39.1--39.6, 39.16, 39.18, 39.24 & Recovery, resources, and value bridges \\
How full history generates measurable prediction and minimal representations & 39.7, 39.11--39.12, 39.69, 39.82--39.88 & Chapters 4, 17, and 27 \\
How continuous noise, stochastic integrals, and change of measure acquire actual sources & 39.9--39.15, 39.17--39.19, 39.22--39.23, 39.26--39.31, 39.34--39.48 & Chapters 9, 19, and 22; apply each section only to its actual content \\
How the same parent, actual births, and population accounts are connected & 39.8, 39.20--39.21, 39.25, 39.32, 39.49--39.70 & Chapters 11, 25, and 26; related decision and continuous branches are interleaved \\
How the book's named sources return to existing consumers & 39.71--39.74 & Finite common-source examples, Perron theory, and spectral semigroups \\
How finite decisions are computed and implemented under unknown source parameters & 39.75--39.88 & Beliefs, wallets, computable representations, and the full comparison class \\
What remains guaranteed when the fault structure is unknown & 39.89--39.95 & Admissible experiments, nonidentifiability, and worst-case regret \\
\end{longtable}

\section{39.1 Conditions still supplied by the model}\label{conditions-still-supplied-by-the-model}

The preceding interfaces state program execution, repair, parent objects, resources, and evaluation separately. This chapter constructs a common nondegenerate source family and derives several bridges that were previously independent. It is not a proof of autopoiesis in an arbitrary world: the executor and interpreter are reliable, the retained program is protected, and copying inherits the program according to the declared rule. Environmental modes and birth labels follow specified independent probability laws. Action permissions, prices, discounting, initial history, and resource supply are also explicitly registered.

Let \(I\) be a finite register set with a designated execution bit \(i\in I\), and let the physical state be \(x\in\{0,1\}^I\). The retained program bit is \(b\). Allowed actions are only idle, write-false, and write-true; \(U_i(x,a)\) either preserves or writes bit \(i\) as specified. The actual interpreter first reads \(x_i\): it idles if \(x_i=b\) and otherwise executes write-\(b\). Interpretation of the finite instructions proves that their action is exactly this policy \(\pi_b\); the axiom of choice is not used to find another program achieving the same objective.

At each step, mode \(R_t\) has repair probability \(r=1-\epsilon\) and fault probability \(\epsilon\in[0,1]\). The physical transition is \[
X_{t+1}=\begin{cases}
U_i(X_t,A_t),&R_t=1,\\
X_t[i\leftarrow\neg b],&R_t=0.
\end{cases}
\] A fault sets the designated register to the opposite of the program target while preserving the program and controller. It does not flip an arbitrary damaged bit again; that alternative could accidentally repair damage and defines a different source. The pushforward of the source-mode PMF equals the physical fault/repair mixture row of P5 and the full physical row of this chapter's control model. Kernel equality is proved from execution definitions, not supplied by the caller.

\section{39.2 Repairability, actual births, and history}\label{repairability-actual-births-and-history}

In the joint organizational state, the retained program exactly characterizes the finite repair carrier. The actual potential and fault envelope are constructed from the transition above, providing a source for P5's execution and fault conditions. Membership in this carrier means that recovery remains possible. A register under a fault need not satisfy the physical target; carrier membership does not imply health at every later moment.

In repair mode, an actual offspring is produced only if the target bit is correct after the action; failure is recorded as None. The offspring literally copies the program and post-action physical state, and its history appends the current label to its parent's history. Actual birth therefore implies that the offspring satisfies the target and inherits the program; this same offspring can seed P5 with the same \(\epsilon\). There is no retry or deletion of failures conditional on success.

For an actual log of repair, damage, and birth instructions, the resulting history is the original history followed by the birth labels occurring in the log, in order. Thus the parent-identity bridge comes from execution history. Repair adds no birth label; birth strictly increases generation. Two unmarked \nolinkurl{physical/program} records may coincide despite different actual histories, so the former alone cannot recover the latter. Label history describes the lineage but does not guarantee globally unique offspring identities across times or concurrent populations.

This also corrects a potential confusion in a finite model. On a finite nonempty table of permanent individuals, if every positive reproductive edge strictly increases generation, then not every row can have reproductive mass one. An individual of maximal generation has no positive outgoing edge, contradicting a unit row sum. A valid finite Price application uses rectangular indices for parent types and offspring-event labels, obtaining a normalized transmission kernel from the actual birth-label law. Parent frequencies, reproductive intensities, and trait assignments remain inputs. Types and event labels are not permanent individuals; a strict individual genealogy cannot silently be replaced by a type chain. Sections 39.20--39.21 derive individual parent identities, a full queue of old parents, and actual offspring frequencies from newborn slots in a separately registered serial population model. They do not retain the impossible combination of a finite, normalized, strictly increasing genealogy.

\section{39.3 Reward and resources on the same record}\label{reward-and-resources-on-the-same-record}

Let a controlled write cost \(c\ge0\), let idle cost zero, and register price \(\lambda\ge0\). The actual record retains action, mode, label, parent, offspring, and cost together. Put \[
B_t=\mathbf1_{\{R_t=1\}}\mathbf1_{\{U_i(X_t,A_t)_i=b\}},\qquad
C_t=c\,\mathbf1_{\{R_t=1\}}\mathbf1_{\{A_t\ne\mathrm{idle}\}}.
\] \(B_t\) is the actual birth count for the transaction. In fault mode there is neither a controlled-write charge nor a birth. Physical damage, code copying, and rebinding are not separately charged in this cost model, so \(C_t\) must not be called total energy expenditure. For \(0<d<1\), the finite evaluation is \[
J_T(\pi)=E\!\left[\sum_{t<T}d^t(B_t-\lambda C_t)\right].
\] The current state in initial policy memory must match the parent's actual physical state, and its retained program must be the registered target; the designated seed of the compiled policy satisfies these conditions automatically. The complete finite-record PMF is generated step by step from the same policy, modes, and labels. Direct proof gives integrability of the record evaluation and its equality to the same policy's truncated value in the physical control model. Policies may use their original causal memory of physical states and actions. Recording modes and labels does not thereby make them policy observations.

Actual payment records also enter the original P6 resource account. For every actual action/mode sequence, \(0\le C_t\le c\) at each step. Thus nonnegative initial reserves and per-step supply \(s\ge c\) imply nonnegative cumulative balances on every path. Supply is a design input in this sufficient condition; it is not created by mean reward or reproductive success. Zero extension beyond an original finite record defines only a finite-horizon cost application. The added infinite interpreter extends the same history event by event and gives almost-sure prefix compatibility. Every finite prefix of its complete infinite-record law is exactly the earlier recordLaw. For this actual infinite cost stream, the same P6 supply condition holds pathwise without zero padding.

\section{39.4 Optimality of the interpreted policy itself}\label{optimality-of-the-interpreted-policy-itself}

Under the sufficient condition \(\lambda c\le1\), \(\pi_b\) is optimal in the original class of all causal policies, using the same program for every \(\epsilon\in[0,1]\). Its healthy and damaged state values are \[
G=\frac{r(1-d\epsilon\lambda c)}{1-d},\qquad H=G-r\lambda c.
\] The proof first averages actual one-step birth minus cost under the same mode law, then checks the Bellman inequality and equality for \(\pi_b\), and finally invokes the original causal-optimality endpoint. The condition is asserted only to be sufficient. The fault objective is freshly derived from actual records, rather than imported as a path integral from the earlier fault-free model.

For a single register, target \(b=\mathrm{true}\), fair birth labels, and \(\epsilon=1/4,c=1,\lambda=d=1/2\), we obtain \(G=45/32,H=33/32\). Starting damaged, the complete two-step tape with write-true then idle, repair then fault, and labels true then false has mass \(3/64\) and discounted net reward \(1/2\) (the actual write costs \(1\), converted by the price to \(1/2\)). The second-step failure remains in the record. The same program's finite-record expectations converge to the optimal control-model value. For the actual infinite interpreter, each reward has absolute value at most \(1+\lambda c\); geometric discounting gives absolute summability and integrability. Dominated convergence and the proved prefix laws then identify the actual infinite-record integral with the infinite value of the same original causal policy. The compiled program's actual integral is therefore the optimum above. Arbitrary Type0 internal memories of original policies are replayed by countable event words, without imposing finite memory. Modes and labels never enter the original policy's information permissions.

\section{39.5 A source for purpose under an unknown fault rate}\label{a-source-for-purpose-under-an-unknown-fault-rate}

Birth readouts for actual write-false and write-true, together with controlled-write cost readouts, give three exact means: \[
\bigl(r\mathbf1_{\{b=\mathrm{false}\}},\ r\mathbf1_{\{b=\mathrm{true}\}},\ rc\bigr).
\] For \(r>0\), the sum of the first two identifies \(r\), the nonzero entry identifies \(b\), and division of the third by \(r\) identifies \(c\). These means are computed from actual post-action births and payments, independently of the parent's current physical state and history; no reset permission is needed. Source-parameter identification then yields identical policy orderings for the complete-record evaluation, supplying the original C6 faithfulness application. Price, discount, and the evaluation contract remain registered inputs: the mechanism alone does not determine ethical evaluation.

At \(\epsilon=1\), all these means vanish and source parameters cannot be identified this way, but all registered control values also vanish. This does not prove nonidentifiability of the wholly indifferent purpose-order class. Access to exact experimental means is not finite-sample estimation, an executed inversion algorithm, or a post-selection risk guarantee. Connecting this mechanism to statistical calibration requires its own common application. The retained-program controller itself does not read the true fault rate.

\section{39.6 Verified sources and subsequent connections}\label{verified-sources-and-subsequent-connections}

The sources above are \nolinkurl{RegisterRepairSource}, \nolinkurl{RegisterRepairHistory}, \nolinkurl{RegisterFaultSource}, \nolinkurl{RegisterResourceSource}, \nolinkurl{RegisterFaultControlSource}, \nolinkurl{RegisterFaultControlModel}, \nolinkurl{RegisterFaultJointRecord}, \nolinkurl{RegisterBirthMaintenance}, \nolinkurl{RegisterFaultPurposeInformation}, and \nolinkurl{RegisterFaultInfiniteRecord.} They invoke the original P5/P6/P12/C6 and causal-control endpoints. Direct source mappings and verification receipts are in the \href{../formalization/ueot-core/docs/compression/theory_completion/unified_closure/register_fault_resources_v1/RESULTS_AND_OPEN_OBLIGATIONS.md}{fault-source package}.

For general prediction, \nolinkurl{MeasurablePredictiveSource/Sensor} derive the almost-everywhere projection of the actual posterior and same-policy recursion from measurable shared-noise updates and an observer that does not read hidden coordinates. Continuous noise and partial noisy observations are allowed. A strict leakage counterexample covers the case where hidden coordinates influence the next visible state and destroy visible autonomy. \nolinkurl{MeasurablePredictiveJointPath} then preserves the complete finite action--observation--posterior record of the same time-varying visible-belief policy. \nolinkurl{MeasurablePredictiveLatentPath} starts with \(X_0\sim b_0\), retains and transitions the actual latent state under the original P/O, and updates the same posterior from actual observations. A one-step disintegration identity gives compatibility of the randomized lift; induction over prefixes proves that the joint law of the complete observation record and terminal latent state disintegrates through the final actual posterior. This yields the finite observation-record law of actual latent execution and extends it to the full countable path by uniqueness on countable prefixes. It neither assumes pointwise lumpability for arbitrary \nolinkurl{belief/latent-state} pairings nor resamples the physical state from the belief at each step. The visible belief comes from the declared source; it is not claimed to be a minimal measurable predictive image for every process.

These results provide a nonempty common-source chain connecting objects, programs, actions, actual histories, births, resources, and registered purposes. The following sections construct sources for minimal prediction under actual finite protocols, full-history policies, timestamped identity, and continuous Gaussian processes. At the source-construction stage recorded in this section, outstanding connections included uncovered protocol and action classes, general continuous stochastic analysis, dependent or adaptive calibration and discovery, formation, lineage and scale transport in uncovered models, and complete long-term resource and global composition results. Sections 39.7--39.95 and the definitive core overview record the subsequent treatment and current scope. This historical list does not collectively replace current conclusions or remove the model and policy conditions that individual results still explicitly retain. Compilation, mathematical validity, and empirical support are reported separately; this chapter does not claim that a universal mechanism of nature has been established.

\section{39.7 Predictive recursion for finite protocols from a general actual source}\label{predictive-recursion-for-finite-protocols-from-a-general-actual-source}

Visible dynamics need not be autonomous. Given any nonempty standard Borel latent space \(Z\), original transition kernel \(P\), and observation kernel \(O\), let \(K_u\) be the full ordered observation law for the admissible action word \(u\): transition by \(P\), observe the next latent state by \(O\) at that step, then execute the remaining word. Repeated actions retain distinct positions, and the empty word returns the unique empty tuple. For prior \(b\), define \[
M_u(b)=K_u\circ b,\qquad C(b)=(M_u(b))_{u\in A^*}.
\] This retains the entire observation law for each finite action word, rather than a few empirical scores. Disintegration of the same actual P/O joint law directly gives \[
M_{a::u}(b)=\operatorname{Obs}_{b,a}\otimes
  \bigl(K_u\circ\operatorname{Posterior}_{b,a}\bigr).
\] The numerator and continuation factor come from actual execution of the longer word; the caller need not supply a separate Bayes factorization. If \(C(b)=C(c)\), the actual observation laws under the same first action agree, and subsequent posterior responses agree almost everywhere for every continuation word. When actions are countable, intersecting over the countable family of finite action words gives one probability-one set. That set is for fixed \(b,c,a\), not a version simultaneous over uncountably many priors.

This recursion of source responses requires only a countably generated observation space. If the observation space is additionally standard Borel, a uniform measurable disintegration of the response coordinates \(M_{a::u}(b)\) constructs a total response-array update. The update reads neither a latent state, source parameters, nor a section back to latent beliefs. With countable actions and measurable singletons, it is jointly measurable in state, action, and observation and yields an actual Markov kernel: pushing the same C1 posterior transition through \(C\) equals the next-response kernel computed from response arrays. Arrays outside the source image receive only a mathematical extension, not a claim of physical realizability. A total kernel does not empirically certify \nolinkurl{zero-observation-probability} branches.

An existing continuous Gaussian source, with explicitly allowed integer drift actions, supplies a nonatomic common application whose first observation from the origin is \(N(0,1)\). The countable-action result does not silently include all real actions, nor does it alone prove arbitrary adaptive-protocol or general minimal-predictive-image results. Its source is \nolinkurl{MeasurablePredictiveResponseRecursion} and its actual application is \nolinkurl{actual_C1_response_kernel_factor.}

\nolinkurl{MeasurablePredictiveCanonicalSource} starts from the already generated actual past and retained terminal latent variable, executes future action words with the same P/O, and derives that the past--future joint law is the actual past law multiplied by the future kernel supplied by the final posterior. This is an explicit future intervention: the future action word may differ from continuation of the policy that generated the past. The source then supplies the original \nolinkurl{PredictionDependent} endpoints for predictive sufficiency, event conditional expectations, and minimality modulo null sets. A competing statistic's ability to predict these actual future laws is the definition of its sufficiency, not a separately supplied realization certificate. Minimality is relative to the actual past law modulo null sets and applies for countable actions. The embedding from probability-response arrays into original Measure arrays concerns all normalized arrays; it neither requires a Borel physically realizable image nor constructs a section to latent beliefs.

\section{39.8 Offspring of the same parent at different times}\label{offspring-of-the-same-parent-at-different-times}

Original Bool labels may repeat. After two successful births by the same parent, the offspring's earlier label histories may be identical. The new source maintains an actual counter outside the original transaction interpreter. Every transaction, including faults and failures, increments the parent counter. An actual offspring identity appends the current counter and label to its parent's identity, and its own local counter starts at zero. The executed prefix then proves that successful births by the same retained parent at distinct transaction times have distinct identities; the parent's identity is recovered as the prefix of the actual offspring identity.

Erasing the added metadata from the entire timestamped record law gives exactly the original complete-record law for the same policy, initial state, fault rate, and label source. Actions, physical states, programs, payments, and failures are preserved. For the damaged single register, true program, fault rate \(1/4\), fair labels, actions \nolinkurl{write-true/idle/write-true}, modes \nolinkurl{repair/fault/repair}, and all three sampled labels true, the complete tape has mass \(9/512\). The intervening failure remains. The two actual offspring identities are \([(0,\mathrm{true})]\) and \([(2,\mathrm{true})]\), although both old Bool histories are \([\mathrm{true}]\).

\nolinkurl{RegisterTimestampedBirthSource} actually supplies the original P8 \nolinkurl{repair/offspring} distinction, P5 maintenance of the same offspring, and P12 parent-identity discovery. The last returns Option(parent ID) after a successful birth and retains None after failure; it does not identify the parent of a failed request. Reliable accessible serial metadata and initial founder \nolinkurl{identities/counters} remain declared conditions. The original controlled-write model does not separately charge metadata. This is neither historical knowledge automatically inferred from an unlabeled physical state nor a solution of founder allocation, global identity, or reproductive-intensity provenance for arbitrary concurrent populations.

\section{39.9 An actual all-time source from finite Gaussian laws}\label{an-actual-all-time-source-from-finite-gaussian-laws}

Continuous-time recovery interfaces also require an actual random source. Given a source with Brownian finite-dimensional laws, the time change \[
\tau(t)=\frac{e^{2kt}-1}{2k},\qquad
X_t=e^{-kt}\bigl(x_0+sB_{\tau(t)}\bigr),\quad k>0,
\] yields the actual pushforward law, deterministic initial condition, square integrability, and \[
E[X_t^2]=e^{-2kt}x_0^2+\frac{s^2}{2k}(1-e^{-2kt}).
\] This expectation is absolutely continuous, with almost-everywhere derivative \(E[-2kX_t^2+s^2]\). The original p\_rec\_02 uses the resulting actual Dynkin certificate. The right time derivative of this quadratic observable also comes from the source; it does not directly establish a general generator domain, conditional transition semigroup, or SDE identity.

The source assumption itself has an actual construction. \nolinkurl{BrownianCountableSource} starts from existing finite Gaussian laws, disintegrates adjacent finite prefixes, and iterates the existing trajectory kernels to construct a probability space for any countable time sequence. Every finite coordinate set has the original Gaussian law, and repeated times share a coordinate. The same dense-time source starts at zero almost everywhere and has unit variance at unit time.

\nolinkurl{BrownianProjectiveSource} extends the construction to all nonnegative real times. Each countable family of cylinder sets and its cylindrical union involves only countably many times. Realize all of them on the same countable source and use its genuine probability union bound to prove \(\sigma\)-subadditivity of cylinder content; the existing Carathéodory extension then applies. The resulting all-time probability law has all the original finite-dimensional Brownian Gaussian laws and actually supplies \nolinkurl{IsPreBrownianReal.} Neither a projective limit, subadditivity, nor a Brownian source is assumed as input. The same source enters the OU recovery construction, giving strictly positive variance at every positive time from every deterministic initial state when noise is positive.

Finite-dimensional laws, path continuity, and joint time measurability are separate proof obligations. The original coordinate process is not presumed almost surely continuous; the next section constructs its continuous modification. General Itô, controlled SDE/HJB, and Girsanov results remain obligations at this stage of the theory. Constructing a mathematical source does not empirically confirm a model of physical noise.

\section{39.10 An actual continuous modification on the same Gaussian source}\label{an-actual-continuous-modification-on-the-same-gaussian-source}

\nolinkurl{BrownianContinuousSource} derives from the source's Gaussian increments that \[
E[|B_t-B_s|^4]=3|t-s|^2.
\] On each compact interval with integer endpoint, use grids of spacing \(32^{-n}\), adjacent-increment thresholds \(2^{-n}\), and Markov's inequality to obtain summable bad-event probabilities. Over all integer intervals and all levels, construct one measurable probability-one event on which grid approximations converge uniformly. Refinement and oscillation estimates yield a continuous limit. Outside that event, use the zero path uniformly, making every path of the constructed process continuous. Limits are measurable at each fixed time, and continuity together with timewise measurability gives joint time measurability.

The relation to the original source must also be proved. Gaussian second moments imply convergence in probability of the fixed-time grid approximations to the original coordinate. Uniqueness of the almost-sure and in-probability limits then makes the new process equal to the original coordinate almost everywhere at each fixed time. No intersection over uncountably many null sets is taken. Finite-dimensional distributions require only finite intersections, so all original Gaussian finite-dimensional laws transfer, giving \nolinkurl{IsBrownianReal} without a Brownian certificate as input. The zero initial state and unit-time variance are preserved; construction has not eliminated nonzero noise.

The same probability law and continuous process feed back into the previous OU construction. Every OU path is continuous, and every deterministic initial state has strictly positive variance at positive times with positive noise. At each fixed time, the OU pushforward law equals that of the original coordinate construction, so the same source still supplies the recovery bound. This does not prove general stochastic integration, Itô's formula, controlled SDEs, or Girsanov change of measure. Conditional increments under the natural filtration and these more general applications retain separate source obligations.

\section{39.11 Original policies read the complete actual history}\label{original-policies-read-the-complete-actual-history}

Let a controller be a measurable Markov kernel on complete finite histories \(h\in(A\times W)^*\), with time recovered from history length. The history measurable space is generated by length and each ordered event coordinate with a missing-entry marker, allowing nonatomic observations; histories are not all assigned a discrete measurable space. The controller reads the same actual observation history, not latent states, posteriors, response arrays, or a model oracle.

\nolinkurl{MeasurablePredictiveHistoryControl} samples a latent state once from the registered initial belief, then retains and advances it by the original P/O at every step. Actual posterior disintegration gives one-step consistency, and induction on complete prefixes gives the same posterior distribution of the terminal latent variable. Converting the current posterior to its full response array while retaining the history unchanged yields the same controller's finite and countable full action--observation history laws. A nonempty standard Borel latent space and countably generated observations suffice for \nolinkurl{source/filter} path compatibility. The response process additionally requires standard Borel observations and countable actions with measurable singletons.

Retaining history beside the response is essential. In an actual Bool source, the first two actions are fixed, while subsequent actions read the first observation. Two length-two histories each have actual probability \(1/4\) and share their length and last action--observation event, yet yield different next-action laws. This counterexample neither selects successful runs nor proves that the response alone implements every controller with memory. The common Gaussian source supplies a nonatomic application for every such integer-action controller reading full real-valued observation histories. Comparing different initial beliefs still requires the same controller and initial memory; distinct controller histories cannot silently be merged.

\section{39.12 Minimality of adaptive prediction from actual history programs}\label{minimality-of-adaptive-prediction-from-actual-history-programs}

Finite-word responses must address two further issues: future actions may be selected after observations, and the past itself may have been generated by a full-history program. Fix the same P/O source. The past program reads only actual action--observation history, and the latent state is sampled once from the registered initial belief and subsequently retained. Prefix induction shows that the joint law of the entire actual past and terminal latent variable disintegrates through the actual past law and the same terminal posterior. The future program executes from that retained actual latent variable, not from a newly created physical state independent of the past.

The actual \nolinkurl{past/retained-latent-state} joint law requires only countably generated observations. Constructing the subsequent measurable fresh/memory prediction kernels from response arrays additionally requires standard Borel observations, countable actions, and measurable singletons. If each new future program explicitly starts with empty control memory, such a prediction kernel for its complete countable action--observation history takes the current full response array as input. The same response, future program, and initial control memory therefore give the same full future law. If inherited memory is a measurable function \(m(h)\) of the past, the predictor retains \((C(h),m(h))\) rather than equating action rules with different memories.

Minimality also requires inclusion of actual experiments. For each finite action word \(u\), construct an actual history program that executes the corresponding allowed action of \(u\) according to current history length, then uses a registered allowed default action after the word ends. On the actual history of length \(|u|\), an explicit ordered observation projection gives exactly the original \(M_u\). Repeated actions retain distinct observation positions, and default readouts for missing positions alter no actual word observation. Consequently, a statistic predicting all these actual fresh adaptive programs necessarily predicts every original finite word. The original canonical minimality theorem then applies, without a separately assumed word-prediction certificate.

More precisely, let \(\mu_n\) be the actual prefix law of any full-history past program. Suppose a statistic \(S\) has a measurable Markov prediction kernel for each fresh adaptive program, giving that program's actual future law almost everywhere under \(\mu_n\). Then \(C\) is, almost everywhere under \(\mu_n\), a measurable factor of \(S\), and \(\sigma(C)\preceq\sigma(S)\) modulo null sets. A common null-set complement is obtained only for the original countable family of action words, not by intersection over the uncountable family of all programs. The original \nolinkurl{finite/countable} canonical protocol convention and observation-space conditions remain; no uncountable version lacking a common measurable structure is asserted.

\nolinkurl{MeasurablePredictiveAdaptiveCanonical}, \nolinkurl{MeasurablePredictiveWordController}, and \nolinkurl{MeasurablePredictiveHistoryCanonical} provide the actual measurable future kernel, physical word projection, and arbitrary actual past source, respectively. The last removes the restriction of the past to visible-belief policies; control programs still read only their permitted observation histories. The actual past--future joint law and these \nolinkurl{sufficiency/minimality} applications use the same source.

\section{39.13 Conditional noise relative to the entire natural past}\label{conditional-noise-relative-to-the-entire-natural-past}

Let \(B\) be the continuous process actually constructed in Section 39.10 and let \(\mathcal F_s=\sigma(B_r:r\le s)\) be its natural filtration. \nolinkurl{BrownianMartingaleSource} derives independence of future increments from the entire natural past using the same source's independent increments, not merely independence from the current coordinate. For \(0\le s\le t\), it obtains \[
E[B_t-B_s\mid\mathcal F_s]=0,\qquad
E[(B_t-B_s)^2\mid\mathcal F_s]=t-s.
\] The required integrability follows from Gaussian moments. Thus \(B_t\) and \(B_t^2-t\) are actual martingales for this natural filtration, rather than fields supplied with the source.

For each fixed real \(u\), the same-source Gaussian MGF also gives \[
Z_t^u=\exp(uB_t-u^2t/2),\qquad E[Z_t^u]=1,
\qquad E[Z_t^u\mid\mathcal F_s]=Z_s^u.
\] Hence the finite-horizon measure \(Q_T^u=Z_T^uP\) is a probability law. For each fixed \(u,T\), its actual Radon--Nikodym derivative is, almost everywhere under \(P\), the same explicit \(Z_T^u\). For \(A\in\mathcal F_s\) and \(s\le t\), \(Q_t^u(A)=Q_s^u(A)\). These are finite-horizon normalization and consistency on past events, not a proof that the drift-shifted process is Brownian after change of measure, a general Girsanov theorem, or infinite-horizon absolute continuity.

\section{39.14 The same time input returns the same stochastic integral}\label{the-same-time-input-returns-the-same-stochastic-integral}

Take a finite nondecreasing deterministic partition \(0=t_0\le\cdots\le t_N=T\). Each real coefficient \(H_i\) is strongly measurable with respect to the interval-start information \(\mathcal F_{t_i}\) and essentially bounded under the same \(P\). Define the actual output and time input by \[
I(H)=\sum_{i<N}H_i(B_{t_{i+1}}-B_{t_i}),\qquad
h(t,\omega)=\sum_{i<N}H_i(\omega)\mathbf1_{(t_i,t_{i+1}]}(t).
\] Conditional moments relative to the original natural past give the one-step zero mean and second moment. A preceding output is measurable at its endpoint, which licenses its use in conditioning at a later interval's start. This proves cross-term orthogonality and \[
E[I(H)^2]=\sum_{i<N}(t_{i+1}-t_i)E[H_i^2]
=\|h\|_{L^2([0,T]\times\Omega,dt\,dP)}^2.
\] Repeated times give zero increments and zero input intervals; empty partitions and \(T=0\) are retained. Anticipating coefficients do not permit this conditional-mean argument.

An energy formula of the same shape for each partition table does not yet define a single integral. For two tables, \nolinkurl{BrownianCommonRefinement} constructs the common refinement of interval intersections and proves pointwise preservation of both original time inputs and execution outputs. Empty intersections uniformly receive zero coefficients, since the other table's coefficient may become observable only later. Positive-length intersections transfer permissions only forward along the filtration. The energy identity for the common-source difference then proves that equal native time-input L² classes give almost-everywhere equal original literal outputs and equal native output classes. All L² distances between different inputs are preserved as well.

\nolinkurl{BrownianIntegralClosure} derives the Gram identity from the proved distances and norms, uses native finite linear combinations to construct a continuous linear isometry on the span of elementary inputs, and extends it continuously to that span's closure in native L². Every original finite input still returns the same class of its original \(I(H)\). Convergence of those inputs in time--source L² implies convergence of their original output classes to the actual integral. The unit coefficient on \((0,T]\) has integral energy \(T\), so extension does not discard nonzero execution noise.

At this stage the integration domain is explicitly only the closure of the elementary span. The next section establishes its relation to the full predictable L² class; that relation needs an independent proof and is not supplied by the word ``closure.'' The time-indexed integral process, its martingale and quadratic-variation properties, and the earlier general Itô/\nolinkurl{Girsanov/controlled-SDE} applications each have separate source obligations.

\section{39.15 The integration domain is the native predictable L² space}\label{the-integration-domain-is-the-native-predictable-luxb2-space}

The relationship between the closure above and predictable inputs now has an independent proof. Start with the predictable σ-algebra of the native natural filtration, generated by the zero-time section and rectangles \((s,\infty)\times A\) with \(A\in\mathcal F_s\), and pull it back to the real time coordinate used in this book. Define the full predictable L² space as the native L² classes admitting an almost-everywhere strongly measurable version for that σ-algebra, without defining it as an elementary closure.

Every actual interval input is predictable, so the elementary closure lies in this full space. The reverse inclusion is also derived from actual inputs. Under the same \([0,T]\) measure, existing finite interval records realize the generating rectangles; rectangles beyond the horizon and the zero-time section give the zero class. The π-system property of the generators, σ-algebra induction, L² convergence of indicators under the same product measure, and native L² strong-measurability induction put every full predictable L² class into the original closure. The two spaces therefore coincide, including \(T=0\).

\nolinkurl{BrownianPredictableInputs} and \nolinkurl{BrownianPredictableDensity} connect both inclusions to the original integral, yielding an actual continuous linear isometric integral on full predictable L². Each original finite input still returns the same literal output and the unit coefficient still has energy \(T\). This constructs the terminal integral; a time-indexed natural-filtration martingale, continuous version, quadratic variation, and the earlier general Itô/Girsanov applications must still be derived separately.

\section{39.16 Purpose certificates from finite actual birth and cost data}\label{purpose-certificates-from-finite-actual-birth-and-cost-data}

Exact-mean access in Section 39.5 is not finite data. Fix a complete budget: first execute \(N\) write-false requests, then \(N\) write-true requests. Bin the cost readout into integer bins with preregistered rational width \(\tau>0\), and register rational price \(\lambda\ge0\). Allowed actions, the original unknown real fault rate, program target, and real costs are unchanged. Training observations contain only birth bits and cost bins. The complete transaction sequence retains the same parent, all failures and costs, and all \(2N\) coordinates. The same product law of source modes and labels yields independence of the actual sequential samples and each observation-row law, without resetting to a successful state or resampling.

If a request for bit \(b\) produces an actual offspring, then \(b\) is the source target and the actual cost is \(c\). Its integer bin \(k\) satisfies \[
k\tau\le c<(k+1)\tau.
\] The executable selector reads only the complete Bool/Nat batch and registered rational parameters. If \(\lambda(k+1)\tau<1\), it returns the corresponding repair controller. If \(1<\lambda k\tau\), then under explicit \(\epsilon<1\) and an actually damaged state, it gives a negative verdict that one-step idle is preferable to repair. Otherwise it returns UNKNOWN. Source parameters, true means, and the calibration-good event are not selector inputs. Interval and target certificates are error-free in this exact-readout model; physical execution is not thereby fault-free.

On a positive output, the learned bit equals the original protected target, and the cost certificate supplies the original sufficient condition \(\lambda c\le1\). The new application starts from the actual final parent and physical state after all probes and executes the same returned controller. One actual infinite-record law retains all original finite prefixes, failures, and costs; its integrable infinite reward equals the returned policy's own value and the optimum over the original class of all causal policies. No alternative existing policy is selected, and execution does not return to the intermediate state at the first success. Training costs remain in the original full record. Optimality evaluates the registered future deployment; it asserts neither optimality of total net training-plus-future reward nor free training.

If an informative bin still cannot decide, then \(|\lambda c-1|\le\lambda\tau\). When the actual bin is separated from the threshold, the original same-source scalar Hoeffding calibration gives, for any analytical radius satisfying \(0\le u<1-\epsilon\), \[
P(\mathrm{UNKNOWN})\le\min\{1,4e^{-2Nu^2}\}.
\] The unknown repair-probability lower bound is not an input to the selector. Under the fixed source conditions \(N=64,\lambda=1/2,\tau=1/4,c=1,\epsilon=1/4\), the analytical radius \(u=1/4\) gives \(4e^{-8}<1\). This is a formal conclusion for the fixed-budget model, and the executed selector fixtures are deterministic checks. The package generates no new real-world sample batch. Zero budget always returns UNKNOWN; at total failure \(\epsilon=1\), UNKNOWN is returned almost everywhere. No positive repair probability is inferred from an imaginary successful tuple of zero mass.

Necessary negative bounds are retained. At fault rate \(1/2\), the common all-failure sample has positive mass \(2^{-2N}\). Sources with price one, target true, and cost zero or two are indistinguishable on this sample but have opposite actual one-step orderings at the same damaged state. A deterministic, data-only, zero-error procedure therefore cannot always give a strict verdict for both sources; no randomized-estimator version is asserted. Bin width two also makes all coarse readouts identical for costs \(1/2\) and \(3/2\), an obstruction relative to that fixed coarse readout.

A negative verdict cannot be extended to other horizons. With \(\epsilon=1/4,\lambda=d=1/2,c=3\), the damaged-state one-step repair advantage is \(-3/8\), but the same compiled policy's own infinite value is \(3/32\), whereas idle has infinite value zero. This uses policy evaluation, not the inapplicable optimality condition \(\lambda c\le1\). Partial purpose certificates from finite data, complete source-parameter identification, and all-policy ordering at arbitrary horizons are different obligations. The sources are \nolinkurl{RegisterFaultFinitePurposeProbe} and \nolinkurl{RegisterFaultFinitePurposeInfinite.} Section 39.18 separately registers a resource-maintenance channel; Section 39.24 adds an actual training ledger under a separately registered complete idle/write protocol. The earlier protocol and evaluations accounting for more costs retain their respective scopes.

\section{39.17 An actual continuous process for the same integral}\label{an-actual-continuous-process-for-the-same-integral}

Fix \(T\ge0\) and continue to use the same native predictable \(L^2\) domain of Section 39.15 and the same \(dt|_{[0,T]}\otimes P\). For input \(f\), define the actual truncated input and output class by \[
f_s=\mathbf1_{(0,\min(s,T)]}f,\qquad I_s(f)=I_T(f_s).
\] The time measure is not changed with \(s\). Multiplication by an indicator preserves the original genuinely predictable domain and is a contraction. For an actual interval input \(H\mathbf1_{(a,b]}\), the output at time \(u\) is \(H(B_{\max(a,\min(b,u))}-B_a)\). When \(u<a\), the conditional expectation vanishes by the tower law through \(\mathcal F_a\), without treating future \(H\) as \(\mathcal F_u\)-known. Finite decompositions and the proved native density extend the conditional-projection identity to every \(f\); only then is a natural-filtration martingale obtained that agrees almost everywhere with the actual truncated output at every fixed time.

Continuity is derived separately. Finite literal paths are continuous. Conditional Jensen for squares and the native maximal inequality give a grid maximum-deviation bound; ordering and continuity from below extend it to a single countable maximum event. Geometrically accurate approximations by original finite inputs and Borel--Cantelli yield one event of uniform convergence over all times. The actual continuous version \(M_s(f)\) always takes limits of finite adapted values at time \(s\), rather than choosing outputs according to that future global event. It therefore preserves the original natural filtration, has joint and progressive measurability, and has continuous paths almost everywhere. At each fixed time, \(M_s(f)\) equals the original \(I_s(f)\) almost everywhere. Each original finite input returns its entire literal path for all times on one common full-measure event. There is neither an uncountable intersection of timewise null-set complements nor one event asserted for all inputs simultaneously. The terminal second-moment energy of the actual unit input remains \(T\), including the zero horizon.

The direct sources are \nolinkurl{BrownianIntegralTimeTruncation}, \nolinkurl{BrownianIntegralMartingale}, \nolinkurl{BrownianIntegralFinitePaths}, \nolinkurl{BrownianIntegralMaximalControl}, and \nolinkurl{BrownianIntegralContinuousVersion.} General Itô, change of measure, and transport to arbitrary filtrations remain separate obligations.

\section{39.18 Feedback resource supply and the closed-reserve obstruction}\label{feedback-resource-supply-and-the-closed-reserve-obstruction}

The per-step supply at least as large as the write cost in Section 39.3 is legitimate but strong. A separate open-environment port is now registered: every idle request receives \(g\) resource units, even if a fault preempts that physical execution; write requests receive zero. This is a new external channel, not energy created within the original closed source. Allowed actions, earlier physical transactions, birth projections, and cost projections are unchanged. If deployment starts with reserve \(R\ge c\) and \(g\ge c\), the compiled repair program's actual state transitions preserve a reserve cushion determined by \nolinkurl{healthy/damaged} status. The step of actual paid repair may receive zero supply; after repair, the same program requests idle next. Every original positive-mass finite record therefore has nonnegative actual balance. Prefix consistency of the complete infinite record further gives P6 pathwise resource viability almost everywhere.

By contrast, in the original closed single-register source, fix birth label false and suppose \(c>0\) and \(0<\epsilon<1\). Starting damaged, the alternating word of actual repair and faulty idle has mass \[
\bigl((1-\epsilon)\epsilon\bigr)^k>0,
\] and actually costs \(kc\). Every finite reserve is therefore exceeded by some positive-mass prefix. The complete infinite law of this particular original compiled policy is not almost-surely resource viable. This is not an almost-sure bankruptcy theorem and does not cover zero costs, degenerate fault rates, or a universal impossibility for arbitrary policies.

Positive finite-data output retains the original returned program and the actual parent after all original probes. The same actual infinite record both attains the registered optimal future reward and satisfies resource P6 under the new port conditions. Reserves here are measured from deployment; all past probe costs are retained. A funding source for training and deployment together still requires proof at this stage. The source is \nolinkurl{RegisterFaultFeedbackResource.}

\section{39.19 Actual squared increments and time energy of the same input}\label{actual-squared-increments-and-time-energy-of-the-same-input}

For a nondecreasing deterministic partition \(0=t_0\le\cdots\le t_n=T\) of the original Brownian source, allowing repeated times, put \(Q=\sum_j(B_{t_{j+1}}-B_{t_j})^2\). Original Gaussian fourth moments and conditional independence from the entire natural past give orthogonality of centered squared increments, hence \[
E[(Q-T)^2]=2\sum_j(t_{j+1}-t_j)^2
\le2\,\operatorname{mesh}(t)\,T.
\] The mesh is computed from the caller's actual finite clock. Every partition family whose actual mesh tends to zero yields mean-square and native \(L^2\) convergence. Mesh convergence is also proved for the equal-spaced clock family itself. Original past-measurable uniformly bounded weights may vary with the partition; the same argument connects the weighted squared sum to the same weighted time sum. The latter is not yet called an unproved time integral. Sources are \nolinkurl{BrownianQuadraticVariationSteps} and \nolinkurl{BrownianQuadraticVariation.}

For any original full predictable \(L^2\) input \(f\), define the actual pathwise time energy \[
A_s(f)=\int_0^{\min(s,T)}|f(t,\omega)|^2\,dt,
\] using the native predictable representative of the same input class. Native predictability and Fubini imply that \(A_s\) is strongly adapted to the original \(\mathcal F_s\) and integrable. One common event of integrable sections gives monotonicity and continuity at all times. The natural space for squared outputs is \(L^1\): \[
\|u^2-v^2\|_1\le\|u-v\|_2(\|u\|_2+\|v\|_2).
\] No fourth-moment condition is silently added for a general \(L^2\) input.

The key conditional second-moment energy is also derived from actual inputs. For \(A\in\mathcal F_s\), restrict to the future window \((s,\min(t,T)]\) and multiply \(f\) by \(\mathbf1_A\) there. Original interval outputs and native density show that its actual integral is exactly \(\mathbf1_A(M_t-M_s)\). Isometry and Fubini for the same time integral yield the second-moment energy relation on every such event. Cross terms of the actual martingale have zero conditional mean. Uniqueness of native conditional expectation then proves that \[
M_s(f)^2-A_s(f)
\] is an actual continuous martingale for the original natural filtration, starting at zero almost everywhere. Neither \(M,A\) nor the compensated martingale is a caller field. Sources are \nolinkurl{BrownianIntegralTimeEnergy}, \nolinkurl{BrownianIntegralSquareL1}, \nolinkurl{BrownianIntegralPastLocalization}, and \nolinkurl{BrownianIntegralCompensator.} This does not by itself establish partition quadratic variation for arbitrary \(f\), a general C² Itô formula, Girsanov, or controlled SDEs.

\section{39.20 Parent identity from actual population scheduling}\label{parent-identity-from-actual-population-scheduling}

Timestamped identity extends to a separately registered serial population model. Distinct slots in a finite founder table represent distinct owners even when program, physical value, or birth label agrees. Each request first reads the actual parent at its address, calls the original transaction and updates that slot, then appends a new optional slot. An actual birth occupies the new slot. Failure or an invalid address leaves it empty and still consumes a clock step. No alternative parent is borrowed, no reset occurs, and records are not discarded by conditioning on success.

The newborn ID is the table length before execution. Successful transactions across parents and generations therefore receive distinct IDs even with equal labels. Actual birth also determines the unique parent that existed at that time, supplying the original P12 parent bridge and P8 \nolinkurl{repair/offspring} distinction. The copied program has a source as well: if the selected parent's program equals the registered mechanism target, the offspring enters original P5. A mixed-program population does not thereby acquire a common P5 guarantee. Scheduling addresses and reliable logs are explicit additional permissions, not ancestry inferred from unlabeled physical states. The module is \nolinkurl{RegisterPopulationHistorySource}; a statistical-discovery good-event algorithm still needs a separate source.

\section{39.21 The original Price theorem uses actual offspring}\label{the-original-price-theorem-uses-actual-offspring}

Freeze the address table before execution and attempt each old address once; newborns do not enter the current generation's schedule. Actual execution proves that the selected parent remains the original object and that the new slot in the terminal table is precisely the offspring produced by this transaction. Parent frequencies count only occupied old slots, and reproductive rates are computed from actual births. Only when the actual offspring count \(B>0\) are normalized transmission rows constructed on the actual offspring. A row with no offspring may be completed using an already existing offspring, but its reproductive rate is zero. Only entries with \(bK>0\) represent actual genealogy edges. The impossible combination refuted in Section 39.2---a finite set of permanent individuals with every row normalized and every positive edge strictly increasing generation---must not be reused.

After normalization, the source-computed \(f bK\) update equals the uniform empirical frequency of actual offspring. Thus the frequencies, reproductive rates, and transmission matrix in the original p\_evo\_01 all come from the same actual births. For the inherited program trait, the transmission term vanishes, and the change in actual offspring mean equals selection covariance divided by mean reproductive rate. In an example with both programs and equal labels, the parent true-type frequency is \(1/2\) but the actual offspring frequency is zero, giving covariance divided by mean reproductive rate \(-1/2\). The all-failure realization of the same schedule produces no offspring and is not normalized.

This is a realization-wise empirical Price identity; it does not replace the expectation of a random ratio by a ratio of expectations. Old parents remain in the terminal table, and the updated frequencies explicitly concern only the offspring census, not the entire terminal population. Nor does the result prove that selection necessarily approaches an optimum. The source is \nolinkurl{RegisterPopulationTypePrice.} Formation, dependent statistical discovery, scale, long-term funding closure, and global composition retain their respective source obligations at this stage.

\section{39.22 Partition quadratic sums and actual weighted tables for the same full integral}\label{partition-quadratic-sums-and-actual-weighted-tables-for-the-same-full-integral}

The square-compensated martingale in Section 39.19 did not by itself prove partition quadratic variation. The new source chain fixes an original full predictable \(L^2\) input \(f\) on real time times the original probability law, using the same continuous integral process \(M_f\) and actual time energy \(A_f(t)=\int_0^{t\wedge T}f_s^2\,ds\). For any declared deterministic finite nondecreasing partitions \(\pi_m\) from \(0\) to \(T\), if the computed maximum interval length tends to zero, then

\[\begin{gathered}
\sum_{[u,v]\in\pi_m}(M_f(v)-M_f(u))^2\longrightarrow A_f(T)
\\\text{in the original source's }L^1\text{, hence in probability.}
\end{gathered}\]

No other existing input is substituted. The proof starts with actual left sampling on the original finite span and extends to the full input through actual time-marginal integrals and a squared-sum stability bound uniform over all partitions. It requires only original \(L^2\), not an added \(L^4\) assumption. Repeated endpoints and empty time domains are included. Random stopping-time partitions are not covered, nor is a single event claimed to give full-sequence almost-sure convergence for all partition families simultaneously.

For two inputs \(f,g\) on the same Brownian source, actual cross-increment sums converge to the actual \(\int_0^T f_sg_s\,ds\), and \(M_fM_g-\int fg\,ds\) is an integrable martingale for the original natural filtration. These inputs are not two independent noises. For an actual uniformly bounded, jointly measurable, adapted coefficient process \(H\) with almost-everywhere continuous paths, original \(F_u\)-simple-function approximation, future-interval localization, and dominated convergence for the original input further give

\[
\begin{aligned}
\sum_{[u,v]\in\pi_m}H_u\,(M_f(v)-M_f(u)) &\longrightarrow I(Hf)_T,\\
\sum_{[u,v]\in\pi_m}H_u\,(M_f(v)-M_f(u))^2 &\longrightarrow \int_0^T H_sf_s^2\,ds.
\end{aligned}
\]

in probability under the original source. Every table uses the same actual \(M_f\) term by term. The target is computed from the original time integral, not supplied as the limit to be proved. These are sources for the first- and second-order stochastic sums needed by Itô's formula. Section 39.26 combines them for the same scalar input using actual Taylor remainders and telescoping sums under explicit global derivative bounds. A bounded \(C^2\) function need not have bounded first and second derivatives; each later application must supply its own regularity and integrability conditions. The present source is the canonical scalar source and does not remove the earlier obligations for general dimensions, filtrations, controlled SDEs, or change of measure.

\section{39.23 Finite-dimensional noise reading the whole joint history and native input identity}\label{finite-dimensional-noise-reading-the-whole-joint-history-and-native-input-identity}

Register a native product of finitely many original scalar Brownian sources, with the genuine Euclidean vector norm. Independence of future vector increments from the entire joint natural history follows from coordinate pair laws and product rearrangement. Independence of each coordinate from its own past is insufficient to authorize the whole joint history. Thus an actual coefficient may read another coordinate's past, for example to construct cosine/sine weights from the second coordinate. If each coordinate coefficient at a deterministic START time is strongly measurable for the entire joint natural history and essentially bounded on the same source, the actual sum of dot-product increments has zero mean and Euclidean isometry derived from the same-source conditional covariance.

The actual input lies in Euclidean-valued \(L^2\) on unnormalized real time times the same vector probability source. Common refinement of distinct deterministic partitions preserves the original START permissions. Collapsed common cells receive zero coefficients; a future-measurable function is not moved to an earlier start. All distances between native inputs and outputs are then equal: the same input class returns the actual output class of the same finite program, with original outputs equal almost everywhere. Native finitely supported linear combinations and continuous linear extension construct the integral on the closure of this span, and every original finite input returns its original output.

This section first constructs the closure of the actual elementary input span. Section 39.27 separately proves from actual generators and density that it equals the full jointly predictable Euclidean-valued \(L^2\) space, and derives actual conditional processes and mean-square time continuity rather than adding these as fields. Zero dimension, zero horizon, and repeated times are included in the exact scope. Section 39.31 separately constructs the sample-continuous version of the vector-noise integral. General vector Itô, control, and Girsanov results do not follow automatically from finite-sum isometry.

\section{39.24 One ledger for all training costs and the actual returned policy}\label{one-ledger-for-all-training-costs-and-the-actual-returned-policy}

The externally protected idle port from Section 39.18 now enters a separately registered complete training protocol: request idle before each original allowed write, and request idle once more at the end. The original \(2N\) writes therefore correspond to \(4N+1\) primitive draws. All faults, births, costs, and idle records remain, without reset, program installation, successful-sample selection, or redrawing. The marginal law of the selected write coordinates equals the original training source, but the full trajectory and actual final parent differ. The old trajectory cannot be claimed to be identical.

Under the declared port assumptions \(g\ge c\) and actual initial reserve \(R_0\ge0\), every actual training prefix has nonnegative balance and final reserve at least \(c\). The same count-based selector reads only actually allowed observations. Under original program matching, registered price, positive readout width, and discount \(0<d<1\), a positive repair verdict returns the original policy, now starting from the new actual final parent state and entering the original future-optimality and P6 resource results. Literally prefixing the full training record to that actual future-record law yields one almost-surely viable complete ledger from the same initial reserve. Optimality is still evaluated from the actual deployment start; total net training-plus-future reward is not claimed optimal. The port is separately registered external supply, not internally generated energy.

\section{39.25 Self-maintenance of mixed programs, actual offspring, and complete parent discovery}\label{self-maintenance-of-mixed-programs-actual-offspring-and-complete-parent-discovery}

Under the original mechanism with a fixed global target, mixed-program populations do not automatically share P5. The new model computes the current target from the selected actual parent's protected program after each lookup, and a fault writes the opposite of that parent's own program. This is an explicitly different environmental semantics. The earlier source and restrictions remain, and no biological action is added. An actual offspring's physical seed matches its inherited program, so each child type enters its own original P5 without separately assuming that the parent program equals a global target.

With the old address queue frozen and each original parent visited once, the complete actual child options of the new model equal the earlier offspring-census projection. The original Price result can therefore reuse actual frequencies, birth counts, and transmission terms. Retained parents after faults and the full terminal table may differ; offspring frequencies are not terminal-table population frequencies. Newborn slots, unique birth IDs, and parents are derived from the same allocator. A three-parent source produces actual offspring of both types, one retained fault, and a later actual paid repair. Its four-step primitive word has original probability \(27/4096\). A newborn's zero local cushion and the declared idle port also supply its actual first-step funding bridge, not a population-wide bank or globally optimal schedule.

Parent discovery separately registers complete noisy-read access to a marked journal. The count-based selector uses only actual ID readouts, a finite registry, and a fixed fallback. Hidden targets, noise parameters, correct marks, and Good are not selector inputs. When \(\eta<1/2\), the actual independent-read source gives correctness probability at least \(1-\eta\) and hence a strict-majority event. Source means and Hoeffding bound its complement by \(\exp[-2N(1/2-\eta)^2]\). The bridge from actual ID classes to parents supplies P12. At \(\eta=1/2\) in an unknown arbitrary confuser class, different actual parents can have identical complete-readout laws; every deterministic estimator has errors on the two sources summing to at least one.

A stronger common application covers the product law of all four population draws and all \(N\) read draws. For each primitive population word, the fourth action is computed from the actual first-three-step state, so the actual key and parent may vary with that word. On the registered source with both population and read noise \(1/4\), joint actual-parent risk is at most \(\exp(-N/8)\) and tends to zero. No conditioning on a positive example, successful birth, or acceptance is used. The \nolinkurl{all-fault/no-birth} word of mass \(1/4096\) remains in the full source. Only on source-derived Good and when the target is actually born does the same returned actual child enter its own P5. None denotes absence of an offspring parent, not the identity of an unsuccessful transaction's attempted parent. Read costs, population-wide funding, and global purpose optimality retain their own source obligations and are not obtained for free from metadata access.

\section{39.26 Actual Itô and Dynkin sources for the same full input}\label{actual-ituxf4-and-dynkin-sources-for-the-same-full-input}

Let \(f\) be any actual predictable square-integrable input on the same original Brownian source of Sections 39.13--39.17, and let \(M\) be its constructed continuous integral. If \(\varphi\in C^2(\mathbb R)\) has globally bounded first and second derivatives, then almost everywhere on the original source,

\[
\begin{aligned}
\varphi(M_T)-\varphi(0)
&=I_T\bigl(\varphi'(M)f\bigr)\\
&\quad+\frac12\int_0^T\varphi''(M_t)f_t^2\,dt.
\end{aligned}
\]

The finite table is the actual telescoping Taylor sum of the same \(M\). The second-order remainder is controlled by uniform continuity of \(\varphi''\) on the actual compact path interval, tightness in probability of the actual squared sums, and uniqueness of subsequential limits. Neither \(f\in L^4\) nor a third derivative is required, and the Itô identity is not input. The function \(\varphi\) itself need not be bounded; conversely, boundedness of a \(C^2\) function does not imply the derivative bounds above.

The constant-coefficient form of the same actual input gives the actual state law of \(X_t=x_0+\sigma B_t\). Source-derived zero means, time Fubini, and the fundamental theorem of calculus construct the original \nolinkurl{RecoveryDynkin} certificate of absolute continuity of expectations and an almost-everywhere generator identity. Displacement second moments for \(\sigma=2\) are \(4t\), and a nonlinear sine test uses the same certificate. Free Brownian motion and a bounded sine test do not satisfy all the conditions for forced recovery. General recovery and controlled-SDE obligations remain, while the earlier OU quadratic recovery result keeps its separate identity.

\section{39.27 The full vector predictable domain and actual conditional processes}\label{the-full-vector-predictable-domain-and-actual-conditional-processes}

The closure domain of Section 39.23 now equals the independently defined native predictable Euclidean-valued \(L^2\) space on the full joint natural filtration. The proof realizes vector indicator inputs for predictable generating rectangles and applies native \(L^p\) induction for density; zero-time generators have zero original time-product measure. The result covers every finite dimension, including zero dimension, zero horizon, and repeated times. The new integral uses the same input class and returns the original finite program's original output.

Native truncation uses the same fixed-horizon \(dt\times P\), rather than a newly selected measure or noise. When observation precedes START, the conditional prefix follows by the tower law through START; a future coefficient cannot be pulled out of an earlier conditional expectation. Actual finite prefixes and continuous linear extension by density yield the full input's joint-past martingale, original-output identity at every fixed time, and Euclidean energy. Nonatomic time boundaries and dominated convergence of squared norms give mean-square time continuity. Mean-square continuity is not a sample-continuous version; Section 39.31 constructs that version for the same integral from actual finite paths and maximal control. Vector quadratic variation and general vector Itô retain independent obligations.

\section{39.28 The original OU conditional kernel and actual constant-control change of measure}\label{the-original-ou-conditional-kernel-and-actual-constant-control-change-of-measure}

For the original time-changed OU source with \(k>0\), the actual conditional filtration is \(\mathcal F_{\tau_k(s)}\), where \(\tau_k(s)=(e^{2ks}-1)/(2k)\). The same Brownian path yields the actual state decomposition and Gaussian innovation independent of that entire past. The current state in the conditional kernel is the actual random state. Native conditional distributions and the tower law derive the transition-operator semigroup for every initial state and bounded Borel test, including discontinuous indicators, zero times, and degenerate noise. This operator semigroup is not a proof of a physical-clock SDE, and the time-changed past is not identified with the unchanged \(\mathcal F_s\).

Constant control \(u\) on finite \(T\) uses the same original noise source's actual exponential density. Completing the square in the actual joint law of finite Gaussian residuals proves that under the resulting actual \(Q\),

\[
W_t=B_t-u\min(t,T)
\]

is Brownian, with independence from the entire original \(\mathcal F_t\) past and the corresponding martingale property. Drift stops accumulating after \(T\). The original predictable integral returns \(u(B_T-B_0)\). The controlled integral, actual time energy \(u^2T\), and drift identity construct the original \nolinkurl{TerminalGirsanovData.} Hence actual full-noise KL is \(u^2T/2\), and the KL of any measurable readout is no greater. The proof includes \(T=0\), \(u=0\), and increments crossing the density horizon. General predictable controls still need sources for their actual exponential martingales, drift, and \(Q\) energy.

\section{39.29 The same noisy measurements, actual inverse, and returned action}\label{the-same-noisy-measurements-actual-inverse-and-returned-action}

C5's deterministic finite-dimensional error result uses the supremum norm, distinct from the Brownian Euclidean norm above. Construct \(\widehat M\) from actual noisy anchor columns; for the same returned matrix \(R\), verify exactly \(R\widehat M=\widehat M R=I\), and compute coefficients using the same noisy query. Error and the original action-gap ordering bound explicitly depend on \(\lVert R\rVert\). Positive anchors can approach degeneracy, amplifying query error; the condition-number factor cannot be dropped. Failure of inversion is not a general impossibility proof.

A separately registered original source has two unknown positive rational diagonal gains and two distinct rational latent responses in \([-1,1]\). Unit anchors and independent read noise are explicit model permissions. The risk result holds for every row size \(N>0\). The registered example uses \(N=32\) per row and \(128\) complete signs, with sensor noise amplitude \(1/10\). These signs and observations are a synthetic reconstruction of preregistered formal words and actual sourceData, not new empirical samples. Fair independence describes the original probability law, not an empirical fact inferred from a fixed word. Signs, budgets, and thresholds are unchanged after registration. The selector reads only actual rational means, computes the same diagonal inverse, condition number, and query scores, and rejects or returns an action under the original threshold. True gains, true ordering, noise signs, and Good are not inputs. Source Hoeffding and the same error event yield

\[
P(\mathrm{Accept}\cap\mathrm{Wrong})\le 8e^{-N/8},
\]

which is \(8e^{-4}<1\) at the original budget. This is neither a conditional error rate nor zero error on every sample. The actually registered word has nonzero second-anchor bias \(1/160\) and returned-coordinate error \(1/321\). All signs and complete observations for two legal parameter sources are independently reconstructed term by term and proved in the kernel equal to actual sourceData. Reversing the source ordering retains the same original sign word but changes query observations; the selector correctly changes the returned action.

A strict limitation uses the same source and selector. With legal gains \((1,9/100)\) and true response \((-1,1)\), another complete word gives observed means \((1,-1/100,-19/100,19/100)\), condition number \(100\), and scores \((19,-19)\), causing acceptance of the wrong action. This word has mass \(2^{-128}>0\) and lies outside Good. It preserves the joint risk bound while refuting samplewise correctness. It is a formal source word, not a newly sampled empirical batch; the original actual A--F materials are unchanged. General \nolinkurl{dependent/adaptive} data, rank deficiency, nonlinearity, and complete dynamic-purpose interfaces retain their own source obligations.

\section{39.30 Full change-of-measure paths for actual past-feedback tables}\label{full-change-of-measure-paths-for-actual-past-feedback-tables}

On the same original Brownian source, take a finite deterministic clock \(0=t_0\le\cdots\le t_n=T\). Each coefficient \(H_i\) may read the entire original \(\mathcal F_{t_i}\) and is essentially bounded on that source. Coefficients, clocks, and the actual returned control table retain their identity. Rowwise conditional Gaussian exponential normalization and past-weighted integrability give a genuine prefix-density martingale, with terminal density

\[
D=\exp\!\left(\sum_{i<n}H_i(B_{t_{i+1}}-B_{t_i})
-\frac12\sum_{i<n}H_i^2(t_{i+1}-t_i)\right),\qquad Q=D\cdot P.
\]

Normalization is not an input field. The original predictable integral returns this same random sum; actual real-time control predictability and time energy are derived from the same table. The law of each original row-start past under final \(Q\) equals its preceding-density law, generally not its law under original \(P\). Conditional tilting retaining the entire past, elimination of future densities, and induction on actual prefix processes prove that

\[
W_t=B_t-\sum_{i<n}H_i\,\bigl((t\wedge t_{i+1})-(t\wedge t_i)\bigr)
\]

is a full Brownian path under actual \(Q\), with future increments independent of the entire original \(\mathcal F_t\). This includes times inside cells, times after \(T\), zero-width cells, and empty tables. For these finite causal tables, equality of the natural filtration of \(W\) with the original filtration is also derived; it cannot be passed without proof to limits of arbitrary path-dependent controls. A single random block allows normalization and path laws for a general original row-start-measurable coefficient; constructing the terminal KL record additionally uses essential boundedness.

The baseline integral of original \nolinkurl{TerminalGirsanovData}, the controlled terminal integral formed from the same \(W\) increments, its zero-mean martingale, and actual time energy all come from the source. The original KL application gives

\[
\mathrm{KL}(Q\Vert P)=\frac12 E_Q\!\left[\sum_{i<n}H_i^2(t_{i+1}-t_i)\right].
\]

Every measurable readout retains the data-processing upper bound. Energy must be evaluated under final \(Q\); it cannot generally be replaced by \(E_P\). An actual example uses, on \((0,1]\), the coefficient \(1\), and on \((1,2]\), the coefficient \(1+\tfrac12\sin(B_1)\), with actual cost in \([5/8,13/8]\). Under final \(Q\), the original \(B_1\) has mean \(1\), so the second row is actually \(1+\tfrac12\sin(W_1+1)\); no different policy has been substituted. Deterministic clocks are not random stopping times, and this terminal controlled-integral record alone establishes neither a general \(Q\) predictable integral nor all SDEs.

\section{39.31 An Actual Continuous Version for the Same Joint Vector Input}\label{an-actual-continuous-version-for-the-same-joint-vector-input}

For an arbitrary native Euclidean-valued predictable \(L^2\) input on the entire joint natural filtration of §39.27, actual finite continuous paths, native density, and control of the maximum over the whole interval provide a continuous version on a common almost-sure set. The maximal estimate is first proved at actual finite ordered times, followed by a single countable union. Path control is not replaced by an endpoint-by-endpoint union bound growing with the grid.

At each time, the output uses limUnder measurable with respect to the original \(\mathcal F_t\) at that time. The global future Good event is used only in the proof, not as a branch condition for selecting program outputs. At each fixed time, the result is tied exactly to the original integral and the energy of the same input. For finite tables and their span, it returns the original path at every time on a common almost-sure set. All-sample continuous, strongly adapted approximants on each horizon and their measurable limit further establish actual progressive measurability. No shortcut from ``almost-sure continuity plus adaptedness'' is taken on an uncompleted filtration.

A two-dimensional example lets one integral coefficient read the discontinuous sign function of the other coordinate at START. Actual return, energy, and the zero path before START are checked from the same joint source. The precise scope also includes zero dimension, zero horizon, repeated times, and finite families of Euclidean outputs. This constructs an actual sample-continuous version. Section 39.35 further supplies scalar-integral QV for the same fixed Euclidean input, and §39.42 separately gives a finite-state Itô formula with explicit derivative bounds. Drift, regularity not covered here, and controlled SDEs retain independent source obligations.

\section{39.32 Execution-Preserving Scales and Funding an Actual Continuing Population}\label{execution-preserving-scales-and-funding-an-actual-continuing-population}

A noninjective projection that retains the actual active register, protected program, and history while forgetting passive coordinates commutes with the actual interpreter, updates, faults, births, and costs. Given a coarse causal policy, the fine execution updates its original control memory using only the same projected past. The same raw-noise word and action weights then give equality of all finite and countable complete-record laws. The original formation--binding composition and operational target-preservation results thus have actual consumers. This does not assert descent of every semantic property or of arbitrary fine-level policies.

The same actual population source also retains timestamps, allocation slots, offspring IDs/parents, and the entire failure log. The original Price identity consumes the actual offspring of the same fixed parent population at old addresses; it does not substitute the whole final table. Hiding an active coordinate yields a strict counterexample with the same coarse observation but different actual birth outcomes and compilation costs. This refutes exact descent of that actual compilation policy, not of every other policy.

Continuing funding uses a separately specified original admissible-action protocol: request one idle supply before each actual work transaction. A shared funding pool and an external supply port are explicit model assumptions. Only an actual Some idle record at a valid address receives \(g\); an invalid lookup creates neither fictitious income nor a charge, while the failure clock is retained. Actual work costs lie in \([0,c]\). The original infinite-product source generates an actual continuing store recursion and a complete-record path. Every \(2N\) prefix agrees exactly with the finite protocol on the same raw word. A None tail after a finite stopping protocol is not used as evidence of continuing execution.

If \(R\ge0\) and \(g\ge c\), the balance at every actual prefix of every registered plan is at least \(R\). Hence the original \nolinkurl{ResourceViable} holds almost surely under its actual path law. Conversely, repeatedly requesting work at an actually occupied initial parent slot admits a prefix of \(N\) consecutive legitimate \nolinkurl{idle-fault/paid-repair} transactions with positive original-source mass \((3/64)^N\), whose balance is exactly \(R+N(g-c)\). When \(0\le g<c\), every finite \(R\) has such an overspending prefix for some length. This is neither almost-sure bankruptcy nor impossibility for each individual policy: from the same nonempty initial state, a plan repeatedly requesting new empty slots has neither income nor expenditure and remains viable with zero supply.

The exact equivalence established by the two directions is: when \(R,g\ge0\) and at least one initial slot is actually occupied, \textbf{every registered plan individually has an almost-sure funding guarantee} if and only if \(g\ge c\). It is not the statement ``each plan is viable if and only if \(g\ge c\),'' nor does it claim a common probability-one event for all plans. Both the actual original founder population and a population with a nonzero active port consume this result. Supply, the aggregate ledger, and the fixed-parent counterexample do not establish endogenous resources, population optimality, or general organizational autonomy.

\section{39.33 Useful Acceptance by the Same Calibrated Selector}\label{useful-acceptance-by-the-same-calibrated-selector}

The original four-row source, data, inverse computation, and returned selector of §39.29 remain unchanged. If both true gains are at least \(1\) and the latent gap between the truly best action and the other action is at least \(1/2\), the original Good event ensures not only correctness of accepted actions but actual acceptance of the best action. Hence, on the same complete source,

\[
P(\mathrm{Select}=\mathrm{Best})\ge1-8e^{-N/8}.
\]

Quantitative separation is a sufficient model assumption. The true gains and ranking are not program inputs, and these conditions have not been proved sharply necessary. Both the original forward-order and reverse-order registered sources actually consume the result. The original positive-mass word on which an incorrect action is accepted remains outside Good. No samples are added, no threshold changes, and joint risk is not renamed conditional error. Related materials, successor protocols for actual adaptive readouts, and general nonlinear sources retain their own obligations.

\section{39.34 Actual Change of Measure for Bounded Original Inputs and Integration under the Constructed Measure}\label{actual-change-of-measure-for-bounded-original-inputs-and-integration-under-the-constructed-measure}

Let \(f\) be a genuinely predictable native \(L^2(dt\times P)\) input on the original Brownian natural filtration, with \(|f|\le K\) almost everywhere under that time--source measure and \(K\ge0\). For the same \(f\), write \(I_P(f)\) and \(A_f\) for its actual integral and energy. Native density, left sampling, and truncation at the same \(K\) produce actual finite tables; no separate certificate of approximation existence or uniform integrability is supplied. The square of a finite-table density is the auxiliary doubled-coefficient density times the exponential of its actual energy, yielding a uniform \(L^2\) bound. Actual convergence and Vitali then give

\[
D_f=\exp(I_P(f)-A_f/2),\qquad E_P D_f=1,\qquad Q=D_f\cdot P.
\]

Normalization is therefore not an unconditional assertion for arbitrary \(L^2\) exponentials. An actual predictable representative and a time mask produce an actual control \(a\), zero outside \((0,T]\). The process \(W_t=B_t-\int_0^t a_s\,ds\) is a complete Brownian process under the same \(Q\). Every event in the original \(\mathcal F_s\) is independent of the entire subsequent \(W\)-increment readout. For a general path-dependent limit, the natural filtration of \(W\) is not assumed equal to the original filtration, nor is a reverse SDE used to remove that obligation.

Conditional moments of these actual \(Q\) increments construct finite sums and an isometry; native closure then extends them to construct \(I_Q\). Its domain is explicitly the closure of the span of basic inputs. This section does not identify that domain with all \(Q\)-predictable \(L^2\) inputs. The original tables for the same bounded \(f\) enter this domain through \(dt\times Q\ll dt\times P\), an actual almost-everywhere subsequence, and bounded convergence. Proving substitution of \(W=B-\int a\) for actual finite tables and comparing the limits of the two actual sides yields

\[
I_Q(f;W)=I_P(f;B)-A_f\quad Q\text{-a.e.}.
\]

The coefficient of energy here is one. The integral \(I_Q\) is first constructed by isometric extension, not defined as the right-hand side. An integration horizon \(U\) may exceed the density horizon \(T\): actual Brownian noise in \(W\) continues after drift stops. A zero density horizon does not make the integral of every subsequent input zero.

Weak \(L^2\) pairing of the actual density and a strong \(L^2\) limit of the same finite controlled integral further give zero terminal mean. The original \(\mathcal F_0\) is not replaced by the bottom σ-algebra; the original time-zero past is constant on a \(Q\)-full event. These results construct the original \nolinkurl{TerminalGirsanovData}, whose returned controlled integral is precisely the constructed \(I_Q(f;W)\). The original consumer then gives

\[
\mathrm{KL}(Q\Vert P)=\tfrac12 E_Q A_f\le K^2T/2.
\]

Every measurable readout satisfies the corresponding data-processing bound. On \((1,2]\), an actual feedback control uses the sign of \(B_1\) to select \(1/2\) or \(3/2\), giving the actual cost interval \([1/8,9/8]\). Zero horizon or zero control gives actual KL equal to zero. This terminal record is not a general all-time \(Q\)-integral process, and it does not complete arbitrary controlled SDEs or all source obligations of the original HJB interface.

\section{39.35 Actual Partition Quadratic Variation for the Same Fixed Joint Input}\label{actual-partition-quadratic-variation-for-the-same-fixed-joint-input}

For any finite noise dimension \(d\), let \(f\) belong to the native Euclidean-valued predictable \(L^2\) space on the entire joint past of §39.27. The output \(M_f\) is the \textbf{scalar} continuous integral of that input against the \(d\)-dimensional noise, with actual energy \(A_f(t)=\int_0^t\|f_s\|^2ds\). For any family of deterministic finite monotone clocks running from zero to the same \(T\), whose actually computed mesh tends to zero,

\[
\sum_j\bigl(M_f(t_{j+1})-M_f(t_j)\bigr)^2
\longrightarrow A_f(T)
\]

in \(L^1\) of the original source and in measure for the actual readout. Every partition uses the same fixed \(f\) and the same \(M_f\); an approximation error for tables changing with the mesh is not presented as a fixed-input limit.

The proof first derives finite-table errors from Gaussian fourth moments conditional on the entire joint past. Products \(H_iH_j\) within a row are not assumed independent. Conditional centering and orthogonality across rows give the mean-square bound \(2d^4K^4\,\mathrm{mesh}\,(t_{\rm last}-t_{\rm first})\). Left traces and the original essential bound for inputs in the finite span give a common full event. Truncated analytical representatives remain exactly in the original input class. The caller's clock samples at START, and actual \(L^2\) convergence gives the result for a fixed span input. Density of the original complete domain and a uniform stability bound for partition error then extend it to a fixed complete-domain \(f\).

The complete input \(f\) need not be bounded or in \(L^4\). For example, a past Gaussian increment of another coordinate can be used directly as a coefficient, with native second-moment integrability constructing the same complete input. Zero dimension, zero horizon, empty partitions, and repeated times are retained. Sections 39.40 and 39.42 further give weights on the same state and a finite-state Itô formula with explicit derivative bounds. Random stopping-time partitions, almost-sure convergence along partitions, and calculus domains not covered here retain separate obligations.

\section{39.36 Submitted Programs and Implementations Returned by Computation}\label{submitted-programs-and-implementations-returned-by-computation}

Instruction on finite Bool registers is a current-state controller issuing the original admissible Option Bool actions. Its role differs from that of the protected inherited Bool repair program. A static checker recursively handles tests, outputs, and nested branches and returns an actual coarse-level AST. Structural induction on syntax proves that the returned program issues the same action as the submitted program at every actual state. Actual updates and projection then give equivalence of the original \nolinkurl{RepairProgramValid} and a consumer of P9 target preservation. Valid here means preservation of the registered target, not repair of every possible damage. An incorrect write may pass the syntax checker while failing Valid.

The static checker is incomplete: it may reject a hidden test even when the two branches have equal semantics. The exact finite checker instead computes the truth table of the same interpreter and is equivalent to equality of actions on the original observation fibers. It computes and returns an original admissible two-branch coarse program. Successful checking proves that \textbf{this returned program} implements \textbf{the submitted program}. Conversely, exact realization by any coarse current-state action function requires the check to succeed. Rejection is therefore evidence of nonrealizability for this all-state action contract, not a universal denial concerning all history-dependent policies or a smaller reachable domain. Canonical representatives used by the checker appear only in computation and proof; they do not become actual resets, preparation of initial states, or additional observations.

Every finite raw fault word repeatedly interprets the same submitted program at the fresh state of the same retained parent. The complete record of actual parent--child relations, faults, failed births, protected \nolinkurl{program/history}, and writing costs commutes word by word, and complete projected laws agree for any correlated finite-word PMF. The actual repair example retains costs \([1,0,1]\) and births \([\mathrm T,\mathrm F,\mathrm T]\). A program genuinely dependent on a hidden test can instead produce different actual births and costs at the same coarse observation, giving a strict counterexample. The wordwise result in this section is restricted to finite words. Section 39.41 separately constructs countable continuing records for the same AST. Costs remain the original write-transaction charges; no internal branch-reading or runtime cost is added silently.

\section{39.37 Actual Adaptive Readouts and a Separately Registered Direction Policy}\label{actual-adaptive-readouts-and-a-separately-registered-direction-policy}

On the same four-row observation source, first read each row \(m>0\) times and then allow \(K\) additional reads. Each read accesses only the next original coordinate of the selected row, recording the row, index, value, and actual cache. The computed \(n_*\) is the shortest cache length and satisfies \(m\le n_*\le m+K\). Actual decoding is precisely this prefix of the original sourceData. A Good event simultaneously covering every prefix in that interval gives, for the original selector,

\[
P(\mathrm{Accept}\cap\mathrm{Wrong})\le8(K+1)e^{-m/8}.
\]

There is no conditioning on \(n_*\) or on acceptance. Unread coordinates occur only in the proof event, not as selector inputs. For the registered \(m=128,K=8\), both complete formal words retain all 544 latent symbols and actually read 520. Their computed shortest caches are 130 and 128, respectively. These are actual execution-reproduction facts; the general decoding and risk results are proved in Lean. Concrete read counts are not presented as another kernel theorem.

Preserving the marginal law of each read is insufficient to apply an independent-sample bound. In a correlated source that repeats one fair sign within each row, the same original selector accepts an incorrect action with probability at least \(1/16\). This counterexample is retained.

A separately registered new policy acts only on actions accepted by the original selector: if the mean of the same observational anchor row is negative, it flips the Bool; otherwise it retains it. Original rejection remains rejection. There are no new actions, samples, observations, or thresholds. Wordwise, errors of the original means are at most \(0.1\). The original acceptance-score gap and anchor lower bound imply that the actual query-mean gap corresponding to the new action exceeds \(0.3>0.2\). The original action-gap consumer and actual positive gains then prove correctness of the new action. Hence the new policy's ``accepted and wrong'' set is empty under arbitrary row/time dependence and arbitrary actual read-routing. This is a result about the new policy and does not rewrite the original policy's positive risk.

Useful acceptance still requires a separate source-separation condition: a legitimate weak-gain source may reject at every sample size. The general routing theorem is uniform for a given history callback; it is not a program sandbox forbidding that callback from capturing external parameters. The original actual sampling batches, thresholds, and failure words remain unchanged.

\section{39.38 Computing Carriers and Physical Cut Sets from Actual Response Tables}\label{computing-carriers-and-physical-cut-sets-from-actual-response-tables}

Given finite source-state, observation, and target tables, compute the actual distinguishing set \(D_{xy}\) for each pair of source rows with different targets. Sufficient reading in the original sense is exactly equivalent to hitting all these sets. A structured frontier adds atoms only from the current distinguishing set and removes nonminimal generated candidates. The general engine requires only decidable equality of atoms; it does not enumerate every subset of the ambient universe. The returned carrier table is exactly the original family of inclusion-minimal carriers, and the actual find decoder returns the same target at every source row. The returned cut-set table is exactly the original blocker.

A newly registered duplicate-pair snapshot source has rows \(\operatorname{Option}(\operatorname{Fin}m)\), observation atoms \(\operatorname{Fin}m\times\mathrm{Bool}\), and a uniform full-support law. Its actual responses yield one distinguishing edge for each duplicate pair. The complete actual carrier output has exactly \(2^m\) members. For any fixed natural numbers \(C,k\), some such source has output size exceeding \(C\) times the \(k\)th power of the total number of cells in its observation and target tables. This is a strict obstruction arising from complete explicit output size. It does not establish the same complexity for finding one carrier or for a compressed answer, and it makes no claim of a polynomial-delay algorithm.

Physical deletion writes None exactly on the deleted set and retains Some of the original readout everywhere else. The frame and failure conditions are derived from this actual update before invoking the original \nolinkurl{minimal-erasure/blocker} consumers. Physical deletion is not replaced by an absence-of-repair assumption. A constant-target zero-duplicate source has one empty carrier and no cut set; a blind-observation source with differing targets has no carrier and a cut-set table consisting of the empty set. These are separately registered snapshot-deletion permissions. They do not silently expand the original \nolinkurl{repair/calibration} actions or prove universal equivalence with death or organizational autonomy.

\section{39.39 Readouts, Infinite Rewards, and Original Control Consumers under a Fixed Actual Measure}\label{readouts-infinite-rewards-and-original-control-consumers-under-a-fixed-actual-measure}

The deployment input \(b\), bound \(K\ge0\), and finite deployment window \(A\) retain their actual measure \(Q_A\). The integral of another bounded original predictable test input \(g\) is first constructed through an actual isometric closure. Limits of finite actual outputs and the cross-drift term then prove

\[
I_P(g)=I_{Q_A}(g;W)+\int g_s a_s\,ds.
\]

Changing the readout horizon \(U\) changes only the native time input and integration endpoint; it does not replace the measure by \(Q_U\). The actual control vanishes after \(A\), while the physical state \(X_t=x_0+\sigma B_t\) continues to diffuse under the same \(Q_A\). It also equals \(x_0+\sigma W_t+\sigma\int_0^t a_sds\) and is not described as a zero-drift Gaussian state.

For a bounded \(C^2\) test function \(\varphi\) with explicit global derivative bounds, the same \(\sigma\varphi'(X)\) is both an actual original predictable input and a member of the above \(Q_A\) integration domain. The original actual scalar Itô formula and the genuinely constructed test integral give, for each \(U\), the following \(Q_A\)-almost-sure path identity:

\[
\varphi(X_U)=\varphi(x_0)+I_{Q_A}\bigl(\sigma\varphi'(X);W\bigr)
+\int_0^U\left(\sigma a_s\varphi'(X_s)+\tfrac12\sigma^2\varphi''(X_s)\right)ds.
\]

The expectation identity follows from zero mean of this same integral, integrability, and actual Fubini. Zero mean alone is not presented as a process-martingale property. Absolute continuity on intervals and the almost-everywhere derivative of the actual expectation construct the original \nolinkurl{DynkinExpectationCertificate}; integration by parts then yields the discounted identity. The state--action pushforward of the actual discounted \(dt\times Q_A\) measure gives the original \textbf{unnormalized finite occupation measure}. Only the terminal pushforward is a probability law.

If the discount rate is strictly positive and the registered reward is measurable and bounded on the actually admissible \(|a|\le K\), actual integrability of the exponential tail, Fubini, and the native improper-integral limit produce the original \nolinkurl{reward_tendsto.} It is not an input field. These ingredients construct the original ItoRun field by field. The model generator identity records only its actual operator structure, not the required Itô or reward-limit conclusion.

One nonconstant common consumer uses \(V(x)=2+\sin x\), \(\sigma=2\), and discount rate \(1\), and registers

\[
r(x,a)=V(x)-\bigl(2a\cos x-2\sin x\bigr)-a^2.
\]

Every declared bounded predictable finite-deployment policy with bound \(K\) has its own actual run. The original \nolinkurl{ControlFamily} and P-CTL-03 show that each policy's actual infinite-horizon reward is at most \(V(x_0)\), while the actual zero-action selector attains \(V(x_0)\) exactly. Reward is an explicit evaluation contract; the provenance of the process, occupation measure, reward limit, and ordering is derived separately. Policies may have different \(A\), without any assertion that the measures \(Q_A\) are projectively consistent across \(A\). This is not a general solution of the HJB equation, nor does it complete infinite deployment, unbounded control, or arbitrary SDEs.

\section{39.40 Actual START Weights and Cross Variation on the Same State}\label{actual-start-weights-and-cross-variation-on-the-same-state}

This section remains on the original joint-noise product source, with inputs \(f,g\) in the same complete predictable \(L^2\) domain. Polarization and the \nolinkurl{addition/subtraction} identities for the same continuous version first derive the original cross energy \(\int\langle f,g\rangle dt\). Weighting by a fixed matrix and weighting by a state-dependent process are different problems.

If \(H\) belongs to the original source \(L^\infty\) and can read the entire original \(\mathcal F_s\), window multiplication of native inputs satisfies

\[
I\bigl(H\,1_{(s,t]} f\bigr)=H\bigl(M_f(t)-M_f(s)\bigr)
\quad\text{a.e.}
\]

An empty window returns zero before coefficient permissions are considered. A readout from \(\mathcal F_s\) cannot be borrowed in advance into the time-zero past. For a jointly measurable, strongly adapted, source-almost-surely continuous, explicitly globally bounded live process \(H\), actual START sampling on the caller's clock converges in native \(L^2\). Its actual integral is exactly the START-weighted sum of increments of the same \(M_f\).

Applying unit windows to this same sampled input further shows that each cell increment is exactly \(H_{\rm START}\Delta M_g\). The original grid cross sum is therefore exactly

\[
\sum_j H_{t_j}\,\Delta M_f(t_j)\,\Delta M_g(t_j)
\longrightarrow\int_0^T H_s\langle f_s,g_s\rangle ds.
\]

Convergence of the native \(L^2\) inputs and a grid-uniform \(L^1\) stability bound give this actual limit. There is no jump from unweighted convergence to a randomly weighted conclusion. Any bounded continuous \(F\) applied to the same actual finite state supplies a genuine state weight. Complete inputs \(f,g\) need neither \(L^4\) membership nor samplewise bounds. Repeated times, zero dimension, and zero horizon are retained; random stopping-time partitions are outside the conclusion.

\section{39.41 Actual Countable Continuing Records of Submitted Programs}\label{actual-countable-continuing-records-of-submitted-programs}

The same external AST is reinterpreted at the state of each actual parentAfter, and the same raw mode/label stream drives the next transaction. For each \(N\), the first \(N\) entries are literally the original \nolinkurl{programRecords} execution on the same raw prefix, not a None tail after stopping. Each coordinate factors through a finite raw prefix, proving actual measurability; the countable-product interface then supplies a probability law on complete records. No measurability assumption is imposed on a map from arbitrary uncountable discrete record tuples to lists.

The same \(q\) actually returned by synthesize gives equality after projection of each coordinate, the entire countable path, and measurable-event probabilities, allowing arbitrary temporal dependence in the raw stream. P9 target preservation is used together with the operational law of that returned program. The protected inherited Bool program retains the same value at every time. Faults, failed births, actual history, and writing costs are all retained. The prefix theorem proves that the default values in the finite adapter are unreachable; they do not enter continuing execution.

These remain continuing records of a retained parent, not an infinite offspring population, and they do not establish permanent legitimacy under all faults. No internal branch-read or interpreter-runtime costs are added. The original \nolinkurl{paid/fault/paid} costs and birth records remain an actual prefix of the same continuing source.

\section{39.42 Actual Taylor and Itô Formulae for the Same Finite State}\label{actual-taylor-and-ituxf4-formulae-for-the-same-finite-state}

Let the state dimension \(m\) and noise dimension \(d\) be arbitrary and finite. Let \(X_j=M_{f_j}\) use the same original joint source and a family of complete predictable inputs, so the initial state is zero. The native second Fréchet derivative of a \(C^2\) function is connected, through actual derivatives along straight lines, to the existing purely scalar Lagrange bound. Uniform continuity of the Hessian on a compact ball containing the same state path supplies an actual remainder modulus. The maximum, over the actual partition cells, of the absolute Taylor remainder divided by the squared norm of the state increment (with zero assigned to a zero increment) tends to zero; together with the original energy bound for the sum of squared state increments, this makes the sum of actual Taylor remainders tend to zero in measure. This remainder conclusion requires neither global derivative bounds nor \(C^3\) regularity nor \(L^4\) integrability.

The actual Itô formula separately retains the global native operator-norm bounds \(\|D\varphi\|\le K_1\) and \(\|D^2\varphi\|\le K_2\), ensuring admissibility of the original \(L^2\) integral input and Hessian weights. The actual gradient of the same state computes a concrete itoInput, represented by \(\sum_j\partial_j\varphi(X_s)f_j(s)\). The actual first-order sum converges to the integral of that input; by §39.40 the Hessian sum converges to the actual time energy; the remainder tends to zero. Finite telescoping and uniqueness of a common subsequential limit yield

\[
\varphi(X_T)=\varphi(0)+M_{\rm itoInput}(T)
+\frac12\int_0^T\sum_{j,k}\partial_{jk}\varphi(X_s)
\langle f_j(s),f_k(s)\rangle ds
\quad\text{a.e.}
\]

No Itô, bracket, or remainder certificate is put back among the assumptions. An independent consumer uses unit inputs on two independent coordinates to generate the full-rank triangular state \((\Delta B_0,\Delta B_0+2\Delta B_1)\), with actual cross-time energy \(T\); the mixed Hessian of \(\sin(x_0+x_1)\) equals \(-1\) at the specified point. Quadratic functions with unbounded gradients verify only the Taylor-remainder conclusion and are not presented as satisfying the global derivative assumptions for the full Itô formula here. Section 39.45 adds an explicit initial state and actual bounded predictable drift. General localization for unbounded \(C^2\) functions and implicit SDEs do not follow automatically from this construction.

\section{39.43 Actual Q-Integral Martingales on the Original Filtration}\label{actual-q-integral-martingales-on-the-original-filtration}

Fix the same deployment measure \(Q_A\). Independence of the actual future of \(W\) from the entire original past, together with the tower property and window splitting, gives conditional expectations at any observation time \(u\). A row-start coefficient \(H\) multiplied by \(W_b-W_a\), conditional on the original \(\mathcal F_u\), returns exactly the same increment truncated at \(\operatorname{clamp}(u,a,b)\). Times before the row, inside it, after it, and after \(A\) are all included. The filtration is not replaced by the natural filtration of \(W\).

Actual finite \(W\) paths consequently give continuous, strongly adapted actual martingales on the original filtration, returning the same \nolinkurl{simpleIntegral} at the integration endpoint. The native indicator continuous linear map on \(dt\times Q_A\), actual finite-window identities, and continuity on a closed space prove that truncation preserves the original explicit integration closure; truncation membership is not an input field. The genuine \(Q\) conditional-expectation projection and a dense linear-map identity further prove that projection equals actual truncation. Thus every same input in this closure has an actual integral martingale on the original \(\mathcal F_t\), actual second-order time energy, and time-zero and terminal identifications.

For general closure inputs, the identities hold almost surely separately at each time. They do not take an arbitrary intersection of uncountably many full events, and a general sample-continuous version has not been proved. Section 39.48 separately proves that the complete original predictable \(Q\) domain \(L^2\) equals this closure. All-sample continuity of finite actual paths is not transferred to general inputs. A consumer with control window \(A=3\) and integration window \(U=6\) uses a nonconstant coefficient from the original past at the second time. The actual increment from \(4\to5\) still has variance one; the corresponding weighted output has second moment at least \(9/16\), and the unit input has energy six. The end of a deployment window is not the end of the noise.

\section{39.44 Deriving Restoration from Checks on Actual Updates}\label{deriving-restoration-from-checks-on-actual-updates}

Target preservation and restoration are different: idle can preserve an already intact target while never restoring a damaged state. The new \nolinkurl{restorationCheck} enumerates finite actual states and computes whether the update after interpreting the submitted program restores the retained port to the target from every state. Passing the check implies that the same complete physical-state kernel equals a reliable target write. The original dynamical realization, program-specific recovery basin, and eventual permanent-target conclusion in the absence of faults therefore receive \textbf{that submitted program}, without an externally supplied \nolinkurl{realization/recovery} certificate.

The same \(q\) returned by synthesize passes this restoration check if and only if the submitted \(p\) passes. From that same pair accepted by the executable checks, the original \nolinkurl{PreservesRepairability} is derived and the actual P9 repairability consumer is invoked. This is a sufficient criterion for one-step restoration, not a complete test for all eventual recovery or history-dependent policies.

When the initial protected program equals the target, the actual fault-exposed retained parent's legitimacy after one step is exactly equivalent to that step's actual mode being repair. The original fault-free permanence result therefore cannot be applied to a source with continuing faults. Constant writing and conditional repair can produce the same physical restoration while paying different actual costs in an intact state; repairability is not thereby identified with optimal purpose or funding viability. No new inherited AST, observation, initial-state reset, or resource is installed here. Section 39.47 adds actual long-term legitimate frequency, purpose, and GOD--GOA. All six P9 semantics, organization, and \nolinkurl{splitting/fusion} remain subject to their original separate obligations.

\section{39.45 Actual Drift Integrals and Itô's Formula for the Same State}\label{actual-drift-integrals-and-ituxf4s-formula-for-the-same-state}

State and noise dimensions are arbitrary and finite. Retain the same original joint source, a complete predictable input family, and an explicit initial state, and define

\[
Y_t=x_0+\int_0^t\beta_s\,ds+M_f(t).
\]

The drift \(\beta\) is strongly predictable with respect to the entire original joint past and has an explicit global samplewise norm bound. Its actual Lebesgue primitive is therefore strongly adapted, jointly measurable, and samplewise continuous; the norm of an increment is at most the bound times the interval length. Its sum of squared increments is controlled by the mesh. Mixed drift--noise sums tend to zero through tightness of the original quadratic variation and a Cauchy bound, with no new \(L^4\) condition.

Actual drift sums use \(\Delta\int\beta ds\). Measurability alone does not ensure that the left endpoint value of \(\beta\) times a time step equals or approximates that increment. The deterministic indicator of rational times is a predictable source; on the actual unit horizon \([0,1]\) its time integral is zero, while its left sums on every rational uniform grid equal one. The correct result retains the original measurable drift class rather than adding time continuity to conceal the counterexample.

For a \(C^2\) function \(\varphi\) with globally bounded native first and second Fréchet derivatives, the gradient along the same \(Y\) produces a concrete original integration input \(g\). Actual START sums, weighted drift sums, Hessian cross sums, Taylor remainders, and finite telescoping jointly give

\[
\begin{aligned}
\varphi(Y_T)&=\varphi(x_0)+M_g(T)+\int_0^T D\varphi(Y_s)[\beta_s]ds\\
&\quad+\frac12\int_0^T\sum_{j,k}\partial_{jk}\varphi(Y_s)
\langle f_j(s),f_k(s)\rangle ds.
\end{aligned}
\qquad\text{a.e.}
\]

This is an actual identity for the constructed process, not an existence theorem for arbitrary implicit SDEs. Globally bounded drift cannot be applied directly to the unbounded drift \(-kX\). The next section derives the OU process separately from its actual explicit mechanism.

\section{39.46 Physical-Clock OU and the Original Recovery Consumer}\label{physical-clock-ou-and-the-original-recovery-consumer}

The earlier time-changed OU construction gives the correct transition law, but its sample paths cannot simply be treated as Brownian-driven paths on the same physical clock. Set \(\bar B_t=B_t-B_0\) and define the actual Volterra state

\[
X_t=e^{-kt}\left(x_0-k\sigma\int_0^t e^{ks}\bar B_sds\right)
+\sigma\bar B_t.
\]

It starts samplewise at \(x_0\) and is adapted and continuous on the original physical filtration. The native fundamental theorem of calculus gives directly

\[
X_t=x_0-k\int_0^t X_sds+\sigma\bar B_t.
\]

On the same full event \(B_0=0\), \(\bar B\) reduces to the actual \(B\). The condition \(B_0=0\) is not imposed on every raw sample. Native deterministic Gaussian integration, actual \nolinkurl{Riemann/integration-by-parts} bounds, and uniqueness of \(L^2\) limits bind the same stochastic integral to the Volterra path above. A measurable limit of future coordinates and independence from the entire past of the original source give the actual Gaussian innovation law.

When \(k>0\), this physical path's own conditional kernel equals the original OU transition kernel exactly. Equality of the actual physical pushforward law and the earlier valid OU law is \textbf{equality in law}, not equality of sample paths. The actual quadratic Dynkin certificate is derived from that law, and the original P-REC-02 receives the physical path's own pushforward law, giving

\[
\mathbb E[X_t^2]\le e^{-2kt}x_0^2+
\frac{\sigma^2}{2k}(1-e^{-2kt}).
\]

The actual second moment here attains the right-hand side. An independent consumer uses \(k=3/2,\sigma=-4,x_0=5\), the original physical past, a nonlinear past readout, and the actual conditional kernel. Time zero, repeated times, and zero noise are retained. When \(k=0\), the new physical process is \(X_t=x_0+\sigma(B_t-B_0)\) with variance \(\sigma^2t\). Its native unit-exponential integral has variance \(t\); the time-change clock that degenerates to zero in the old formula cannot be used. General SDEs and localization for unbounded functions remain separate from this explicit OU source.

\section{39.47 Long-Term Behavior, Purpose, and GOD--GOA for the Same Actual Program}\label{long-term-behavior-purpose-and-godgoa-for-the-same-actual-program}

When the initial protected program equals the target and the submitted AST passes the actual restoration check, legitimacy after one step is exactly the indicator of the actual repair mode. For the original registered IID mode/label source, let the fault probability be \(\varepsilon\). The actual legitimate frequency is then almost surely \(1-\varepsilon\). If \(\varepsilon<1\), legitimate states are revisited infinitely often; if \(\varepsilon>0\), the original ``eventually permanently legitimate'' event has probability zero. Only these mode indicators have been proved IID. Complete action, cost, and history records are not claimed to be IID. Paying costs also does not automatically provide resources.

The same actual Organism, submitted program \(p\), and computed return \(q\) continue into the actual infinite-record law, integrable discounted reward, and original control consumer. Prices, costs, and discounting are explicit evaluation contracts. The relationship from actual interpretation, updates, births, and costs to that reward is separately derived from the source. The original native Bellman result compares all declared fine-state history-dependent causal policies, not merely the Bool compiler class. The actual record-reward identity and this comparison yield purpose transport and P9 consumption for the same \(p/q\).

The original contract with fault probability \(1/4\), unit price and cost, and discount \(1/2\) still gives conditional-repair value \(21/16\) and constant-paid-writing value zero, although the physical recovery paths can coincide. Another explicitly registered zero-price contract has value \(3/2\) and fine-level optimal actions depending on hidden state that do not factor through a coarse current-action representation. It does not change the original contract or rule out the existence of a correct coarse optimal policy.

GOD is the original set-valued Bellman correspondence. Optimal actions of an actual program may differ from the tie-breaking choice of a canonical greedy selector. Actual retained-parent recursion and the original IID source derive the positive-time physical marginals, the actual selector's transition rows, and an invariant law, which then enter the original \nolinkurl{InvariantLawGOA.} This joint optimality consumer specifically uses the original conditional-repair program \(p=\operatorname{compile}(i,\mathrm{goal})\) and its computed return \(q\), retaining \(0<d<1\), price times write cost at most one, and an initial protected program equal to the target. The same returned \(q\) therefore has native optimal reward, GOD, and GOA simultaneously. Passing an arbitrary restoration check does not itself imply optimality. The coarse constant-row kernel has a unique invariant law, while passive coordinates retained at the fine level permit different invariant laws.

A constant-paid policy can still have the same physical GOA while having zero value. Under all faults, a GOA can also lie entirely on states of zero legitimacy probability. Weak optimality does not imply recovery; stationarity does not imply being alive. Actual IID marginal provenance and native value adaptation are bound separately here, without an additional claim that entire adapter paths agree or that arbitrary temporally dependent raw streams retain the same optimality. All six P9 properties, general organization, and \nolinkurl{splitting/fusion} retain their separate original obligations.

\section{39.48 The Complete Original Predictable Q Domain and Genuinely Unbounded Coefficients}\label{the-complete-original-predictable-q-domain-and-genuinely-unbounded-coefficients}

Fix the \(Q_A\) produced by the same bounded deployment. First, on the original predictable \(\sigma\)-algebra, define independently its native \(L^2(dt\times Q_A)\) subspace. Successive truncations of an actual predictable representative produce admissible bounded original \(P\) inputs. The existing target-input theorem places each truncation in the constructed \(Q_A\) integration closure. Under the actual \(Q_A\), \(L^2\) dominated convergence and closedness give the reverse inclusion; predictability of the original actual finite tables gives the forward inclusion. Equality of the complete domain and the original closure is thus a proved conclusion, not a renaming by definition.

The identity-class isometry for the same native input yields the actual full-domain integral, an actual martingale on the original filtration, the conditional process, and truncated time energy. The caller no longer supplies closure membership externally. The filtration remains the entire original past; it is not assumed to equal the natural filtration of \(W\).

An independent nonconstant-control consumer, with deployment window \(A=2\), retains the same \(Q_A\) and, over integration window \(U=5\), uses \(H=W_1\). Under the actual \(Q_A\), this coefficient lies in \(L^2\) but not in \(L^\infty\), and

\[
I_{Q_A}\bigl(W_1\,1_{(1,4]}\bigr)=W_1(W_4-W_1),
\qquad \mathbb E_{Q_A}[I^2]=3,\qquad\mathbb E_{Q_A}[I]=0.
\]

Truncation and limits of actual finite records prove this product formula for the same integral. An older lemma restricted to essentially bounded coefficients cannot be applied directly. Row starts, repeated windows, zero integration horizons, and the presence of noise even when \(A=0\) are checked separately. Process identities for general inputs still hold almost surely at each fixed time; they do not assert arbitrary uncountable common full events, a general sample-continuous version, or every localized Itô formula.

\section{39.49 An Actual Time-Varying Body and Returns of the Same Parent}\label{an-actual-time-varying-body-and-returns-of-the-same-parent}

One registered round services all ports in the order of \(\mathrm{Fin}\,n\). Each step remains an original Option Bool instruction with actual mode/label input and a literal transition. The round clock, actual costs, \nolinkurl{failed/actual} births, and retained parent are all in the complete record. If the initial protected program equals the target, BodyGoal---every port matching its target at round end---is determined by the mode of the last service of each port. It differs from the earlier current-port Legitimate predicate.

For the original registered IID mode/label source, let the fault rate be \(\varepsilon\). A round succeeds completely with probability \(p=(1-\varepsilon)^n\). When \(n>0\), actual infinite-record execution performs continuing nonempty rounds, without None padding. Equality of finite prefixes with actual sequential execution yields measurability and a probability law for countable records. For any starting round \(N\) and number of rounds \(R\), the probability of no return is exactly \((1-p)^R\). If \(\varepsilon<1\), the same retained parent returns to BodyGoal infinitely often almost surely. If \(\varepsilon=1,n>0\), the body is never entirely intact after service. The case \(n=0\) is the empty-body boundary. Overlapping health indicators at the microstep scale are not claimed to be IID.

The actual conditional expectation of the zero--one defect from the same initial state, given the original past containing all previous costs, faults, and child objects, supplies the drift with \(\kappa=0,\eta=1-p\) required by the original P-REC-01, and hence its mean and tail bounds. This is not yet the expected hitting time of the original P-REC-04. An independent execution with three ports and fault rate \(1/3\) has success probability \(8/27\), while its complete records still contain actual charges and failures.

\section{39.50 Identity and Executed Repair of Actual Child Objects}\label{identity-and-executed-repair-of-actual-child-objects}

The earlier local birth gate checks only the current port. An actually born child can therefore remain damaged at other ports; the birth condition cannot be changed retrospectively to claim that its entire body was intact at birth. Source definitions derive the actual child's copied program, the parentAfter parent--child relation, and a history dropLast bridge. Only execution of a complete repair-mode service round on that same child restores every port, while preserving program and history.

This actual provenance consequently enters the original P8, stagewise P5, and SafeControl. It proves continuation of the same child object, not another birth. No old action at a fixed port can repair a different damaged coordinate, so \nolinkurl{RepairProgramRecoverableAt} for the earlier single stationary kernel is not silently replaced by a time-varying sequence of kernels. The correct consumer uses the actual finite service round and stagewise contracts. In an independent four-port example, birth costs are \([0,0,5,0]\), and subsequent repair costs for that same partially damaged child are \([5,5,0,5]\). The subsequent stochastic law is an explicitly registered new continuation source; no claim is made about a joint conditional law at a random birth time or automatic funding.

\section{39.51 Dynamic Purpose at All Prices and the Actual Returned Program}\label{dynamic-purpose-at-all-prices-and-the-actual-returned-program}

Retain the original single-port mechanism, actual state-history causal policies, and registered discounted evaluation. Set \(r=1-\varepsilon\), \(\kappa=\mathrm{price}\cdot c\), and \(0<d<1\), keeping actual cost \(c\) separate from its evaluation price. The original transition, cost, and Bellman rows directly yield the threshold

\[
\kappa(1-dr)=1.
\]

When the left-hand side is at most one, the original conditional-repair program is optimal among all declared fine-state causal policies. When it is at least one, the actual idle program is optimal, with bad/good state values \(0\) and \(r/(1-dr)\), respectively. Both selectors are allowed at the boundary. If \(r>0\), the corresponding GOD conditions have converses as well. If \(r=0\), every action is weakly optimal and the value is zero; neither the nondegenerate converse nor a recovery conclusion can be carried over. An evaluation price of zero can coexist with a positive actual cost. The set-valued GOD does not impose a canonical tie-breaking choice.

The same actually selected \(p\), the \(q\) actually returned by semantic synthesis, the complete infinite records, and their integrable reward jointly enter the original P9 and native optimal value. The original counterexample with \(c=3,\mathrm{price}=1/2,\varepsilon=1/4,d=1/2\) is retained: the one-step gain is \(-3/8\), but infinite discounted repair value from the bad state is \(3/32\). With \(c=4\), idle is optimal; \(c=16/5\) is the boundary with both GOD actions. This is an exact classification for known parameters; it does not itself supply true parameters to a data-driven selector. The earlier product-price sufficient condition of §39.47 remains valid, while this section extends the applicable price domain.

\section{39.52 A Finite-Resource Obstruction on the Same Native Policy Law}\label{a-finite-resource-obstruction-on-the-same-native-policy-law}

Fix an arbitrary original Type0 fine-state causal policy \(\pi\) and its actual history, with the initial parent's physical state matching the current state of that history and its protected program equal to the target. Suppose the actual zero-income resource condition holds almost surely on the same \nolinkurl{recordPathLaw}: actual accumulated cost at every prefix is at most a finite reserve \(R\). This is neither a discounted nor an average budget, and it does not assert that an arbitrary policy is automatically affordable.

When \(\varepsilon>0,c>0\), a native supersolution for the source transition rows gives a finite discounted-prefix bound. Letting the discount tend to one only within that same finite polynomial yields a uniform bound on partial sums of actual expected health. Genuine summability is proved first, then transferred to the actual binary health sequence using native \(L^1\)/Tonelli and a common full event. Let \(h_0\in\{0,1\}\) be initial health and \(H\) the total health occupation after all actual transactions. Then \(H\) is integrable and

\[
\varepsilon\bigl(h_0+\mathbb E H\bigr)
\le h_0+R/c\le 1+R/c.
\]

The actual subsequent transaction states of the same policy are almost surely eventually never Legitimate. The auxiliary proof price \(1/(\varepsilon c)\) does not alter the actual policy, costs, or evaluation. The conclusion bounds total health occupation, not the expected last healthy time or a calendar-time first hit.

The actual clock-visible policy ``write once, then idle'' derives reserve \(R=c\) from its own literal costs. An independent consumer with four ports, cost \(7/3\), fault rate \(3/5\), and a bad initial state gives \(\mathbb E H\le5/3\). Boundary cases excluded by the assumptions still use the same native law: zero-cost constant writing with fault rate \(\varepsilon<1\) has positive expected health at every period; one write with zero faults has expected health one at every period. These expectations are not summable. Independent consumers additionally use bounded dominated convergence to refute almost-sure eventual death. No other IID law is substituted, and almost-sure infinite return is not inferred merely from marginal means.

\section{39.53 Formation Transport for Actual Injected Blocks}\label{formation-transport-for-actual-injected-blocks}

Let an injection \(e:\mathrm{Fin}\,k\to\mathrm{Fin}\,n\) retain every port actually visited. Original literal steps, actual child \nolinkurl{objects/programs/histories}, sequential execution, and complete finite traces all commute with this restriction. Native calculations on independent coordinates derive the restriction law of the actual product source. Hence the record law of the embedded \(k\)-step scan equals one round of the earlier \(k\)-port record law; it is not the law of a complete \(n\)-port scan.

The fine-level full BodyGoal holds exactly when service at every visited port uses repair and all unvisited coordinates initially meet the target. Block health cannot replace whole-body health. Actual child formation, binding at different next ports, and parental identity thereby enter the original C4/FBT/P8 consumers. An independent example with \(n=5,k=3\), reversed ports \([4,2,0]\), and a false target checks the complete trace and product law. The same block observation with \nolinkurl{intact/damaged} unvisited coordinates yields full-body \nolinkurl{success/failure}, respectively. The repeated list \([2,0,2,1]\) uses the same original input for each port, so repeated visits cannot be claimed as fresh independent samples. The boundaries \(k=0,n=0\) are treated separately. This section uses only finite transport; infinite block laws and transport of hidden-history policies retain separate source obligations.

\section{39.54 Data-Selected Dynamic Programs and Joint Risk from the Same Source}\label{data-selected-dynamic-programs-and-joint-risk-from-the-same-source}

The rational-data selector reads only the target bit in the original first-birth record, actual paid-cost bins of the same declared width, and the target-row birth frequency \(\widehat r\) in the complete original Batch. Source records derive the target bit and cost interval, while the original calibration supplies the unique common Good event. The selector reads neither the true target, true fault rate, nor true cost, and no additional cost upper bound is required. Let actual binning give the price--cost interval \([\underline\kappa,\overline\kappa]\), and let the originally registered radius be \(u\). Then

\[
\begin{aligned}
\overline\kappa\bigl(1-d(\widehat r-u)\bigr)<1
&\Longrightarrow\text{select conditional repair},\\
1<\underline\kappa\bigl(1-d(\widehat r+u)\bigr)
&\Longrightarrow\text{select idle}.
\end{aligned}
\]

Source-derived margin inequalities and the actual all-price theorem of §39.51 ensure that accepted programs are genuinely GOD/optimal on Good. Incorrect acceptance, and UNKNOWN under explicit separation and \(u<r\), are both controlled by the complement of that same Good event, with original bound \(4\exp(-2Nu^2)\). No birth, an exact boundary, or inadequate separation can honestly remain UNKNOWN. Zero samples and all-fault boundaries are retained.

The same actual \(p\) and synthesized return \(q\) execute on the actual \nolinkurl{postProbeParent}, preserving physical state, protected code, passive coordinates, all training costs, and history. The actual \nolinkurl{training-times-deployment} Markov joint law has the original training marginal and the same joint risk of incorrect acceptance. Actual infinite reward and the original P9 refer to this very policy. They do not assert reward maximization on each trajectory or optimal net training reward, and funding is not derived from reward.

A preregistered formal source example fixes \(N=64,u=1/4\), all 128 original inputs, price \(1/2\), width \(1/100\), and discount \(1/2\). The source gives \(\widehat r=3/4\) and cost bins 220/600. With \(c=11/5\), the result is repair, whereas the earlier one-step choice is idle. The actual terminal parent is damaged and the same \(p/q\) has value \(15/32\). With \(c=6\), the result is idle. The joint bound \(4e^{-8}<1\) is checked. These are complete mathematical source examples, not newly collected empirical data.

The batch with original width \(1/4\) has not been changed to finer bins. An independent consumer retaining that width obtains bins 8/24 and dynamic repair/idle, but its earlier low-cost choice is UNKNOWN; it cannot borrow the earlier idle choice of the fine-width example. Although the original \(c=3\) has positive infinite repair value, the new choice remains UNKNOWN at fixed \(N=64,u=1/4\) for both widths, without extra samples or changed thresholds. A separately registered formal example with \(N=96\) and width \(1/20\) checks all 192 actual readouts. Its actual post-training parent is intact, the same \nolinkurl{low/high-cost} \(p/q\) values are \(9/5\) and \(3/2\), and total training costs are \(2592/5\) and \(1728\). Together, the different terminal states and nonzero costs test the connection without resets or free training.

\section{39.55 Actual Payment under the Same Training--Deployment Joint Law}\label{actual-payment-under-the-same-trainingdeployment-joint-law}

The original \(2N\) training transactions and actual post-training parent of §39.54 remain unchanged. The AST returned by the same data then continues to execute under the native infinite-deployment law. Every original literal transaction costs an amount in \([0,c]\), including failures, faults, UNKNOWN, and incorrect selections. The source derives this cost bound; the caller is not asked to supply the desired path-viability certificate.

Now use the already registered open-environment permission: an independently known upper cost bound \(C\ge c\), external supply \(g\ge C\) arriving before payment for each actual transaction, and initial reserve \(R_0\ge0\). Let \(e_t\) be the actual cost in the same complete training-plus-deployment record. Then

\[
B_{t+1}=B_t+g-e_t,\qquad B_0=R_0.
\]

Almost surely under the same actual joint law, every prefix satisfies the original P6 \nolinkurl{ResourceViable}, and each charge is affordable before payment: \(e_t\le B_t+g\). The case \(R_0=0\) is allowed because supply arrives before the first charge. A different supply order requires a different contract. No earlier waiting action is inserted, no different protocol of length \(4N+1\) substitutes for the original \(2N\) trajectory, and no post-training state is reset. The measurable augmentation by a capital coordinate gives a probability law whose erasure returns exactly the original \nolinkurl{training-times-deployment} jointLaw.

The supply condition is a sufficient environmental permission, not a necessary threshold for all policies or a conclusion about endogenous energy. The selector does not read \(C,g\) or the true cost. GOD/P9 and infinite value for the actual \(p/q\), and the original unique calibrated risk, retain the same provenance. Funding holds almost surely without adding another statistical failure budget, but optimal net training reward has not been proved.

The complete 128-input example with original width \(1/4\), \(N=64\), and actual cost \(11/5\) retains, under \(C=g=3,R_0=0\), the actual damaged terminal parent and value \(15/32\). An independent 160-input formal example with \(N=80\) has actual training cost 312, a bad terminal parent, and the same \(p/q\) has value \(321/400\); every charge is payable under \(C=4,g=9/2,R_0=1/2\). In the same example with \(g=1/2\), the first balance is \(-8/5\) on positive joint mass. Mathematical source examples do not replace the original empirical data. For the original single-bit compiled policy, the zero-income no-AE obstruction remains when \(c>0\), \(0<\varepsilon<1\), and the given initial reserve is finite. Under all faults, actual costs can be zero, so this result cannot be generalized into a resource obstruction for every positive fault rate or every policy.

\section{39.56 Countable Block-Record Laws and the Information Obstruction for the Complete Body}\label{countable-block-record-laws-and-the-information-obstruction-for-the-complete-body}

The actual injected block of §39.53 now continues indefinitely, still with exactly \(k\) original transactions per round. The actual retained parent and every finite sequential-record prefix commute with port restriction. The native infinite-product pushforward theorem and the proved actual single-round source restriction derive the law of the entire input stream. Measurability of each record coordinate follows from actual finite-source factorization. This gives equality of pushforward laws for the same complete countable record, permitting transport of every measurable record event; a finite-prefix result is not merely renamed an infinite result.

The original \(k\)-port body-return and P10 consumers therefore apply to the same block target in this fine-level \(n\)-port source. The success probability is \(p=(1-\varepsilon)^k\), and the probability of no return over \(R\) rounds after any starting round is exactly \((1-p)^R\). When \(\varepsilon<1\), the block target is revisited infinitely often almost surely. The complete fine-body target is still equivalent to the block target together with all unvisited coordinates initially intact. Without PassiveReady, the same actual stream cannot make the complete body intact at any clock time.

Two fine parents differing only in a hidden passive initial state have identical entire projected infinite records, while for \(\varepsilon<1\) their positive masses of complete-body success differ. Consequently, no decoder reading these complete block records can determine complete-body success almost surely correctly under both actual initial-state sources. The proof uses actual positive mass, not a zero-probability all-success word. An independent consumer with \(n=5,k=3\), ports \([4,1,3]\), fault rate \(2/5\), and cost \(7/3\) checks six actual transactions with total cost 7, transport of countable tail events, failure rate \(98/125\), and the original P10. Empty blocks, empty bodies, and all-fault boundaries are treated separately. This is not a complete \(n\)-port scan and does not automatically transport arbitrary fine-history control/GOD, resources, or discounted time units.

\section{39.57 The Original Hitting-Time Consumer on the Same Actual Path}\label{the-original-hitting-time-consumer-on-the-same-actual-path}

Each new round-end physical state is generated from the current original input independently of the \(\sigma\)-algebra of all preceding original inputs. Together with the actual retained-parent recursion, this derives the actual physical path's own homogeneous transition kernel, its actual Dirac initial law, and the complete path law of the original \nolinkurl{homTrajMeasure.} On this path, the original hittingAfter defines the first hit \(\tau\) of the registered body target, including time zero. It is a measurable stopping time for the original complete past. Registered execution continues; no new stop or restart is performed.

Actual survival-tail counts are bound to the original nonnegative \nolinkurl{hittingValue.} Nonnegative tail integration on the same source gives \(\mathbb E\tau=0\) if the initial state is already in the target, and otherwise

\[
\mathbb E\tau=\bigl((1-\varepsilon)^k\bigr)^{-1}.
\]

The values are extended nonnegative numbers: the reciprocal of zero success probability is infinity. The actual bad-state indicator potential is nonnegative, and integrating the actual transition row gives decrease \(p=(1-\varepsilon)^k>0\) outside the target. The original P-REC-04 therefore consumes the same kernel, initial state, and drift to give \(\mathbb E\tau\le1/p\); the exact tail formula gives equality for a bad initial state. Source proofs cover zero for an initially good state, one for a bad initial state with zero faults, zero for an empty body, almost-sure nonhitting from a bad initial state with positive body dimension under all faults, and almost-sure finiteness of the actual stopping time when faults are not certain.

This \(\tau\) observes the target only after the \(k\) transactions of each round, not at the earliest microstep hit. An independent example with two ports, a true target, fault rate \(1/3\), and cost \(5/2\) gives \(p=4/9\), \(\mathbb E\tau=9/4\), and the same upper bound from the original P-REC-04. In its actual word, the body is intact after transaction three and faulty again after transaction four, but the registered first hit is at the end of round three: six transactions and actual total cost 5. Round-end time cannot be identified with microstep first-hit time or finite-resource lifetime, nor can the discount clock be changed. The section proves a link from an actual path to the original consumer; it does not supply externally a geometric variable with the desired mean.

\section{39.58 Actual Returned Repair and Parent History in the Same Complete Record}\label{actual-returned-repair-and-parent-history-in-the-same-complete-record}

Fix the original model family whose initial protected program equals the mechanism target. When the original complete Batch returns repair, the actual first-birth record has already proved that the returned target bit equals the mechanism target; calibration Good is not needed for that fact. All restoration branches of that same returned AST are therefore checked, instead of choosing a different repair program known merely to exist. The actual \nolinkurl{postProbeParent} may already be damaged. It enters the original recoverable P5 carrier and P12 maintained-parent endpoint directly, without resetting to an intact seed. The endpoint invoked is precisely parent maintenance. An \nolinkurl{organizational-representation} clock is not treated as a literal transaction clock, nor is the entire strict finite-lifecycle endpoint claimed to be realized.

Action support and recursion for the native retained parent prove that, for any actual deployment program, parent history before and after each transaction retains the input history. The actual original probe recursion retains it as well. Thus, under the same original \nolinkurl{training-times-native-deployment} jointLaw, parent history before and after every transaction in wholeRecords equals the original parent history. This holds on the same full-measure event as the original P6/prepay of §39.55, including all training costs and \nolinkurl{incorrect/UNKNOWN} outputs.

Preserving actual parent history does not mean identifying it from Batch. Two parents with the same initial physical state and program but different actual histories produce exactly the same original visible Batch and original input law. For every parent-history estimator reading the complete Batch, the error probabilities under the two actual sources sum to at least one. The old Batch does not thereby acquire permission to read a global parent ID. Explicitly registered identity channels remain separate from this nonidentifiability conclusion.

An independent 144-input formal example with \(N=72\) retains the same previously frozen input word, width, bad terminal parent, and actual repair/idle results. Low/high actual costs are 4/20 and total training costs 432/2160. The actual P12 carrier at low cost 4, parent history under the same joint law, and \nolinkurl{funding/prepay} are checked with \(C=9/2,g=5,R_0=1/3\). High cost 20 gives an idle, nonrepair counterexample and does not borrow the lower-cost bound or supply certificate. The original impossibility of normalized strict generational increase on a finite permanent set of individuals is not withdrawn, and cyclic histories are not used to fabricate actual multiple generations.

\section{39.59 Native Complete-Program Laws and Actual Joint Long-Term Maintenance}\label{native-complete-program-laws-and-actual-joint-long-term-maintenance}

The original IID interpreter records of an AST and the native causal-deployment records previously had separate definitions. They cannot be identified in law merely because their physical effects agree. Now, for any original AST \(p\) and actual initial state, the proof starts from event marginals of the native growing \nolinkurl{action/mode/label} words. Strong lumping to the last primitive, path naturality of the original \nolinkurl{homTrajMeasure}, and the actual shift removing the time-zero default yield an actual IID primitive innovation stream. Action support of the same \(p\) and the actual retained-parent recursion reconstruct every complete record, giving

\[
\operatorname{Law}_{\rm native}(\text{actual full record};p,o)
=\operatorname{Law}_{\rm primitive}(\text{the same actual full record};p,o).
\]

This equality in law does not receive ``source matching'' as an external input, nor does it require \(p\) to be synthesizable or to pass the restoration check. Programs depending on an unobserved passive port, failing synthesis, and not restoring the target are checked independently as well. It does not assert equal records for arbitrary policies with the same physical restoration: compiled repair idles in a good state, while constant writing may still incur a cost. The cost sequence is not treated as IID.

If the initial protected program equals the target and the same \(p\) passes the actual restoration check, this complete-law identity transfers the original long-term theorem to that native deployment law: legitimate frequency is almost surely \(1-\varepsilon\); when \(\varepsilon<1\), legitimate states are revisited infinitely often; when \(\varepsilon>0\), legitimacy is not eventually permanent. The repair program \(p\) actually returned from data and the \(q\) actually returned by synthesis each obtain their own native complete-record law and corresponding frequency. Initial states remain the actual \nolinkurl{postProbeParent} and its actual projection.

Under the original \nolinkurl{training-times-native-deployment} jointLaw, one full event simultaneously gives parent history for complete training plus deployment, P6/prepay, and the future frequency, return, and nonpermanence conclusions for every actual repair output. These execution conclusions no longer require Good externally: the actual birth record already ensures correctness of the returned bit. Good remains the statistical evidence for optimal purpose/GOD, and the original unique calibration budget is not counted twice. A high-cost branch returning idle is not rewritten as repair or a guarantee of life.

An independent consumer reuses the same \(N=72\), 144 inputs, and actual terminal state. Exact enumeration of every \nolinkurl{primitive/native} event for two-step prefixes compares all fields of the complete records; the full infinite-law identity is proved separately in Lean. For a program passing the actual restoration check with matching initial protected code, zero faults make every actual post-transaction state legitimate, all faults make every such state bad, and idle from a bad initial state stays bad forever. Zero cost can still permit positive health frequency, so the positive-cost resource-exhaustion result cannot be applied. This removes the \nolinkurl{native/primitive} complete-record-law assumption left in §39.58. The same-tick source for P5 organization is treated separately in §39.64; the information condition for parent identity from Batch remains.

\section{39.60 Finite Linear Rank and Existence of a Two-Dimensional Factorization}\label{finite-linear-rank-and-existence-of-a-two-dimensional-factorization}

The sufficient direction required by the original A5 differs from GL uniqueness between two already-given factorizations. For a finite-dimensional real input space \(X\), any real vector space \(Y\), and a linear response \(J:X\to Y\), the following is now proved:

\[
\dim\operatorname{range}J\le2
\quad\Longleftrightarrow\quad
\exists H:X\to\mathbb R^2,\ G:\mathbb R^2\to Y,
\quad J=G\circ H.
\]

The arrows here denote linear maps. Actual \(H,G\) are constructed from the actual range embedding and a native linear left inverse; ranks zero and one permit zero padding. If the response has rank exactly two, every given actual two-dimensional factorization has surjective \(H\) and injective \(G\). Once these conditions have been derived from dimension, the existing nondegenerate GL-uniqueness and independently registered unit-intervention-anchor consumers are actually invoked. The matrix-specific equivalence is bound to the matrix of the same actual \(J\).

The independent mixed matrix

\[
\begin{pmatrix}
2&1&1&-3\\
4&-1&2&3\\
0&5&0&-15
\end{pmatrix}
\]

is bound entry by entry to an actual four-input, three-output linear map and has rank exactly two. A nontrivial shear with coefficient \(-2/3\) changes the encoder while preserving every response; the original unit-anchor consumer removes its coordinate freedom. Degenerate \nolinkurl{rank-zero/rank-one} factors and the rank-three obstruction to a two-dimensional factorization for the three-dimensional \(2I\) are rigorously checked. This supplies finite linear existence; low rank does not determine the physical names Π/Φ, the real existence of independent interventions, exact rank under noisy measurement, or a common global nonlinear factor. The original chain-rule necessity for a general differentiable common bottleneck and the existing noisy calibration are reused within their respective scopes.

\section{39.61 Six actual properties of the same data-selected program}\label{six-actual-properties-of-the-same-data-selected-program}

The original P9 endpoint requires all six properties to belong simultaneously to one actual object, program, and scale. The original timestamped probe now proves that erasing the time registration gives the original probeTrace, that the terminal parent's physical state, program, and history are those of the actual \nolinkurl{postProbeParent}, that the owner is preserved, and that the clock advances by the original \(2N\) training transactions. The initial owner and starting clock are explicitly registered source information; the selector does not infer them from the old Batch.

For this actual clocked parent, the original data-returned AST, the genuine idle comparison program, and the original single-port scale are used to derive, one by one, the actual registered identity, the actual infinite-horizon value ordering, the original Bellman GOD, the GOA of the actual native marginals, actual recovery without changing the goal, and a positive long-run legitimate frequency. The original six-property endpoint is then invoked. The ordering compares actual repair with actual idle, rather than a program with itself; the six conclusions are not stored in a record left to be supplied. Legitimacy here remains Legitimate for the originally registered active port; it does not mean that every physical port satisfies BodyGoal.

Execution recovery requires the actual repair output and an initial protected program matching the target. Positive long-run frequency also requires \(\varepsilon<1\), which is derived almost everywhere from the complete actual source. Optimal ordering and GOD still use the original Good calibration evidence, together with a nonnegative price, positive observation width, and discount \(0<d<1\). A single full event under the same original \nolinkurl{training-times-native-deployment} jointLaw simultaneously supplies parental history, original P6, \nolinkurl{transaction-by-transaction} prepayment, actual long-run frequency and return, and, on Good, the six properties of this very object and scale. The supply condition remains an explicitly external income per transaction sufficient to cover the preregistered cost cap; the mechanism does not generate its own energy.

No additional observation or threshold is introduced to exclude the all-fault source: under the actual training law, that source returns UNKNOWN almost everywhere, so an actual repair output implies \(\varepsilon<1\) almost everywhere. Even if a zero-mass input word computationally returns repair, this does not establish that the true noise is less than one. An independent formal consumer with \(N=8\) uses a different owner, starting clock, and cost, and all 16 original inputs, to verify the original endpoint. It is not a new empirical batch, and the statistical upper bound at this small budget is not described as a nontrivial high-confidence guarantee.

\section{39.62 Two-generation dispatch by the actual data-returned AST}\label{two-generation-dispatch-by-the-actual-data-returned-ast}

For the same complete training data, take its actual returned AST and the actual terminal parent's physical state, protected program, and Bool history, and register serial slot0 under an explicit new protocol. This registration starts with the root mark as its ID and clock zero. Registration transport is still needed to relate it to the source owner and starting clock retained in the preceding section; equality of physical state and program alone does not establish identity. There is only one actual founder.

Each request looks up its parent in the actual store and then executes this single AST; an actual child retains the same program. The original allocator, actual P8 parent bridge, and the respective P5 results act on actual records, without supplying a separate birth satisfying the desired conditions. On a branch with actual repair output and initial protected code matching the target, the two normal steps of dispatch plan \([0,1]\) actually produce two generations \(0\to1\to2\). A fault in the first step remains an actual failure record wrapped in Some, with None in slot1. If the second step cannot find a parent, there is no transaction, no child, and no fee. Actual fault transactions that do exist are retained; None cannot stand in for a successful record or receive income.

The complete original training law times the two-step primitive input law retains all these branches. The original low-cost reference has costs \([11/5,0]\); the previously frozen independent formal example with \(N=72\) and 144 inputs gives \([4,0]\) at the different cost of four. Actual data are not resampled. The first actual birth cohort already enters the original TypePrice, but a single parent with a constant trait has covariance zero; this is not an adaptive advantage. Nor is this finite dispatch the original infinite deployment law retaining a fixed parent: its funding and global GOD are not borrowed from that law.

\section{39.63 Price terms from actual histories and the failure boundary}\label{price-terms-from-actual-histories-and-the-failure-boundary}

First establish child1 by the single founder's actual first birth, and then dispatch to the actual old cohort \([0,1]\). The complete plan is \([0,0,1]\). The original complete training remains present through inputLaw and the probe; \nolinkurl{experimentRecords} outputs the subsequent three Option records, without claiming that the training trajectory has already been prepended to this output.

On the branch just described, with actual repair output and initial code matching the target, the first actual child under mode word TFT is UID1. The second step fails and leaves an empty slot at UID2; the third step produces UID3 from the actual UID1. The old parents have history depths one and two, and the new child has depth three. The original birth count \(B=1\) counts only the last two dispatches to the old cohort, excluding the first birth that establishes the cohort. Applying the original Price identity to the actual incidence and fertility gives selection term \(1/2\), transmission term one, and mean increment \(3/2\). The costs are \([11/5,0,0]\); failure costs have not been changed into success costs.

On the same source, FFF has positive mass but \(B=0\), so no normalization is performed; the final request may yield None. An independent consumer, using the same earlier formal example with \(N=72\), 144 inputs, and the full three-step input law, also verifies a positive-mass TTF branch: selection term \(-1/2\), transmission term one, and increment \(1/2\). A positive realization therefore gives neither a universal nor an expected selection advantage. The branch in which the first step fails and the root gives birth in the second also shows that \(B>0\) does not imply a child in the third step. Successful outcomes have not been selected to hide this boundary.

This source with strictly finitely many generations preserves the earlier impossibility of normalized strict generational increase on a finite permanent set of individuals. It uses a rectangular parent--offspring channel of actual child events, rather than cyclic identities to restore the permanent finite strict lineage already ruled out by the counterexample.

\section{39.64 One organizational step per literal transaction}\label{one-organizational-step-per-literal-transaction}

The previous parent-maintenance result already used an organizational carrier, but its representation clock had not yet been identified with literal transactions. The actual \nolinkurl{parentBefore.physical} and protected \nolinkurl{parentBefore.program} are now read directly from complete native records and connected to the original mode-free \nolinkurl{safeRepairKernel.} When this actual AST passes \nolinkurl{restorationCheck} and the initial protected program equals the target, the original physical shadow, the same actual initial state, and one-step kernel commutation derive the complete homTraj organizational path law, including time zero, and every marginal at the same \(t\). One literal transaction is one organizational update here; no two-step clock that first restores code and then executes an action is used.

The program is read from the actual record, not replaced by the mechanism's target. Native parent reconstruction is first established for an arbitrary AST. An actual repair output then implies the restoration check; only after combining this with equality of the initial program and target may the corresponding organizational path law and P12 maintained-parent result be invoked. P5 and P12 on this native source share a full event with parental history, \nolinkurl{P6/prepayment}, and long-run life under the same complete training--deployment jointLaw. The organizational carrier allows recovery after a bad state; it does not assert health at every time.

The independent consumer retains the original \(N=72\) and 144 inputs, and checks time zero, one step, and arbitrary times. Incorrect protected code remains incorrect. A constant-write program and compiled repair can have the same organizational law while paying different fees, four and zero. The organizational projection therefore loses the actual AST, fee, label, and child; it is not an invertible isomorphism of complete records. The existing cap \(9/2\), income five, and balance \(1/3\) for the low cost of four do not automatically pay a high cost of twenty.

This section supplies the same-tick organizational source missing in Sections 39.58--39.59. At the source stage recorded here, the registration transport between the source owner and clock in Section 39.61 and the serial population in Sections 39.62--39.63, the actual resources of the complete finite-generation process, and original P12 parent-identity discovery still had to be connected to the same process. Registration views do not turn population dispatch into infinite deployment of a fixed parent. The extra permissions and fees for reading marks must be explicitly registered, and the old Batch still cannot identify global parent identity. The OPEN / \nolinkurl{NOT_ESTABLISHED} assessment belongs to that source stage. Subsequent Sections 39.65--39.95 and the core overview record its later resolution and the scope of the current \nolinkurl{CLASSIFIED_CORE_CLOSED} status; the model restrictions just stated remain in force.

\section{39.65 Registration transport for source identity and actual dispatch clocks}\label{registration-transport-for-source-identity-and-actual-dispatch-clocks}

The actual postClocked in Section 39.61 preserves the source owner and post-training starting clock, whereas the actual population uses a serial root mark and zero clock. An explicit registration view is now constructed: the root mark is replaced by the actual founder's history prefix, and subsequent positive phase marks are offset by a registered epoch. Thus the initial serial root represents exactly the actual postClocked after the complete original training. Its parameters come from explicit source-registration permissions; they are not global identities estimated from the old Batch, and no second founder is added.

The epoch cannot be added uniformly to all clocks. An actual newborn non-root child has clock zero, so a uniform epoch shift gives the wrong result; a strict counterexample already exists within the original model. The correct view preserves the newborn's zero clock and offsets only actual positive dispatch phases. Actual lookup, THIS AST requests, and the interpreter establish commutation of step, child, run, the complete finite trace, and its law, without inserting the desired parent equality into a field. Injectivity of reachable identities, reflection of SameObject, strict parenthood, and original P8 are thereby preserved.

Actual history depth changes by a fixed constant, while the original incidence and fertility remain unchanged; hence the Price selection term, transmission term, and mean increment are all preserved. The independent earlier formal example with \(N=72\) and 144 inputs uses a different founder prefix and starting clock 17, with epoch 161. Actual dispatch clocks are 162, 163, 164, or None; newborn clocks remain zero. Exact replay of all 64 three-step primitive words is not new empirical sampling.

Store.length also increases for an invalid None request, so phase is generally not the number of actual transactions. Registration transport does not remove this strict boundary. Nor does registered time turn finite dispatch to children into the original infinite evaluation protocol retaining the parent. P9/GOD retains its original protocol type: at the original cost \(11/5\), price \(1/2\), and fault rate \(1/4\), if only the final step of a bad parent is evaluated, compiled repair has actual expected immediate payoff \(-3/40\), while idle has payoff zero. This verification does not refute the original infinite-horizon value optimality, and it is not an optimal solver for the whole finite population.

\section{39.66 Executable discovery of an actual birth's parent mark}\label{executable-discovery-of-an-actual-births-parent-mark}

The true parent mark comes from the actual three-step episode of this same data-returned AST, not from a separately supplied parent-identity answer. The last request in the explicit protocol family may select root0 or the actual child1; candidate1 is established only by an actual first birth. An additional sensor reads the actual birth's parent mark, with allowed observations None, Some0, and Some1. The majority algorithm directly returns this operational serial mark. It does not call actualParent, the hidden request address, or a world owner lookup to complete the answer. These may serve as semantic proof bridges, but must not silently be added to the algorithm's inputs.

For any one preregistered event slot, fixed mark noise \(\eta=1/4\), an arbitrary unknown confuser, and all \(B\) readings, the joint law of the original complete training, the three-step future, and the readings gives

\[
\Pr(\text{returned tag differs from the actual event's parent tag})
\le \exp(-B/8).
\]

The actual ParentBridge and the discovery endpoint of original P12 are derived from actual birth histories and invoked. Both None branches without a child and false positives remain present; an event slot is not necessarily an actual successful transaction. This is a risk bound for one preregistered target, not a free simultaneous guarantee for arbitrarily many targets.

The independent consumer retains the original \(N=64\) and the full formal word of 128 inputs, takes a nonuniform unknown confuser, and exactly handles every count class of the fixed 64 readings. With \(B=64\), the upper bound is \(\exp(-8)\), and the error risk for the actual finite set of readings is also proved strictly positive. Each reading costs \(1/20\), for a total of \(16/5\) for 64 readings, payable even when the actual child is absent or a reading or estimate is wrong. A fixed budget limits the number \(B\) of readings; sending the error to zero as \(B\) grows requires a separately specified family of resources growing with \(B\).

There are two distinct information boundaries. Under positive execution noise, two actual episodes have a common positive-support atom of bare observations but different true parents. Any decoder based only on those bare observations therefore has a sum of error risks across the two experiments at least equal to that common positive atom's mass; at least one experiment has strictly positive risk. This does not assert equality of the complete bare laws: actions or fees at an intermediate fault may distinguish them. The swapped-confuser counterexample at \(\eta=1/2\) instead proves only that, for different fixed true marks, the laws of all mark-read sequences coincide. This is nonidentifiability for that read channel; it does not prove equality of the complete world laws of training, dispatch, and readings. Reading a complete ancestral ID for free changes the parent-information conditions, so the conclusion of the bare-observation counterexample cannot simply be carried over.

\section{39.67 The original P6 ledger for all training and actual finite dispatch}\label{the-original-p6-ledger-for-all-training-and-actual-finite-dispatch}

Funding is now based transaction by transaction on the original interpreter: retain all \(2N\) original training transactions, followed by the finite population's Option records from the same actual AST. Each actual fee is nonnegative and at most the original writeCost. Fault transactions remain in the record; an invalid None has physical cost zero. If the preregistered finite horizon \(H\), cost bound cap, and initial capital reserve satisfy

\[
\text{writeCost}\le\mathrm{cap},\qquad
(2N+H)\mathrm{cap}\le\mathrm{reserve}
\]

then original P6 \nolinkurl{ResourceViable} and \nolinkurl{transaction-by-transaction} prepayment follow directly for that same complete finite cost stream, with income identically zero. The zero-cost extension after the finite endpoint is merely a ledger convention; it does not fabricate infinite None execution. No idle-before-work, income, free second founder, allowed action, or primitive sample is added.

The formal reference preregisters \(N=64,H=3,\mathrm{cap}=9/4,\mathrm{reserve}=300\). Actual training costs \(1056/5\) in total, and the positive-support TFT branch has total physical cost \(1067/5\). Initial capital covering only the training amount is strictly insufficient on TFT from the same source, so it cannot fund this complete protocol almost everywhere. The independent earlier formal word with \(N=72\) and 144 inputs uses cap \(9/2\) and reserve 665: actual training costs 432 and the entire TFT branch costs 436. Training-only capital 432 gives a negative balance of four; FFF, B0, None, and branches with different fees are retained. This refutation concerns this protocol and ledger, not every strategy or the originally registered external supply under another contract.

This section still pays only controlled-write fees. The actual consumer that adds the full read fee \(B r\) to the same initial wallet is given in Section 39.68; None and incorrect readings are not exempt. At the source stage recorded here, the correct rectangular finite-generation replacement for original P12 still required an actual connection among the owner and clock, the same returned program, actual parenthood, P5/P6, Price, the originally registered P9 evaluation, and necessary parent-identity discovery. Registration, evaluation, and sensor permissions are explicit model conditions, not synonyms for those internal derivations. The OPEN / \nolinkurl{NOT_ESTABLISHED} assessment belongs to that source stage. Subsequent Sections 39.68--39.95 and the core overview record the later connections and the scope of the current \nolinkurl{CLASSIFIED_CORE_CLOSED} status; the explicit model conditions remain unchanged.

\section{39.68 Joint consumption from one wallet and the actual first birth's parent}\label{joint-consumption-from-one-wallet-and-the-actual-first-births-parent}

The physical write-fee ledger of Section 39.67 is now genuinely combined with all mark readings in one initial wallet. Retain all \(2N\) original training transactions, the three dispatch positions of plan \([0,0,h]\) executed by the same actual returned AST, and all \(B\) readings. Preregistering

\[
(2N+3)\mathrm{cap}+B r\le\mathrm{reserve},
\qquad \text{writeCost}\le\mathrm{cap}
\]

then derives original P6 and prepayment for every transaction from the actual fees, with zero income during the finite stage. Readings follow the original complete write-fee trajectory, and their actual bill is exactly the original readCharge. None, failure of the first birth, incorrect readings, and false positives do not waive \(B r\); there is no second observation account.

The main same-parent consumer fixes the preregistered first event \(e=0\). Its request starts from the actual postClocked founder. If an actual child exists, the actual birth establishes original P8 and SameObject between the reregistered transaction actor and that founder. The clock is updated by the actual dispatch; complete equality of the Clocked tuples is not asserted. In contrast, the parent of the last event \(e=2,h=1\) is child1 and cannot be called the founder merely because its program agrees. The mark algorithm still returns only an operational mark; world lookup merely interprets the actual meaning of the returned mark, rather than providing the algorithm with new directory access.

Original P12 marked discovery, the original maintained-parent organizational family, and the six properties of original P9 are invoked for this actual first birth's parent. Optimal purpose uses the original Good event. Under the original complete joint input law, actual repair implies \(\varepsilon<1\) almost everywhere, without giving the selector access to the true fault rate. The original bad training-calibration event and actual parent-mark error are each counted once, giving under the complete \nolinkurl{training/future/readings} joint law

\[
\Pr(\text{bad calibration or incorrect parent tag})
\le \min\{1,4\exp(-2Nu^2)\}+\exp(-B/8).
\]

This bounds the union of two specified error events. It is neither conditional risk after acceptance nor a guarantee of birth or success of the entire process. For the original \(N=64,u=1/4,B=64\), the bound is \(5\exp(-8)\); first-step faults, UNKNOWN, and None retain their full probabilities. The asymptotic discovery result concerns a family of experiments paying separately for the corresponding readings, not infinitely many free readings from a fixed wallet.

In the original formal example, the same wallet of 300 pays total physical and read fees \(1083/5\) on TFT, leaving \(417/5\); on FFF it pays \(1072/5\), leaving \(428/5\). The independent original formal word with \(N=72\) and 144 inputs uses the same wallet of 665, with all 64 readings costing \(16/5\). Both an actually supported Good reading atom and an FFF false-positive atom have positive mass. On FFF the returned mark may be Some0 while the actual parent is None, and the bill is still paid. Successful samples have not been selected and failure records have not been overwritten.

The precise remaining boundaries must also be retained. The native maintained-parent contract here starts at the actual post-training state. It is another protocol retaining the parent of the same object, rather than the actual continuation after the three-step population dispatch, and the finite zero-income wallet does not fund infinite maintenance. The actual episode changes the founder. The next source must therefore begin with the actual terminal root's physical state, code, identity, and clock, retain its current balance, and explicitly register subsequent supply. The original rectangular Price result already has a source, but its role for a general positive birth count of the same actual \(h=1\) cohort still has to be consumed on that common source; zero births remain unnormalized. At this stage, the entire finite-generation chain has therefore not yet been closed.

\section{39.69 Canonical prediction and the actual program from the same Register source}\label{canonical-prediction-and-the-actual-program-from-the-same-register-source}

Existing active-port execution and value transport do not automatically mean that the canonical predictive endpoint has been applied to this source. The mass of the finite stochastic kernel is now taken directly from the original actual Model.transition, and proved equal to the original \nolinkurl{physicalSource} and the pushforward of the actual modeLaw. For the originally allowed Option Bool actions and pre-action readout \(q(x)=x_i\), induction on words of arbitrary length using actual primitive expectations gives

\[
\begin{aligned}
\operatorname{Future}_{K,q}(x)=\operatorname{Future}_{K,q}(y)
&\Longleftrightarrow x_i=y_i\\
&\Longleftrightarrow
\operatorname{projectState}_i(x)=\operatorname{projectState}_i(y).
\end{aligned}
\]

Here Future uses the pre-action emission convention of the original \nolinkurl{controlledOutputTrace}; the empty word alone cannot distinguish states. The two actual one-letter tests \([(\mathrm{none},\mathrm{false})]\) and \([(\mathrm{none},\mathrm{true})]\) reconstruct the indicator of each actual predictive class, deriving hresolve, strong lumpability, and the unique normalized quotient kernel instead of taking them as certificate inputs.

The macro reward is the same mechanism's actual \nolinkurl{expectedModePayoff.} It is first proved constant on classes and independent of the representative, after which the original \nolinkurl{predictiveFutureExactControlQuotient} is actually invoked. The entire constructed micro Model---transition, reward, bound, and discount---equals the original actual Register model. Original P-QUO exact optimal value and macro greedy lifting thus apply to this source and dominate all original fine causal competitors; this does not mean that every controller reading hidden information can be factored through the quotient.

For the synthesized pair \(p,q\) actually returned by the checker, macro actions are defined by this very \(q\), and their lift executes precisely this very \(p\) at every fine state. No other strategy is selected by choice. The mathematical definition of a quotient class may choose a representative, but actual actions have been proved independent of that representative. This canonical value commutes with the existing active-port value at every nonnegative price; GOD for the same actual \(p/q\) and original P9 scale reduction also connect to this exact micro Model.

The independent original nonzero-noise formal example with \(N=64\) retains all 128 training inputs and actually invokes repair, high-cost idle, the actual returned \(q\), the six original properties, and canonical control. States differing in hidden coordinates may be physically distinct but have the same predictive signature; an arbitrary program reading hidden bits does not automatically descend. When \(\varepsilon=1\), post-action transition rows may merge while pre-action readings still distinguish the current active bit, so the timing condition cannot be deleted.

This is a state quotient only within the same fixed mechanism and read protocol. At \(\varepsilon=0\), two different targets can have identical physical rows but different actual birth rewards and opposite strategy rankings. Physical prediction therefore neither identifies purpose across mechanisms nor identifies ancestry or owner. The original birth/cost Batch still supplies the actual program source; the selector obtains no true-parameter oracle. The named missing arrow connecting canonical prediction to actual control and purpose has been supplied. Funded continuation of the actual post-birth parent and the complete common Price source are constructed separately in Section 39.70.

\section{39.70 Funded continuation of the actual terminal parent and a common source for the original roles}\label{funded-continuation-of-the-actual-terminal-parent-and-a-common-source-for-the-original-roles}

The source in Section 39.68 cannot substitute the post-training state for the state after the population experiment. Starting from the original complete \nolinkurl{training/future/readings} inputs, this section actually executes the original requests \([0,0,h]\) and uses lookup in the actual terminal store to obtain the still-living root0. Its program and history are proved preserved; its identity, protected target, and actual stored clock come from that store, while its physical state is the actual terminal state. Even after an intermediate failed birth, it is not reset to healthy or replaced by a newborn object.

The same returned AST is interpreted on this actual \nolinkurl{currentParent.} New native fault primitives define a conditional kernel, which is joined with the original complete finite input law. Projection back to the finite stage gives exactly the original law: training, false positives, None, faults, and all readings are retained. The conditional native-word law equals the source law of the original \(p\) and the same returned \(q\) executed from the actual current state. Together with Section 39.69, this links the complete canonical micro Model, the specified strategy, and the actual process.

The original formal example's post-training clock 139 is not the current parent's clock. The actual root has clock 142 and is damaged on TFT and FFF; on TTT it has clock 142 and is healthy. The fee-conversion index 195 differs again from the stored clock. The general theorem connects \(t,t+1\) from the actual terminal root clock. It neither inserts extra physical faults because of read billing or scheduling phases nor uniformly shifts newborn clock zero.

Funding continues from the same initial wallet. Write \(F(z)\) for the sum of all actual finite write and read fees; the actual remaining balance is \(R-F(z)\), not a separately supplied budget certificate. The finite fee schedule is zero after \(S=2N+3+B\), and every continuation fee equals the actual native expenditure. A separately registered external income \(g\ge\mathrm{cap}\) arrives before each continuation tick. From \(R\ge F(z)\) and the actual write-fee bound, original P6 and \nolinkurl{transaction-by-transaction} prepayment follow for the complete schedule. This is a sufficient condition, not a claimed necessary one; nor does it assert that a finite zero-income wallet funds infinite life. The original actual TFT and FFF balances remain \(417/5,428/5\), respectively.

Original P5 and the native organizational source retain their respective general parameter assumptions. The final quantitative common consumer specifically uses the original registered parameters with \(N=64\). For every complete finite input \(z\) returning repair, the actual organizational marginal at each \(t\) under the conditional native law equals the original \nolinkurl{homeostaticMarginal} from the actual current-parent initial state \(\mu_0(z)\). This is not equality of the entire joint law to a marginal with one fixed initial state. This conditional law actually invokes the original P5 mean theorem and finite-horizon bound on bad-state occupation,

\[
\kappa\sum_{t<K}\Pr(\mathrm{bad}_t)
\le V_0+K\lambda.
\]

The mean bound comes from the same actual organizational marginals and the original drift proof, not by working backwards from an almost-everywhere frequency statement. Under the complete joint law of the same \(N=64\) instance, each tick's actual fee is measurable, bounded, and integrable. For any finite prefix, interchanging a finite sum and an integral gives original P6 for expected costs in the same \(R,g\) ledger. This does not assert a bound on infinite total cost. Expected costs here are still the actual write, read, and native-continuation fees, not an independent replacement cost process.

Native same-tick organizational records reconstruct the same parent object, program, and identity. If an actual next birth occurs, it satisfies the original parent relation and has newborn clock zero. The native almost-everywhere frequency and recovery conclusions for repair with \(\varepsilon<1\) continue to hold; optimality of purpose separately uses the original calibration event Good. The current physical state enters the original value integral, and the actual returned \(p/q\) connects to the original six P9 properties, canonical quotient GOD, and the declared infinite discounted evaluation retaining the parent. It has not become a finite-population endpoint optimum: the original finite-endpoint counterexample, repair value \(-3/40\) versus idle value 0, remains valid.

History depth, the selection term, and actual birth count for the same actual \(h=1\) cohort are produced by execution. With positive births, the original rectangular Price response is the selection term plus one; with zero births, only the unnormalized formula is retained. The Price identity is not adaptive improvement or long-run convergence. When the actual mark equals the true first-birth parent's mark and the child exists, the returned mark is bound to the store lookup actually executed, yielding original P8. This does not require the mistaken inference from ``the mark is correct'' to the unobservable marked Good event.

The original bad calibration event and actual parent-mark error are still counted only once under the same joint law, with the same upper bound \(5\exp(-8)\). Native funding, the organizational source, and exact transport are pointwise or almost-everywhere facts and do not consume additional statistical failure budget. This bound is neither a birth success rate, an entire-process reliability rate, nor conditional risk after acceptance.

The independent complete formal word with the original \(N=72\), a different owner, actual port, and target retains all training and all readings. It is checked exactly, item by item, across every future word of the original fault source and the following two native steps. Both a correct \nolinkurl{birth/readout} atom and an FFF false-positive atom have positive mass. This is a witness of nonemptiness for the formal model; no real samples were generated or selected. Initial protected-program agreement with the target, trusted interpretation semantics, actual owner and mark-read permissions, the cost cap, initial balance, subsequent income, and evaluation contract are all explicitly registered. Program execution, historical identity, conditional kernels, ordering, and organizational sources are derived by proof.

\section{39.71 Actual proofs and consumers for the original book's named source obligations}\label{actual-proofs-and-consumers-for-the-original-books-named-source-obligations}

An overall acceptance review previously found that Sections 21.7, 35.4, and 35.7 still had only written derivations, and therefore rejected the then-current draft claim of overall closure. That failed review is retained. The subsequently added \nolinkurl{ConflictParametricSource/Control/Stationary/QSD/Path/Infinite/CommonConsumer}, \nolinkurl{ConflictOrganizationalSource}, and \nolinkurl{WaitingRepairControl} now provide actual Lean proofs and independent invocations item by item, without deleting the original problems or exempting them as future directions.

The new common consumer strictly distinguishes the original M from the actually restricted R. For M, the original \nolinkurl{GreedyInvariantGOASet} is exactly μ1 outside K. For R, the complete GOA consists of all μθ, but the original \nolinkurl{CesaroLimitGOA} from the registered safe δ initial state contains only μ0. For θ\textless1, the conditional terminal distribution at any positive horizon from any safe initial state is that QSD. Zero steps retain the initial state; θ=1 gives zero survival. Conditioning of the actual complete path, finite cylinders of the same infinite experiment, and AE eventual failure when θ\textgreater0 are verified separately. Terminal damage θ, failure at some point in two steps 1−(1−θ)², and actual return θ(1−θ) arise from the same complete records and cannot be interchanged.

The independent \nolinkurl{waiting/repair} model with cost 8/5 retains its own source, rewards, and complete action domain. Original-interface consumers are also available for terminal splitting in P-ALG on a source with nonzero rewards, actual posterior support, and two-tick organizational recovery. Sources, exact quantifiers, counterexamples, actual public builds, and independent reviewer identities are recorded in the \href{../formalization/ueot-core/docs/compression/theory_completion/unified_closure/conflict_family_book_source_v1/RESULTS_AND_OPEN_OBLIGATIONS.md}{named source package}. These new proofs discharge the named written obligations listed here. A later overall review of the original scope again found missing derivations: the complete Perron source in Section 25.4 and the actual reversible spectral expansion and remainder source in Section 10.4. It therefore again rejected the then-current inference of overall completion. The earlier negative judgments are retained. The subsequent general constructions and original consumers appear in Sections 39.72--39.74 and their corresponding independent reviews. Item-by-item source settlement for the current mathematical contract is recorded separately from release adoption; a local pass or a book build cannot replace overall acceptance.

\section{39.72 A complete Perron source for every finite primitive matrix}\label{a-complete-perron-source-for-every-finite-primitive-matrix}

Original Section 25.4 starts from the same nonempty finite primitive nonnegative matrix M. \nolinkurl{PrimitivePerronSource} maximizes a compact set of feasible growth rates on the nonnegative unit simplex. A positive primitive power smooths a nonzero candidate into a strictly positive vector. If the eigenvector residual were still nonzero, the same positive power would send it to a vector with all coordinates strictly positive, permitting a strict increase of the growth rate and contradicting maximality. This constructs the original \nolinkurl{RightEigenData} rather than assuming existence of a principal root.

\nolinkurl{PrimitivePerronDoob} proves transport for every power of the same original matrix: \[
(H^n)_{ij}=\frac{(M^n)_{ij}r_j}{R^n r_i},\qquad
H_{ij}=\frac{M_{ij}r_j}{Rr_i}.
\] Stochasticity of H follows from the actual eigenvector equation, and primitivity from the same positive power. The original Cesàro existence theorem produces an actual stationary probability row for H; the left vector of the original matrix and both normalization conditions are derived from this row. For any primitive stochastic matrix, \nolinkurl{PrimitiveMarkovMixing} first proves strict Dobrushin contraction at a positive power. Nonincrease of distance at every step and a cofinal subsequence then give convergence at all original times, with the limit bound to that same constructed stationary row. Transport back to M yields \(R^{-n}M^n\to rl\).

Uniqueness of the positive eigendirection is proved using finite ratio extrema and the actual positive power. Strict dominance over the other complex spectral values follows from the actual power limit above. \nolinkurl{PrimitivePerronSimplicity} separately proves the dimensions of the full eigenspace and generalized eigenspace at the principal root, so that the root has multiplicity one in the characteristic polynomial of both the original real matrix and its complexification. No diagonalizability assumption is imposed on the whole matrix; nonprincipal modes may have Jordan structure.

The terminal \nolinkurl{certificateOfPrimitive} takes only M and a proof of primitivity and returns the original complete KPF01Certificate. Its return value for the actual nonsymmetric matrix \[
M=\begin{pmatrix}2&1\\2&3\end{pmatrix}
\] is proved to have \(R=4\), \(l=(1/2,1/2)\), and \(r=(2/3,4/3)\), and is actually used by the original growth and composition theorems for nonzero counts. Independent examples with a different matrix, a primitive matrix containing zeros, and nonprincipal Jordan structure also verify the same general interface. Lineage transmission still requires actual identity conditions; primitivity does not replace birth records. Sources and settlement of the original scope are recorded in the \href{../formalization/ueot-core/docs/compression/theory_completion/unified_closure/general_perron_compact_spectral_source_v1/RESULTS_AND_OPEN_OBLIGATIONS.md}{general Perron and actual compact-spectral source}.

\section{39.73 An actual compact resolvent generates the spectral expansion of the same semigroup}\label{an-actual-compact-resolvent-generates-the-spectral-expansion-of-the-same-semigroup}

Fix the actual evolution \(S_t\) and a compact self-adjoint resolvent R on the original real Hilbert space. The source contract requires \(S_0=I\) and, on the actual generator domain, \[
\frac{d}{dt}S_t(Rv)=S_t(Rv-v),\qquad t\ge0.
\] This expresses the execution semantics of \(R=(I-L_D)^{-1}\). Actual contractivity, the original first mode, orthogonal spectral bounds, a second mode, and a strict gap are also explicitly registered.

For each nonzero resolvent mode \(Rv=av\), \nolinkurl{CompactSpectralSemigroupSource} derives from this differential relation \[
S_tv=\exp\!\bigl((1-a^{-1})t\bigr)v.
\] The complete orthogonal Hilbert sum for the compact self-adjoint spectrum comes from the original mathlib theorem. The generator relation at time zero makes R injective, and actual contractivity makes the nonzero resolvent eigenvalues positive. The coordinate action and HasSum expansion of the complete S are obtained by applying the actual continuous linear map to that sum. The original orthogonal Rayleigh bound then gives Parseval decay for the actual remainder.

\nolinkurl{CompactKilledSpectralBinding} takes the same source's initial probability density \(g\) and an implementation by the actual killed evolution, and constructs the original \nolinkurl{SpectralData.} Both \(a_1=\langle g,\phi_1\rangle\) and the actual remainder \(S_t(g-a_1\phi_1)\) are computed directly from the source. Positivity of the leading coefficient, positivity of the principal mode's total mass, the time-zero identity, the integral squared norm, and both original spectral outputs are proved. Original P-QSD-04 is then invoked, transporting the conditional law and survival law exactly to this killed object.

The general implementation contract binds the actual evolution from the specified g; a single initial-state trajectory does not identify the entire kernel. The actual origin of the generator and resolvent, compact self-adjointness, the original spectral conditions, and source-operator semantics are explicit assumptions. The construction does not claim to establish them for every SDE. The common instance in the next section separately provides a kernel implementation for the whole operator.

\section{39.74 A nonzero spectral remainder and conditional law from the same actual kernel}\label{a-nonzero-spectral-remainder-and-conditional-law-from-the-same-actual-kernel}

Take uniform reference measure \(m\) on Bool, with \(\phi=1\) and \(\psi=(1,-1)\). The actual killed kernel is \[
k_t(x,y)=\frac12\begin{cases}
 e^{-t}+e^{-2t},&x=y,\\
 e^{-t}-e^{-2t},&x\ne y.
\end{cases}
\] Nonnegativity, the subprobability property, the identity kernel at time zero, and semigroup composition are proved directly from this table. R/S on the same weighted \(L^2(m)\) satisfy the actual generator relation of the preceding section, and the full operator identity \[
(S_tv)(y)=\sum_x k_t(x,y)v(x),
\] is proved for every vector v, together with the actual kernel-pushforward identity for nonnegative densities. This finite reversible killed instance witnesses the spectral operator class; it is not called a Brownian diffusion.

The actual initial state is \((3/4,1/4)\), so \(g=1+\frac12\psi\). Its leading coefficient is one and its initial remainder norm is \(1/2\). The same source gives \[
S_tg=e^{-t}\phi+\tfrac12e^{-2t}\psi,\qquad
\|R_t\|_2=\tfrac12e^{-2t},
\] \[
\Pr(\tau>t)=e^{-t},\qquad
D_{\rm TV}(\mu_t^c,m)=\tfrac14e^{-t}>0\quad(t<\infty).
\] An independent verification registers initial state \((5/8,3/8)\) while retaining exactly the same kernel, R, and S. It obtains initial remainder norm \(1/4\) and actual TV distance \(e^{-t}/8\), and actually invokes the same original P-QSD-04. Both are formal-model verifications retaining a nonzero remainder; no real-world samples are redrawn.

Item-by-item settlement of this edition's original mathematical contract, classification of assumptions, compatible common sources, and strict counterexamples is recorded in the \href{../formalization/ueot-core/docs/compression/theory_completion/unified_closure/general_perron_compact_spectral_source_v1/independent/ueot-original-contract-general-gaps-discharge-review-20261103.json}{review discharging the original contract gaps}. Mathematical-scope verification and actual release adoption are recorded separately. Final adoption is confirmed by release receipts matching the source version, both books, actual data bindings, and GI/index. External empirical support remains UNVERIFIED.

\section{39.75 A joint decision contract with one unknown source and a finite wallet}\label{a-joint-decision-contract-with-one-unknown-source-and-a-finite-wallet}

This section begins a new conditional mathematical contract. The completion judgments for Sections 39.70--39.74 and the original eight acceptance gates remain valid. The question here is how information, repair, dispatch, and stopping paid intervention jointly determine net payoff within the same finite process. Original samples, failures, and counterexamples are retained; new formal conclusions add no external empirical support.

The finite parameter \(\theta\sim\pi_0\) is drawn only once, at initialization. Conditional on it, each step uses the same registered primitive to generate an execution mode and an independent fair birth label; the probability of normal mode is \(r(\theta)\in[0,1]\cap\mathbb Q\). Actions, visible receipts, and noise alphabets are finite, while priors, fees, initialization, and discount have exact rational representations. The general algorithm does not require the hidden world to be globally finite: finite initial support and finite noise already ensure finite posterior support at every finite history.

The actual Register consumer has a public target bit \(g\), an inherited protected program equal to \(g\), and one actually damaged founder as its initial state. The unknown quantity is the persistent \(r(\theta)\); this family does not claim to identify an unknown program. The main menu is stop, or an original idle, write-false, or write-true request to address 0 or 1. Address 1 must pass the original Store lookup and may be absent. Each original transaction appends an actual child or None. Subsequent birth addresses are allocated by the original Store length; addresses 2 onward remain in the population records, but this menu does not acquire free dispatch permission for them.

The newly registered observation permission is a receipt at the end of every step, \[o_t=(\text{actual birth},\ \text{write fee actually paid})\in\{0,1\}^2.\] The write fee \(c>0\) is public, so the second bit is equivalent to that step's literal fee, \(0\) or \(c\). A receipt is a projection of the complete execution record. It does not return the mode, label, protected program, owner, actual parent identity, or complete population state. Receipt timing is finer than in the old batch protocol; this is an explicit new permission and cannot be combined for free with the permissions of the old physical controller or parent-mark sensor. Writing itself can provide information, while all its actual physical effects, fees, and failures are retained. The main menu contains no separate free parent-mark reading.

There is only one initial wallet \(B\ge0\), with no income, subsidy, or reset. Any write request must have balance at least \(c\) before execution and lookup, even if a subsequent fault or empty address means that no fee is charged. A normal write is charged even when no birth occurs; fault and None fees retain their original source definitions. Terminal value is zero, and evaluation is fixed in advance as \[
J_H(\pi)=\mathbb E_{\pi_0,\pi}\!\left[
\sum_{t=0}^{H-1}d^t\bigl(N_t-\lambda C_t\bigr)\right],
\qquad 0\le d\le1,\quad\lambda\ge0,
\] where \(N_t\) is the birth count of the original transaction and \(C_t\) its actual fee; discounting uses the protocol's decision index. Hard prepayment and the price penalty constrain, respectively, whether payment is possible and whether the action is worthwhile. The solver's computing-hardware cost is outside this physical-wallet model; execution time is reported separately.

stop commits the current and all remaining steps to root idle, still subject to the original environment, faults, births, and clocks. It does not freeze the world. One protocol step is one original Store transaction. Undispatched slots retain their state under the original primitive; the entire population does not age simultaneously. This permission is the actual source of the backup value below. The outer finite-table scheduler carries history and wallet information; the old memoryless AST is not tacitly assumed to store a posterior.

\section{39.76 Filtering, interpretation, and payment from the same joint source}\label{filtering-interpretation-and-payment-from-the-same-joint-source}

Write \(w\) for the hidden world containing persistent \(\theta\), \(z\) for the current noise, \(K(w,z)\) for its source weight, and \(T(w,a,z)\) and \(O(w,a,z)\) for the actual next world and allowed receipt, respectively. For a finitely supported belief \(b\), first construct the same joint weight \[
q_b^a(w,z,o)=b(w)K(w,z)\mathbf1_{\{O(w,a,z)=o\}},
\qquad m_b^a(o)=\sum_{w,z}q_b^a(w,z,o).
\] When \(m_b^a(o)>0\), pushforward by \(T(w,a,z)\) followed by normalization gives the new posterior; a zero-mass receipt uses an explicit default. The proof derives total expectation and identities for test functions of complete finite traces from these joint weights. It does not independently multiply the laws of training, modes, births, and post-selection deployment.

The hidden compressed Register Store retains only the optional physical bit at each slot, but the complete native Store still retains programs, histories, owners, parents, clocks, and labels. The invariant equating inherited programs with the public target makes the original transactions, receipts, fees, and evaluation commute with the compressed transition. The finite table receives only allowed histories. World-side compression is used for the proof and posterior; the actual hidden slot table is not supplied to the selector. The same actions and noise drive the native and compressed interpreters, with equality of transaction-record projections, final states, and all fees. The probabilities under the original modeLaw and fair labels are also proved equal item by item.

Let \(b_t\) and \(B_t\) be the scheduler's posterior and balance obtained from allowed histories. Feasible actions are filtered before execution by \(c\,\operatorname{cap}(a)\le B_t\), where a write request has cap 1 and every other request cap 0. Every prefix with positive source probability satisfies \[
0\le C_t\le c\operatorname{cap}(a_t)\le B_t,
\qquad B_{t+1}=B_t-C_t,
\qquad B_t=B_0-\sum_{j<t}C_j\ge0.
\] The implementation defaults to truncating balances at zero on zero-mass branches. On actually supported paths, this truncation is proved inactive. Positive initial atoms and positive primitive-noise paths establish positive receipt mass, after which the same prefix ledger is connected to the literal fees of the original Store. Average costs therefore do not replace payment on each prefix.

\section{39.77 A computably returned table and the complete comparison class}\label{a-computably-returned-table-and-the-complete-comparison-class}

The information state is \((b,B)\). At each finite node, enumerate the currently feasible actions and use rational comparisons to compute \[
V_0(b,B)=0,\qquad
V_{h+1}(b,B)=\max_{a\in A(b,B)}
\sum_o m_b^a(o)\left[R(a,o)+dV_h(b^{a,o},B-C(o))\right].
\] The implementation also returns a finite action table containing all receipt branches, resolves ties in a fixed order, and proves that lookup for every history after serialization returns the same action. It does not merely output a symbolic argmax, an existence claim, or a printed objective value.

\textbf{Return and optimality theorem.} For every registered finite horizon, the actual returned table has expected \(J\) in the original interpreter equal to the recursive value and is prepayable on every positive-probability execution prefix. The comparison class includes causal strategies based on complete allowed action--receipt histories and using arbitrary real randomization weights. Only histories reached with positive probability are required to be legal. Zero-probability counterfactual rows may be arbitrary; no legality requirement is imposed on these rows to restrict the comparison class. The proof first expands the actual primitives, then links them level by level to the recursion using conditional expectation under the same posterior. Any randomized mixture of actions is bounded above by the rational maximum action value.

This is Bayes optimality for a specified prior. Changing the prior or information permissions changes the problem. The theorem does not assert minimax optimality, omniscient optimality, infinite-horizon optimality, or scalable computational complexity. The finite table may be large; engineering time required by an exact algorithm is not mathematical impossibility. The current consumer reuses the original transaction and birth-identity mechanisms, and explicitly registers the new outer table scheduler as a finite-memory implementation.

\section{39.78 Actual source-parameter boundaries: continuation with two objects after repair}\label{actual-source-parameter-boundaries-continuation-with-two-objects-after-repair}

In what follows, write the total horizon as \(H+1\), fix the wallet at \(B=c\), write the priced write fee as \(p=\lambda c\), and define \[S_H(d)=\sum_{t=0}^{H-1}d^t.\] This parameter subfamily has a substantive converse boundary. It is not presented as a closed-form solution for every price and wallet; general inputs are still handled by the exact algorithm of the preceding section.

For continuation without further writes, let \(U_1(r,d,n)\) and \(U_2(r,d,n)\) denote the optimal remaining birth payoff with one and two dispatchable healthy objects, respectively. The actual source gives \[
\begin{aligned}
U_1(r,d,0)&=U_2(r,d,0)=0,\\
U_1(r,d,n+1)&=r\bigl(1+dU_1(r,d,n)\bigr),\\
U_2(r,d,n+1)&=r\bigl(1+dU_2(r,d,n)\bigr)
 +(1-r)dU_1(r,d,n).
\end{aligned}
\] These bounds are attained by choosing a known healthy object that may be dispatched and idling it: continue with it after a birth, and switch to the other surviving object after no birth. The allowed receipts suffice to make this choice. The original primitive preserves undispatched objects, so one fault does not simultaneously damage the other object. This is not a theorem of free redundancy in arbitrary whole-population environments.

\textbf{Exact formula for a general finite horizon.} If \(p\ge S_H(d)\), then \[
\boxed{\quad
V_{H+1}^*=
\max\left\{0,\ \mathbb E_{\pi_0}\!\left[
r\bigl(1-p+dU_2(r,d,H)\bigr)\right]\right\}.
\quad}
\] An attaining strategy first writes the correct bit to root when the expectation in braces is positive. After a successful birth it uses the receipt-based schedule above, and after no birth it stops paid intervention. When the expectation is nonpositive, stop attains zero; ties may admit several optimal tables.

The proof has five steps. First, from an entirely damaged population there are no births before the first charged write. If the first charge occurs at step \(t\) of the remaining process, all subsequent discounted birth payoff is at most \(d^tS_{H-t}\le d^tS_H\le d^tp\), already covered by that equally discounted fee; later fees are nonnegative. Thus, when \(p\ge S_H\), net payoff cannot be positive on any path, and without any charge the process cannot repair itself. Second, an initial action other than the correct write gives no birth and leaves a damaged population covered by this impossibility result. Third, the first correct write succeeds with the same probability \(r\), immediately yields \(1-p\), produces an actual root and child1, and leaves wallet zero. Failure retains the actual Store and the unspent wallet. Fourth, a zero wallet forbids writes by any legal subsequent action table; \(U_2\) is a source-derived upper bound for each fixed \(r\) and is attained by the same receipt policy. Fifth, take expectations over the complete process conditional on fixed \(\theta\), and compare initial repair with stop. This establishes both an upper bound and attainment, not merely containment of one strategy class in another.

Let \(\mu_1=\mathbb E r>0\) and \(M_H=\mathbb E[rU_2(r,d,H)]\). In the same high-price region, the two-direction boundary is \[
V_{H+1}^*>0
\quad\Longleftrightarrow\quad
p<1+d\frac{M_H}{\mu_1}.
\] At or above the right-hand threshold, the optimal net value is zero. A separate resource impossibility does not require the high-price condition: if every population in the positive initial support is entirely damaged and \(B<c\), the optimal value is zero at every finite horizon. This follows from the pre-execution prohibition on writing and cannot be evaded by arguing that a fault might make the transaction free.

\section{39.79 A persistent unknown parameter, actual posteriors, and short-horizon checks}\label{a-persistent-unknown-parameter-actual-posteriors-and-short-horizon-checks}

Write \(\mu_k=\mathbb E r^k\). For total horizons 2, 3, and 4, the preceding subfamily reduces, respectively, to \[
\begin{aligned}
V_2^*&=\max\{0,\mu_1(1-p)+d\mu_2\},&&p\ge1,\\
V_3^*&=\max\{0,\mu_1(1-p)+(d+d^2)\mu_2\},&&p\ge1+d,\\
V_4^*&=\max\{0,\mu_1(1-p)+(d+d^2)\mu_2
             +d^3(2\mu_3-\mu_4)\},&&p\ge1+d+d^2.
\end{aligned}
\] All these formulas use \(B=c\) and the same finite-terminal contract. With equal weights on \(r\in\{1/4,3/4\}\), \(d=1/2\), and \(p=13/10\), the two-step value is \(1/160\). Replacing each step by the mean kernel gives repair value \(-1/40\). The difference, \(d\operatorname{Var}(r)\), comes from correlation induced by the same persistent parameter throughout the complete process. It is not a profit from free information and cannot be represented by a process that redraws \(\theta\).

Actual failure receipts also affect choice. For an initial correct write, receipt \((0,0)\) is the uncharged fault branch, whose posterior moments are \[
\mu'_k=\frac{\mu_k-\mu_{k+1}}{1-\mu_1}
\qquad(\mu_1<1).
\] For a nondegenerate two-point prior with \(0<w<1\) and \(0<r_0<r_1<1\), \[\mu_1\mu_3-\mu_2^2=w(1-w)r_0r_1(r_0-r_1)^2>0.\] Consequently, when \(d>0\), the repair-price threshold for the remaining two steps strictly falls from \(1+d\mu_2/\mu_1\) to \[1+d\frac{\mu_2-\mu_3}{\mu_1-\mu_2}.\] For the two-point source above, the thresholds are \(21/16\) and \(5/4\). The price \(p=13/10\) lies between them: retaining the old prior leads to a retry, whereas the actual fault posterior gives retry value \(-3/160\). After a fault, the weight on low \(r\) changes from \(1/2\) to \(3/4\).

At the same parameters and total horizon 3, the actual general solver returns value \(27/320\). A policy that incorrectly freezes the prior and retries has value \(51/640\), a difference of \(3/640\). Here \(p<1+d\), so the preceding high-price closed-form theorem is not substituted for the actual computation. The posterior formula, row-by-row optimality verification, and complete replay on the same source supply their respective evidence. All failure records remain in the expectation; neither the source nor the object is redrawn or reset after failure.

\section{39.80 Comparisons under the same conditions, zero value of information, and native populations}\label{comparisons-under-the-same-conditions-zero-value-of-information-and-native-populations}

The three-step example uses the same equal-weight prior on \(r\in\{1/4,3/4\}\), with \(d=3/4\) and \(p=7/4\), and has optimal net value \(9/256\). The fixed procedure ``write root, idle root, idle child1'' already attains this value. Joint decision-making therefore has zero gain over this fixed procedure under the same menu. If only root may be dispatched, the initial-repair candidate has value \(-9/512\). That comparison changes dispatch permissions and cannot be presented as a learning advantage under equal permissions.

At four steps, after successful initial repair exhausts the wallet, the best birth value over predetermined sequences of three continuation requests is \(r+dr+d^2r^2\), attained, for example, by alternating between the two objects. With receipt feedback, the two-object value is \(r+dr+d^2(2r^2-r^3)\). Under initial repair and the high-price conditions, the nonnegative raw feedback difference is \[d^3\mathbb E[r^3(1-r)].\] It is strictly positive only if \(d>0\) and \(\Pr(0<r<1)>0\); it is zero if \(d=0\) or the source takes only values \(0,1\). If stop truncates either negative candidate to zero, the final optimal difference must account for the corresponding positive parts. The raw difference is not unconditionally the final gain. With equal weights on \(r\in\{3/4,7/8\}\), \(d=3/4\), and \(p=37/16\), both candidates are positive: \[
V_4^*=\frac{39525}{524288},\qquad
V_{4,\mathrm{fixed}}^*=\frac{2325}{65536},\qquad
\Delta=\frac{20925}{524288}.
\] The comparison uses the same prior, receipts, single wallet, available addresses, and complete fees. The fixed procedure differs only in not adapting requests to receipts. From these damaged initial states with \(p>1\), a myopic policy selects a zero-value action and misses potentially positive continuation payoff. Agreement with an exact full-history reference under the same observations is a correctness check, not a claim that ``UEOT outperforms general POMDPs.''

The four-step difference comes from knowing which dispatchable body remains healthy and switching accordingly. It does not establish a strict benefit from paying extra to identify \(\theta\). After success leaves a zero wallet, freeSchedule simultaneously attains, for every \(r\), the omniscient upper bound for the allowed actions. In that continuation domain, therefore, revealing \(\theta\) for free has exactly zero additional value. A positively priced reading that adds only this information while leaving the rest of the process unchanged cannot recover its cost. Parent-mark readings outside the main menu have their own timing, error, and fee contracts; this result does not make every read protocol useless.

The actual population is more than a birth count. A concrete legal trajectory of the original Store may first produce child1 from root, append None on the next transaction when root faults, and then dispatch child1 to produce a grandchild at address 3. The original transaction determines parenthood, owner, and newborn clock zero. If address 1 is absent, dispatch still yields None. The selector sees only receipts; complete identities remain in world records and consumers. Formal proofs connect native execution for all noise words and all tables. Exact enumeration retains all original branches, including faults, wrong writes, empty addresses, and no birth.

Net payoff, budget safety, and event success are reported separately. In the two-step and three-step high-price instances above, the probability of at least one actual birth is \(1/2\); in the four-step instance it is \(13/16\). This does not make net payoff nonnegative on every path or imply that birth probability itself is maximized. Balance safety on each positive-source path follows from the prepayment theorem, not from these event probabilities.

\section{39.81 Domains of the new finite and old long-run conclusions}\label{domains-of-the-new-finite-and-old-long-run-conclusions}

The earlier long-run results on maintenance, supply, and purpose with a retained parent use their own infinite records, source costs, and evaluations. The present contract instead solves the complete decision problem with a finite zero-valued terminal, no income, and one wallet. Both problems may have correct but different optimal actions. The original finite-process counterexample with value \(-3/40\) is retained; no new discount, objective, menu, or parameter overwrites it. Nor is the earlier fixed-sample-size risk bound used without proof as an adaptive-stopping guarantee.

This family assumes a public target and positive fixed write fee. It does not solve unknown protected programs, unknown charges, endogenous cognitive hardware, arbitrary sensor permissions, arbitrary dispatch addresses, simultaneous global aging, or infinite-population planning. The mathematical closure consists in computable optimality and both positive and negative boundaries, derived jointly from the same actual source, allowed information, actual table, native population, and wallet under these explicit conditions. Normative reasons for changing the objective and external empirical applicability still require independent support.

The implementation entry point is the \href{../formalization/ueot-core/docs/compression/theory_completion/unified_closure/joint_decision_closure_v1/README.md}{joint decision delivery package}. It retains the frozen contract, Lean module index, THIS returned table, complete rational replay, failed attempts, and reviews by nonauthors. Delivery receipts record the final source--output binding, affected integrations, and adoption in both books. Core parameter endpoints appear in \href{../formalization/ueot-core/UEOT/V3/Compression/TheoryCompletion/UnifiedClosure/RegisterJointDecisionOptimality.lean}{general-horizon optimality}, \href{../formalization/ueot-core/UEOT/V3/Compression/TheoryCompletion/UnifiedClosure/RegisterJointDecisionSource.lean}{the actual source}, and \href{../formalization/ueot-core/UEOT/V3/Compression/TheoryCompletion/UnifiedClosure/RegisterJointDecisionBoundary.lean}{posterior boundaries}.

\section{39.82 A decision representation derived from permitted histories}\label{a-decision-representation-derived-from-permitted-histories}

The source, observation, resource, and evaluation contract of Sections 39.75--39.81 remains unchanged. This section further answers the concrete implementation question of Section 29.4: which historical distinctions must be retained, and which may be discarded without changing future legal decisions? The original complete conditional mathematical contract and the preceding finite joint decision contract remain complete. The new contribution derives executable statistics and a shared policy artifact from the original Register transactions. Existing general predictive-quotient, C1, Belief, and Bellman theorems are reused and are not counted again as new sources.

The fixed public inputs are a finite prior \(\pi_0\), rates \(r_\theta\in[0,1]\cap\mathbb Q\), target bit \(g\), write fee \(c>0\), initial wallet \(B_0\ge0\), price \(\lambda\ge0\), discount \(0\le d\le1\), and total horizon \(H\). The initial state remains the original founder, known to be damaged, with its protected program aligned with \(g\). The scheduler's only new data are its own actions and the post-transaction receipts \((\mathrm{birth},\mathrm{paidWrite})\). It has no Store query, actual parent-identity, hidden-label, or true-\(\theta\) interface.

Write the executable local observer as \[q=(x_0,x_1,\mathrm{stopped},\mathrm{fresh}),\] where \(x_0\in\{0,1\}\) is the root's actual physical bit and \(x_1\in\{\mathrm{None},0,1\}\) records the existence and physical bit of child1. fresh holds only at the initial length of one. The actual append at the first tick determines whether child1 is a newborn object or None; later appends never fill address1. Under the fixed initialization, fresh can also be derived from whether elapsed time is zero. The implementation retains this redundant bit to make the local update definition and source proof clear.

For an existing selected slot, post-execution birth is true exactly when the new bit equals \(g\); if birth is false, the new bit is \(\neg g\). Unselected slots retain their bits, and native tail slots are appended normally. If the current request is stop, or a previous request has already stopped, the effective request is root idle. Environment, births, and clocks continue. The pre-execution request cap is still calculated from the original request; an already-stopped state cannot relax the original menu for free.

\texttt{update} and \texttt{updateStat} take only the public target, previous summary, the scheduler's own request, and the new receipt. Source-side \texttt{projectLocal} occurs only in the proof: projecting the native Store through the old PackedStore and then locally gives the same result as the observer computes from receipts. This commutation holds for each original transaction and each noise input, hence for every actual finite execution history. Complete identity, lineage, and None records remain in the native interpreter.

\newpage

\section{39.83 Informative zero receipts and the counting posterior}\label{informative-zero-receipts-and-the-counting-posterior}

Not every zero receipt is a fault sample. Let \(o_N(q,a)\) be the normal-mode receipt determined by the known local state and effective action; the fault-mode receipt is \((0,0)\). Summing over the original mode and fair hidden label gives \[L_\theta(o\mid q,a)=r_\theta\mathbf1_{o=o_N(q,a)}+(1-r_\theta)\mathbf1_{o=(0,0)}.\] If both modes produce the same output, a zero receipt has likelihood one and must not increase the fault count.

\begin{longtable}[]{@{}
  >{\raggedright\arraybackslash}p{(\columnwidth - 4\tabcolsep) * \real{0.3333}}
  >{\raggedright\arraybackslash}p{(\columnwidth - 4\tabcolsep) * \real{0.3333}}
  >{\raggedright\arraybackslash}p{(\columnwidth - 4\tabcolsep) * \real{0.3333}}@{}}
\toprule\noalign{}
\begin{minipage}[b]{\linewidth}\raggedright
Selected slot and effective action
\end{minipage} & \begin{minipage}[b]{\linewidth}\raggedright
Normal receipt
\end{minipage} & \begin{minipage}[b]{\linewidth}\raggedright
Does a zero receipt identify a fault?
\end{minipage} \\
\midrule\noalign{}
\endhead
\bottomrule\noalign{}
\endlastfoot
Existing slot, write target & \((1,1)\) & Yes \\
Existing slot, write nontarget & \((0,1)\) & Yes; a normal unsuccessful write is still charged \\
Healthy slot, idle & \((1,0)\) & Yes \\
Damaged slot, idle & \((0,0)\) & No \\
Absent slot, any effective request & \((0,0)\) & No \\
\end{longtable}

The observer counts a sample only when \(o_N\ne(0,0)\): an actual nonzero receipt increments \(s\), and an actual zero receipt increments \(f\). Define \[w_\theta(s,f)=\pi_0(\theta)r_\theta^s(1-r_\theta)^f,\qquad Z(s,f)=\sum_\theta w_\theta(s,f).\] For every permitted history \(h\) of positive mass, the actual conditional posterior is \[\boxed{\displaystyle\pi_h(\theta)=\frac{w_\theta(s(h),f(h))}{Z(s(h),f(h))}.}\]

This formula is not restricted to the new deterministic controller. For any real-randomized strategy using complete permitted histories, expansion of the actual primitive probabilities for a fixed receipt-compatible history \(h\) gives \[\mathbb P_\pi(\theta,h)=\alpha_\pi(h)\,\pi_0(\theta)r_\theta^{s(h)}(1-r_\theta)^{f(h)},\] whereas an incompatible history has actual mass zero. Here \(\alpha_\pi(h)\) is the product of the real action probabilities along that history. The strategy updates only from the same initially permitted memory, its own actions, and receipts, so this factor is the same for every \(\theta\). Positive actual history mass implies receipt compatibility, \(\alpha_\pi(h)>0\), and \(Z(s,f)>0\); only then may the action factor be cancelled in conditioning.

The proof first sums over the actual four-branch primitive. Hidden labels and uninformative modes are integrated out rather than counted as extra samples. Zero prior weights, \(r=0\), \(r=1\), and \(s=f=0\) all lie within the theorem's scope; \(0^0=1\) corresponds to the absence of samples of that type. A zero prior weight remains zero. At \(Z=0\), the computational function returns the original prior as an explicit default making the function total. Positive-support execution first proves that the denominator is strictly positive, so actual conditional probabilities are not interpreted through that default.

\section{39.84 Joint sufficiency for all actions}\label{joint-sufficiency-for-all-actions}

The actual statistic is \(\sigma=(t,q,s,f,k)\), where \(k\) counts actual paidWrite events, and the graph key is \((n,\sigma)\). Keys reachable from initialization satisfy \(n+t=H\), and the wallet is given by \(B_0-kc\). The total-function implementation uses \(\max(0,B_0-kc)\) on impossible branches. On every feasible source-supported prefix, the prepayment proof gives the untruncated literal balance.

\textbf{Source sufficiency theorem.} Two positive-mass permitted histories with the same public inputs, remaining horizon, and statistic have the same legal menu. For every action in that menu, their receipt, next statistic, actual charge, birth, and next balance have the same joint law. This conclusion is derived from the original primitive, local-observer commutation, and the counting Bayes recursion; sufficiency is not an externally supplied field.

The proof establishes a world-side relation between the original Belief and the statistic. Every nonzero belief atom retains root and projects to the same \(q\); its \(\theta\) marginal equals the counting posterior; wallets agree; and the counting denominator is positive. The actual damaged founder establishes the initial relation, and updates on receipts of positive conditional mass preserve it. The joint test-function identity for each action includes birth, fees, and balance; a product of three marginal distributions cannot replace it. Equality of native fees and receipts reuses the completed original transaction bridge.

Recursion along supported successors gives the same finite receipt-trace law and the same value for every common table policy. Comparing every legal action at each level then establishes equality of optimal values in the original complete Belief model and the new statistical model. The already-proved upper bound for the original model covers arbitrary legal real-randomized full-history strategies, so compression does not narrow the comparison class to strategies that read only the summary. The statistic preserves the distinctions required by the current menu and evaluation; it is not claimed to be the unique or globally smallest encoding.

\section{39.85 A shared graph, the actual returned artifact, and native execution}\label{a-shared-graph-the-actual-returned-artifact-and-native-execution}

On each call with public inputs, the candidate generator creates a new empty HashMap keyed by \((n,\sigma)\). It completes positive-probability successors for all legal actions in decreasing remaining time. Children are written to an Array first, and each parent key stores only a node index; histories reaching the same key reuse that index. The graph contains alternative-action nodes needed for the upper bound, so its size must be counted separately from the subgraph reachable when only optimal actions are executed. It does not call the old \texttt{solve}, \texttt{fastSolve}, or quaternary-tree expansion.

The acceptor checks uniqueness of every key in the actual Array, the root horizon and initial statistic, and, for every nonterminal node, all legal actions, every positive-mass receipt, the exact next statistic, and remaining time. It recomputes each action's rational Q value from child-node values. Selection must follow the original menu's first-occurring tie rule, and terminal value must be zero. Only successful checking returns that same graph with a proof. Omitting a positive successor of even an unselected action causes rejection; checking only the optimal path is insufficient for this certificate.

The return type is explicitly Option. The theorem proves correctness of every accepted artifact; it does not separately prove that the general candidate generator is accepted on every input. All five main examples and eight boundary examples in this work were actually accepted. Certification before export and direct execution consume the same graph in memory. This work adds no Lean theorem for JSON decoding or round-trip correctness; exported bytes are bound by actual output identity and independent item-by-item checking.

During execution, the controller reads only the current node index, statistic, and receipt. The native interpreter directly calls the original transaction and retains complete records; it does not first materialize a virtual full tree. \texttt{denote} occurs only in the proof connecting this direct interpreter to the old finite-table theorem. Consequently, the actual expected \(J\) of THIS returned graph in the native world equals the reported root-node value and dominates every legal real-randomized history strategy under the same information, wallet, and complete fees.

For any positive initial atom and positive primitive-noise tape, every request satisfies original cap prepayment before its outcome is observed. At each prefix \(m\), the balance updated from actual receipts and the fees in complete native records satisfy \[B_m=B_0-\sum_{j<m}C_j\ge0,\qquad 0\le\sum_{j<m}C_j\le B_0.\] The proof includes unsuccessful writes, faults, None, and stop; actual payment is not replaced by average expenditure. An unknown index or zero-probability receipt enters an explicit safe default action while the world continues to run. Complete execution preserves \(H\) native records, \(H\) action--receipt pairs, and \(H\) newly allocated slots, including None slots without a birth.

\section{39.86 A source-specific size bound on the actual artifact}\label{a-source-specific-size-bound-on-the-actual-artifact}

Each update increases the informative count by at most one and incurs at most one actual payment. Hence \[s+f\le t,\qquad k\le t,\qquad n+t=H.\] The local-state encoding conservatively has \(2\times3\times2\times2=24\) possibilities. For the actual Array checked by the acceptor, inject each unique key into \[\{0,\ldots,H\}^4\times Q,\qquad(n,\sigma)\longmapsto(t,s,f,k,q),\] and recover remaining time from \(H-t\). This yields the bound actually proved here, \[\boxed{N_H=\mathrm{Graph.size}\le24(H+1)^4.}\] \texttt{boundCheck} checks every actually stored row, rather than merely counting an ideal deduplicated set. Key uniqueness excludes storing duplicate graph rows while claiming a quotient-size bound. Charged prefixes also imply \(kc\le B_0\), but the present result uses the conservative bound above and does not claim formalization of a sharper budget-dependent summation constant.

This is only a node-count bound for the known-target family with fixed dispatch addresses 0 and 1. The size of finite \(\Theta\), rational bit lengths, arithmetic costs, the checker's traversal method, temporary arrays and hash tables, serialization, and complete native population records require separate accounting. It proves neither polynomial time for general POMDPs nor a total-memory bound, and it does not treat external computer expenditure as cognitive energy use already internalized within the object.

\section{39.87 Historical distinctions that must be retained and those that can be removed}\label{historical-distinctions-that-must-be-retained-and-those-that-can-be-removed}

Take equal prior weights on \(r\in\{1/4,3/4\}\), \(B_0=c\), and an initially bad root. Both histories request ``write target to root, idle root, idle root'':

\begin{longtable}[]{@{}
  >{\raggedright\arraybackslash}p{(\columnwidth - 6\tabcolsep) * \real{0.2500}}
  >{\raggedright\arraybackslash}p{(\columnwidth - 6\tabcolsep) * \real{0.2500}}
  >{\raggedright\arraybackslash}p{(\columnwidth - 6\tabcolsep) * \real{0.2500}}
  >{\raggedright\arraybackslash}p{(\columnwidth - 6\tabcolsep) * \real{0.2500}}@{}}
\toprule\noalign{}
\begin{minipage}[b]{\linewidth}\raggedright
History
\end{minipage} & \begin{minipage}[b]{\linewidth}\raggedright
Actual three-step receipts
\end{minipage} & \begin{minipage}[b]{\linewidth}\raggedright
Primitive likelihood
\end{minipage} & \begin{minipage}[b]{\linewidth}\raggedright
Posterior: low \(r\), high \(r\)
\end{minipage} \\
\midrule\noalign{}
\endhead
\bottomrule\noalign{}
\endlastfoot
A & \((1,1),(1,0),(0,0)\) & \(r^2(1-r)\) & \((1/4,3/4)\) \\
B & \((1,1),(0,0),(0,0)\) & \(r(1-r)\) & \((1/2,1/2)\) \\
\end{longtable}

Both reach the same time, zero wallet, non-stopped status, bad root, and good child1. Yet idling child1 next gives birth probabilities \(5/8\) and \(1/2\), respectively. Physical state, wallet, and time alone are therefore insufficient. The third step of B idles a bad root, so its zero receipt does not provide another fault sample. Positive mass, likelihood, and the next-birth joint probability for both histories are obtained by expanding the original primitive and connecting to the native Store; unreachable hand-specified worlds are not substituted.

Other histories can actually be merged. Both first perform a normal write to root. A then has a fault on root idle followed by normal child1 idle, whereas B first has normal child1 idle followed by a fault on root idle. At the same time they have the same complete statistic, with \((s,f)=(2,1)\), but their actual tail slots are, respectively, \[A=(\mathrm{bad},\mathrm{good},\mathrm{None},\mathrm{good}),\qquad
B=(\mathrm{bad},\mathrm{good},\mathrm{good},\mathrm{None}).\] Under the original address-0/1 menu, the future law relevant to decisions is the same for every allowed action. The shared artifact stores this key at the same actual node index. Native execution still retains the different histories, owners, and parent--offspring records. Past births occurred at different times, so realized discounted payoff may also differ. What is shared is the conditional future and continuation value; actual total payoff is still accumulated from each complete record.

If idle at address2 is separately allowed, that slot is empty in A and cannot give birth, but exists and is healthy in B and gives birth with probability \(r_\theta\). This boundary is witnessed by the two histories above, both reachable under the old menu. Enlarging future action permissions requires refinement of the representation. Decision equivalence neither erases historical identity in the world nor supplies a new dispatch capability for free from the old protocol.

\section{39.88 Finite consumers, verification scope, and the stopping boundary}\label{finite-consumers-verification-scope-and-the-stopping-boundary}

The actual run reuses four original successful outputs for exact comparison and preregisters, before obtaining results, a longer consumer with \(H=8\), \(B_0=2\), \(c=1\), equal weights on \(r\in\{1/4,3/4\}\), \(d=3/4\), \(\lambda=13/10\), and \(g=\mathrm{true}\). All outputs are generated through the new Lean shared entry point, checked for all actions, and exported. An independent checker recomputes every alternative branch of the original complete menu and executes the artifact index by index using the native Store.

\begin{longtable}[]{@{}
  >{\raggedright\arraybackslash}p{(\columnwidth - 6\tabcolsep) * \real{0.2000}}
  >{\raggedleft\arraybackslash}p{(\columnwidth - 6\tabcolsep) * \real{0.2667}}
  >{\raggedleft\arraybackslash}p{(\columnwidth - 6\tabcolsep) * \real{0.2667}}
  >{\raggedleft\arraybackslash}p{(\columnwidth - 6\tabcolsep) * \real{0.2667}}@{}}
\toprule\noalign{}
\begin{minipage}[b]{\linewidth}\raggedright
Consumer
\end{minipage} & \begin{minipage}[b]{\linewidth}\raggedleft
Exact optimal net value
\end{minipage} & \begin{minipage}[b]{\linewidth}\raggedleft
Certificate nodes
\end{minipage} & \begin{minipage}[b]{\linewidth}\raggedleft
Execution nodes
\end{minipage} \\
\midrule\noalign{}
\endhead
\bottomrule\noalign{}
\endlastfoot
Two steps, B=1 & \(\frac{1}{160}\) & 21 & 6 \\
Three-step backup, B=1 & \(\frac{9}{256}\) & 52 & 11 \\
Three-step posterior, B=1 & \(\frac{27}{320}\) & 52 & 11 \\
Four-step feedback, B=1 & \(\frac{39525}{524288}\) & 104 & 16 \\
Eight steps, B=2 & \(\frac{884450457}{1342177280}\) & 1661 & 151 \\
\end{longtable}

Node counts include terminal rows. The eight-step certificate has 8,221 positive-probability edges over all actions, and the actual optimal-execution subgraph has 182 edges; its JSON output is 960,155 bytes. The policy has probability \(215/256\) of at least one birth, probability \(1233/4096\) of actually paying twice, expected total births \(57313/16384\), and expected actual total write fees \(4673/4096\). These totals are not conflated with discounted net value.

For the original Mac entry point in that run, eight-step generation took 852 milliseconds, Lean certificate checking 1,740 milliseconds, node preparation for JSON 931 milliseconds, and encoding 6 milliseconds. Complete wall time, including loading and entry-point processing, was 7.250 seconds. Peak RSS for the entire Lean process was 5,812,731,904 bytes, including imports and runtime; the node-count bound is not a guarantee of low total memory use. The independent Linux checker took approximately 0.286 seconds for all-action verification and 0.304 seconds for the complete check, with peak RSS 27,398,144 bytes. Environments and workloads differ, so no cross-environment speedup ratio is calculated.

The four earlier examples reuse previously successful full trees, with action-by-action agreement at 26 positive-support decision histories. Complete native replays cover 32, 128, 128, and 512 runs, respectively, totaling 800. The original 28-minute entry point was not rerun. The longer example checks 32 fixed distinguishing native paths and two further Lean direct-graph execution outputs. Exact optimality is supplied by the all-action certificate and general native-execution theorem; finite path instances only check identity, fees, and controller wiring. A further positive-mass legal native prefix contains two consecutive normal writes, incurs two actual charges, and takes the balance from 2 to 0. Whether the optimal policy uses the second unit of budget is reported separately from its output; rewards are not changed to manufacture repeated payment.

The run record retains compilation results and entry-point wall time including loading, and separately reports generation, serialization, Lean checking, independent checking, actual graph nodes, edges, bytes, and memory metrics. Pure loading time was not isolated. The existing Python reference already used memoization and remains a stronger exact baseline. Its reduction to representative equivalent actions after stop is explicitly recorded, while the new Lean certificate still checks all legal requests of the original menu. The original four-step time of 1713.890130 seconds was unsegmented wall time for the earlier entry point; it cannot justify a fixed speedup ratio for a particular pure solving phase.

Required boundary cases include \(H=0\), a zero wallet, insufficient prepayment, zero prior weights, endpoint rates, a charged wrong write, absence of a child, first allocation, stop, and zero-probability default branches. Initial compilation failures and correction logs are retained; earlier failures, negative examples, and original samples were not redrawn. This work changed no QM, physics, or GI implementation and did not disable CI status.

The new connection runs from permitted histories to concrete statistics and a shared policy, then back to the same actual world. Its inputs remain a conditional model and a preselected evaluation; external empirical support remains UNVERIFIED. Unknown targets, arbitrary addresses, common aging, infinite horizons, approximate regret, and internal hardware belong to different contracts. This contract ends on completing its six exit conditions and does not automatically expand into those investigations. Source code, actual inputs and outputs, checks, limitations, and failures are recorded in the \href{../formalization/ueot-core/docs/compression/theory_completion/unified_closure/decision_representation_closure_v1/README.md}{decision representation delivery package}.

\section{39.89 Fault structure is itself a mechanism condition to be tested}\label{fault-structure-is-itself-a-mechanism-condition-to-be-tested}

This section separately registers two fixed mechanisms while preserving the earlier local-fault source and completion conclusions of Sections 39.75--39.88. Both sources initialize a Store containing only one bad root; the short-horizon and budget bounds below use this common initial state. S is the original serial transaction: a fault damages only the existing individual actually dispatched in that transaction. G first executes the same transaction and, if that tick faults, sets the relevant physical coordinate of every living individual in the resulting complete Store to \(!g\). This includes requests to absent addresses and already-stopped runs. G changes no other physical coordinate, protected program, identity, lineage, slot order, None record, or individual clock determined by the original transaction. Undispatched individuals do not acquire an extra clock increment. A fault produces no child; successful children still come from the original transaction and allocator.

The mechanism \(m\in\{S,G\}\) is fixed before the run and remains unchanged throughout. Both mechanisms use the same public fixed normal rate \(r\), with conditionally independent tick modes. Native birth labels are independent and fair and remain in world records. The controller reads only its own actions, \nolinkurl{birth/payment} receipts, and actual wallet, and may use private randomness independent of the mechanism and all world noise. Mechanism, mode, labels, complete Store, and actual parent identity are not controller inputs.

The seven actions remain stop, idle at root0 or child1, and either write value at either address. Every write request is checked against the tariff-\(c\) cap before its outcome is observed; only then is the actual transaction fee deducted. A zero fee caused by failure or an absent slot does not legalize a request that could not originally be prepaid. stop permanently executes actual root idle from the current tick onward, while the world advances tick by tick. All probing, repair, birth, and stopping use the same initial wallet \(B\), horizon \(H\), and evaluation fixed in advance, \[J_m(\pi)=\mathbb E_m^\pi\!\left[\sum_{t=0}^{H-1}d^t\bigl(\mathrm{birth}_t-\lambda\,\mathrm{actualFee}_t\bigr)\right],\] with terminal value zero. G is a new environmental mechanism, not a free damage action granted to the controller, and it does not change the meaning of the original S.

\section{39.90 Identical receipts no longer determine the same body}\label{identical-receipts-no-longer-determine-the-same-body}

Start from a bad root. First write the correct bit to root and receive \((1,1)\), then idle root and receive \((0,0)\). For \(0<r<1\), this prefix has probability \(r(1-r)>0\) under both mechanisms and spends one write fee. Under S, root is now bad but child1 remains good; under G, both are bad. The old S observer retains a good child1 after that receipt and therefore cannot be transferred to G without proof. This is an explicit counterexample at the boundary of the new mechanism's scope, not an error in the old source proof.

The new information state retains the \textbf{joint} conditional law of mechanism and physical state determined by permitted history. If the initial mechanism weights are \(w_m\), their marginal remains \(w_m\) after this prefix, but physical state is correlated with mechanism. Multiplying the mechanism marginal by the physical-state marginal would create impossible worlds. Complete-history filtering is derived from the same primitive source. Action factors for real-randomized controllers depend only on permitted history and may be cancelled on conditioning at positive mass. Zero-weight mechanisms are not silently added, and zero-mass off-path branches follow the total-function convention.

This finite joint representation directly enters the existing complete-history Bayes recursion; it does not require a new counting compression or shared-graph platform. The native Store retains actual identities and None slots. World projections used in the proof help check semantics but do not become additional sensors.

\section{39.91 Exact distinguishing power of permitted experiments}\label{exact-distinguishing-power-of-permitted-experiments}

Let \(P_m^\pi\) be the law of complete \nolinkurl{action/receipt} histories under mechanism \(m\) for the same legal experimental strategy. Private randomness is independent of mechanism and world, and the classifier may retain its own private random information. All comparisons use the same strategy; separate mechanism-dependent actions for S and G must not leak which mechanism is present. Write \(\Delta_{H,B}\) for the supremum of attainable total variation.

The formal entry point permits history-dependent action kernels with arbitrary real probabilities. If the strategy and terminal discriminator also share an arbitrary private random source, fixing one realization of that source determines a complete finite action tree and a binary decision at each terminal permitted history. Pushing the private random law forward to this pair of finite tables preserves decisions and error probabilities separately under each mechanism. Keeping only the action tree could lose other private randomness used by the discriminator; the decision table is therefore retained as well.

The following zero bounds come from path couplings of the actual source: experiments dispatching only root are indistinguishable for every finite \(H\); experiments with all seven actions are indistinguishable when \(H\le2\); they are indistinguishable from the specified bad founder when \(B<c\); and every finite experiment is indistinguishable when \(r=0\) or \(r=1\). These results equate only the laws of permitted records, not the hidden complete Stores. The low-budget conclusion depends on this initialization and pre-request prepayment; it does not apply to a world supplied with an additional free healthy backup.

Within three steps, a coupling using the same private randomness, modes, and labels can produce its first receipt difference only on the mode word normal, fault, normal. Its mass is \(q=r^2(1-r)\), so every adaptive real-randomized experiment has TV at most \(q\). If \(B\ge c\), the word \[\text{root writes }g\ ;\quad\text{root idle}\ ;\quad\text{child1 idle}\] within the same process attains this bound: receipt sequence \((1,1)/(0,0)/(1,0)\) has mass exactly \(q\) under S and zero under G. Therefore \[\Delta_{3,B}=r^2(1-r)\quad(B\ge c),\qquad\Delta_{3,B}=0\quad(B<c).\] The middle idle cannot be replaced by committed stop. The upper bound over all strategies uses the same-source coupling and the implementation theorem for finite complete response tables; it does not restrict competitors to this fixed action word.

At \(r=1/2\), maximum TV is \(1/8\), and the minimum binary classification error under equal priors is \(7/16\). Positive finite separation can therefore coexist with a high error rate. More generally, for every finite \(H\) and \(0<r<1\), the two mechanisms share all-fault mass \((1-r)^H\). Every total binary classifier has equal-prior error at least \((1-r)^H/2>0\). This does not rule out one-sided certificates allowing UNKNOWN, nor does it make the classification-error lower bound a decision-regret bound of the same numerical size.

Indistinguishability need not cause action loss. For example, retaining \(r=1/2,B=2,c=1,\lambda=1/4,d=1\) and taking \(H=2\), one legal strategy attains the nonzero optimal value \(13/16\) under both mechanisms. Complete mechanism identification is not necessary for correct action in this short task.

\section{39.92 A worst-case regret certificate for a fixed mechanism family}\label{a-worst-case-regret-certificate-for-a-fixed-mechanism-family}

For a preregistered finite mechanism family, \(V_m^*\) is the optimal value when mechanism \(m\) is known but the same receipt permissions, wallet, and horizon must still be respected. It gives no foreknowledge of future modes or the complete hidden world. For one shared legal strategy, define \[R_m(\pi)=V_m^*-J_m(\pi),\quad R(\pi)=\max_mR_m(\pi),\quad R^*=\inf_\pi R(\pi).\] The mechanism must be fixed before the private seed is realized; it cannot be selected according to that seed or replaced by the worst mechanism at each tick.

If \(w_m\ge0,\sum_mw_m=1\) and the existing Bayes upper bound for real-probability strategies on all permitted histories gives \(\sum_mw_mJ_m(\pi)\le V_w^*\), then \[L(w)=\sum_m w_m V_m^*-V_w^*\le R^*\le U(\rho)=\max_m\{V_m^*-J_m(\rho)\}.\] The left inequality follows because maximum regret is at least weighted regret; the right is supplied by an actual legal attaining strategy \(\rho\). The weights \(w\) are only a lower-bound witness, not a claim to have identified a natural prior. Equality of the endpoints certifies both the concrete returned strategy and the minimum. The actual S/G consumer below uses this general certificate; oracle values and Bayes numbers are not unproved assumptions.

For a specified strategy, zero regret is equivalent to optimality under every mechanism. Without a separate attainment proof, a zero infimum over an arbitrary strategy class cannot be rewritten as existence of a common optimal strategy. The exact nonzero result here has an explicit attaining strategy.

\section{39.93 Exact loss with an unknown mechanism over three steps}\label{exact-loss-with-an-unknown-mechanism-over-three-steps}

Fix \(r=1/2,H=3,B=2,c=1,\lambda=1/4,d=1,g=\mathrm{true}\). The original complete-action recursion on the new joint source gives \[V_S^*=\frac{41}{32},\qquad V_G^*=\frac{40}{32}.\] Write the two actual returned tables as \(\pi_S,\pi_G\). Among visited decision prefixes, they differ only at the ambiguous prefix discussed above: \(\pi_S\) finally idles child1, whereas \(\pi_G\) finally writes the correct bit to root.

\begin{longtable}[]{@{}
  >{\raggedright\arraybackslash}p{(\columnwidth - 4\tabcolsep) * \real{0.2727}}
  >{\raggedleft\arraybackslash}p{(\columnwidth - 4\tabcolsep) * \real{0.3636}}
  >{\raggedleft\arraybackslash}p{(\columnwidth - 4\tabcolsep) * \real{0.3636}}@{}}
\toprule\noalign{}
\begin{minipage}[b]{\linewidth}\raggedright
Returned strategy
\end{minipage} & \begin{minipage}[b]{\linewidth}\raggedleft
Full-horizon \(J\) under true S
\end{minipage} & \begin{minipage}[b]{\linewidth}\raggedleft
Full-horizon \(J\) under true G
\end{minipage} \\
\midrule\noalign{}
\endhead
\bottomrule\noalign{}
\endlastfoot
\(\pi_S\) & \(41/32\) & \(37/32\) \\
\(\pi_G\) & \(40/32\) & \(40/32\) \\
\(\rho=\tfrac14\pi_S+\tfrac34\pi_G\) & \(161/128\) & \(157/128\) \\
\end{longtable}

The ambiguous prefix has probability \(1/4\) under both mechanisms, with one tick and wallet one remaining. The conditional continuation value of child idle is \(1/2\) under S and 0 under G; the conditional net value of a correct root write is \(3/8\) under either. This shows that the role of the backup depends on fault structure and explains the payoff conflict between the two tables.

The lower bound that actually excludes all other legal strategies uses \(w=(3/4,1/4)\). The Bayes value with the complete seven-action menu is \(5/4\), so \[L(w)=\frac34\frac{41}{32}+\frac14\frac{40}{32}-\frac54=\frac3{128}.\] The actual \(\rho\) has regret \(3/128\) under each mechanism, hence \[R^*=R(\rho)=\frac3{128}.\] This covers the original full-history real-randomized competitors; merely comparing the two strategies in the table is not a lower-bound proof. Dual weights and deployment-mixture weights are also different and must not be interchanged.

\section{39.94 Private randomness, the deterministic boundary, and actual payment}\label{private-randomness-the-deterministic-boundary-and-actual-payment}

Before execution, the actual controller obtains an independent uniform four-point private seed: one outcome selects \(\pi_S\) and three select \(\pi_G\). It then interprets that same complete returned table using only received receipts. The world mechanism neither reads the seed nor is reselected according to it; native labels do not provide the controller's random source. Original randomized history kernels can be implemented by distributions on finite complete response tables, without assuming rational action probabilities.

At these fixed parameters, every deterministic permitted-history strategy has eight equiprobable mode paths. Labels affect neither permitted control nor payoff, and each path's net payoff lies in \(\tfrac14\mathbb Z\). Thus \(J_m\) and regret lie in \(\tfrac1{32}\mathbb Z\). The strictly positive lower bound \(3/128\) excludes zero, so deterministic worst-case regret is at least \(1/32\), attained by \(\pi_G\). The deterministic optimum is therefore exactly \(1/32\), while independent randomization strictly reduces it to \(3/128\). This lattice argument applies only to the fixed rate, horizon, discount, and price; it does not directly extend to arbitrary parameters.

The same returned object has a complete Store interpretation. On every positive-mass primitive path, for every private seed and every prefix, the request cap is checked first, the actual charge comes from the original transaction, and the balance equals \(B\) minus actual cumulative fees and is nonnegative. Newborn source programs, actual parent identities, lineage extensions, None ticks, and clocks are retained. Only the registered two-bit receipt is supplied to the controller. The source bridge preserves full-horizon native \(J\) for each deterministic component and for their private-seed mixture.

\section{39.95 Verification levels and the stopping boundary of this work}\label{verification-levels-and-the-stopping-boundary-of-this-work}

Verification in this work binds the same input, source code, Lean returned tables, and four-point private seed. Two mechanisms, four private seeds, and all three-step Bool×Bool primitive tapes give \(2\times4\times4^3=512\) short native runs. Complete enumeration covers birth labels and all 1536 payment prefixes. Actual Lean packed execution is compared path by path with an independent complete Store replay. A separate Fraction checker that does not import the solver recursion being checked verifies all-action oracle values, the dual value, the attaining strategy, and fees. Execution enumeration checks wiring and identity; universal optimality is supplied by the formal proof.

The three-step and two-step closed-form numerical equalities use Lean \nolinkurl{native_decide} finite computational certificates. They respectively depend on the two computer-generated native-decision certificates listed in the axiom audit and on correctness of Lean's compiled execution. General source bridges, upper bounds over all strategies, and payment derivations are separately established by logical proofs. The original attempt at direct kernel reduction, stopped because of computational cost, is retained; this engineering interruption is not interpreted as a mathematical counterexample. Independent exact checking reduces implementation-wiring risk and is not presented as sampling evidence replacing formal proof.

Source code and principal endpoints, axiom dependencies, actual artifacts, compilation failures and corrections, independent checks, and release status are recorded in the \href{../formalization/ueot-core/docs/compression/theory_completion/unified_closure/mechanism_decision_closure_v1/README.md}{mechanism decision delivery package}. Completion judgments for the original eight conditional mathematical contracts and PR441 and PR442 are retained. The new conclusions connect organizational fault structure, permitted evidence, and achievable action loss; they do not claim to discover real mechanisms or derive normative objectives.

S/G are synthetic mathematical mechanisms, and empirical support remains \textbf{UNVERIFIED}. The distinguishing H3 receipt occurs at the terminal time; no result proves or requires this information to improve later decisions in the same H3 task. The outer randomized controller's reliable random source and computing resources are registered conditions, not internalized object hardware. Work stops once the six contract conditions are complete; it does not automatically increase H, add a third mechanism or permissions, or introduce a compression platform or performance gate.

\chapter{Appendix A. Index of Symbols and Types}\label{appendix-a}

\textbf{How to use this appendix.} This appendix identifies the spaces, units, and roles attached to symbols. Where a letter is reused in a local model, its explicit local definition applies. Shared notation is not evidence that two sources, policies, or metrics are identical.

The same character may denote a local variable in different chapters. The table gives its principal type. Every use must be interpreted according to the definition at that location.

\begin{longtable}[]{@{}
  >{\raggedright\arraybackslash}p{(\columnwidth - 4\tabcolsep) * \real{0.3333}}
  >{\raggedright\arraybackslash}p{(\columnwidth - 4\tabcolsep) * \real{0.3333}}
  >{\raggedright\arraybackslash}p{(\columnwidth - 4\tabcolsep) * \real{0.3333}}@{}}
\toprule\noalign{}
\begin{minipage}[b]{\linewidth}\raggedright
Symbol/name
\end{minipage} & \begin{minipage}[b]{\linewidth}\raggedright
Mathematical type and meaning
\end{minipage} & \begin{minipage}[b]{\linewidth}\raggedright
Entry point
\end{minipage} \\
\midrule\noalign{}
\endhead
\bottomrule\noalign{}
\endlastfoot
\(X_t,H_t\) & State and history up to time \(t\) & Chapter 2 \\
\(P^u,K_i\) & Protocol path law and conditional response kernel & Chapters 2, 4 \\
\(f,\sigma(f)\) & Measurable readout and its embedded factor & Chapter 3 \\
\(Y_f^+\) & Future relative to a candidate factor & Chapter 4 \\
\(C^f,C\) & Canonical predictive state taking values in a family of probability kernels & Chapter 4 \\
FFIPS & Full predictive sufficient state for a fixed future and protocol family; the state domain must be specified in each application & Chapters 4, 16 \\
\(M,U\) & Internal representation and environmental interface; neither means matter or energy by default & Chapter 6 \\
\(I(\cdot;\cdot\mid\cdot)\) & Conditional mutual information, measured in nats & Chapters 1, 5 \\
\(\epsilon,\eta,\alpha\) & Model defect, estimation radius, and confidence failure probability, recorded separately & Chapter 16 \\
\(\mathcal V,\tau_{\mathcal V}\) & Persistence domain and exit time & Chapter 8 \\
\(V_A,\mathcal L\) & Mean return-time potential and generator & Chapter 9 \\
\(q,h,\varpi\) & QSD, positive right eigenfunction, and invariant law of the Doob transform & Chapter 10 \\
\(\lambda,\gamma\) & Exponential absorption rate and conditional mixing rate & Chapter 10 \\
\(\Gamma,\Theta,\Psi\) & History transport, structural descriptor, and developmental mechanism & Chapter 11 \\
TCIC / DTRC & Frozen-kernel commutation / developmental transport consistency & Chapter 11 \\
\(\Omega\) & Causal integrity structure with a specified certificate; no geometric ring is presupposed & Chapter 12 \\
\(V,\mathcal M_0,b_V(\mathcal M_0)\) & Carrier universe, family of minimal sufficient sets, and family of minimal hitting sets & Chapter 13 \\
\(e(S),r(S)\) & Response-diameter defect and optimal decoding radius & Chapter 13 \\
\(\mathcal L_s,\operatorname{cl}_s\) & Closure system at a resolution and its closure operator & Chapter 14 \\
\(\pi_*,J^-,J^+\) & Coarse image of a minimal family and lower/upper endpoints of the realization interval & Chapter 14 \\
\(\Delta_P\) & Conditional predictive entropy contributed by the parent beyond the child cores & Chapter 15 \\
\(G_N,\mathcal I\) & Design Gram matrix and Fisher information matrix & Chapter 17 \\
\(J,\Pi,\beta\) & Task evaluation, policy class, and discount factor & Chapter 18 \\
\(V^*,Q^*,\mathcal A^*\) & Optimal value, action value, and optimal-action correspondence & Chapter 19 \\
GOD / GOA & \nolinkurl{Optimal-action/generative} mechanism and its family of long-term structures & Chapters 19, 21 \\
TVRC & Teleological realization certificate relative to the family of optimal path laws for a specified task & Chapter 18 \\
\(\operatorname{sp}(v)\) & \(\sup v-\inf v\), the span of a function & Chapters 1, 20 \\
\(d_\mu^\pi\) & Normalized discounted state--action occupancy distribution & Chapter 20 \\
Dual drives \(\Pi,\Phi\) & Independently anchored \nolinkurl{resource/information} coordinates; do not confuse this notation with the policy class & Chapter 23 \\
\(B,z(B),S\) & Budget, common low-dimensional response variable, and stacked Jacobian & Chapter 23 \\
\(M=D_bK\) & Mean reproduction operator, reproduction intensity, and transmission kernel & Chapter 25 \\
\(R,l,r\) & Principal growth factor, long-term composition, and reproductive value & Chapter 25 \\
ETI & The four-component certificate of evolutionary transition in individuality used in this book & Chapter 26 \\
\(b_t(\theta,x)\) & Posterior belief over parameter and state & Chapter 27 \\
\end{longtable}

\chapter{Appendix B. Register of Proved Results}\label{appendix-b}

\textbf{How to use this appendix.} This register connects stable proposition identifiers to their locations in the text. It is a retrieval aid: hypotheses, standard dependencies, and scope must still be checked in the corresponding statements. Later chapter-local propositions are not counted among the original 106 P-IDs.

Chapters 0--31 retain the 106 P-numbered results of the original v3.0, each with its original proof or an explicit assembly of results. An identifier implies neither originality nor unconditional validity. The present Core has a separate formalization trace for these frozen source endpoints. The precise Lean hypotheses and the standard mathematical assumptions in this text must still be read through the chapter-by-chapter correspondence. The added M-numbered results appear separately in Appendix F and do not alter the following stable index of 106 items. Each identifier links to its statement.

\begin{longtable}[]{@{}
  >{\raggedright\arraybackslash}p{(\columnwidth - 4\tabcolsep) * \real{0.3000}}
  >{\raggedleft\arraybackslash}p{(\columnwidth - 4\tabcolsep) * \real{0.4000}}
  >{\raggedright\arraybackslash}p{(\columnwidth - 4\tabcolsep) * \real{0.3000}}@{}}
\toprule\noalign{}
\begin{minipage}[b]{\linewidth}\raggedright
Identifier
\end{minipage} & \begin{minipage}[b]{\linewidth}\raggedleft
Chapter
\end{minipage} & \begin{minipage}[b]{\linewidth}\raggedright
\nolinkurl{Result/section}
\end{minipage} \\
\midrule\noalign{}
\endhead
\bottomrule\noalign{}
\endlastfoot
\hyperref[p-met-01]{P-MET-01} & 1 & Data processing for kernels and readouts \\
\hyperref[p-met-02]{P-MET-02} & 1 & Span estimate \\
\hyperref[p-proc-01]{P-PROC-01} & 2 & Making history into a state \\
\hyperref[p-fac-01]{P-FAC-01} & 3 & 3.2 Covariance under representation changes \\
\hyperref[p-pred-01]{P-PRED-01} & 4 & Canonical minimality for a common protocol family \\
\hyperref[p-pred-02]{P-PRED-02} & 4 & Target pushforward and protocol refinement \\
\hyperref[p-pred-03]{P-PRED-03} & 4 & Recursive update for finite-alphabet processes \\
\hyperref[p-info-01]{P-INFO-01} & 5 & 5.1 Information retention identity \\
\hyperref[p-info-02]{P-INFO-02} & 5 & 5.2 From information residual to prediction error \\
\hyperref[p-info-03]{P-INFO-03} & 5 & Zero-distortion complexity \\
\hyperref[p-info-04]{P-INFO-04} & 5 & Fano-type lower bound for identity \\
\hyperref[p-info-05]{P-INFO-05} & 5 & 5.5 Omission and packing lower bounds \\
\hyperref[p-int-01]{P-INT-01} & 6 & Necessary and sufficient conditions for structured sufficiency \\
\hyperref[p-int-02]{P-INT-02} & 6 & 6.2 Least internal--environmental factorization of a finite structural kernel \\
\hyperref[p-int-03]{P-INT-03} & 6 & 6.5 Conditions for uniqueness of a minimal Markov boundary \\
\hyperref[p-dyn-01]{P-DYN-01} & 7 & Strong lumpability \\
\hyperref[p-dyn-02]{P-DYN-02} & 7 & 7.2 Continuous-time generator criterion \\
\hyperref[p-dyn-03]{P-DYN-03} & 7 & Path error under the same policy \\
\hyperref[p-dyn-04]{P-DYN-04} & 7 & 7.4 Quotients across scales of the same underlying process \\
\hyperref[p-per-01]{P-PER-01} & 8 & 8.2 A necessary invariant core for deterministic persistence \\
\hyperref[p-per-02]{P-PER-02} & 8 & 8.4 Stochastic persistence and an occupation-based invariant measure \\
\hyperref[p-per-03]{P-PER-03} & 8 & 8.5 Finite-state viability kernel \\
\hyperref[p-per-04]{P-PER-04} & 8 & Tangency as a necessary condition \\
\hyperref[p-rec-01]{P-REC-01} & 9 & 9.1 Discrete recovery tube \\
\hyperref[p-rec-02]{P-REC-02} & 9 & 9.2 Continuous stochastic recovery \\
\hyperref[p-rec-03]{P-REC-03} & 9 & 9.3 Recovery time generates a canonical potential \\
\hyperref[p-rec-04]{P-REC-04} & 9 & Drift implies a recovery-time bound \\
\hyperref[p-qsd-01]{P-QSD-01} & 10 & Converse result for conditional stabilization \\
\hyperref[p-qsd-02]{P-QSD-02} & 10 & 10.2 Finite Perron representation \\
\hyperref[p-qsd-03]{P-QSD-03} & 10 & 10.3 Persistence window from the same initial state \\
\hyperref[p-qsd-04]{P-QSD-04} & 10 & 10.4 Spectral version for reversible diffusions \\
\hyperref[p-id-01]{P-ID-01} & 11 & Accumulation of transport defects \\
\hyperref[p-id-02]{P-ID-02} & 11 & Developmental tube \\
\hyperref[p-omg-01]{P-OMG-01} & 12 & Reconstruction from minimal destructive sets \\
\hyperref[p-omg-02]{P-OMG-02} & 12 & 12.3 Integrity margin \\
\hyperref[p-car-01]{P-CAR-01} & 13 & 13.1 Monotone access system \\
\hyperref[p-car-02]{P-CAR-02} & 13 & 13.2 Access erasure and blockers \\
\hyperref[p-car-03]{P-CAR-03} & 13 & Double blocker \\
\hyperref[p-car-04]{P-CAR-04} & 13 & 13.4 Defect diameter and decoding radius \\
\hyperref[p-res-01]{P-RES-01} & 14 & Coarse-graining of minimal properties \\
\hyperref[p-res-02]{P-RES-02} & 14 & Composition law \\
\hyperref[p-res-03]{P-RES-03} & 14 & Naturality of blockers across scales \\
\hyperref[p-res-04]{P-RES-04} & 14 & Adjoint triple \\
\hyperref[p-res-05]{P-RES-05} & 14 & 14.4 Exact microscopic uncertainty interval \\
\hyperref[p-res-06]{P-RES-06} & 14 & Object-dependent uniqueness \\
\hyperref[p-comp-01]{P-COMP-01} & 15 & 15.1 Three forms of commonality \\
\hyperref[p-comp-02]{P-COMP-02} & 15 & 15.2 Interventional irreducibility \\
\hyperref[p-comp-03]{P-COMP-03} & 15 & 15.3 Composition margin \\
\hyperref[p-comp-04]{P-COMP-04} & 15 & 15.4 Additional predictive information at the parent level \\
\hyperref[p-comp-05]{P-COMP-05} & 15 & 15.6 Canonical Booleanization of overlap geometry \\
\hyperref[p-comp-06]{P-COMP-06} & 15 & Lifting physical carriers to coalitions \\
\hyperref[p-comp-07]{P-COMP-07} & 15 & 15.8 Composition window preserving individuality \\
\hyperref[p-stat-01]{P-STAT-01} & 16 & Simultaneous finite-alphabet TV bound \\
\hyperref[p-stat-02]{P-STAT-02} & 16 & 16.2 Defect error has no additional dependence on the number of candidate carriers \\
\hyperref[p-stat-03]{P-STAT-03} & 16 & Recovery of a fixed-tolerance family \\
\hyperref[p-stat-04]{P-STAT-04} & 16 & Exact family recovery \\
\hyperref[p-stat-05]{P-STAT-05} & 16 & 16.5 Recovery of canonical predictive classes \\
\hyperref[p-stat-06]{P-STAT-06} & 16 & Simultaneous embedding error \\
\hyperref[p-stat-07]{P-STAT-07} & 16 & Covariance of kernel transport \\
\hyperref[p-stat-08]{P-STAT-08} & 16 & Finite-candidate discovery \\
\hyperref[p-stat-09]{P-STAT-09} & 16 & Average-error transport \\
\hyperref[p-inv-01]{P-INV-01} & 17 & Binary testing lower bound \\
\hyperref[p-inv-02]{P-INV-02} & 17 & 17.2 Gauge null directions of Fisher information \\
\hyperref[p-inv-03]{P-INV-03} & 17 & 17.2 Gauge null directions of Fisher information \\
\hyperref[p-inv-04]{P-INV-04} & 17 & 17.3 Design identification of developmental laws \\
\hyperref[p-inv-05]{P-INV-05} & 17 & 17.4 OLS bound for predictable designs \\
\hyperref[p-tel-01]{P-TEL-01} & 18 & Discounted potential-function equivalence \\
\hyperref[p-ctl-01]{P-CTL-01} & 19 & Bellman existence and optimality \\
\hyperref[p-ctl-02]{P-CTL-02} & 19 & 19.3 Compact Feller control \\
\hyperref[p-ctl-03]{P-CTL-03} & 19 & 19.4 Continuous-time HJB verification \\
\hyperref[p-quo-01]{P-QUO-01} & 20 & 20.1 Exact control quotient \\
\hyperref[p-quo-02]{P-QUO-02} & 20 & 20.2 Span-sensitive approximate quotient bound \\
\hyperref[p-quo-03]{P-QUO-03} & 20 & 20.3 Action gap and near-optimality error \\
\hyperref[p-quo-04]{P-QUO-04} & 20 & 20.4 Exact occupancy-weighted identity \\
\hyperref[p-quo-05]{P-QUO-05} & 20 & Occupancy-weighted regret of a quotient policy \\
\hyperref[p-goa-01]{P-GOA-01} & 21 & 21.2 Occupancy limit of a finite chain \\
\hyperref[p-goa-02]{P-GOA-02} & 21 & 21.3 Dobrushin mixing margin \\
\hyperref[p-goa-03]{P-GOA-03} & 21 & 21.4 Stability of multiple structures preserving the recurrent decomposition \\
\hyperref[p-goa-04]{P-GOA-04} & 21 & 21.5 Spectral stability of a symmetric killed kernel \\
\hyperref[p-kl-01]{P-KL-01} & 22 & 22.1 Information lower bound for a survival event \\
\hyperref[p-kl-02]{P-KL-02} & 22 & 22.2 Exact event I-projection \\
\hyperref[p-kl-03]{P-KL-03} & 22 & 22.3 KL divergence of discrete Markov paths \\
\hyperref[p-kl-04]{P-KL-04} & 22 & 22.4 Explicit CTMC formula \\
\hyperref[p-kl-05]{P-KL-05} & 22 & Observability-corrected control-energy relation \\
\hyperref[p-ddh-01]{P-DDH-01} & 23 & 23.2 Gauge freedom \\
\hyperref[p-ddh-02]{P-DDH-02} & 23 & 23.3 Finite exponential families and constraint duality \\
\hyperref[p-ddh-03]{P-DDH-03} & 23 & Variational minimality under moment constraints \\
\hyperref[p-ddh-04]{P-DDH-04} & 23 & 23.4 Common two-dimensional bottleneck \\
\hyperref[p-ddh-05]{P-DDH-05} & 23 & Effective-dimension threshold \\
\hyperref[p-ali-01]{P-ALI-01} & 24 & Global integrability criterion \\
\hyperref[p-ali-02]{P-ALI-02} & 24 & 24.3 Child-level update and parent-level value \\
\hyperref[p-ali-03]{P-ALI-03} & 24 & 24.4 Coordination threshold \\
\hyperref[p-evo-01]{P-EVO-01} & 25 & Exact decomposition of selection and transmission \\
\hyperref[p-evo-02]{P-EVO-02} & 25 & Irrelevant labels can inflate hereditary information \\
\hyperref[p-evo-03]{P-EVO-03} & 25 & Mean growth and reproductive value \\
\hyperref[p-evo-04]{P-EVO-04} & 25 & Reproductive-value martingale \\
\hyperref[p-brg-01]{P-BRG-01} & 26 & Concentration under fixed fitness without mutation \\
\hyperref[p-brg-02]{P-BRG-02} & 26 & Selection cannot distinguish behaviorally equivalent types \\
\hyperref[p-ref-01]{P-REF-01} & 27 & Reflexive state augmentation \\
\hyperref[p-ref-02]{P-REF-02} & 27 & Belief sufficiency \\
\hyperref[p-ref-03]{P-REF-03} & 27 & Nonnegative decision value of free information \\
\hyperref[p-ref-04]{P-REF-04} & 27 & Regret bound from uniform objective error \\
\hyperref[p-ref-05]{P-REF-05} & 27 & Monotonicity under expansion of feasible choices \\
\hyperref[p-api-01]{P-API-01} & 28 & Composition of exact and approximate process interfaces \\
\hyperref[p-alg-01]{P-ALG-01} & 28 & Correctness and coarseness of the finite stable partition \\
\hyperref[p-core-01]{P-CORE-01} & 31 & Finite object--realization--control certification assembly \\
\end{longtable}

\chapter{Appendix C. Sources, Improvements Returned to the Core, and Standard Dependencies}\label{appendix-c}

\textbf{How to use this appendix.} This appendix separates the original conceptual sources, standard mathematical dependencies, and the specific connections made in this book. Citations identify what is used and prevent existing theorems from being reassigned to UEOT as original discoveries.

\section{C.1 Roles of the Sources}\label{c.1-roles-of-the-sources}

The theoretical starting point of this edition is the original paper and early mathematical files supplied by the author. Later specialized branches do not replace that philosophical starting point. Later material is used to recover proofs, counterexamples, and interface corrections that can return to the general Core. Research intentions expressed in historical conversations explain the scope of the task; they do not replace mathematical evidence. Historical content that is inaccessible or has not been checked is not used as a proved premise.

The following genealogy lists sources that directly shaped this edition. ``Incorporated'' means that the relevant mathematics has been extracted and restated in the text. It does not mean that every claim in an earlier version is accepted, or that every historical file and conversation has received a word-for-word independent audit.

\begin{longtable}[]{@{}
  >{\raggedright\arraybackslash}p{(\columnwidth - 4\tabcolsep) * \real{0.3333}}
  >{\raggedright\arraybackslash}p{(\columnwidth - 4\tabcolsep) * \real{0.3333}}
  >{\raggedright\arraybackslash}p{(\columnwidth - 4\tabcolsep) * \real{0.3333}}@{}}
\toprule\noalign{}
\begin{minipage}[b]{\linewidth}\raggedright
Source
\end{minipage} & \begin{minipage}[b]{\linewidth}\raggedright
Role in this edition
\end{minipage} & \begin{minipage}[b]{\linewidth}\raggedright
Chapters incorporating it
\end{minipage} \\
\midrule\noalign{}
\endhead
\bottomrule\noalign{}
\endlastfoot
\nolinkurl{UEOT_v3_top.pdf} & Original problem architecture of \nolinkurl{Being/Becoming}, persistent identity, hierarchy, Π/Φ, and GOD/GOA & 0, 8, 12, 18--27 \\
UEOT\_v1.0\_核心理论正式版.md & Development of Ω, structured openness, object hierarchies, adaptation, and meaning & 0, 6, 11--12, 24--27 \\
UEOT\_数学基础\_统一正式版\_v1.1(1).md & Prediction, information, persistence, composition, control, and reflexive structure in the early unified mathematics & 1--12, 15, 18--27 \\
UEOT\_数学基础\_完整系统正式版\_v2.1.md & Convergence of the I--IV main lines, object-relative futures, rigorous reformulation, and span-based control bounds & 3--27 \\
UEOT\_数学扩展\_v2.1\_五大开放主线系统推进.md & Developmental identification, continuous responses, GOA, and multilevel objectives & 10, 16--17, 21, 24 \\
UEOT\_数学基础\nolinkurl{IV_Dynamic-Composition}与\nolinkurl{Multi-GOA_v0.4.md} & Dynamic composition and the distinction between multiple long-term structures & 11, 15, 21 \\
UEOT\_数学基础\nolinkurl{IV_Metastable-Object-GOA}与\nolinkurl{Teleology-Bridge_v0.5.md} & Separation of conditional stability, QSD, and the teleological bridge & 10, 18, 21 \\
UEOT\_数学基础\nolinkurl{IV_Persistence-IProjection}与\nolinkurl{Anchored-Dual-Drives_v0.6.md} & Survival information projection and independently anchored dual drives & 22--23 \\
UEOT\_数学基础\nolinkurl{IV_Lineage-IProjection}与\nolinkurl{Positive-Operator-Evolution_v0.7.md} & Lineage evaluation, positive operators, and the growth bridge & 25--26 \\
UEOT\_数学基础\nolinkurl{IV_Realized-Lineage_Heredity_Object-Selection_v0.8.md} & Realized lineages, hereditary sampling, and levels of object selection & 25--26 \\
UEOT\_数学基础\nolinkurl{IV_Endogenous-Teleological-Bridge_v0.9.md} & Conditions under which an explicit selection mechanism induces purpose & 26 \\
UEOT\_数学基础\nolinkurl{IV_Causal-Omega_Teleonomy_Variational-Realization_v0.10.md} & Causal Ω and variational realization & 12, 22 \\
UEOT\_数学基础\nolinkurl{IV_Protocol-Saturated-Objecthood_v0.11.md} & Protocol-relative completeness and the boundary of saturation & 4, 7, 16 \\
UEOT\_数学基础\nolinkurl{IV_Factor-Ontology}与\nolinkurl{Directed-Identity_v0.12.md} & Embedded factors and directed historical identity & 3, 11 \\
\nolinkurl{UEOT-QOD_Interface_v1.0_}通用对象实现-分辨率-组合接口\_严格版.tex & General carriers, blockers, resolution adjunctions, and proofs for overlapping composition & 13--15, 28 \\
\nolinkurl{UEOT_Core_v2.2_QOD-G}融合\_双主线收敛核心\_严格版.tex & Interface between the object-realization and control lines & 13--21, 28, 31 \\
Earlier v2.3 Core and independent checking work & References for corrected conditions, finite examples, and rigorous statements; its shorter exposition does not replace this complete text & Overall structure and verification \\
\end{longtable}

The sources are ordered by conceptual dependency. Absolute dates of composition are not inferred from version numbers in filenames. Version v3.0 denotes a complete architectural reorganization and integration of general proofs; it does not declare that all earlier branches had converged.

\section{C.2 Connections Explicitly Incorporated and Corrected}\label{c.2-connections-explicitly-incorporated-and-corrected}

\begin{longtable}[]{@{}
  >{\raggedright\arraybackslash}p{(\columnwidth - 4\tabcolsep) * \real{0.3333}}
  >{\raggedright\arraybackslash}p{(\columnwidth - 4\tabcolsep) * \real{0.3333}}
  >{\raggedright\arraybackslash}p{(\columnwidth - 4\tabcolsep) * \real{0.3333}}@{}}
\toprule\noalign{}
\begin{minipage}[b]{\linewidth}\raggedright
Item
\end{minipage} & \begin{minipage}[b]{\linewidth}\raggedright
Treatment in this edition
\end{minipage} & \begin{minipage}[b]{\linewidth}\raggedright
Nature of the change
\end{minipage} \\
\midrule\noalign{}
\endhead
\bottomrule\noalign{}
\endlastfoot
Predictive identity & Fix the object-relative future, protocols, and common conditional versions & Retention and unification of assumptions \\
Internal--external closure & Replace two separate residuals by the joint residual; provide an XOR counterexample & Logical correction \\
Historical identity & Separate directed instance continuation, undirected genealogy, and similarity of type & Return of later proofs to the Core \\
Recovery and persistence & Distinguish deterministic invariant cores, stochastic occupancy, recovery potentials, and QSD & Restoration of mathematical breadth \\
Carrier defect & Distinguish response diameter from optimal decoding radius and prove \(e/2\le r\le e\) & Interface refinement \\
Realization and destruction & Separate access blockers from failure of a physical mechanism & Correction of operational semantics \\
Resolution & Closure, Boolean naturality, adjoint triples, and exact microscopic intervals & Return of general proofs to the Core \\
Statistical recovery & A simultaneous response event covers all carriers; exact structure requires a gap & Unification of the error chain \\
Kernel embeddings & A characteristic kernel does not imply control of TV error & Metric correction \\
Quotient control & Retain the existing span bound and add an occupancy-weighted regret interface & Retention and strengthening \\
Teleological realization & Measure distance to the full family of optimal path laws & Correction for nonunique optima \\
Girsanov & Separate full-noise KL from observed-state path KL & Correction of the assumptions for equality \\
Dual drives & Independent anchors, a common budget, stacked rank, and effective dimension & Making claims testable \\
Evolution & Separate reproduction intensity, transmission kernels, stochastic populations, and mean models & Type correction \\
Finite algorithms & Prove termination, correctness, and coarseness of the coarsest controlled stable partition & Completion of the executable interface \\
\end{longtable}

These revisions improve the validity of general connections. They do not reclassify standard mathematics as original. A P identifier means only that a proof is given in this text.

\section{C.3 Foundational Mathematics Invoked Directly}\label{c.3-foundational-mathematics-invoked-directly}

The book assumes the usual foundations of measure-theoretic probability, linear algebra, calculus, and conditional expectation. The following substantial results are not reproved here. Their roles and applicability conditions must be retained.

\begin{longtable}[]{@{}
  >{\raggedright\arraybackslash}p{(\columnwidth - 4\tabcolsep) * \real{0.3333}}
  >{\raggedright\arraybackslash}p{(\columnwidth - 4\tabcolsep) * \real{0.3333}}
  >{\raggedright\arraybackslash}p{(\columnwidth - 4\tabcolsep) * \real{0.3333}}@{}}
\toprule\noalign{}
\begin{minipage}[b]{\linewidth}\raggedright
Standard result
\end{minipage} & \begin{minipage}[b]{\linewidth}\raggedright
\nolinkurl{Conditions/content} required here
\end{minipage} & \begin{minipage}[b]{\linewidth}\raggedright
Use
\end{minipage} \\
\midrule\noalign{}
\endhead
\bottomrule\noalign{}
\endlastfoot
Regular conditional probabilities and Doob--Dynkin & Standard Borel targets and the relevant measurability; statements modulo null sets use a common completed space & Conditional kernels and minimal sufficient representations \\
Ionescu--Tulcea & Measurable probability kernels along a countable time sequence and an initial law & Consistent existence of a path law \\
KL chain rule and Pinsker & Extended nonnegative KL; expressions involving subtraction require additional finiteness & From information to prediction and path costs \\
Prokhorov and Portmanteau & Tight families of probabilities on Polish spaces and weak convergence & Invariant probabilities from persistence occupancy \\
Marchaud viability / Nagumo & Closed domains and the specified differential inclusion, or a locally Lipschitz nonexplosive ODE & Tangency conditions and viability \\
Spectral decomposition for a self-adjoint operator with compact resolvent & The conditions in §10.4, a positive principal eigenfunction, and a strict principal spectral gap & Spectral proof of diffusion QSD \\
Perron--Frobenius & Finite primitive nonnegative matrices; the corresponding simple principal-root version in the symmetric irreducible case & QSD, Doob transforms, and reproductive growth \\
Measurable \nolinkurl{maximum/selection} theorems & Compact metric action spaces, continuous rewards, weakly continuous kernels, and the other conditions in §19.3 & Existence of optimal policies in continuous state spaces \\
Itô/Dynkin and Girsanov & Domains, integrability, nonexplosion, a true exponential martingale, and observability are specified separately & Continuous recovery, HJB, and energy bounds \\
Hoeffding and McDiarmid & Independent samples and the corresponding bounded-difference conditions & Certification from finite data \\
Ergodic theorems and absolute-regularity coupling & Stationarity and ergodicity for consistency; additional mixing rates for rates of convergence & Dependent data \\
Hodge decomposition & A compact oriented manifold without boundary, a Riemannian metric, and the required smoothness & Nongradient components of local objective fields \\
Nonnegative martingale convergence & Adaptedness, integrability, nonnegativity, and the martingale property & Limits of stochastic reproductive value \\
\end{longtable}

Using foundational mathematics does not turn it into a new physical law of UEOT. A complete application must additionally verify that the particular model satisfies its conditions.

\section{C.4 Original Research Entry Points Retained from v3.0}\label{c.4-original-research-entry-points-retained-from-v3.0}

The following entries retain and correct the standard-module sources of v3.0. For v4.1, accessible author, institutional, or publisher pages were checked for titles, authors, and the stated role. The reference audit records the level of verification and access limitations for each entry. These works do not supply every part of this edition and do not replace the assumptions in the text.

\begin{itemize}
\tightlist
\item
  Predictive equivalence and computational mechanics: the original paper by Shalizi and Crutchfield studies optimal predictive states and their structure, providing a comparison for this edition's protocol-relative extension. \href{https://arxiv.org/abs/cond-mat/9907176}{Original paper}
\item
  Viability: Jean-Pierre Aubin and Arrigo Cellina, \emph{Differential Inclusions: Set-Valued Maps and Viability Theory}, Springer, 1984, particularly its chapters on viability and control systems, provide standard background for the closed-domain and differential-inclusion modules. \href{https://doi.org/10.1007/978-3-642-69512-4}{Publisher and DOI}. The previous WP-81-076 link actually identifies a paper on Lipschitz stability of solutions to convex optimization problems; it is not a direct source for the viability results described in this entry. This edition corrects the citation without changing the frozen mathematical source or proofs.
\item
  QSD: Champagnat and Villemonais study uniform exponential convergence to a quasi-stationary distribution and the Q-process. This edition proves the results it needs only under explicit spectral or finite-state conditions. \href{https://arxiv.org/abs/1404.1349}{Original paper}
\item
  Kernel two-sample testing and mean embeddings: the original work by Gretton and colleagues supplies statistical background for MMD. This book does not convert MMD into TV without additional assumptions. \href{https://www.jmlr.org/papers/volume13/gretton12a/gretton12a.pdf}{Original paper}
\item
  Adaptive linear estimation: Abbasi-Yadkori and colleagues provide tools for self-normalized processes. The text uses a separately proved conservative fixed-time bound. \href{https://arxiv.org/abs/1102.2670}{Original paper}
\item
  Approximate MDP homomorphisms: the work of Ravindran and Barto is established theoretical background for action and state abstraction. \href{https://all.cs.umass.edu/pubs/2004/ravindran_b_KBCS04.pdf}{Original paper}
\item
  Finite-entropy Girsanov theory: Léonard's work supplies standard background for changes of measure under relative-entropy conditions. Observation maps remain subject to data processing. \href{https://leonard.perso.math.cnrs.fr/papers/2012-Leonard-Girsanov\%20theory\%20under\%20a\%20finite\%20entropy\%20condition.pdf}{Author's text}
\end{itemize}

\section{C.5 Bibliographic Records and Entry Points for Standard Results}\label{c.5-bibliographic-records-and-entry-points-for-standard-results}

The following records supply authors, titles, editions or years for the standard mathematics in C.3. The role column identifies dependencies in this book; it does not claim that a cited work proves the entire UEOT assembly. First publication, copyright year, reprint date, and electronic publication date may differ. The table uses the bibliographic year of the stated edition. This revision checked bibliographic records and subject correspondence on accessible primary pages; access limited to an abstract or table of contents is not represented as a page-by-page review of every theorem in a book.

\begin{longtable}[]{@{}
  >{\raggedright\arraybackslash}p{(\columnwidth - 2\tabcolsep) * \real{0.5000}}
  >{\raggedright\arraybackslash}p{(\columnwidth - 2\tabcolsep) * \real{0.5000}}@{}}
\toprule\noalign{}
\begin{minipage}[b]{\linewidth}\raggedright
Bibliographic record
\end{minipage} & \begin{minipage}[b]{\linewidth}\raggedright
Role in this book
\end{minipage} \\
\midrule\noalign{}
\endhead
\bottomrule\noalign{}
\endlastfoot
Cosma Rohilla Shalizi and James P. Crutchfield. \emph{Computational Mechanics: Pattern and Prediction, Structure and Simplicity}. Journal of Statistical Physics 104, 2001; preprint 1999. \href{https://arxiv.org/abs/cond-mat/9907176}{Original preprint} & Chapters 4--5: established background on predictive states; the protocol extension is subject to this book's own conditions. \\
Jean-Pierre Aubin and Arrigo Cellina. \emph{Differential Inclusions: Set-Valued Maps and Viability Theory}. Springer, 1984. \href{https://doi.org/10.1007/978-3-642-69512-4}{DOI} & Chapter 8: differential inclusions, tangency, and viability. \\
Nicolas Champagnat and Denis Villemonais. \emph{Exponential convergence to quasi-stationary distribution and Q-process}. arXiv:1404.1349, 2014. \href{https://arxiv.org/abs/1404.1349}{Author preprint} & Chapters 10, 21: standard background on QSD and conditional convergence. \\
Arthur Gretton, Karsten M. Borgwardt, Malte J. Rasch, Bernhard Schölkopf, and Alexander Smola. \emph{A Kernel Two-Sample Test}. Journal of Machine Learning Research 13(25):723--773, 2012. \href{https://www.jmlr.org/papers/v13/gretton12a.html}{Journal} & Chapter 16: MMD and kernel statistics; no unconditional equivalence to TV. \\
Yasin Abbasi-Yadkori, Dávid Pál, and Csaba Szepesvári. \emph{Online Least Squares Estimation with Self-Normalized Processes: An Application to Bandit Problems}. arXiv:1102.2670, 2011. \href{https://arxiv.org/abs/1102.2670}{Preprint} & Chapter 17: background for estimation under adaptive designs. \\
Balaraman Ravindran and Andrew G. Barto. \emph{Approximate Homomorphisms: A framework for non-exact minimization in Markov Decision Processes}. 2004. \href{https://all.cs.umass.edu/pubs/2004/ravindran_b_KBCS04.pdf}{Author's institutional text} & Chapters 20, 28: prior work on state and action abstraction. \\
Christian Léonard. \emph{Girsanov theory under a finite entropy condition}. Author file dated 2013-12-06; its URL directory contains 2012. \href{https://leonard.perso.math.cnrs.fr/papers/2012-Leonard-Girsanov\%20theory\%20under\%20a\%20finite\%20entropy\%20condition.pdf}{Author text} & Chapter 22: finite-entropy change of measure; it does not remove information loss due to observation. \\
Olav Kallenberg. \emph{Foundations of Modern Probability}. 3rd ed., Springer, 2021. \href{https://doi.org/10.1007/978-3-030-61871-1}{DOI} & Chapters 1--4, 8, 27: conditional probabilities, kernels, convergence of probabilities, and process foundations. \\
Thomas M. Cover and Joy A. Thomas. \emph{Elements of Information Theory}. 2nd ed., Wiley, 2006. \href{https://doi.org/10.1002/047174882X}{DOI} & Chapters 5, 17, 22: entropy, data processing, and information bounds. \\
Martin L. Puterman. \emph{Markov Decision Processes: Discrete Stochastic Dynamic Programming}. Wiley, 1994. \href{https://doi.org/10.1002/9780470316887}{DOI} & Chapters 19--21: discrete control, Bellman equations, and policy comparison. \\
Onésimo Hernández-Lerma and Jean Bernard Lasserre. \emph{Discrete-Time Markov Control Processes: Basic Optimality Criteria}. Springer, 1996. \href{https://doi.org/10.1007/978-1-4612-0729-0}{DOI} & Sections 19.3 and Chapter 27: background on Borel states and actions, discounted control, and belief states. \\
E. Seneta. \emph{Non-negative Matrices and Markov Chains}. Springer, photographic reprint of the 1981 edition. \href{https://doi.org/10.1007/0-387-32792-4}{DOI} & Chapters 10, 21, 25: nonnegative matrices, Perron structure, and Markov chains. \\
Daniel Revuz and Marc Yor. \emph{Continuous Martingales and Brownian Motion}. 3rd ed., Springer, 1999. \href{https://doi.org/10.1007/978-3-662-06400-9}{DOI} & Chapters 9, 19, 22, 39: martingales, stochastic integration, Itô calculus, and change of measure. \\
E. B. Davies. \emph{Heat Kernels and Spectral Theory}. Cambridge University Press, 1989. \href{https://www.cambridge.org/core/books/heat-kernels-and-spectral-theory/39DC7EE2A5DCCE9DEF581D2F52A73E1F}{Publisher} & Sections 10.4 and 39.74: background on self-adjoint elliptic operators and spectral methods; the specific semigroup conditions remain in force. \\
Frank W. Warner. \emph{Foundations of Differentiable Manifolds and Lie Groups}. Springer, 1983; original edition 1971. \href{https://doi.org/10.1007/978-1-4757-1799-0}{DOI} & Chapter 24: differential forms and the Hodge theorem, especially Chapter 6 of the cited book. \\
\end{longtable}

\chapter{Appendix D. Complete Executable Verification from v3.0 and Historical Output}\label{appendix-d}

\textbf{How to use this appendix.} This appendix retains executable code and its original run record for reproducibility of example calculations. Historical output establishes only the finite checks associated with its identified code and environment. The present publication build and any new checks have separate receipts.

\section{D.1 Execution and Code Identity}\label{d.1-execution-and-code-identity}

Save the Python code below as \nolinkurl{verify_ueot_core_v3.py} and run it in Python 3 with NumPy installed:

\begin{verbatim}
python3 verify_ueot_core_v3.py
\end{verbatim}

No access to original data or other documents is required. The synthetic experiment uses the fixed seed 30905. Final floating-point digits may differ across platforms; the assertions include appropriate numerical tolerances. Exact equalities in the algorithm apply only to the explicitly identical finite parameters of this example and are not a proof rule for arbitrary floating-point models.

Code SHA-256: \nolinkurl{4d4072377cbae485f595eddc239ea6962a7bdc5cf7ad400e8d75086ce485c5f7}

\section{D.2 Complete Code}\label{d.2-complete-code}

\begin{verbatim}
"""UEOT Core v3.0: finite-model checks, not a proof assistant.
Run: python3 verify_ueot_core_v3.py
Dependency: NumPy. No network, external datasets, or quantum modules.
"""
import itertools
import json
import math
import numpy as np

results = {}
checks = 0

def check(condition, message):
    global checks
    checks += 1
    if not bool(condition):
        raise AssertionError(message)

def tv(p, q):
    return float(np.abs(np.asarray(p) - np.asarray(q)).sum() / 2)

def minimal(family):
    return frozenset(s for s in family if not any(t < s for t in family))

def powerset(n):
    return [frozenset(i for i in range(n) if mask & (1 << i))
            for mask in range(1 << n)]

def blocker(family, n):
    return minimal([s for s in powerset(n)
                    if all(bool(s & edge) for edge in family)])

# All hypergraphs on three vertices, reduced to antichains.
universe = powerset(3)
antichains = set()
for mask in range(1 << len(universe)):
    antichains.add(minimal([universe[i] for i in range(len(universe))
                           if mask & (1 << i)]))
for family in antichains:
    check(blocker(blocker(family, 3), 3) == family, "blocker involution")
results["blocker_antichains_checked"] = len(antichains)

# Exhaustive resolution tests: all 3-to-2 surjections and both antichain spaces.
coarse_universe = powerset(2)
coarse_antichains = set()
for mask in range(1 << len(coarse_universe)):
    coarse_antichains.add(minimal([coarse_universe[i]
        for i in range(len(coarse_universe)) if mask & (1 << i)]))
surjections = [p for p in itertools.product(range(2), repeat=3) if set(p)=={0,1}]
resolution_checks_before = checks
for mapping in surjections:
    def image(s):
        return frozenset(mapping[i] for i in s)
    def preimage(s):
        return frozenset(i for i in range(3) if mapping[i] in s)
    def push(family):
        return minimal([image(s) for s in family])
    for family in antichains:
        check(blocker(push(family), 2) == push(blocker(family, 3)),
              "blocker quotient naturality")
        up_fine = frozenset(s for s in universe if any(e <= s for e in family))
        restriction = frozenset(s for s in coarse_universe
                                if preimage(s) in up_fine)
        for coarse in coarse_antichains:
            up_coarse = frozenset(s for s in coarse_universe
                                 if any(e <= s for e in coarse))
            ext_minus = frozenset(s for s in universe
                if any(preimage(b) <= s for b in up_coarse))
            ext_plus = frozenset(s for s in universe if image(s) in up_coarse)
            check((ext_minus <= up_fine) == (up_coarse <= restriction),
                  "left resolution adjunction")
            check((restriction <= up_coarse) == (up_fine <= ext_plus),
                  "right resolution adjunction")
            check((push(family) == coarse) ==
                  (ext_minus <= up_fine and up_fine <= ext_plus),
                  "exact realization interval")
    for coarse in coarse_antichains:
        up_coarse = frozenset(s for s in coarse_universe
                             if any(e <= s for e in coarse))
        ext_minus = frozenset(s for s in universe
                             if any(preimage(b) <= s for b in up_coarse))
        ext_plus = frozenset(s for s in universe if image(s) in up_coarse)
        active = frozenset().union(*coarse)
        singleton_active = all(len(preimage({w})) == 1 for w in active)
        check((ext_minus == ext_plus) == singleton_active,
              "active-fiber uniqueness")
results["resolution"] = dict(surjections=len(surjections),
    fine_antichains=len(antichains),coarse_antichains=len(coarse_antichains),
    assertions=checks-resolution_checks_before)

# Redundant representation and its response deficiency.
states = list(itertools.product(range(2), repeat=2))
p_one = np.array([[0.15, 0.85], [0.90, 0.85]])
readout = [(s, s, z) for s, z in states]
response = np.array([p_one[s] for s, z in states])
def deficiency(support, table):
    pairs = [(i, j) for i in range(4) for j in range(4)
             if all(readout[i][k] == readout[j][k] for k in support)]
    return max(float(np.max(np.abs(table[i] - table[j]))) for i, j in pairs)
exact_family = minimal([s for s in universe if deficiency(s, response) == 0])
check(exact_family == frozenset([frozenset([0]), frozenset([1])]),
      "minimal redundant carriers")
check(blocker(exact_family, 3) == frozenset([frozenset([0, 1])]),
      "redundancy blocker")
# Two Bernoulli rows: minimax decoding radius is half the diameter.
a, b = 0.2, 0.8
check(abs(max(abs(a - 0.5), abs(b - 0.5)) - abs(a-b)/2) < 1e-14,
      "radius versus diameter")

# Fresh independent repeated-sampling experiment.
rng = np.random.default_rng(30905)
N, L, alpha = 20000, 8, 0.05
emp = rng.binomial(N, response) / N
eta = math.sqrt(math.log(2 * L / alpha) / (2 * N))
observed = float(np.max(np.abs(emp-response)))
check(observed <= eta, "this run's simultaneous event")
for s in universe:
    check(abs(deficiency(s, emp)-deficiency(s, response)) <= 2*eta + 1e-14,
          "all carrier deficiencies")
tau = 0.2
emp_family = minimal([s for s in universe if deficiency(s, emp) <= tau])
check(emp_family == exact_family, "gap-certified carrier recovery")
results["sampling"] = dict(seed=30905, samples=N*L, eta=eta,
                           observed_max_error=observed,
                           empirical_response=emp.tolist(),
                           threshold=tau, nonzero_gap=0.75)

# Bellman solvers and exact lifting.
beta = 0.9
macro_P = np.stack([1-p_one, p_one], axis=-1)  # state, action, next
reward = np.array([[s-0.12*a for a in range(2)] for s in range(2)])
def solve(P, r):
    v = np.zeros(len(r))
    for _ in range(10000):
        q = r + beta*np.einsum("san,n->sa", P, v)
        new = q.max(axis=1)
        if np.max(np.abs(new-v)) < 1e-13:
            v = new
            break
        v = new
    else:
        raise AssertionError("value iteration did not converge")
    q = r + beta*np.einsum("san,n->sa", P, v)
    policy = q.argmax(axis=1)
    value = np.linalg.solve(np.eye(len(r))-beta*P[np.arange(len(r)), policy],
                            r[np.arange(len(r)), policy])
    return value, policy

true_v, true_policy = solve(macro_P, reward)
micro_P = np.array([[[((p_one[s,a] if sn else 1-p_one[s,a])/2)
                      for sn, zn in states] for a in range(2)]
                    for s, z in states])
micro_r = np.array([reward[s] for s, z in states])
micro_v, micro_policy = solve(micro_P, micro_r)
check(np.max(np.abs(micro_v-np.repeat(true_v, 2))) < 1e-11,
      "exact value lifting")
check(true_policy.tolist() == [1, 0], "true optimal policy")
# Learn macro rows at representative z=0, retaining uniform micro comparison.
hat_one = emp[[0, 2]]
hat_P = np.stack([1-hat_one, hat_one], axis=-1)
hat_v, hat_policy = solve(hat_P, reward)
epsilon_p = float(np.max(np.abs(response-hat_one[[0,0,1,1]])))
bound_D = beta*epsilon_p*float(np.ptp(hat_v))/(1-beta)
confidence_D = beta*eta*float(np.ptp(hat_v))/(1-beta)
value_error = float(np.max(np.abs(true_v-hat_v)))
check(value_error <= bound_D+1e-12, "span value bound")
closed_true = macro_P[np.arange(2), hat_policy]
lift_value = np.linalg.solve(np.eye(2)-beta*closed_true,
                             reward[np.arange(2), hat_policy])
regret = float(np.max(true_v-lift_value))
check(regret <= 2*bound_D+1e-12, "lifted policy regret bound")
def stationary(P):
    A = P.T-np.eye(len(P))
    A[-1] = 1
    bvec = np.zeros(len(P)); bvec[-1] = 1
    return np.linalg.solve(A, bvec)
closed_hat = hat_P[np.arange(2), hat_policy]
stationary_true, stationary_hat = stationary(closed_true), stationary(closed_hat)
mix = tv(closed_true[0], closed_true[1])
stat_error = tv(stationary_true, stationary_hat)
check(stat_error <= epsilon_p/(1-mix)+1e-12, "stationary perturbation bound")
hat_q = reward+beta*np.einsum("san,n->sa",hat_P,hat_v)
action_gaps = np.abs(hat_q[:,0]-hat_q[:,1])
results["control"] = dict(true_value=true_v.tolist(), estimated_value=hat_v.tolist(),
    true_policy=true_policy.tolist(), estimated_policy=hat_policy.tolist(),
    epsilon_p_actual=epsilon_p, value_error=value_error, span_bound_actual=bound_D,
    span_bound_confidence=confidence_D, policy_regret=regret,
    estimated_action_gaps=action_gaps.tolist(),
    certified_actions_at_confidence_radius=(action_gaps>2*confidence_D).tolist(),
    stationary_true=stationary_true.tolist(),stationary_estimated=stationary_hat.tolist(),
    stationary_error=stat_error,mixing_coefficient=mix,
    stationary_confidence_bound=eta/(1-mix))

# Coarsest finite controlled partition.
partition = [[0,1,2,3]]
def refine(blocks):
    refined = []
    for block in blocks:
        groups = {}
        for x in block:
            sig = tuple(micro_r[x]) + tuple(
                float(micro_P[x,a,B].sum()) for a in range(2) for B in blocks)
            groups.setdefault(sig, []).append(x)
        refined.extend(groups.values())
    return refined
for _ in range(5):
    refined = refine(partition)
    if refined == partition:
        break
    partition = refined
check(partition == [[0,1],[2,3]], "coarsest stable partition")
results["stable_partition"] = partition

# Path TV coupling bound, horizon eight, one fixed initial macro state.
T = 8
path_true, path_hat = [], []
for future in itertools.product(range(2), repeat=T):
    x = 0
    pt = ph = 1.0
    for xn in future:
        pt *= closed_true[x,xn]
        ph *= closed_hat[x,xn]
        x = xn
    path_true.append(pt); path_hat.append(ph)
path_error = tv(path_true, path_hat)
check(path_error <= 1-(1-epsilon_p)**T+1e-12, "path coupling bound")
results["path"] = dict(horizon=T, observed_tv=path_error,
                       actual_defect_bound=1-(1-epsilon_p)**T)

# Joint residual XOR: individual mutual information zero, joint ln 2.
def entropy(probabilities):
    return -sum(p*math.log(p) for p in probabilities if p>0)
hy = entropy([.5,.5])
conditional_single = .5*entropy([.5,.5])+.5*entropy([.5,.5])
check(abs(hy-conditional_single) < 1e-14, "XOR marginal information")
check(abs(hy-math.log(2)) < 1e-14, "XOR joint information")
results["xor"] = dict(I_Y_X=0.0,I_Y_E=0.0,I_Y_XE=hy)

# Continuous-time QSD, h-transform, and q versus q*h.
K = np.array([[-2.,1.],[2.,-4.]])
eigs, right = np.linalg.eig(K)
idx = int(np.argmax(eigs.real))
lam = float(-eigs[idx].real)
h = right[:,idx].real
if h.sum()<0: h=-h
leigs,left = np.linalg.eig(K.T)
q = left[:,int(np.argmax(leigs.real))].real
q = q/q.sum()
h = h/float(q@h)
varpi = q*h
K_up = np.diag(1/h)@K@np.diag(h)+lam*np.eye(2)
check(np.max(np.abs(q@K+lam*q)) < 1e-12, "QSD eigenlaw")
check(np.max(np.abs(K_up.sum(axis=1))) < 1e-12, "Doob generator")
check(np.max(np.abs(varpi@K_up)) < 1e-12, "Doob stationary law")
check(tv(q,varpi)>0.01, "QSD is not Doob equilibrium")
results["qsd"] = dict(exit_rate=lam,q=q.tolist(),doob_stationary=varpi.tolist())

# Price decomposition and Perron values.
M = np.array([[1.8,.2],[.4,.8]])
p = np.array([.4,.6]); z = np.array([0.,1.])
b = M.sum(axis=1); transmission = M/b[:,None]; mean_b=float(p@b)
pnew = p@M/mean_b
selection = float(p@(b*z)-mean_b*(p@z))/mean_b
change = float((p*b)@(transmission@z-z))/mean_b
check(abs(float(pnew@z-p@z)-selection-change)<1e-14, "Price identity")
ev,rv=np.linalg.eig(M); k=int(np.argmax(ev.real)); growth=float(ev[k].real)
results["evolution"] = dict(next_frequency=pnew.tolist(),selection=selection,
                            transmission=change,growth_rate=growth)

# Potential shaping, alignment, and a strict recovery counterexample.
psi=np.array([.3,-.2])
shaped_r=reward+beta*np.einsum("san,n->sa",macro_P,psi)-psi[:,None]
shaped_v,shaped_policy=solve(macro_P,shaped_r)
check(np.max(np.abs(shaped_v-(true_v-psi)))<1e-11, "potential shaping")
check(np.array_equal(shaped_policy,true_policy), "shaping policy equivalence")
g=np.array([1.,0.]); G=np.array([-1.,1.])
threshold=-float(g@G)/(float(g@g)-float(g@G))
check(abs(threshold-.5)<1e-14, "alignment threshold")
check(float(g@((1-.6)*G+.6*g))>0, "alignment improves")
x=0.0
for _ in range(8): x=(x+1)/2
check(0<x<1, "contractive recursion need not converge to zero")
# Degenerate observation loses the orthogonal drift exactly.
sigma=np.array([[1.,0.]])
u=np.array([0.,1.])
check(float(np.linalg.norm(sigma@u))==0, "unobserved Girsanov drift")
results["girsanov"] = dict(time=1,full_noise_relative_entropy=.5,
                           observed_state_relative_entropy=0.0)

# Rank-two per-environment responses do not ensure pooled common rank two.
J1=np.diag([1.,1.,0.]); J2=np.diag([0.,1.,1.])
check(np.linalg.matrix_rank(J1)==2 and np.linalg.matrix_rank(J2)==2,
      "separate response ranks")
check(np.linalg.matrix_rank(np.vstack([J1,J2]))==3, "pooled rank counterexample")
results["summary"] = dict(checks=checks,status="all_passed",
    meaning="Finite instances and counterexamples only; not formal proof verification.",
    numpy_version=np.__version__)
print(json.dumps(results,ensure_ascii=False,indent=2))
\end{verbatim}

\section{D.3 Historical Execution Output Retained from the Original v3.0}\label{d.3-historical-execution-output-retained-from-the-original-v3.0}

The following JSON is retained verbatim from the uploaded source text. This publication revision has not rerun that historical Python program. Evidence for the Lean compilation added at the earlier maintenance stage is listed separately in §38.8. The 2,371 numerical assertions are neither 2,371 general theorems nor a count of new proofs in this edition.

\begin{verbatim}
{
  "blocker_antichains_checked": 20,
  "resolution": {
    "surjections": 6,
    "fine_antichains": 20,
    "coarse_antichains": 6,
    "assertions": 2316
  },
  "sampling": {
    "seed": 30905,
    "samples": 160000,
    "eta": 0.01200866457583208,
    "observed_max_error": 0.005449999999999955,
    "empirical_response": [
      [
        0.14565,
        0.8509
      ],
      [
        0.15085,
        0.85295
      ],
      [
        0.89955,
        0.84455
      ],
      [
        0.89675,
        0.851
      ]
    ],
    "threshold": 0.2,
    "nonzero_gap": 0.75
  },
  "control": {
    "true_value": [
      7.771727748691103,
      8.944502617801051
    ],
    "estimated_value": [
      7.769815365791166,
      8.941100066407659
    ],
    "true_policy": [
      1,
      0
    ],
    "estimated_policy": [
      1,
      0
    ],
    "epsilon_p_actual": 0.005449999999999955,
    "value_error": 0.0034025513933926277,
    "span_bound_actual": 0.05745151456523852,
    "span_bound_confidence": 0.12659008583256628,
    "policy_regret": 0.0,
    "estimated_action_gaps": [
      0.623443681598804,
      0.17797859268051575
    ],
    "certified_actions_at_confidence_radius": [
      true,
      false
    ],
    "stationary_true": [
      0.10526315789473684,
      0.8947368421052632
    ],
    "stationary_estimated": [
      0.10558679770851953,
      0.8944132022914805
    ],
    "stationary_error": 0.0003236398137826946,
    "mixing_coefficient": 0.050000000000000044,
    "stationary_confidence_bound": 0.012640699553507453
  },
  "stable_partition": [
    [
      0,
      1
    ],
    [
      2,
      3
    ]
  ],
  "path": {
    "horizon": 8,
    "observed_tv": 0.0020629076678654185,
    "actual_defect_bound": 0.042777333714830834
  },
  "xor": {
    "I_Y_X": 0.0,
    "I_Y_E": 0.0,
    "I_Y_XE": 0.6931471805599453
  },
  "qsd": {
    "exit_rate": 1.2679491924311228,
    "q": [
      0.7320508075688773,
      0.2679491924311227
    ],
    "doob_stationary": [
      0.7886751345948128,
      0.2113248654051871
    ]
  },
  "evolution": {
    "next_frequency": [
      0.6315789473684211,
      0.368421052631579
    ],
    "selection": -0.1263157894736841,
    "transmission": -0.10526315789473686,
    "growth_rate": 1.874456264653803
  },
  "girsanov": {
    "time": 1,
    "full_noise_relative_entropy": 0.5,
    "observed_state_relative_entropy": 0.0
  },
  "summary": {
    "checks": 2371,
    "status": "all_passed",
    "meaning": "Finite instances and counterexamples only; not formal proof verification.",
    "numpy_version": "2.3.5"
  }
}
\end{verbatim}

\chapter{Appendix E. Conventions for Maintenance, Extension, and Reproduction}\label{appendix-e}

\textbf{How to use this appendix.} This appendix specifies how the monograph is modified and checked. Editorial changes, changes to mathematical proofs, and source-code changes carry different responsibilities. Reproduction starts by identifying the exact files and deciding whether an earlier receipt still applies.

\section{E.1 Maintaining a Single Main Text}\label{e.1-maintaining-a-single-main-text}

This file uses UTF-8 Markdown and standard LaTeX math delimiters, with stable anchors for chapters and theorems. Subsequent work can locate results by identifier without reparsing a PDF or assembling a theory from historical patches. The code in Appendix D can be copied and run independently; its models and results are embedded in the appendix.

Every substantive new conclusion must update its definitions or conditions, complete derivation, failure cases, dependencies, application interface, and appropriate verification together. Corrections to an existing identifier retain that identifier and record the change. A fundamentally different proposition receives a new identifier, so that one identifier does not acquire incompatible meanings.

\section{E.2 Requirements for Strengthening a Claim}\label{e.2-requirements-for-strengthening-a-claim}

\begin{itemize}
\tightlist
\item
  From D to P: a definition does not generate external facts automatically; explicit assumptions and a proof are required.
\item
  From H to E: state the experimental design, samples, controls, errors, and competing models. Internal consistency of a simulation is not evidence from nature.
\item
  From E to a universal conclusion: finite experiments do not complete this transition. The model class must be restricted or an additional theoretical mechanism supplied.
\item
  From local to global: supply the relevant compactness, global measurability, support, reachability, or boundary conditions separately.
\item
  From average to worst-case: add coverage, density, or uniformity assumptions; changing notation is insufficient.
\item
  From a finite horizon to the long term: add mixing, spectral, recurrence, or nonstationary-control conditions.
\item
  From a domain-specific result to the Core: remove the specialized assumptions, prove the result again, and check that nontrivial content remains.
\end{itemize}

\section{E.3 Suggested Order of Independent Verification}\label{e.3-suggested-order-of-independent-verification}

First check that the objects and quantifiers are fully specified, then check for circular evidential dependencies. Next construct degenerate examples, redundant examples, hidden common causes, examples outside the support, nonunique optima, and examples near bifurcations. For claims that admit finite models, use exhaustive enumeration and numerical counterexample searches. Only then assess generality, computational efficiency, and explanatory power in applications.

Every claim of completeness should specify its model class and protocol boundary. Completeness in this edition means a self-contained exposition of a single general mathematical architecture, not a declaration of finality immune to refutation.

\section{E.4 Version Delivery Record}\label{e.4-version-delivery-record}

The complete v3.0 text contains Chapters 0--31, Appendices A--E, the source genealogy, proofs, interface contracts, application procedures, and complete executable checks. Quantum-specific structure is not imposed as a foundational or application premise.

The historical v3.0 verification used synthetic data and finite models. Its original record reports 2,371 successful assertions, including 2,316 from finite enumeration of resolution structures. This does not mean that every theorem was machine-formally verified. Maintenance edition 3.1 retained the original code and output and added separately compiled formal connections and witnesses. If data, protocols, implementations, or evaluations change, regenerate the relevant certificates rather than carrying forward historical pass labels.

\chapter{Appendix F. Source Correspondences and Evidence Levels of the Maintained Edition}\label{appendix-f}

\textbf{How to use this appendix.} This appendix connects natural-language claims to formal declarations and evidence levels. Check the source type, invocation hypotheses, and build identity together. A similarly named Lean declaration does not establish that every interpretation in the prose has been formalized.

\section{F.1 Base Text, Stable Identifiers, and Reading Rules}\label{f.1-base-text-stable-identifiers-and-reading-rules}

The complete base text is the \href{source/core-v3.0/UEOT_Core_Mathematics_v3.0_Complete.md}{original v3.0 mathematical text}, comprising 183,440 bytes, with SHA-256:

\begin{verbatim}
ed00dd102157cdafe3a79c45506e86dc574d6cba65feb2df8686e63ce2726303
\end{verbatim}

The maintained edition retains all 106 stable P identifiers, the full index in Appendix B, and the original numerical program and historical output. The added M identifiers designate connecting propositions in this book; they do not enter the frozen P-ID set.

The accompanying \href{UEOT_MATHEMATICS_SOURCES.md}{complete source index} lists all 106 source results, the fully qualified names of the 163 registered source endpoints, links to the actual Lean files, and the 11/89/6 registered dispositions. This appendix supplies a reading route through the whole book and direct correspondences for the added connections. The registration sources are \href{../formalization/ueot-core/docs/compression/COMPRESSION_LEDGER.yaml}{\nolinkurl{COMPRESSION_LEDGER}} and the \href{../formalization/ueot-core/docs/CORE_COMPRESSION_THEOREM_INDEX.csv}{original index of 106 results by chapter}. The source SHA in the ledger agrees exactly with the uploaded base text above.

Some entries in the old PID\_STATUS retain earlier checkpoints and must not be mixed with later compression-ledger counts. This text uses 11/89/6 to describe the registered compression structure. A status table is not a mathematical proof, and a module count does not establish theoretical superiority.

The following abbreviations avoid repeating long prefixes. Concatenation gives the full Lean name:

\begin{longtable}[]{@{}
  >{\raggedright\arraybackslash}p{(\columnwidth - 2\tabcolsep) * \real{0.5000}}
  >{\raggedright\arraybackslash}p{(\columnwidth - 2\tabcolsep) * \real{0.5000}}@{}}
\toprule\noalign{}
\begin{minipage}[b]{\linewidth}\raggedright
Abbreviation
\end{minipage} & \begin{minipage}[b]{\linewidth}\raggedright
Full namespace
\end{minipage} \\
\midrule\noalign{}
\endhead
\bottomrule\noalign{}
\endlastfoot
V3 & UEOT.V3 \\
CMP & UEOT.V3.Compression \\
TC & \nolinkurl{UEOT.V3.Compression.TheoryCompletion} \\
MC & \nolinkurl{UEOT.V3.Compression.TheoryCompletion.MinCore} \\
UC & \nolinkurl{UEOT.V3.Compression.TheoryCompletion.UnifiedClosure} \\
SC & \nolinkurl{UEOT.V3.Compression.TheoryCompletion.ScientificClosure} \\
SV & \nolinkurl{UEOT.V3.Compression.TheoryCompletion.UnifiedClosure.SupportViabilityClosure} \\
POST & \nolinkurl{UEOT.V3.Compression.TheoryCompletion.UnifiedClosure.PosteriorSupportCoverageCheck} \\
STRICT & \nolinkurl{UEOT.V3.Compression.TheoryCompletion.UnifiedClosure.PosteriorSupportStrictWitness} \\
REG & \nolinkurl{UEOT.V3.Compression.TheoryCompletion.UnifiedClosure.RegisteredGoodJointCheck} \\
SELECT & \nolinkurl{UEOT.V3.Compression.TheoryCompletion.UnifiedClosure.RegisteredSelectionWitness} \\
VIABLE-JOINT & \nolinkurl{UEOT.V3.Compression.TheoryCompletion.UnifiedClosure.ViableRegisteredJointCheck} \\
\end{longtable}

For example, \nolinkurl{SV.initial_viability_iff_exists_zero_uncovered_policy} expands to \nolinkurl{UEOT.V3.Compression.TheoryCompletion.UnifiedClosure.SupportViabilityClosure.initial_viability_iff_exists_zero_uncovered_policy.} A filename need not equal its namespace. The module split at that maintenance stage retained the original public mathematical names; it did not create new proof endpoints by duplicating declarations.

\section{F.2 Reading and Assumption Map for the Entire Core}\label{f-core-map}

\begin{longtable}[]{@{}
  >{\raggedright\arraybackslash}p{(\columnwidth - 4\tabcolsep) * \real{0.3333}}
  >{\raggedright\arraybackslash}p{(\columnwidth - 4\tabcolsep) * \real{0.3333}}
  >{\raggedright\arraybackslash}p{(\columnwidth - 4\tabcolsep) * \real{0.3333}}@{}}
\toprule\noalign{}
\begin{minipage}[b]{\linewidth}\raggedright
Book location
\end{minipage} & \begin{minipage}[b]{\linewidth}\raggedright
Direct source entry
\end{minipage} & \begin{minipage}[b]{\linewidth}\raggedright
Conditions to retain with the theorem
\end{minipage} \\
\midrule\noalign{}
\endhead
\bottomrule\noalign{}
\endlastfoot
Chapters 1--4, 32: factors, prediction, kernel descent & \href{../formalization/ueot-core/UEOT/V3/Compression/QuotientDescent.lean}{\nolinkurl{QuotientDescent}}, \href{../formalization/ueot-core/UEOT/V3/PredictionUpdate.lean}{\nolinkurl{PredictionUpdate}} & Set-theoretic descent and measurable descent are separate; recursive Bayes retains the same given joint kernel and the zero-denominator branch. \\
Chapter 5: information and compression bounds & \href{../formalization/ueot-core/UEOT/V3/InformationPacking.lean}{\nolinkurl{InformationPacking}}, \href{../formalization/ueot-core/UEOT/V3/InformationEntropy.lean}{\nolinkurl{InformationEntropy}} & Finite packing, a common decoder, and uniform-input entropy are distinct subclaims. \\
Chapters 6--7, 11, 32: interfaces and error transport & \href{../formalization/ueot-core/UEOT/V3/Compression/TransportCertificate.lean}{\nolinkurl{TransportCertificate}}, \href{../formalization/ueot-core/UEOT/V3/PathError.lean}{PathError} & Intermediate objects in a transport chain must agree; common mass of actual kernels and abstract recurrence are proved separately. \\
Chapters 8, 35: survival fixed points & \href{../formalization/ueot-core/UEOT/V3/ViabilitySource.lean}{\nolinkurl{ViabilitySource}} & Finite \nolinkurl{states/actions}, one fixed source, and an initial distribution supported on the viability kernel; the kernel may be empty. \\
Chapters 9, 22: recovery and path information & \href{../formalization/ueot-core/UEOT/V3/RecoveryDynkin.lean}{\nolinkurl{RecoveryDynkin}}, \href{../formalization/ueot-core/UEOT/V3/GirsanovPathKL.lean}{\nolinkurl{GirsanovPathKL}} & Dynkin expectation certificates, Girsanov densities, and integrability are inputs; equality for observed KL requires further conditions. \\
Chapters 10, 21, 25, 32: invariant and positive spectral structures & \href{../formalization/ueot-core/UEOT/V3/Compression/OccupationLimitInvariance.lean}{\nolinkurl{OccupationLimitInvariance}}, \href{../formalization/ueot-core/UEOT/V3/Compression/PositiveEigenstructure.lean}{\nolinkurl{PositiveEigenstructure}} & \nolinkurl{Limit/compactness} assumptions and positive eigenstructure do not follow from normalization formulas. \\
Chapters 12--13, 34: objects, carriers, and recovery & \href{../formalization/ueot-core/UEOT/V3/Compression/TheoryCompletion/OperationalObject.lean}{\nolinkurl{OperationalObject}}, \href{../formalization/ueot-core/UEOT/V3/Compression/TheoryCompletion/AutopoiesisClosure.lean}{\nolinkurl{AutopoiesisClosure}} & Access protocols, physical destruction, instance identity, program provenance, and expected resource budgets are retained separately. \\
Chapters 14, 32, 34: scales and resolution & \href{../formalization/ueot-core/docs/compression/hierarchy/H0_HIERARCHY_INVENTORY.md}{H0 inventory}, \href{../formalization/ueot-core/UEOT/V3/Compression/TheoryCompletion/ObjectScaleCalculus.lean}{\nolinkurl{ObjectScaleCalculus}} & Eight types of hierarchy are not conflated; the current RG-flow record has no compatibility equation with actual scale maps. \\
Chapters 15, 34: parent formation and binding & \href{../formalization/ueot-core/UEOT/V3/Compression/TheoryCompletion/ScientificClosure/C4FBT.lean}{C4FBT} & The two transports and the existence of parent formation are separate inputs; formation alone does not imply binding. \\
Chapters 16, 36--38: statistics, paths, and actual coverage & \href{../formalization/ueot-core/UEOT/V3/Compression/TheoryCompletion/UnifiedClosure/StatisticalRiskCalibration.lean}{\nolinkurl{StatisticalRiskCalibration}}, \href{../formalization/ueot-core/UEOT/V3/Compression/TheoryCompletion/UnifiedClosure/DeploymentCoverage.lean}{\nolinkurl{DeploymentCoverage}} & Registered sampling laws, parameters from the same source, a common failure event, and actual occupancy cannot be replaced by a training average. \\
Chapters 17--18, 23, 26, 34: sources of purpose and gauges & \href{../formalization/ueot-core/UEOT/V3/Compression/TheoryCompletion/TeleologicalGauge.lean}{\nolinkurl{TeleologicalGauge}}, \href{../formalization/ueot-core/UEOT/V3/Compression/TheoryCompletion/UnifiedClosure/ProgramPurposeInformationBoundary.lean}{\nolinkurl{ProgramPurposeInformationBoundary}} & Existence of a representation, numerical gauge, the ordering class of purpose, and identification of its source are different questions. \\
Chapters 19, 22: continuous-time HJB & \href{../formalization/ueot-core/UEOT/V3/DiffusionHJBVerification.lean}{\nolinkurl{DiffusionHJBVerification}} & \nolinkurl{V3.DiffusionHJBVerification.Model.ItoRun} supplies actual occupancy, terminal data, the Itô expectation identity, and limit obligations. \\
Chapters 20, 28, 31, 35: organizational control quotients from the same source & \href{../formalization/ueot-core/UEOT/V3/Compression/TheoryCompletion/UnifiedClosure/OrganizationalCoreAssembly.lean}{\nolinkurl{OrganizationalCoreAssembly}}, \href{../formalization/ueot-core/UEOT/V3/Compression/TheoryCompletion/UnifiedClosure/UnifiedCoreClosureCertificate.lean}{\nolinkurl{UnifiedCoreClosureCertificate}} & Registered properties and rewards actually enter the stable partition; physical validity of the labels requires separate evidence. \\
Chapters 23, 33: quantitative defects and dual drives & \href{../formalization/ueot-core/UEOT/V3/Compression/TheoryCompletion/MinCore/RelationDefectNoGo.lean}{\nolinkurl{RelationDefectNoGo}}, \href{../formalization/ueot-core/UEOT/V3/Compression/TheoryCompletion/MinCore/CalibratedDefectComparison.lean}{\nolinkurl{CalibratedDefectComparison}}, \href{../formalization/ueot-core/UEOT/V3/Compression/TheoryCompletion/MinCore/DefectGaugeCovariance.lean}{\nolinkurl{DefectGaugeCovariance}}, \href{../formalization/ueot-core/UEOT/V3/Compression/TheoryCompletion/UnifiedClosure/DualDriveGaugeCompleteness.lean}{\nolinkurl{DualDriveGaugeCompleteness}} & The same zero relation does not determine numerical values; two-sided calibration carries units; uniqueness of drives requires independent anchors. \\
Chapters 24--26, 34: values, parent--offspring relations, and lineages & \href{../formalization/ueot-core/UEOT/V3/Compression/TheoryCompletion/LineageEvolution.lean}{\nolinkurl{LineageEvolution}} & SameObject and Offspring are distinct; the Perron bridge consumes the eigen-certificate of the same reproduction operator. \\
Chapters 27, 34: parametric beliefs and scientific interfaces & \href{../formalization/ueot-core/UEOT/V3/Compression/TheoryCompletion/ScientificClosure/C1ParametricBelief.lean}{\nolinkurl{C1ParametricBelief}}, \href{../formalization/ueot-core/UEOT/V3/Compression/TheoryCompletion/ScientificClosure/C5DualDriveTest.lean}{\nolinkurl{C5DualDriveTest}}, \href{../formalization/ueot-core/UEOT/V3/Compression/TheoryCompletion/ScientificClosure/C6PurposeMechanism.lean}{\nolinkurl{C6PurposeMechanism}} & Joint measurability requires regularity; rank gives only a necessary condition; identification of purpose requires faithfulness. \\
Chapters 29, 34: evidence from the world & \href{../formalization/ueot-core/docs/compression/theory_completion/scientific_closure/SCIENTIFIC_CLOSURE_STATUS.json}{Scientific Closure status} & Package closure does not close every domain interface; Lean compilation supplies no empirical support. \\
Chapter 35: safe optimality and repair & \href{../formalization/ueot-core/UEOT/V3/Compression/TheoryCompletion/UnifiedClosure/CoreSourceAllTimesAssembly.lean}{\nolinkurl{CoreSourceAllTimesAssembly}}, \href{../formalization/ueot-core/UEOT/V3/Compression/TheoryCompletion/UnifiedClosure/SafeOptimalControlFromViability.lean}{\nolinkurl{SafeOptimalControlFromViability}}, \href{../formalization/ueot-core/UEOT/V3/Compression/TheoryCompletion/UnifiedClosure/RepairProgramSafeControlBridge.lean}{\nolinkurl{RepairProgramSafeControlBridge}} & The original source kernel is unchanged; initial viability, restriction to safe actions, and a trusted program substrate are retained; safety does not imply optimality. \\
\end{longtable}

The paths in this table are mathematical source entry points. The accompanying index gives the complete fully qualified names for all 106 results. For the P12 status in Chapter 34, also consult the \href{../formalization/ueot-core/docs/compression/theory_completion/P12_FINAL_AUTOPOIESIS_AUDIT.md}{P12 conditional-boundary audit}. This book preserves the substantive meaning of PARTIAL there; a CLOSED work package is not rewritten as an unconditional principle of autopoiesis.

\section{F.3 Exact Endpoints of Added Chapter Propositions}\label{f.3-exact-endpoints-of-added-chapter-propositions}

The list gives the principal endpoints. Auxiliary theorems cited in each proposition explain the essential proof steps. An M identifier need not correspond to a single Lean declaration. The text explicitly distinguishes written combinations of conditional results from the hand-worked examples in the book.

\begin{longtable}[]{@{}
  >{\raggedright\arraybackslash}p{(\columnwidth - 2\tabcolsep) * \real{0.5000}}
  >{\raggedright\arraybackslash}p{(\columnwidth - 2\tabcolsep) * \real{0.5000}}@{}}
\toprule\noalign{}
\begin{minipage}[b]{\linewidth}\raggedright
Chapter proposition
\end{minipage} & \begin{minipage}[b]{\linewidth}\raggedright
Exact namespace and declaration
\end{minipage} \\
\midrule\noalign{}
\endhead
\bottomrule\noalign{}
\endlastfoot
M-32.1 & \nolinkurl{CMP.QuotientDescent.fiberCompatible_iff_existsUnique_descend} \\
M-32.2 & \nolinkurl{CMP.TransportCertificate.weighted_chain_bound} \\
M-32.3 & \nolinkurl{CMP.OccupationLimitInvariance.invariant_of_continuous_observable_residual} \\
M-32.4 & \nolinkurl{CMP.PositiveEigenstructure.qsdDoob_bundle}; \nolinkurl{CMP.PositiveEigenstructure.stochasticRightEigen_bundle} \\
M-33.1 & \nolinkurl{MC.same_relation_arbitrary_positive_scale}; \nolinkurl{MC.no_relation_only_quantitative_recovery}; \nolinkurl{MC.relationDefectPresentation_does_not_force_triangle} \\
M-33.2 & \nolinkurl{MC.DefectCalibration.zero_defect_iff}; \nolinkurl{MC.DefectCalibration.forward_budget}; \nolinkurl{MC.DefectCalibration.reverse_budget}; \nolinkurl{MC.DefectCalibration.tendsto_zero_iff} \\
M-33.3 & \nolinkurl{MC.DefectCalibration.triangle_quasi}; \nolinkurl{MC.calibrated_cross_stage_contraction} \\
M-33.4 & \nolinkurl{MC.rescale_contraction_transformed}; \nolinkurl{MC.rescale_local_budget_transformed}; \nolinkurl{MC.rescaled_weighted_chain_bound} \\
M-33.5 & \nolinkurl{UC.same_objective_iff_gauge_related}; \nolinkurl{UC.pi_identified_of_phi_anchor}; \nolinkurl{UC.phi_identified_of_pi_anchor} \\
M-34.1 & \nolinkurl{TC.incomparableBoolContract_has_no_numerical_representation}; \nolinkurl{SC.mechanismObservation_identifies_teleologicalClass}; \nolinkurl{SC.equalMechanismObservation_does_not_determine_purpose} \\
M-34.2 & \nolinkurl{SC.approximate_fbt_realized_bound}; \nolinkurl{SC.certified_formed_structural_continuation}; \nolinkurl{SC.formationTransport_does_not_imply_bindingTransport} \\
M-34.3 & \nolinkurl{TC.no_physicalOnly_synthesizer_for_samePhysical_distinctSeeds}; \nolinkurl{UC.perfect_program_decoder_iff_injective_information} \\
M-35.1 & \nolinkurl{V3.ViabilitySource.p_per_03}; \nolinkurl{UC.single_source_terminal_object_viability_all_times} \\
M-35.2 & \nolinkurl{UC.viable_restricted_greedy_preserves_source_kernel}; \nolinkurl{UC.viable_restricted_optimal_dominates_all_causal_policies} \\
M-35.3 & \nolinkurl{UC.reconstructed_program_source_kernel_agrees_with_safe_action}; \nolinkurl{UC.reconstructed_program_preserves_source_viability}; \nolinkurl{UC.reconstructed_program_joint_and_source_safety} \\
M-36.1 & \nolinkurl{UC.risk_map_mul_fiberPosterior}; \nolinkurl{UC.risk_conditionedRow_mass}; \nolinkurl{UC.risk_map_extend_eq} \\
M-36.2 & \nolinkurl{UC.coverage_tv_extend_same_initial}; \nolinkurl{UC.coverage_tv_extend_same_kernel} \\
M-36.3 & \nolinkurl{UC.coverage_projected_causalLaw_error}; \nolinkurl{UC.coverage_rowLaw_expectation}; \nolinkurl{UC.coverage_same_policy_path_error} \\
M-36.4 & \nolinkurl{UC.coverage_expected_error_le}; \nolinkurl{UC.coverage_same_policy_microscopic_failure_risk} \\
M-36.5 & \nolinkurl{POST.supported_history_actions_bound_actual_coverage}; \nolinkurl{POST.posterior_supported_actions_bound_actual_coverage} \\
M-37.1 & \nolinkurl{SV.mem_traceSupport_iff_realized}; \nolinkurl{SV.mem_coarseLaw_succ_support_iff}; \nolinkurl{SV.traceSupport_eq_fiberTerminalPosterior_support}; \nolinkurl{SV.successorSupport_eq_fiberTerminalPosterior_support} \\
M-37.2 & \nolinkurl{SV.mem_deploymentRowLaw_support_iff}; \nolinkurl{SV.zero_uncovered_iff_support_certified} \\
M-37.3 & \nolinkurl{SV.viablePolicy_zero_uncovered}; \nolinkurl{SV.supported_viable_of_zero_uncovered}; \nolinkurl{SV.initial_viability_iff_exists_zero_uncovered_policy} \\
M-37.4 & \nolinkurl{SV.finiteViabilityCheck_correct}; \nolinkurl{SV.finiteRelationViabilityCheck_correct} \\
M-37.5 & \nolinkurl{STRICT.two_positive_coarse_histories}; \nolinkurl{STRICT.posterior_supports_are_singletons}; \nolinkurl{STRICT.actual_uncovered_probability_zero}; \nolinkurl{STRICT.whole_current_fiber_has_no_common_action}; \nolinkurl{STRICT.null_history_boundary} \\
M-37.6 & \nolinkurl{SV.SafePrefixDeadend.a_complete_zero_uncovered_policy_exists}; \nolinkurl{SV.SafePrefixDeadend.deadend_policy_second_step_not_covered}; \nolinkurl{SV.SafePrefixDeadend.viable_source_with_a_covered_deadend} \\
M-38.1 & \nolinkurl{UC.selectedTrainingDeployment_failure_tower}; \nolinkurl{UC.selectedTrainingDeployment_failure_integral} \\
M-38.2 & \nolinkurl{UC.selectedTrainingDeployment_data_dependent_bound} \\
M-38.3 & \nolinkurl{UC.selectedTrainingDeployment_failure_le}; \nolinkurl{REG.registered_selected_joint_failure_of_cap_on_good} \\
M-38.4 & \nolinkurl{VIABLE-JOINT.viable_registered_selected_joint_failure_on_good} \\
Complete example in §38.6 & \nolinkurl{SELECT.registered_selection_joint_witness}; \nolinkurl{SELECT.actual_joint_failure_exact}; \nolinkurl{SELECT.allTrue_positive_probability} \\
Sharp example and fixed law in §38.7 & \nolinkurl{UC.TrainingDeploymentWitness.sharp_mixture_is_attained}; \nolinkurl{UC.statistical_fixed_policy_training_deployment_witness} \\
\end{longtable}

The sources added in Chapters 36--38 have the following single-module ownership:

\begin{longtable}[]{@{}
  >{\raggedright\arraybackslash}p{(\columnwidth - 2\tabcolsep) * \real{0.5000}}
  >{\raggedright\arraybackslash}p{(\columnwidth - 2\tabcolsep) * \real{0.5000}}@{}}
\toprule\noalign{}
\begin{minipage}[b]{\linewidth}\raggedright
Module
\end{minipage} & \begin{minipage}[b]{\linewidth}\raggedright
Principal content
\end{minipage} \\
\midrule\noalign{}
\endhead
\bottomrule\noalign{}
\endlastfoot
\href{../formalization/ueot-core/UEOT/V3/Compression/TheoryCompletion/UnifiedClosure/ObservableActionCoverage.lean}{\nolinkurl{ObservableActionCoverage}} & Certification of the whole current observation fiber; from action escape to actual coverage; budget substitution. \\
\href{../formalization/ueot-core/UEOT/V3/Compression/TheoryCompletion/UnifiedClosure/PosteriorSupportCoverage.lean}{\nolinkurl{PosteriorSupportCoverage}} & Reuse of the actual posterior; actual coverage bounds from actions supported by positive histories. \\
\href{../formalization/ueot-core/UEOT/V3/Compression/TheoryCompletion/UnifiedClosure/ReachablePolicyExtension.lean}{\nolinkurl{ReachablePolicyExtension}} & Compatibility of prefix policies and noncircular sufficient conditions for stepwise extension. \\
\href{../formalization/ueot-core/UEOT/V3/Compression/TheoryCompletion/UnifiedClosure/FiniteSupportViability.lean}{\nolinkurl{FiniteSupportViability}} & Policy-independent trace support, the actual-probability bridge, backward recursion, and both directions of the policy theorem. \\
\href{../formalization/ueot-core/UEOT/V3/Compression/TheoryCompletion/UnifiedClosure/FiniteSupportViabilityCheck.lean}{\nolinkurl{FiniteSupportViabilityCheck}} & Execution of finite relation tables and correctness with respect to the actual source. \\
\href{../formalization/ueot-core/UEOT/V3/Compression/TheoryCompletion/UnifiedClosure/ControlCoverageWitness.lean}{\nolinkurl{ControlCoverageWitness}} & Insufficiency of representatives, strict relaxation by posterior support, and dead ends in viable problems. \\
\href{../formalization/ueot-core/UEOT/V3/Compression/TheoryCompletion/UnifiedClosure/RegisteredDeploymentRisk.lean}{\nolinkurl{RegisteredDeploymentRisk}} & The canonical Good event and post-selection joint risk requiring assumptions only on that event. \\
\href{../formalization/ueot-core/UEOT/V3/Compression/TheoryCompletion/UnifiedClosure/RegisteredDeploymentRiskWitness.lean}{\nolinkurl{RegisteredDeploymentRiskWitness}} & Two registered rows, actual product sampling, nonconstant selection, and exact joint risk. \\
\href{../formalization/ueot-core/UEOT/V3/Compression/TheoryCompletion/UnifiedClosure/ViableRegisteredJoint.lean}{\nolinkurl{ViableRegisteredJoint}} & Use of the same support-viable policy to discharge a separately supplied actual-coverage input. \\
\end{longtable}

Module ownership is a lookup address and does not alter the mathematical hypotheses of accepted endpoints. Checks at the manuscript stage concern the fixed assembled sources listed below. Build and release receipts for the final production modules must be taken from the actual records accompanying the release.

\subsection{F.3.1 Precise Levels of Previously Added Material}\label{f.3.1-precise-levels-of-previously-added-material}

The following table supplements the added sections without changing the endpoints above, the 106 P-IDs, or the historical evidence.

\begin{longtable}[]{@{}
  >{\raggedright\arraybackslash}p{(\columnwidth - 4\tabcolsep) * \real{0.3333}}
  >{\raggedright\arraybackslash}p{(\columnwidth - 4\tabcolsep) * \real{0.3333}}
  >{\raggedright\arraybackslash}p{(\columnwidth - 4\tabcolsep) * \real{0.3333}}@{}}
\toprule\noalign{}
\begin{minipage}[b]{\linewidth}\raggedright
Material
\end{minipage} & \begin{minipage}[b]{\linewidth}\raggedright
Formal correspondence
\end{minipage} & \begin{minipage}[b]{\linewidth}\raggedright
Present scope
\end{minipage} \\
\midrule\noalign{}
\endhead
\bottomrule\noalign{}
\endlastfoot
M-21.1--M-21.2 & \nolinkurl{UC.ConflictParametricStationary.actual_one_step}, \nolinkurl{actual_invariant_unique}, \nolinkurl{dobrushin_zero}, cesaro\_exact, \nolinkurl{cesaro_tv_exact}, \nolinkurl{cesaro_residual}; \nolinkurl{UC.ConflictParametricCommonConsumer.original_cesaro_limit_GOA} & The same original source, for every θ∈{[}0,1{]}; from any initial law, one step gives μθ; exact N=n+1 averages and TV for the registered δ initial law; consumes the original \nolinkurl{CesaroLimitGOA.} \\
M-21.3--M-21.4 & \nolinkurl{UC.ConflictParametricQSD.actual_qsd_goa}, \nolinkurl{actual_doob_fair}, \nolinkurl{actual_doob_invariant}, \nolinkurl{actual_conditioned_terminal}, \nolinkurl{actual_conditioned_time_zero} & The actual killed kernel on the safe subtype, BiEigenData, and original GOA; any safe initial law, positive T, and θ\textless1; zero steps and θ=1 are proved separately. \\
M-21.5 / full conditioned path / infinite experiment & \nolinkurl{UC.ConflictParametricPath.full_conditioned_path_eq_zero_policy}; \nolinkurl{UC.ConflictParametricCommonConsumer.terminal_damage_probability}, \nolinkurl{two_step_return_probability}; \nolinkurl{UC.ConflictParametricInfinite.finite_record_cylinder}, \nolinkurl{actual_eventual_failure_ae} & The full state--action record comes from the original causal PMF; all finite cylinders belong to the same infinite experiment; terminal risk differs from the probability of ever failing. \\
M-21.6--M-21.7 & \nolinkurl{UC.ConflictParametricControl}; \nolinkurl{UC.ConflictParametricCommonConsumer.original_unrestricted_GOA_iff}, \nolinkurl{original_restricted_GOA_all_rows}, \nolinkurl{original_restricted_cesaro_GOA_iff} & All causal controls and actual greedy policies of the original M; the actual R retains both actions outside the domain; the full original GOA consists of all mixtures of the two closed classes, while the Cesàro law from safe δ is only μ0. \\
§35.4 / §35.7 & \nolinkurl{UC.WaitingRepairControl.healthy_values_differ}, \nolinkurl{actual_macro_optimal}; \nolinkurl{UC.ConflictOrganizationalSource.terminal_relation}, \nolinkurl{Posterior.actual_posterior_support}, \nolinkurl{two_tick_recovered} & The separate 8/5 source; original P-ALG, actual posterior, and two-tick organizational recovery for the same source with nonzero reward; physical clocks are not conflated and past failure is not erased. \\
M-36.6 & \nolinkurl{UC.FiniteObservationEscapeCheck.tvDist_atomEncode_eq}; \nolinkurl{tvDist_embedded_eq_encoded}; \nolinkurl{tvDist_eq_of_live_reference_masses} & The production interface in this edition, with one side supported on a finite carrier and with the entire escape mass retained; no global isometry for arbitrary pairs of reference laws is claimed. \\
M-36.7 & \nolinkurl{UC.PathObservationEscapeCheck.escape_path_tv}; \nolinkurl{UC.ArbitraryObservationRiskCheck.coverage_arbitrary_observation_path_error} and its event-risk endpoints & The production interface in this edition preserves the actual source, policy, complete occupancy, and original Q(F); Y has measurable singletons and F is measurable. \\
§36.8 positive-escape boundary & \nolinkurl{UC.FiniteObservationEscapeBoundaryCheck.infinite_ambient_positive_escape_boundary} & The production interface in this edition: TV from the fixed finite source is 5/12; escape mass and TV between the two original references are 1/4; TV between the encoded references is zero. \\
M-38.5 & \nolinkurl{UC.VariableRegisteredSamples.registered_row_calibration}; \nolinkurl{product_experiment_calibration} & The production interface in this edition uses actual rowwise N j; independence between rows is not required in general. \\
M-38.6 & \nolinkurl{UC.VariableRegisteredJointCheck.registered_selected_joint_failure_of_cap_on_good} & The arbitrary-Y interface in this edition places \nolinkurl{hfit/hcap/hF/henvelope/hchecked} on the same canonical Good event. \\
M-38.7 & \nolinkurl{UC.ViableVariableRegisteredJointCheck.viable_registered_selected_joint_failure_on_good} & The arbitrary-Y interface in this edition proves zero occupancy for the same \nolinkurl{viableFamilyPolicy} inside the same Good event; hF also holds on that event. \\
\end{longtable}

The original \#354 research prototype and the first finite-Y version are retained as historical records. The seven modules described here have entered the production interface. Historical checks must still be distinguished from checks of the present source. Complete signatures and the acceptance scope appear in \href{LEAN_API_SCOPE.md}{\nolinkurl{LEAN_API_SCOPE}} and the \href{../formalization/ueot-core/docs/compression/theory_completion/unified_closure/REGISTERED_INTERFACES_CLOSURE_V1.md}{interface closure document}. The historical check records in the next section are retained unchanged.

\section{F.4 Locatable Identities of Historical Lean Checks}\label{f.4-locatable-identities-of-historical-lean-checks}

The original checks remain in the project implementation evidence. They are identified below by full source hashes. Failed stages retain their original status; acceptance of a third draft does not overwrite failure of the first.

\begin{longtable}[]{@{}
  >{\raggedright\arraybackslash}p{(\columnwidth - 4\tabcolsep) * \real{0.3333}}
  >{\raggedright\arraybackslash}p{(\columnwidth - 4\tabcolsep) * \real{0.3333}}
  >{\raggedright\arraybackslash}p{(\columnwidth - 4\tabcolsep) * \real{0.3333}}@{}}
\toprule\noalign{}
\begin{minipage}[b]{\linewidth}\raggedright
Check and source
\end{minipage} & \begin{minipage}[b]{\linewidth}\raggedright
Source SHA-256
\end{minipage} & \begin{minipage}[b]{\linewidth}\raggedright
Actual result
\end{minipage} \\
\midrule\noalign{}
\endhead
\bottomrule\noalign{}
\endlastfoot
check-03 / \nolinkurl{SupportViabilityPolicyCheck_v1.lean} & \nolinkurl{ceed43d8aa3e4c6b140c0f998105d53ca2bd770b3ae29519d9f3b80562da581f} & Exit 0; 12 standard-axiom reports \\
check-06 / \nolinkurl{SupportViabilityFiniteCheck_v2.lean} & \nolinkurl{59964729e9b415ebe6f856f41fe0b9359ca1dda4318b184f5f500385610d08fc} & Exit 0; 17 standard-axiom reports; actual output true, false, true \\
check-06 / \nolinkurl{ViableRegisteredJointCheck_v1.lean} & \nolinkurl{c876dbfaefeac67c1ac01fd19fc9c6d54c818b545f4707b3d93eb4c96e3ae6d5} & Exit 0; printed endpoints use only the standard three axioms \\
check-08 / \nolinkurl{RegisteredSelectionWitnessCheck_v3.lean} & \nolinkurl{6bbdb1ee39e6d65960123bff1410182b2a1a0096a55d5e74d4f0f7afb0f666bf} & Exit 0; 12 standard-axiom reports; no warnings \\
\end{longtable}

Here ``the standard three axioms'' means a subset of propext, Classical.choice, and Quot.sound. The older assembled sources in check-03 and check-06 had observed \nolinkurl{unused-variable/simp} warnings and are therefore not described as warning-free. The complete check-04 batch and an earlier mixed \nolinkurl{strict-posterior/extension} batch included failures in other sources; §38.8 uses only the individual checks that actually passed.

Success of a historical check, a successful build of present production sources, acceptance of an actual Git candidate, a remote merge, and empirical support are distinct facts. A mathematical book must state which fact it cites; it need not turn the entire development history into a philosophical argument. Conditional mathematical conclusions and explicit counterexamples are borne first by the text; compilation evidence provides additional verification of their particular implementation.

\section{F.5 Added Correspondences for Execution from the Same Source and Measurable Prediction}\label{f.5-added-correspondences-for-execution-from-the-same-source-and-measurable-prediction}

\textbf{Current mathematical contract.} Sections 39.70--39.74 and the \href{../formalization/ueot-core/docs/compression/theory_completion/unified_closure/whole_theory_formal_closure_v1/FINAL_ACCEPTANCE.md}{whole-theory acceptance table} record the proved sources, revisions, admissible premises, and conditions for adoption in a release. The paragraphs below retain OPEN/pending statements in their historical integration order. They do not declare that the final edition retains those same obligations. The frozen 106, endpoint counts, a CLOSED individual package, or successful builds of two PDFs cannot by themselves establish completion of the whole theory.

Source responsibilities for Chapter 39 are as follows:

\begin{itemize}
\tightlist
\item
  Program execution and history: \nolinkurl{RegisterRepairSource}, \nolinkurl{RegisterRepairHistory}.
\item
  Faults and resources: \nolinkurl{RegisterFaultSource}, \nolinkurl{RegisterResourceSource}.
\item
  Maintenance of actual offspring: \nolinkurl{RegisterBirthMaintenance}.
\item
  Control and records from the same source: \nolinkurl{RegisterFaultControlSource}, \nolinkurl{RegisterFaultControlModel}, \nolinkurl{RegisterFaultJointRecord}.
\item
  Actual infinite execution, value, and cost: \nolinkurl{RegisterFaultInfiniteRecord}.
\end{itemize}

\nolinkurl{UC.RegisterFaultPurposeInformation.observed_source_identifies_actual_purpose} derives the original C6 faithfulness property from source means. The fault-free \nolinkurl{UC.RegisterProgramOptimality} and \nolinkurl{UC.RegisterJointRecord} retain their original contracts for a different source and cannot replace the faulty-record law.

The C1 source and sensor correspond to \nolinkurl{UC.MeasurablePredictiveSource} and \nolinkurl{UC.MeasurablePredictiveSensor.} \nolinkurl{UC.MeasurablePredictiveJointPath.actual_timevarying_joint_path}, \nolinkurl{actual_finite_record_event}, and \nolinkurl{actual_finite_record_lintegral} preserve the complete finite joint record of the same time-varying visible-belief policy. The added \nolinkurl{UC.MeasurablePredictiveLatentPath.actual_finite_prefix_consistency} and \nolinkurl{actual_source_timevarying_joint_path} prove conditional consistency of the recursive filter and a countable record law from the same source, using the actual transition that retains the latent state. Specific verification, independent review, and source identities appear in the \href{../formalization/ueot-core/docs/compression/theory_completion/unified_closure/actual_infinite_sources_v1/RESULTS_AND_OPEN_OBLIGATIONS.md}{source-obligation balance}, \href{../formalization/ueot-core/docs/compression/theory_completion/unified_closure/register_fault_resources_v1/RESULTS_AND_OPEN_OBLIGATIONS.md}{fault package}, and \href{../formalization/ueot-core/docs/compression/theory_completion/unified_closure/c1_joint_path_v1/RESULTS_AND_OPEN_OBLIGATIONS.md}{C1 path package}. Compilation of the two books does not alter their mathematical scope. At that source stage, the whole-theory status remained OPEN / \nolinkurl{NOT_ESTABLISHED.}

The actual endpoints added in §§39.7--39.11 are \nolinkurl{UC.MeasurablePredictiveResponseRecursion.actual_C1_response_kernel_factor}; \nolinkurl{UC.MeasurablePredictiveCanonicalSource.actual_past_future_factorization}, \nolinkurl{actual_response_ae_minimal}, and \nolinkurl{actual_canonical_recursion}; the same-record erasure, clock, and actual-birth identity consumers of \nolinkurl{UC.RegisterTimestampedBirthSource}; the actual all-time finite-dimensional Gaussian source of \nolinkurl{UC.BrownianProjectiveSource}; \nolinkurl{UC.BrownianContinuousSource.process_isBrownian}, \nolinkurl{measurable_joint_process}, and \nolinkurl{actual_ou_recovery}; and \nolinkurl{UC.MeasurablePredictiveHistoryControl.actual_action_observation_history_path}, \nolinkurl{actual_finite_action_observation_history_path}, and \nolinkurl{actual_positive_memory_histories.} Full hypotheses and source identities appear respectively in the \href{../formalization/ueot-core/docs/compression/theory_completion/unified_closure/predictive_genealogy_source_v1/RESULTS_AND_OPEN_OBLIGATIONS.md}{general-response and timestamp package}, \href{../formalization/ueot-core/docs/compression/theory_completion/unified_closure/canonical_brownian_source_v1/RESULTS_AND_OPEN_OBLIGATIONS.md}{canonical-prediction and all-time Gaussian package}, and \href{../formalization/ueot-core/docs/compression/theory_completion/unified_closure/continuous_history_source_v1/RESULTS_AND_OPEN_OBLIGATIONS.md}{continuous-source and history-control package}. Added sources do not recount the frozen result index or its registered dispositions. Section 39.12 further proves minimality for arbitrary actual historical pasts and all fresh adaptive futures in the declared protocol class. Broader protocols remain subject to §4.3. General Itô/Girsanov results and a global assembly do not follow automatically from these local sources.

Section 39.12 corresponds to \nolinkurl{actual_past_future_factorization}, \nolinkurl{actual_response_adaptive_ae_minimal}, and \nolinkurl{futureKernel_response_memory_factor} in \nolinkurl{UC.MeasurablePredictiveHistoryCanonical.} Section 39.13 corresponds to \nolinkurl{UC.BrownianMartingaleSource.} The input representation, original output, and full predictable domain in §§39.14--39.15 correspond to \nolinkurl{UC.BrownianIntegralRepresentation}, \nolinkurl{UC.BrownianIntegralClosure.integral_returns_literal}, and \nolinkurl{UC.BrownianPredictableDensity.predictableIntegral_returns_literal.} Section 39.16 corresponds to \nolinkurl{UC.RegisterFaultFinitePurposeProbe} and \nolinkurl{UC.RegisterFaultFinitePurposeInfinite.returned_positive_actual_infinite_record.} Complete types and obligations appear in the \href{../formalization/ueot-core/docs/compression/theory_completion/unified_closure/adaptive_stochastic_sources_v1/RESULTS_AND_OPEN_OBLIGATIONS.md}{\nolinkurl{adaptive/conditional-noise} package}, \href{../formalization/ueot-core/docs/compression/theory_completion/unified_closure/history_integral_source_v1/RESULTS_AND_OPEN_OBLIGATIONS.md}{\nolinkurl{actual-past/integral-representation} package}, and \href{../formalization/ueot-core/docs/compression/theory_completion/unified_closure/predictable_finite_purpose_v1/RESULTS_AND_OPEN_OBLIGATIONS.md}{\nolinkurl{full-predictable/finite-purpose-certificate} package}. At that stage, the whole theory remained OPEN / \nolinkurl{NOT_ESTABLISHED.}

Sections 39.17--39.19 correspond to \nolinkurl{UC.BrownianIntegralContinuousVersion}, \nolinkurl{UC.RegisterFaultFeedbackResource}, \nolinkurl{UC.BrownianQuadraticVariation}, \nolinkurl{UC.BrownianIntegralPastLocalization}, and \nolinkurl{UC.BrownianIntegralCompensator}; §§39.20--39.21 correspond to \nolinkurl{UC.RegisterPopulationHistorySource} and \nolinkurl{UC.RegisterPopulationTypePrice.} Direct obligations and exact types appear in \href{../formalization/ueot-core/docs/compression/theory_completion/unified_closure/continuous_resource_source_v1/RESULTS_AND_OPEN_OBLIGATIONS.md}{continuous \nolinkurl{integrals/feedback} resources} and \href{../formalization/ueot-core/docs/compression/theory_completion/unified_closure/compensator_population_source_v1/RESULTS_AND_OPEN_OBLIGATIONS.md}{actual second-order \nolinkurl{energy/population} sources}. A compensator martingale is not a general \nolinkurl{partition-QV/It}ô theorem; an open resource interface is not a closed source, and an offspring census is not the entire terminal table. At that stage, the whole theory remained OPEN / \nolinkurl{NOT_ESTABLISHED.}

Sections 39.22--39.25 use the integrated UC sources for full scalar partition QV, \nolinkurl{cross/weighted} integrals, Euclidean \nolinkurl{product/vector-closure} integrals, funded training, and \nolinkurl{common-population/log} learning. Actual public consumers, exact premises, and the remaining total obligations appear in \href{../formalization/ueot-core/docs/compression/theory_completion/unified_closure/full_qv_vector_training_source_v1/RESULTS_AND_OPEN_OBLIGATIONS.md}{full \nolinkurl{QV/training/discovery} sources} and \href{../formalization/ueot-core/docs/compression/theory_completion/unified_closure/vector_weighted_population_source_v1/RESULTS_AND_OPEN_OBLIGATIONS.md}{vector \nolinkurl{representation/weighted} \nolinkurl{tables/common-population} sources}. Sections 39.26--39.29 use the integrated \nolinkurl{UC.BrownianScalarItoSource}, \nolinkurl{UC.BrownianBoundedC2DynkinSource}, \nolinkurl{UC.OUSemigroupSource}, \nolinkurl{UC.BrownianVectorIntegralL2TimeContinuity}, \nolinkurl{UC.BrownianConstantGirsanovSource}, and \nolinkurl{UC.DualDriveFiniteNoisyCalibration.} Original sources, actual returned values, independent reviews, and boundaries appear in \href{../formalization/ueot-core/docs/compression/theory_completion/unified_closure/ito_vector_calibration_source_v1/RESULTS_AND_OPEN_OBLIGATIONS.md}{same-input \nolinkurl{calculus/full} vector \nolinkurl{domain/calibration} sources}. Sections 39.30--39.33 further use the integrated \#414/\#415: \nolinkurl{UC.BrownianPastStepGirsanovSource}, \nolinkurl{UC.BrownianFiniteStepShiftedSource}, \nolinkurl{UC.BrownianVectorIntegralContinuousVersion}, \nolinkurl{UC.RegisterActuatorPopulationScale}, \nolinkurl{UC.RegisterPopulationFundingObstruction}, and \nolinkurl{UC.DualDriveFiniteCalibrationCoverage.} Actual public consumers, independent reviews, and precise premises appear in \href{../formalization/ueot-core/docs/compression/theory_completion/unified_closure/past_feedback_scale_source_v1/RESULTS_AND_OPEN_OBLIGATIONS.md}{past \nolinkurl{feedback/actual} \nolinkurl{scales/acceptance}} and \href{../formalization/ueot-core/docs/compression/theory_completion/unified_closure/table_continuous_funding_source_v1/RESULTS_AND_OPEN_OBLIGATIONS.md}{actual table \nolinkurl{paths/continuous} \nolinkurl{vectors/sustained} funding}. This paragraph records the historical sources of \#414/\#415; subsequently integrated bounded predictable control, adaptive calibration, and remaining carrier-algorithm obligations are recorded below. At that stage, the whole theory remained OPEN / \nolinkurl{NOT_ESTABLISHED.}

Sections 39.34--39.38 use the integrated \#417/\#418. Actual bounded control, adaptive caches, policies for new directions, and \nolinkurl{carrier/physical} deletion appear in \href{../formalization/ueot-core/docs/compression/theory_completion/unified_closure/bounded_control_carrier_source_v1/RESULTS_AND_OPEN_OBLIGATIONS.md}{bounded-control and algorithm sources}. The constructed Q integral and original terminal KL, QV for the fixed complete joint input, and \nolinkurl{submitted/returned} instructions correspond to \nolinkurl{UC.BrownianOwnQIntegralShift}, \nolinkurl{UC.BrownianBoundedTerminalKLSource}, \nolinkurl{UC.BrownianVectorPartitionQV}, and \nolinkurl{UC.RegisterSemanticInstructionScale}; see \href{../formalization/ueot-core/docs/compression/theory_completion/unified_closure/q_integral_vector_qv_program_source_v1/RESULTS_AND_OPEN_OBLIGATIONS.md}{actual-integral and program sources}. This paragraph records the historical sources of \#417/\#418; subsequent adoption is recorded below. At that stage, the whole theory remained OPEN / \nolinkurl{NOT_ESTABLISHED.}

Sections 39.39--39.44 use the integrated \#420/\#421. \href{../formalization/ueot-core/docs/compression/theory_completion/unified_closure/fixed_q_run_stateweighted_source_v1/RESULTS_AND_OPEN_OBLIGATIONS.md}{Actual original \nolinkurl{control/weights/persistent} programs} corresponds to the original ItoRun/P-CTL, actual START weights at the same state, and the complete countable law of the same AST. \href{../formalization/ueot-core/docs/compression/theory_completion/unified_closure/state_ito_q_martingale_recovery_source_v1/RESULTS_AND_OPEN_OBLIGATIONS.md}{State calculus/Q \nolinkurl{martingales/recovery}} corresponds to \nolinkurl{UC.BrownianFiniteStateItoSource}, \nolinkurl{UC.BrownianOwnQIntegralMartingale}, and \nolinkurl{UC.RegisterInstructionRecoveryScale.} Each package determines the actual source, public consumers, independent review, and boundaries. Successors to these sources now appear in §§39.45--39.48. At the recorded stage, the whole theory remained OPEN / \nolinkurl{NOT_ESTABLISHED.}

Sections 39.45--39.48 correspond to \nolinkurl{UC.BrownianFiniteStateDriftItoSource}, \nolinkurl{UC.PredictableDriftPointSamplingCounterexample}, \nolinkurl{UC.OUPhysicalRecoverySource}, \nolinkurl{UC.RegisterInstructionHomeostasis/PurposeScale/GODScale/GOASource}, and \nolinkurl{UC.BrownianOwnQPredictableDensity/FullIntegralSource/UnboundedIntervalSource.} Actual sources, public consumers, independent reviews, and quantifiers appear in \href{../formalization/ueot-core/docs/compression/theory_completion/unified_closure/drift_ou_instruction_qfull_source_v1/RESULTS_AND_OPEN_OBLIGATIONS.md}{physical sources and actual-program obligations}. Rotating bodies, actual child objects, purpose for all prices, finite resources, finite-block transport, and dynamic selection from original data now appear in §§39.49--39.54. At the recorded stage, the whole theory remained OPEN / \nolinkurl{NOT_ESTABLISHED.}

Sections 39.49--39.54 correspond to the integrated \#425/\#426 modules \nolinkurl{UC.RegisterRotatingBodySource/Returns/P10}, \nolinkurl{UC.RegisterRotatingChildRepair}, \nolinkurl{UC.RegisterFaultDynamicPrice/Program}, \nolinkurl{UC.RegisterFaultFiniteResourceMaintenance/Witness/Boundary}, \nolinkurl{UC.RegisterBlockFormationTransport}, and \nolinkurl{UC.RegisterFaultDynamicProbe/Infinite.} Exact sources, independent public consumers, and conditions appear in \href{../formalization/ueot-core/docs/compression/theory_completion/unified_closure/rotating_body_dynamic_price_source_v1/RESULTS_AND_OPEN_OBLIGATIONS.md}{actual time-varying bodies and purpose for all prices} and \href{../formalization/ueot-core/docs/compression/theory_completion/unified_closure/resource_block_dynamic_data_source_v1/RESULTS_AND_OPEN_OBLIGATIONS.md}{\nolinkurl{finite-resource/block/data-dynamic} sources}. Successors concerning funding, countable blocks, and actual end-of-round hitting now appear in §§39.55--39.57. \nolinkurl{Native/primitive} full-law connections now appear in §§39.58--39.59. Complete P12, the P9 common consumer, and population sources retain their separate responsibilities. At the recorded stage, the whole theory remained OPEN / \nolinkurl{NOT_ESTABLISHED.}

Sections 39.55--39.57 correspond to \nolinkurl{UC.RegisterFaultDynamicProbeResource}, \nolinkurl{UC.RegisterBlockInfiniteTransport/ReturnTransport}, and \nolinkurl{UC.RegisterBodyCanonicalPath/HittingSource.} Exact sources, public consumers, independent reviews, and applicability boundaries appear in \href{../formalization/ueot-core/docs/compression/theory_completion/unified_closure/funded_countable_block_source_v1/RESULTS_AND_OPEN_OBLIGATIONS.md}{same-source \nolinkurl{funding/countable} \nolinkurl{blocks/actual} hitting sources}. The original finite strict-reproduction impossibility remains. Successors for Lifecycle, the complete native program law, and finite linear-rank existence now appear in §§39.58--39.60; the original P9 common terminal and actual-population successors still bear separate source responsibilities. At the recorded stage, the whole theory remained OPEN / \nolinkurl{NOT_ESTABLISHED.}

Sections 39.58--39.60 correspond to \nolinkurl{UC.RegisterFaultDynamicProbeLifecycle}, \nolinkurl{UC.RegisterInstructionNativeSource/RegisterFaultDynamicProbeNativeHomeostasis}, and \nolinkurl{UC.finite_rank_le_two_iff_linear_factorization}, \nolinkurl{UC.rank_two_factorizations_unique_gauge}, and \nolinkurl{UC.rank_two_registered_units_identify} in \nolinkurl{DualDriveFiniteRankFactorization.} Exact source code, axioms, independent public consumers, and boundaries appear in \href{../formalization/ueot-core/docs/compression/theory_completion/unified_closure/native_lifecycle_rank_source_v1/RESULTS_AND_OPEN_OBLIGATIONS.md}{\nolinkurl{actual-parent/native-full-lifecycle/finite-rank} sources}. Sections 39.61--39.64 further use sources for the actual six properties, finite descendants and Price, and same-tick organization. At the recorded stage, the whole theory remained OPEN / \nolinkurl{NOT_ESTABLISHED.}

Sections 39.61--39.64 correspond to \nolinkurl{UC.RegisterFaultDynamicProbeClocked/RegisterFaultDynamicProbeP9Terminal}, \nolinkurl{RegisterFaultDynamicProbePopulation/RegisterFaultDynamicProbeHistoryPrice}, and \nolinkurl{RegisterFaultNativeOrganization.} Exact source code, independent public consumers, types, and boundaries appear in \href{../formalization/ueot-core/docs/compression/theory_completion/unified_closure/p9_lineage_organization_source_v1/RESULTS_AND_OPEN_OBLIGATIONS.md}{\nolinkurl{actual-six-property/finite-descendant-Price/same-tick-organization} sources}. Sections 39.65--39.67 add registered transport, actual parent tags, and full finite write costs. At that source stage, a common wallet including read costs and a common finite-generation consumer had not yet entered the edition, and the whole theory remained OPEN / \nolinkurl{NOT_ESTABLISHED.}

Sections 39.65--39.67 correspond to \nolinkurl{UC.RegisterPopulationOwnerClockTransport}, \nolinkurl{RegisterFaultDynamicProbeParentDiscovery}, and \nolinkurl{RegisterPopulationFiniteEpisodeResource.} Exact source code, independent public consumers, types, and boundaries appear in \href{../formalization/ueot-core/docs/compression/theory_completion/unified_closure/owner_discovery_resource_source_v1/RESULTS_AND_OPEN_OBLIGATIONS.md}{actual \nolinkurl{registration/birth-parent} \nolinkurl{tags/finite-funding} sources}. Sections 39.68--39.69 further give the \nolinkurl{same-wallet/same-first-birth} parent role and actual canonical predictive control; see \href{../formalization/ueot-core/docs/compression/theory_completion/unified_closure/finite_predictive_source_v1/RESULTS_AND_OPEN_OBLIGATIONS.md}{same-wallet common roles/actual canonical-prediction sources}. At that source stage, funded continuation of the actual end-of-episode parent and general common Price still awaited source connections, and the whole theory remained OPEN / \nolinkurl{NOT_ESTABLISHED.}

Sections 39.68--39.69 correspond to \nolinkurl{UC.RegisterFiniteGenerationCommonConsumer} and \nolinkurl{RegisterPredictiveSource.} Exact sources, independent public consumers, types, and scope appear in \href{../formalization/ueot-core/docs/compression/theory_completion/unified_closure/finite_predictive_source_v1/RESULTS_AND_OPEN_OBLIGATIONS.md}{same-wallet common roles/actual canonical-prediction sources}. The checks of these two sources retain their precise historical scope. Actual end-parent continuation and quantitative common roles now appear in §39.70; the overall verdict and release conditions appear in §39.71.

Sections 39.70--39.71 use \nolinkurl{UC.RegisterPopulationCurrentParentContinuation}, the supplemental public consumer for the original P5 \nolinkurl{mean/expected-resource} interface, and the final independent cross-interface review. Exact declarations, standard axioms, and bindings to actual inputs/logs appear in the \href{../formalization/ueot-core/docs/compression/theory_completion/unified_closure/finite_predictive_source_v1/SOURCE_VALIDATION.json}{three-source common-source record}.

Sections 39.72--39.74 correspond to \nolinkurl{UC.PrimitivePerronClosure.certificateOfPrimitive}, \nolinkurl{actual_p_evo_03}, and \nolinkurl{actual_root_simple}, together with \nolinkurl{UC.CompactSpectralSemigroupSource.actual_principal_expansion_and_remainder} and \nolinkurl{UC.CompactKilledSpectralBinding.ofActualSource} and \nolinkurl{actual_p_qsd_04.} Complete matrix and actual-semigroup sources, independent common instances, and precise semantic boundaries appear in the \href{../formalization/ueot-core/docs/compression/theory_completion/unified_closure/general_perron_compact_spectral_source_v1/SOURCE_VALIDATION.json}{general-source validation record}. These factories supply the internal outputs of the original interfaces. They neither modify the frozen P-IDs nor measure theoretical completeness by counts.

For Sections 39.75--39.81, the original actual source, finite-horizon optimality, and parameter boundaries correspond, respectively, to \nolinkurl{UC.RegisterJointDecisionSource}, \nolinkurl{UC.RegisterJointDecisionOptimality}, and \nolinkurl{UC.RegisterJointDecisionBoundary.} The source, returned action table, complete comparison class of real-randomized strategies on permitted histories, actual single-wallet prepayment, complete rational replay, and independent review are recorded in the \href{../formalization/ueot-core/docs/compression/theory_completion/unified_closure/joint_decision_closure_v1/README.md}{joint decision delivery package}. This contract retains the same persistent parameter, explicit action and receipt permissions, a finite zero-income wallet, and a terminal evaluation fixed in advance. Finite-horizon Bayes optimality is not replaced by an infinite-horizon or unknown-mechanism minimax conclusion.

Sections 39.82--39.88 are adopted through the public entry \nolinkurl{UC.RegisterDecisionRepresentation.} \nolinkurl{RegisterDecisionObserver/Source} and \nolinkurl{CountPosterior/History} supply admissible-history statistics; \nolinkurl{RepresentationModel} supplies the joint refinement derived from the source; \nolinkurl{FiniteJointDecisionDAG/Generator} supplies full-action checking and shared artifacts; \nolinkurl{RepresentationNative/Payment} supplies direct native execution and a literal prefix ledger; \nolinkurl{Bounds/Witnesses} supplies actual array bounds and positive and negative boundary cases. Exact inputs, outputs, sources, axioms, and execution records appear in the \href{../formalization/ueot-core/docs/compression/theory_completion/unified_closure/decision_representation_closure_v1/README.md}{decision-representation delivery package}.

Sections 39.89--39.95 use the public entry \nolinkurl{UC.RegisterMechanismDecision.} The new native S/G source and joint-history filtering, all-policy experimental coupling, finite fixed-mechanism regret certificates, the same privately randomized returned program, and the actual prefix ledger jointly support the conclusions. Source endpoints, actual inputs and outputs, and independent acceptance appear in the \href{../formalization/ueot-core/docs/compression/theory_completion/unified_closure/mechanism_decision_closure_v1/README.md}{mechanism-decision delivery package}.

\backmatter
\end{document}
